\documentclass[11pt,letterpaper]{article}
\usepackage[T1]{fontenc}
\usepackage[utf8]{inputenc}
\usepackage[margin=.96in,top=.87in,bottom=.95in,headheight=15pt,headsep=17pt]{geometry}
\usepackage{amsmath,amssymb,amsthm,mathtools,mathpazo,microtype}
\usepackage{array,booktabs,tabularx,longtable,fancyhdr,needspace,xcolor}
\usepackage[backend=biber,style=numeric,sorting=nyt,sortcites=true,giveninits=true,maxbibnames=99]{biblatex}
\addbibresource{references.bib}
\definecolor{inklink}{RGB}{20,78,89}
\usepackage[bookmarksnumbered,hypertexnames=false,colorlinks,linkcolor=black,urlcolor=inklink,citecolor=inklink,pdfencoding=auto,psdextra]{hyperref}
\usepackage{bookmark}
\hypersetup{pdftitle={The Alpöge–Claude Argument for the Riemann Hypothesis},pdfauthor={Jimmie Adriazola},pdfsubject={Conditional reconstruction of the argument of Levent Alpöge with Claude}}
\pagestyle{fancy}\fancyhf{}
\fancyhead[L]{\small\itshape The Alp\"oge--Claude Argument for the Riemann Hypothesis}
\fancyhead[R]{\small\thepage}\renewcommand{\headrulewidth}{.2pt}
\newtheorem{theorem}{Theorem}[section]
\newtheorem{proposition}[theorem]{Proposition}
\newtheorem{lemma}[theorem]{Lemma}
\newtheorem{corollary}[theorem]{Corollary}
\newtheorem{chatAtheorem}[theorem]{Theorem}
\newtheorem{chatAlemma}[theorem]{Lemma}
\newtheorem{chatAproposition}[theorem]{Proposition}
\theoremstyle{definition}
\newtheorem{scalehypothesis}{Hypothesis}
\renewcommand{\thescalehypothesis}{S}
\numberwithin{equation}{section}
\newcommand{\Tot}{\mathcal T}
\newcommand{\QQ}{\mathbb Q}
\newcommand{\RR}{\mathbb R}
\newcommand{\CC}{\mathbb C}
\newcommand{\src}[2]{\href{https://github.com/lsernau/grh-lean/blob/72fead102f663d5c7a407964a13dc0cf9d977207/grh/src/#1.lean}{#2}}
\newcommand{\lineSrc}[3]{\href{https://github.com/lsernau/grh-lean/blob/72fead102f663d5c7a407964a13dc0cf9d977207/grh/src/#1.lean\#L#2}{#3}}
\let\ChatAFile\src
\let\srcline\lineSrc
\setlength{\emergencystretch}{2em}
\setcounter{tocdepth}{1}
\setcounter{secnumdepth}{2}
\allowdisplaybreaks[1]
\raggedbottom

\title{\bfseries The Alp\"oge--Claude Argument for the Riemann Hypothesis}
\author{Jimmie Adriazola$^{1,2}$\thanks{\texttt{jimmie.adriazola@asu.edu}}\\[6pt]
\small $^1$School of Mathematical and Statistical Sciences, Arizona State University,\\[-2pt]
\small Tempe, AZ 85287, USA\\[3pt]
\small $^2$Simon A. Levin Mathematical, Computational and Modeling Sciences Center,\\[-2pt]
\small Arizona State University, Tempe, AZ 85287, USA}
\date{November 26, 2026}

\begin{document}
\maketitle
\begin{abstract}
This manuscript reconstructs the determinant argument of Levent Alp\"oge and Claude
from the Lean~4 refactor by Luke Sernau, his colleagues at Anthropic, and Claude.
Alp\"oge takes his starting point from Stepanov, who bounded the number of rational points on a curve over a finite
field by constructing an auxiliary function with a zero of high multiplicity at each point.
Alp\"oge sought an arithmetic analogue of that multiplicity over the integers.
Successive theta lifts provide that analogue through increasing divisibility at inert primes.
For a sum of length \(x\), Claude proposed a chain of about \(\sqrt{x}\) lifts.
An upper bound for the Gram determinant of translates of short integral forms
and a lower bound from completing them to a lattice basis give a one-sided
square-root bound for the smoothed prime sum.
For each putative zero to the right of the critical line, Alp\"oge chooses a
prime \(d\) so that twisting the Dirichlet characters by the Legendre symbol
\((\cdot/d)\) preserves a pole at that zero in the transform of the sum,
which Landau's theorem rules out. The functional equation for primitive Dirichlet \(L\)-functions reflects any zero left of the critical line in the critical strip to a zero on the right for the conjugate character, so the zero-free half-plane gives the generalized Riemann hypothesis (GRH). This working manuscript gives a written proof of the conditional reduction to the determinant criterion, and work continues on translating the remaining Lean arguments into a proof that the required arithmetic hypotheses hold simultaneously for the theta lifts at each sufficiently large cutoff.
\end{abstract}

\noindent\textbf{2020 Mathematics Subject Classification.} 11M26, 11M06, 11G05, 11F27, 11F70, 11F80, 14G40.

\section{Introduction}\label{sec:p1intro}
Alp\"oge traces the construction to Stepanov's method~\cite{stepanov} and credits Claude with the choice of a chain of length about \(\sqrt{x}\) for a von Mangoldt sum of length \(x\) \cite{alpo}. Zhang describes the accumulation of divisibility along the chain \cite{zhang}. Slijep\v{c}evi\'c discusses the slope inequality in relation to Stepanov's auxiliary function \cite{slijepcevic}, and Rubinstein gives the reduction to a weighted prime-power estimate for the zeta function \cite{rubinstein}.

Fix a Dirichlet character \(\chi\) and suppose that \(L(\rho,\chi)=0\) with \(1/2<\Re\rho<1\).
The logarithmic derivative has a pole at \(\rho\), but the coefficients of its prime sum have no fixed sign.
The determinant counts prime powers \(n\equiv1\pmod{q_0}\) with nonnegative weight \(1-(n/d)\), where \(d\) is an auxiliary prime and \((n/d)\) is the Legendre symbol.
By character orthogonality the corresponding Dirichlet series is a sum of differences of logarithmic derivatives.
A zero of a quadratic twist can cancel the contribution at \(\rho\). We choose \(d\) after fixing \(\rho\) to exclude this cancellation.

Let \(q_0\) be a positive multiple of \(18480\) and \(d\) a prime congruent to \(q_0/2-1\pmod{q_0}\). Call \((q_0,d)\) permissible.
For an admissible datum \(D\) of primary cutoff \(x\) (\S\ref{sz:setup}), let \(m\) be the length of its chain of lifts, \(b\) its integral weight, and \(f\) the number of infinite places of its CM field. Put \(a=xb\) and let \(\Tot_D\) be the signed excess defined in \eqref{eq:total}. The determinant bound is
\begin{equation}
                m\Tot_D\le fa(\log b)^5. \tag{D}\label{eq:det}
\end{equation}

\begin{scalehypothesis}[Determinant bound at each primary cutoff]\label{hyp:scale}
For every permissible pair \((q_0,d)\) and every \(0<\eta<1/100\), there exists \(X_0=X_0(q_0,d,\eta)\) such that for every integer \(x\ge X_0\), at least one datum in the family constructed in Theorem~\ref{thm:p1-data} at that \(x\) and \(\eta\) satisfies \eqref{eq:det}.
\end{scalehypothesis}

\begin{theorem}[Analytic implication]\label{thm:main}
Assume Hypothesis~\ref{hyp:scale}. For each positive modulus \(q\) and each Dirichlet character \(\chi\bmod q\), \(L(s,\chi)\) is nonzero for \(\Re s>1/2\), with the principal pole at \(1\) excepted. Every zero in \(0<\Re s<1\) has real part \(1/2\). Every zero of \(\zeta(s)\) outside the negative even integers has real part \(1/2\).
\end{theorem}

By Proposition~\ref{prop:count}, the signed excess contains one smoothed remainder at the primary cutoff \(x\) and two at the secondary cutoff \(x_1\), of moduli \(q_0\) and \(1\). The prime number theorem in progressions controls the secondary cutoff. At the primary cutoff, division by \(m\ge\sqrt{x}\) gives the one-sided bound. The prime \(d\) chosen for \(\rho\) also gives the sign of the real residues required by Landau's theorem \cite{MV,widder}. The threshold \(X_0\) may depend on this \(d\), so the argument requires no uniform estimate in the height of \(\rho\). Conversely, under GRH every admissible datum at the constructed scales satisfies \(m|\Tot_D|\le fa(\log b)^5\) for all sufficiently large cutoffs at fixed \((q_0,d,\eta)\) (Corollary~\ref{cor:scale-equivalence}).

To obtain \eqref{eq:det} arithmetically, select an independent family from translates of short integral forms. The upper and lower determinant bounds for this family contain the same main height. Theorem~\ref{sz:conclusion} cancels it under the hypotheses of Sections~\ref{A:section}, \ref{B:frontier}, and~\ref{c3:frontier}. If a single datum satisfies all these hypotheses at each required cutoff, Corollary~\ref{cor:construction-grh} gives Hypothesis~\ref{hyp:scale}. Appendix~\ref{sz:register} records the estimates still assumed.

The cost of completing the family determines the precision needed in its two moment estimates. Under the auxiliary-column bound and the other hypotheses of Proposition~\ref{prop:weakerprecision}, normalized errors of order \(b^{-1}(\log b)^{-A}\) for the same map suffice when \(A>202/\eta\).
Sections~\ref{sec:unitrow-frame}--\ref{sec:other-witnesses} construct the frame and derive the local Taylor witnesses from normalized curve identities in each of the three splitting cases. Section~\ref{sec:odd-cocycles-reduction} gives a single integer bounding the odd Tate-twist cocycles uniformly in the torsion level. For the simultaneous construction of the twisted Selmer polynomials, the representation, its scalar inertia, and its local conditions are fixed before the auxiliary primes are chosen (\S\ref{sec:selmer-polynomial-descent} and \S\ref{sec:terminal-selmer}). Verifying these properties for the Hecke constituents of the theta lifts remains part of hypothesis R in \S\ref{c3:frontier}.

\clearpage
\tableofcontents
\clearpage
\part{From the determinant estimate to Dirichlet zeros}\label{part:analytic}
\section{The excess measured by the determinant}\label{sec:excess}
Fix a permissible pair $(q_0,d)$ as in Hypothesis~\ref{hyp:scale}. Thus $q_0$ is a positive multiple of $18480$, and $d$ is a prime congruent to $q_0/2-1\pmod {q_0}$. In particular $d\equiv7\pmod 8$ and $(q_0,d)=1$.
For a modulus $q\ge1$ coprime to $d$, write $\chi_d(n)=(n/d)$ and define
\begin{align}
 \lambda_{d,q}(n)&=\mathbf 1_{n\equiv1\bmod q}\Lambda(n)(1-\chi_d(n)),\notag\\
 R_{d,q}(y)&=\sum_{1\le n\le y}\lambda_{d,q}(n)(1/n-1/y)
             -\frac{\log y}{\varphi(q)}-c_{d,q},\qquad y\ge1.\label{eq:R}
\end{align}
Take $c_{d,q}=G_{d,q}(1)$, where $G_{d,q}$ is the logarithmic derivative average in \eqref{eq:p1G}. This value is real and is the unique choice for which $R_{d,q}(y)\to0$ (Lemma~\ref{lem:pnt}). For $1\le z\le a$, put $v_z(n)=(a/n-a/z)_+$. Then
\begin{equation}
 \sum_{n\le a}\lambda_{d,q}(n)v_z(n)
 =a\left[R_{d,q}(z)+\frac{\log z}{\varphi(q)}+c_{d,q}\right].
 \label{eq:scaled}
\end{equation}
Here $w_+=\max(w,0)$, and sums in $n$ start at $1$.

\subsection{Fixing the centering constant}
Let $Q\ge1$ and let $d$ be an odd prime coprime to $Q$. Write $L_\eta(s)=L(s,\eta)$, with characters extended by zero on nonunits and the Euler product taken at the indicated level. By classical continuation and nonvanishing for Dirichlet $L$-functions, the only singularity in $\Re s\ge1$ is the pole at $1$ for the principal character, and there are no zeros there. To remove the pole, set
\[
 D_\eta(s)=
 \begin{cases}
 (s-1)L(s,\eta),&\eta\text{ principal},\\
 L(s,\eta),&\eta\text{ nonprincipal}.
 \end{cases}
\]
Each $D_\eta$ is entire and $D_\eta(1)\ne0$. Each twist $\eta\chi_d$ is nonprincipal at level $Qd$. By the Chinese remainder theorem it takes the value $-1$ on an integer congruent to $1$ modulo $Q$ and to a quadratic nonresidue modulo $d$.


With the weights in \eqref{eq:R}, put
\begin{equation}\label{eq:p1G}
 G_{d,Q}(s)=\frac1{\varphi(Q)}\left[
  \sum_{\eta\bmod Q}\left(\frac{D'_{\eta\chi_d}}{D_{\eta\chi_d}}
                             -\frac{D'_\eta}{D_\eta}\right)(s)-1\right].
\end{equation}
For $\Re s>1$, orthogonality and the absolutely convergent logarithmic derivatives give
\begin{align}
 \sum_{n\ge1}\lambda_{d,Q}(n)n^{-s}
 &=\frac1{\varphi(Q)}\sum_{\eta\bmod Q}
       \left(\frac{L'_{\eta\chi_d}}{L_{\eta\chi_d}}-
             \frac{L'_\eta}{L_\eta}\right)(s)\notag\\
 &=G_{d,Q}(s)+\frac{s}{\varphi(Q)(s-1)}.\label{eq:p1Dir}
\end{align}
For the principal character, $D'_1/D_1=L'_1/L_1+1/(s-1)$. The constant $-1$ in \eqref{eq:p1G} supplies the remaining term in $s/(s-1)$. All logarithmic derivatives in \eqref{eq:p1G} are regular at $1$, and \eqref{eq:p1Dir} is real on $(1,\infty)$. Passing to the limit gives $c_{d,Q}=\Re G_{d,Q}(1)=G_{d,Q}(1)$.


\subsection{Counting the places}
In an admissible datum the CM field is $E=\QQ(\sqrt{-d},\zeta_{q_1})$, where $q_0\mid q_1$ and $q_0,q_1$ have the same prime divisors. The infinite places of $E$ are identified with $\Phi=(\mathbb Z/q_1)^\times$ by choosing the representatives extending a fixed embedding of $\QQ(\sqrt{-d})$. Denote by $\Phi_a$ the full fibre of $1\bmod q_0$ and by $\Phi_b$ the prescribed sparse subset in the fibre of $-1\bmod q_0$. Write
\[
 A=|\Phi_a|,\quad B=|\Phi_b|,\quad f=|\Phi|=A\varphi(q_0),\quad
 L=\log b,\quad A\le BL.
\]
With $x\le x_1\le a=xb$, the weight $\kappa(\tau)$ equals $b=a/x$ on $\Phi_b$ and $a/x_1$ elsewhere. For an integer $t$ let $u(t)$ be its class if it is a unit modulo $q_1$, and the identity of $\Phi$ otherwise. For $0\le y\le a$ put
\begin{align}
 S(y,\tau)&=\sum_{n\le a}\lambda_{d,1}(n)
       \big(y/n-\kappa(u(-n)\tau)\big)_+,\notag\\
 S_0&=a\left(\log a-\frac1f\sum_{\tau\in\Phi}\log\kappa(\tau)+c_{d,1}\right),
 &\Tot_D&=\sum_{\tau\in\Phi_a}\big(S(a,\tau)-S_0\big).\label{eq:total}
\end{align}

\begin{proposition}[Two-cutoff identity]\label{prop:count}
For every admissible datum,
\begin{equation}
 \Tot_D=a\left\{B\big[R_{d,q_0}(x)-R_{d,q_0}(x_1)\big]+A R_{d,1}(x_1)\right\}.
 \label{eq:count}
\end{equation}
\end{proposition}
\begin{proof}
Fix $n$. If $n\equiv1\bmod q_0$, it is a unit modulo $q_1$ and multiplication by $u(-n)$ maps $\Phi_a$ bijectively to the fibre of $-1$. Exactly $B$ images belong to $\Phi_b$. If $n$ is a unit in any other class, the image fibre misses $\Phi_b$. If $n$ is a nonunit, $u(-n)=1$ and $\Phi_a$ is disjoint from $\Phi_b$. Hence the sum over places of $S(a,\tau)$ is
\[
 \sum_{n\le a}\left[A\lambda_{d,1}(n)v_{x_1}(n)
    +B\lambda_{d,q_0}(n)\big(v_x(n)-v_{x_1}(n)\big)\right].
\]
Apply \eqref{eq:scaled} at the two cutoffs. The two copies of $c_{d,q_0}$ cancel. The remaining main terms sum to
\[
 a\left[A(\log x_1+c_{d,1})
       +\frac{B}{\varphi(q_0)}\log(x/x_1)\right].
\]
This equals $AS_0$ since $\sum_{\tau}\log\kappa(\tau)=f\log(a/x_1)-B\log(x/x_1)$ and $f=A\varphi(q_0)$. Subtracting $AS_0$ gives \eqref{eq:count}.
\end{proof}

\begin{corollary}[One datum]\label{cor:one}
For each admissible datum $D$ satisfying the pointwise estimate \eqref{eq:det}, one has
\begin{equation}
 R_{d,q_0}(x)\le\frac{\varphi(q_0)L^6}{m}
               +|R_{d,q_0}(x_1)|+L|R_{d,1}(x_1)|.\label{eq:one}
\end{equation}
\end{corollary}
\begin{proof}
Solve \eqref{eq:count} for $R_{d,q_0}(x)$ to obtain
$\Tot_D/(aB)+R_{d,q_0}(x_1)-(A/B)R_{d,1}(x_1)$. Since $aB>0$, \eqref{eq:det} bounds the first term from above by $(A/B)\varphi(q_0)L^5/m$. Use $A/B\le L$ and take absolute values of the two secondary-cutoff errors.
\end{proof}

The saving $1/m$ at the primary cutoff comes from \eqref{eq:det}. Section~\ref{sec:comparison} derives \eqref{eq:det} under the arithmetic hypotheses by comparing determinants on the same family of lifts.

\section{Admissible data at every large scale}\label{sec:data-supply}

Fix $q_0>0$ divisible by $18480$ and a prime
$d\equiv q_0/2-1\pmod {q_0}$. Thus $d\equiv7\pmod8$ and
$(q_0,d)=1$. Put $K=\mathbb Q(\sqrt{-d})$.
We construct the ranks, place weights, and levels of an admissible datum for each sufficiently large prescribed length $x$.

\begin{theorem}\label{thm:p1-data}
Fix $0<\eta<1/100$. Every sufficiently large integer $x$ occurs as the
length of an admissible datum for $(q_0,d)$. With $L=\log b$, one may take
\begin{equation}\label{eq:p1-data-scale}
 x_1=\lfloor\exp(x^{\eta/6})\rfloor,\qquad
 x^\eta<L\le x^\eta+1,\qquad \sqrt{x}\le m<2\sqrt{x}.
\end{equation}
The place counts satisfy $f=A\varphi(q_0)$ and $0<A/B\le L$.
In particular $\log x_1\ge x^{\eta/12}$ eventually. The threshold for $x$
may depend on $q_0,d,\eta$.
\end{theorem}

\paragraph{Ranks, moduli and weights.}
Write $Y=\log x$, $W=x^{\eta/6}$. Choose $m=2^u$ and $r=3^v$ to be the
least powers at least $\sqrt{x}$ and $x^5$, respectively. Thus
$r<3x^5$. Set
\[
 s=22u,\quad o=2^s=m^{22},\quad
 \ell_0=3^{\lfloor\log_3s\rfloor+4},\quad
 A=3\ell_0 o,\quad q_1=q_0 A,\quad q_2=(q_0d)^{s+6}.
\]
Here $27s<\ell_0\le81s$ and $s<22+11Y/\log 2\le32Y$ eventually, so $o,q_1,q_2^2$ are
polynomially bounded in $x$ with $(q_0,d)$ fixed. Take
$E=K(\zeta_{q_1})$. Its infinite places, represented by embeddings
extending one fixed embedding of $K$, are indexed by
$\Phi=(\mathbb Z/q_1)^\times$. Indeed, $K$ has conductor $d$ and the
cyclotomic conductors $d,q_1$ are coprime. Reduction modulo $q_0$ is
surjective and all its fibres have size $A$, since $q_0$ and $q_1$
have the same prime divisors. The places of class $1\bmod q_0$ form
$\Phi_a$, of size $A$, and $f=|\Phi|=A\varphi(q_0)$.
See \src{New/Reduction/Data/FieldsAndPlaces}{the tower and its infinite places}.

For $x_1=\lfloor e^W\rfloor$, let $b$ be the least multiple of $q_1x_1$
strictly greater than $e^{W^6}$, and set $a=xb$. Since
$\log(q_1x_1)\le W+O_{q_0,d}(\log x)$, the rounding costs at most
$e^{W^6}$ for large $x$. Consequently
\[
 W-1\le\log x_1\le W,\qquad W^6< L\le W^6+1.
\]
These bounds give
$10^6Y\le\log x_1$, $10^6\log x_1\le L$, $(\log x_1)^6\le L$,
and $100\log L\le Y$. The last inequality uses $100\eta<1$.
Also $\log L$ eventually exceeds $10^{12}$ times the fixed parameter
constant in the determinant data.
Give the places in $\Phi_b$ weight $b$ and the remaining places weight $a/x_1$. Both weights are positive integers, and $a/x_1\le b$ for large $x$.

\paragraph{A large sparse fibre subset.}
Put $k=r+1-m$, $t=5r$, and $D_*=kt$. The determinant estimate requires a set $\Phi_b$ of at least $A/L$ places in the fibre of $-1\bmod q_0$ with the following property. For all primes $p\le D_*$ and integers
$0<i<j<\ell_0$, there do not exist three members $\sigma,\tau,\rho$
with $\tau/\sigma\equiv\pm p^i$ and $\rho/\sigma\equiv\pm p^j$
modulo the largest divisor of $q_1$ prime to $p$. The signs are independent
and repetitions among the three places are allowed. Write $B=|\Phi_b|$.

Here $D_*\ll x^{10}$, $o\ge x^{11}$, so $2\ell_0(D_*+1)<o$ eventually.
Choose an integer $n\ge1$, fixed in $x$ and sufficiently large that
$o,q_2^2,m\le W^n$ eventually. Since $W$ eventually dominates every
fixed power of $Y$, we also have
\[
 W\le o,\qquad(2\ell_0)^n\le W,\qquad
 3\ell_0 n2^n\le W^3 .
\]
Let $H$ be the largest power of $2$ satisfying $H^n\le o/W$.
Then $H\le W$, $o<W(2H)^n$, and $(2\ell_0H)^n\le o$.

Encode the digit vectors in $\{0,\ldots,H-1\}^n$ in base $2\ell_0H$ (cf.~\cite{behrend}).
At least $H^n/(nH^2)$ of them have the same squared Euclidean norm.
Let $S\subset[0,o)$ be the corresponding integers. Consider a congruence with positive integer coefficients $c,c'$ of the form
\[
 (c+c')e_2\equiv ce_1+c'e_3\pmod o,\qquad
 e_1,e_2,e_3\in S,\quad c,c'>0,\quad c+c'\le2\ell_0 ,
\]
where each digit on either side is at most $(c+c')(H-1)<2\ell_0H$. Thus no carries occur. Both sides are less than $(2\ell_0H)^n\le o$. The congruence is therefore an equality in each digit. The middle vector is a convex combination, with both coefficients positive, of the other two vectors on the same sphere. Strict convexity of the squared norm forces $e_1=e_2=e_3$.

Set $g=1+3\ell_0q_0$ and take $\Phi_b=\{-g^e\bmod q_1:e\in S\}$.
If $v_2(q_0)=v$, then $v\ge4$ and $g$ has order $2^s$ modulo
$2^{s+v}$ because successive squaring increases $v_2(g^{2^i}-1)$ from $v$
to $v+i$. Thus the places are distinct and in the prescribed fibre. Moreover
\[
 A=3\ell_0o
 <\frac{H^n}{nH^2}\,(3\ell_0 n2^n)WH^2
 \le B W^6\le BL .
\]
For odd $p$ the two forbidden ratios, squared to remove their signs
and raised to powers $j,i$, imply
$2j e_2\equiv2(j-i)e_1+2i e_3\pmod o$.
Hence all exponents agree. Since $2^{s+v}$ divides the part of $q_1$ prime to $p$, the first squared ratio would give $p^{2i}\equiv1\bmod 2^{s+1}$.
The identity
$v_2(p^{2i}-1)=v_2(p-1)+v_2(p+1)+v_2(i)$ for odd $p$ implies
$2^s\le (p+1)2i$, contrary to the choice of $o$. For $p=2$, every place
has the same residue modulo $3\ell_0$, so a forbidden first ratio gives
$4^i\equiv1\bmod 3\ell_0$. But $4=1+3$ has order $\ell_0$ there,
whereas $0<i<\ell_0$. This proves sparseness. These arguments specialize
the digit-sphere and order calculations in
\src{New/Reduction/Data/Existence}{the construction of sparse places}.

\paragraph{Two principal-level primes.}
All three level primes must be $1\bmod q_2^2$. We use the following
consequence of the classical Siegel--Walfisz theorem. For a fixed $C>0$
and all sufficiently large $W$, if $1\le Q\le W^C$ and
$W\le V\le W^C$, some prime $p\equiv1\bmod Q$ has
$V\le\log p\le V+W$. To see this, put $Z=e^{V+W}$. Then $Q\le(\log Z)^C$,
so the weighted prime count
$\sum_{p\le Z,\ p\equiv1\,(Q)}\log p$ has main term $Z/\varphi(Q)$
and error $O_C(Z e^{-c_C\sqrt{\log Z}})$. The error is
$o(Z/\varphi(Q))$ uniformly. Removing primes below $e^V$ costs
at most $Ve^V=o(Z/\varphi(Q))$, as $V,Q$ are polynomial in $W$.
We use the prime-weighted Siegel--Walfisz estimate of Montgomery and Vaughan \cite[Corollary~11.21]{MV}. Corollary~11.19 gives the von Mangoldt form. The readable edition works with
$Q=(q_0d)^{2(s+6)}$. Real characters at these levels are induced from the fixed modulus $(q_0d)^2$. The zero-free region and Newman's contour
argument then give \src{Lib/Analytic/Primes}{the restricted prime window}.
The classical estimate used here allows ineffective thresholds.

Let $h$ be the class number of $K$. For $i=0,\dots,h$ place logarithmic
windows of width $W$ at $10L+2iW$. They are disjoint and contained in
$[10L,11L]$ since $(2h+1)W\le L$ eventually. The preceding estimate
with $Q=q_2^2$ supplies $h+1$ distinct primes, all split in $K$
because they are $1\bmod d$. Select a prime ideal above each. Two have
the same ideal class. For their underlying primes $p_0,p_1$, the
product of the first ideal and the conjugate of the second is principal.
Its integral generator $\alpha$ has $\alpha\bar\alpha=M=p_0p_1$.
In particular $b^{20}\le M\le b^{22}$.

\paragraph{The level matched to the weights.}
Let $\kappa_*=\exp(f^{-1}\sum_{\tau\in\Phi}\log\kappa(\tau))$ and
$\delta_E$ be the root discriminant. After choosing $M$ put
\[
 V_*=m\log\kappa_*+(m-1)\log(M\delta_E/(2\pi))
                    -\tfrac{15}{2}L .
\]
We choose the third level prime using $V_*$, which now incorporates the mean weight of the selected places. Since
$E\subseteq\mathbb Q(\zeta_{q_1d})$, the root-discriminant tower
inequality and the cyclotomic discriminant formula give
$1\le\delta_E\le q_1d\le b$ for all large $x$. Also
$1\le\kappa_*\le b$. Hence
$19L\le\log(M\delta_E/(2\pi))\le23L$ eventually, and
\[
 mL\le V_*\le24mL
\]
since $m\ge2$. The modulus $q_2^2$ and $24mL$ are bounded by fixed
powers of $W$, and $V_*\ge W$. Apply the prime-window estimate once
more to choose a prime $N\equiv1\bmod q_2^2$ with
$V_*\le\log N\le V_*+W\le V_*+\sqrt L$.
Since $m\to\infty$, eventually $\log N\ge mL>11L$, so $N$ differs from $p_0$ and $p_1$. Finally $N\equiv1\bmod4$ is a sum of two integral squares and
$\sqrt{-1}\in E$. Thus $N=\beta\bar\beta$ with
$\beta\in\mathcal O_E$. This gives the required level and norm generators $\alpha$ and $\beta$. The bound $\log x_1\ge W-1\ge\sqrt W=x^{\eta/12}$ for large $W$ completes the proof of Theorem~\ref{thm:p1-data}.
The readable edition gives the \src{New/Reduction/Data/AdmissibleData}{definition of admissibility and
bounds on $V_*$} and
proves \src{New/Reduction/Data/Existence}{eventual existence with prescribed
integer length}.

\section{A one-sided prime number theorem}\label{sec:from-data}
Both secondary remainders in \eqref{eq:one} are evaluated at \(x_1\), with
\(q_0\) and \(d\) fixed as \(x\to\infty\). Since \(\log x_1\ge x^{\eta/12}\)
eventually, a sufficiently large inverse power of \(\log x_1\) absorbs both errors.

\begin{lemma}[Centering and variation]\label{lem:pnt}
Fix an odd prime $d$ and a positive integer $q$ coprime to $d$.
For every integer $K\ge1$ and all sufficiently large $y$,
\begin{equation}
 |R_{d,q}(y)|\le (\log y)^{-K}.\label{eq:PNT}
\end{equation}
For all $1\le u\le v$ one has
\begin{equation}
 |R_{d,q}(v)-R_{d,q}(u)|\le 13(v-u)/u.\label{eq:variation}
\end{equation}
\end{lemma}
\begin{proof}
Write $A_{d,q}(t)=\sum_{n\le t}\lambda_{d,q}(n)$ and $e(t)=A_{d,q}(t)-t/\varphi(q)$.
Integrating the step function gives
\begin{equation}
 R_{d,q}(y)=\int_1^y\frac{e(t)}{t^2}\,dt-c_{d,q}.\label{eq:integralR}
\end{equation}
Each \(n\le y\) contributes \(\lambda_{d,q}(n)\int_n^y t^{-2}\,dt\).
By \(0\le\lambda_{d,q}\le2\Lambda\) and Chebyshev's bound
\(\sum_{n\le t}\Lambda(n)\le6t\), we have \(|e(t)|\le13t\).
Integration over \([u,v]\) gives \eqref{eq:variation}.

For \eqref{eq:PNT}, apply the prime number theorem in progressions modulo
the fixed integer \(qd\)~\cite[Corollary~11.19]{MV}. On the reduced classes congruent to \(1\bmod q\),
the weights \(1-\chi_d\) sum to \(\varphi(d)\) by the Chinese remainder
theorem. Among nonunits modulo \(qd\), only powers of \(d\) can contribute,
with total weight \(O_d(\log(2t))\). Thus for each fixed \(B>1\),
\[
 A_{d,q}(t)=\frac{t}{\varphi(q)}+O_{q,d,B}\big(t/(\log t)^B\big)\qquad(t\ge2).
\]
Hence \(e(t)t^{-2}\) is integrable at infinity. For real \(s>1\), partial
summation and \eqref{eq:p1Dir} give
\[
 s\int_1^\infty e(t)t^{-s-1}\,dt
 =\sum_{n\ge1}\lambda_{d,q}(n)n^{-s}-\frac{s}{\varphi(q)(s-1)}
 =G_{d,q}(s).
\]
Letting \(s\downarrow1\) by dominated convergence gives
\(\int_1^\infty e(t)t^{-2}\,dt=G_{d,q}(1)=c_{d,q}\). Thus \(R_{d,q}(y)\)
is the negative of the tail integral. With \(B=K+2\) this is
\(O_{q,d,K}((\log y)^{-K-1})\), giving \eqref{eq:PNT} with coefficient one
for sufficiently large \(y\).
\end{proof}
The formal argument obtains \eqref{eq:PNT} from a quantitative Tauberian
theorem for the regularized logarithmic derivative. The fixed-modulus
prime number theorem used here follows from the classical zero-free estimate
near \(\Re s=1\).

\begin{theorem}[All real cutoffs]\label{thm:all-y}
Assume Hypothesis~\ref{hyp:scale}. For every $q_0$ divisible by $18480$, $q_0>0$, every prime $d\equiv q_0/2-1\pmod{q_0}$ and every $\varepsilon>0$ there is $Y=Y(q_0,d,\varepsilon)$ such that
\begin{equation}
 R_{d,q_0}(y)\le y^{-1/2+\varepsilon}\qquad(y\ge Y).\label{eq:real}
\end{equation}
\end{theorem}
\begin{proof}
Choose \(0<\eta<1/100\) with \(7\eta\le\varepsilon\) and an integer
\(K\ge24/\eta\). For every sufficiently large integer \(x\),
Hypothesis~\ref{hyp:scale} gives a datum of Theorem~\ref{thm:p1-data}
at \(x\) satisfying \eqref{eq:det}. On this datum \eqref{eq:p1-data-scale}
gives \(m\ge\sqrt{x}\), \(L\le2x^\eta\), and \(\log x_1\ge x^{\eta/12}\).
For \(q=q_0,1\), \eqref{eq:PNT} now gives \(|R_{d,q}(x_1)|\le x^{-2}\).
Substitution in \eqref{eq:one} yields
\begin{align*}
 R_{d,q_0}(x)&\le64\varphi(q_0)x^{-1/2+6\eta}
                 +(1+2x^\eta)x^{-2}\\
              &\le C x^t,\qquad C=64\varphi(q_0)+3,\quad t=-1/2+6\eta .
\end{align*}
The factor \(x^{-1/2}\) comes from \(m^{-1}\). For real \(y\ge2\) put
\(x=\lfloor y\rfloor\). Since \(x\ge y/2\) and \(-1\le t\le0\), we have
\(x^t\le2y^t\), and \eqref{eq:variation} bounds the increment by \(26/y\).
Hence for sufficiently large \(y\),
\[
 R_{d,q_0}(y)\le(2C+26)y^t\le y^{t+\eta}\le y^{-1/2+\varepsilon}.
\]
\end{proof}


\section{The transform and its real residues}\label{sec:p1polar}
Let \(d\) be an odd prime coprime to \(Q\). We use the entire functions
\(D_\eta\) and their regularized logarithmic derivative average \(G_{d,Q}\)
from \eqref{eq:p1G}. Under \(y=e^v\) the Mellin transform of the smoothed
remainder becomes a Laplace transform.
\begin{proposition}[Transform and residues]\label{prop:p1polar}
Put $E(v)=R_{d,Q}(e^v)$ and initially, for $\Re s>2$, set
\[
 F(s)=\int_0^\infty E(v)e^{(1-s)v}\,dv.
\]
Then
\begin{equation}\label{eq:p1F}
 F(s)=\frac{G_{d,Q}(s)/s-c_{d,Q}}{s-1}.
\end{equation}
Its meromorphic continuation has at most simple poles in $\Re s>0$ and is regular at $1$. Write
\[
 n_Q(z)=\sum_{\eta\bmod Q}\operatorname{ord}_z D_\eta,
 \qquad n_{Q,d}(z)=\sum_{\eta\bmod Q}\operatorname{ord}_z D_{\eta\chi_d}.
\]
With residue zero at regular points, for $\Re z>0$, $z\ne1$ one has
\begin{equation}\label{eq:p1res}
 \operatorname{Res}_{s=z}F(s)
     =\frac{n_Q(z)-n_{Q,d}(z)}{\varphi(Q) z(1-z)} .
\end{equation}
\end{proposition}
\begin{proof}
For the contribution of $n$ the integral starts at $\log n$. Direct integration gives
\[
 \int_{\log n}^\infty (n^{-1}-e^{-v})e^{(1-s)v}\,dv
       =\frac{n^{-s}}{s(s-1)}.
\]
The same integral at \(\Re s\) gives absolute convergence after summing
against \(\lambda_{d,Q}(n)\ge0\). The transforms of \(v/\varphi(Q)\) and
\(c_{d,Q}\) are \(1/[\varphi(Q)(s-1)^2]\) and \(c_{d,Q}/(s-1)\),
respectively. Substitution of \eqref{eq:p1Dir} cancels the double pole and
gives \eqref{eq:p1F}. Since \(c_{d,Q}=G_{d,Q}(1)\), the singularity at
\(1\) is removable. Elsewhere in \(\Re z>0\), each logarithmic derivative
contributes the order of its entire factor. The twisted orders enter
\(G_{d,Q}\) with positive sign, the untwisted with negative sign.
Dividing by \(z(z-1)\) gives \eqref{eq:p1res}.
\end{proof}
The remainder \(E\) is real and measurable. Since
\(0\le\lambda_{d,Q}(n)\le2\log n\), comparison with
\(\sum_{n\le y}2\log(n)/n\) bounds the smoothed sum by \(O(1+v^2)\)
at \(y=e^v\). Thus \(|E(v)|\ll_{d,Q}e^v\) for \(v\ge0\).
The corresponding formal constructions are
\src{Lib/Analytic/EntireLFunction}{the entire \(L\)-functions} and
\src{New/Reduction/Zeros/PrimeError}{the polar prime remainder}.

\subsection{Landau's argument on the real axis}
To deduce holomorphy from an upper bound on \(E\), we need the sign of the
residue at a possible real boundary pole. The following form of Landau's
singularity theorem makes the sign explicit \cite[Lemma~15.1]{MV}\cite[Chapter~II, Theorem~5b]{widder}.

\begin{lemma}[One-sided Landau criterion]\label{lem:p1landau}
Let $E$ be real measurable on $[0,\infty)$, with $|E(v)|\le Ce^v$. Suppose its transform $F$ as above continues meromorphically to $\Re s>0$ with at most simple poles there. Fix $\beta>0$. If $E(v)\le e^{(\beta-1)v}$ eventually and
\[
 \Re\operatorname{Res}_{s=u}F(s)\ge0
 \quad\text{for each real }u>\beta,
\]
then $F$ is holomorphic for $\Re s>\beta$.
\end{lemma}
\begin{proof}
The eventual upper bound and \(|E(v)|\le Ce^v\) on the initial bounded
interval give a constant \(K\ge1\) with \(E(v)e^v\le K e^{\beta v}\) for
every \(v\ge0\). Define the positive measure
\[
 d\mu(v)=(K e^{\beta v}-E(v)e^v)\,dv
 \qquad(v\ge0).
\]
Its Laplace transform $P(s)=\int e^{-sv}\,d\mu(v)$ converges for $\Re s>\max(2,\beta)$, where
\[
 P(s)=\frac K{s-\beta}-F(s).
\]
The right side continues \(P\) meromorphically to \(\Re s>0\).
Let \(\alpha\) be the abscissa of convergence of the positive integral,
allowing \(\alpha=-\infty\). Local absolute convergence makes \(P\)
holomorphic on \(\Re s>\alpha\).

Suppose $\alpha>\beta$. The continuation of $P$ has at most a simple pole at $\alpha$. Its residue has nonpositive real part by hypothesis, whereas
\[
 \lim_{u\downarrow\alpha}(u-\alpha)P(u)
\]
is a nonnegative real number because $P(u)$ is a positive integral there. The limit equals the residue, so the residue vanishes and the continuation is regular at $\alpha$.

Choose \(r>0\) so that the continuation is holomorphic on the disc of
radius \(3r\) about \(\alpha\), and set \(u=\alpha+r/2\). Differentiation
in the half-plane of convergence gives
\((-1)^kP^{(k)}(u)=\int v^k e^{-uv}\,d\mu(v)\).
The Taylor series at \(u\) converges at \(u-r<\alpha\). Its terms at this
point are nonnegative, so monotone convergence gives
\[
 \sum_{k\ge0}\frac{r^k}{k!}\int v^k e^{-uv}\,d\mu(v)
   =\int e^{-(u-r)v}\,d\mu(v).
\]
This contradicts the definition of $\alpha$. Hence $\alpha\le\beta$, and uniqueness of meromorphic continuation identifies $F=K/(s-\beta)-P$ throughout $\Re s>\beta$. Both terms are holomorphic there.
\end{proof}
The twist in the next section must therefore control residues on the real
axis as well as the residue at the proposed zero.

\section{A quadratic twist free of zeros in a prescribed box}
\label{sec:p1-good-twist}

Let $Q$ be a positive multiple of $8$ and $a_0\bmod Q$ a unit class with
$a_0\equiv7\pmod8$. For each prime $d\equiv a_0\pmod Q$ put
$\chi_d(n)=(n/d)$. We shall choose $d$ after fixing the height $H$ and
left edge $\beta$ of the box containing the proposed zero.

\begin{proposition}\label{prop:p1-good-twist}
For $H\ge0$ and $\beta>1/2$ there is a prime $d\equiv a_0\pmod Q$ such that
\[
 L(s,\xi\chi_d)\ne0
 \qquad (\xi\bmod Q,\ \Re s\ge\beta,\ |\Im s|\le H).
\]
\end{proposition}

It suffices to treat $1/2<\beta<1$, since the box with $\beta=3/4$
contains every box of the same height with $\beta\ge1$. Put
$\delta=(\beta-1/2)/9$. A zero forces a large theta integral. Its mean
over quadratic symbols bounds the number of exceptional primes.

\subsection{Detection by a truncated Mellin transform}
For $j=0,1$ and a sequence $c(n)$ of modulus at most one, put
\[
 w_j(y)=y^j e^{-\pi y^2},\qquad
 \theta_{j,c}(u)=u^{1/2}\sum_{n\ge1}c(n)w_j(nu),\qquad
 G_j(s)=\tfrac12\pi^{-(s+j)/2}\Gamma((s+j)/2).
\]
Thus $G_j$ is the Mellin transform of $w_j$. Let $\chi$ be primitive and
nonprincipal of conductor $q$ and parity $j$. Suppress $j$ in $\theta$,
write $\mu$ for the M\"obius function, and for $T\ge\sqrt q$ put
\[
 J(\chi,T)=\int_{1/T}^{\infty}
       (|\theta_\chi(u)|+|\theta_{\mu\chi}(u)|)\frac{du}{u},\qquad
 M_T(s)=\int_{1/T}^{\infty}\theta_{\mu\chi}(u)u^{s-1/2}\frac{du}{u}.
\]
\begin{lemma}\label{lem:p1-twist-detect}
For $T\ge\sqrt q$ above a threshold depending only on $\delta,H$, a zero $s$ of
$L(s,\chi)$ in the box of Proposition~\ref{prop:p1-good-twist} implies
$J(\chi,T)>T^{2\delta}$.
\end{lemma}
\begin{proof}
For $\Re s>1$, termwise integration gives the Mellin transform
$G_j(s)L(s,\chi)$ of $\theta_\chi$ at $s-1/2$. Poisson summation for
the primitive character gives
$|\theta_\chi(u)|=|\theta_{\bar\chi}(1/(qu))|$.
The weight $w_j$ is real, so $|\theta_{\bar\chi}|=|\theta_\chi|$.
Gaussian decay at infinity therefore gives rapid decay at zero and
continues the Mellin integral to every $s$. Splitting at $q^{-1/2}$
and substituting $u\mapsto1/(qu)$ in the lower integral gives,
on $\Re s=1/2$,
\[
 |G_j(s)L(s,\chi)|\le 2\int_{q^{-1/2}}^\infty|\theta_\chi(u)|\frac{du}{u}
 \le2J(\chi,T),\qquad |M_T(s)|\le J(\chi,T).
\]
On $\Re s=2$, absolute convergence identifies the full Mellin transform
for $\mu\chi$ with $G_j(s)/L(s,\chi)$. For $0<u\le1$ the Gaussian
sum gives $|\theta_{\mu\chi}(u)|\ll u^{-1/2}$, so truncation at $1/T$
costs $O(T^{-1})$ uniformly in $\chi$ and $\Im s$. The Dirichlet
series bounds $L(2+it,\chi)$ uniformly, and the Mellin integral for
$w_j$ bounds $G_j$ throughout $1/2\le\Re s\le2$. Thus
\[
 F_T(s)=G_j(s)L(s,\chi)M_T(s)-G_j(s)^2
\]
is $O(T^{-1})$ on $\Re s=2$. If $J(\chi,T)\le T^{2\delta}$, the
bounds on $\Re s=1/2$ give $F_T(s)\ll T^{4\delta}$ on that line. For fixed
$\chi,T$ the Mellin integrals give a bounded holomorphic function on
the strip, continuous on its boundary. The three-lines theorem~\cite[Lemma~4.6]{brown} now gives,
for $\beta\le\Re s\le1$,
\[
 |F_T(s)|\ll T^{4\delta(1-\omega)-\omega}\le T^{-2\delta},\qquad
 \omega=\frac{\Re s-1/2}{3/2}\ge6\delta .
\]
At a zero of $L(s,\chi)$ we have $|G_j(s)|^2\ll T^{-2\delta}$.
For $j=0,1$, the gamma factor has a positive lower bound in modulus
on $1/2\le\Re s\le1$, $|\Im s|\le H$, a contradiction for large $T$.
The Euler product excludes zeros for $\Re s>1$.
\end{proof}

\subsection{Mean over quadratic symbols}
For $n\ge1$ let $n^\circ$ be its odd part and put
$b_v(n)=\bigl(\frac{-v}{n^\circ}\bigr)$, the Jacobi symbol.
For either parity, any $|\alpha(n)|\le1$, and $u\ge1/T$,
\begin{equation}\label{eq:p1-twist-second}
 \sum_{1\le v\le T^2}|\theta_{j,\alpha b_v}(u)|^2
 \ll_{\delta} T^2u^{-2\delta}.
\end{equation}
Gaussian decay permits termwise expansion at fixed $u$. For a pair
$n,m$ the character sum is $\sum_{v\le T^2}b_v(nm)$. When $(nm)^\circ$
is square, its modulus is at most $T^2$. Using $|w_j(y)|\ll_A y^{-A}$
with $A=1/2+\delta$ bounds the pair's contribution by
$O_\delta(T^2u^{-2\delta}(nm)^{-1/2-\delta})$. Write
$n=2^h t k^2$ and $m=2^i t l^2$, with $t$ odd squarefree and $k,l$ odd.
Summation in $h,i$ gives geometric series. The exponents of $t,k,l$
are all $-1-2\delta$, so the remaining sums converge.

When $(nm)^\circ$ is nonsquare, $b_v(nm)$ is a fixed sign times a
nonprincipal character of $v$ modulo $(nm)^\circ$. P\'olya--Vinogradov
\cite[Theorem~9.18]{MV} gives an interval sum
$\ll ((nm)^\circ)^{1/2}(1+\log((nm)^\circ))\ll_\delta(nm)^{1/2+\delta/2}$.
Taking $A=3/2+\delta$ gives
$O_\delta(u^{-2-2\delta}(nm)^{-1-\delta/2})$ per pair. Summing in
$n,m$ and using $u^{-2}\le T^2$ proves \eqref{eq:p1-twist-second}.

Write $J_j(c,T)$ for the integral defining $J$ with coefficients $c$
and parity $j$. Cauchy--Schwarz in $v$ applied to
\eqref{eq:p1-twist-second} bounds the sum of absolute values by
$O_\delta(T^2u^{-\delta})$ for both $\alpha$ and $\mu\alpha$.
Integration gives
\begin{equation}\label{eq:p1-twist-mean}
 \sum_{v\le T^2,\ v\ge1} J_j(\alpha b_v,T)
 \ll_\delta T^2\int_{1/T}^\infty u^{-\delta-1}\,du
 \ll_\delta T^{2+\delta}.
\end{equation}

\begin{proof}[Proof of Proposition~\ref{prop:p1-good-twist}]
Let $\psi$ be the primitive character of conductor $M\mid Q$ inducing
$\xi\bmod Q$. Every prime $d\equiv a_0\pmod Q$ is coprime to $Q$,
so $\psi\chi_d$ is primitive of conductor $Md$ and nonprincipal. The
finite Euler factors relating $L(s,\xi\chi_d)$ to $L(s,\psi\chi_d)$
are nonzero for $\Re s>0$. Since $d\equiv7\pmod8$, we have
$\chi_d(2)=1$. Quadratic reciprocity gives $\chi_d(n)=b_d(n)$ for
all $n\ge1$, including multiples of $d$. The parity $j$ of $\psi\chi_d$
is opposite to that of $\psi$ for every $d$ in the progression.

Take $X$ sufficiently large and set $T=\sqrt{QX}$. For $d\le X$ the
conductor $Md$ is at most $T^2$. A zero of $L(s,\xi\chi_d)$ in the box
therefore gives $J_j(\psi b_d,T)>T^{2\delta}$ by
Lemma~\ref{lem:p1-twist-detect}. Equation~\eqref{eq:p1-twist-mean}
bounds the number of such primes by $O_{\delta,H,Q}(T^{2-\delta})$.
Let $B(X)$ count primes $d\le X$ in the progression for which some
$\xi\bmod Q$ has a twist with a zero in the box. Summing over the
$\varphi(Q)$ characters gives
\[
 B(X)\ll_{Q,H,\delta} X^{1-\delta/2}.
\]
Dyadic summation shows that the reciprocals of the exceptional primes
have finite sum. By Dirichlet's theorem in its analytic form~\cite[Chapter~4]{MV}, the
reciprocals of the primes in the reduced class have divergent sum.
Infinitely many primes in the class satisfy the proposition.
\end{proof}

The formal argument uses the weight $(\log d)/d$ in the last step.
The exceptional sum still converges by the same dyadic bound, and the
sum over the progression diverges. The Gaussian normalization in the
Poisson transformation, with Gauss sum of absolute value $\sqrt q$,
is given in \src{Lib/Analytic/DirichletTheta}{Dirichlet theta transforms}.
The weighted selection argument is in
\src{Lib/Analytic/QuadraticTwist}{quadratic twist selection}.

\section{Nonvanishing and reflection}\label{sec:descent}
\subsection{Choosing the twist at the height of the proposed zero}
\begin{theorem}[Family criterion]\label{thm:p1criterion}
Let $8\mid Q$, $Q>0$, and let $d_0$ be prime to $Q$ with $d_0\equiv7\bmod 8$. Fix $\beta>1/2$. Suppose that for each prime $d\equiv d_0\bmod Q$ there exists $Y_d$ such that
\[
 R_{d,Q}(y)\le y^{\beta-1}\qquad(y\ge Y_d).
\]
Then for every character $\psi$ modulo $Q$ and every $\rho\ne1$ with $\Re\rho>\beta$, one has $L(\rho,\psi)\ne0$.
\end{theorem}
\begin{proof}
Fix $\rho$ with $\Re\rho>\beta$, $\rho\ne1$. Apply Proposition~\ref{prop:p1-good-twist} with $H=|\Im\rho|$. It gives a prime $d\equiv d_0\bmod Q$ for which $n_{Q,d}$ vanishes throughout $\Re s\ge\beta$, $|\Im s|\le H$. In particular $(d,Q)=1$ and $d\equiv7\bmod 8$. For this $d$ the assumed bound gives $E(v)=R_{d,Q}(e^v)\le e^{(\beta-1)v}$ whenever $e^v\ge Y_d$.

For real $\beta<u<1$, formula~\eqref{eq:p1res} now gives residue $n_Q(u)/[\varphi(Q)u(1-u)]\ge0$. The transform is regular at $1$. For $u>\max(1,\beta)$ the untwisted factors are also nonzero, so the residue is zero. Lemma~\ref{lem:p1landau} therefore gives holomorphy in $\Re s>\beta$. At $\rho$ the residue vanishes, and $n_{Q,d}(\rho)=0$ by the choice of $H$. Thus \eqref{eq:p1res} gives $n_Q(\rho)=0$. This is a sum of nonnegative orders, so $D_\psi(\rho)\ne0$ for every $\psi\bmod Q$. As $\rho\ne1$, the same holds for $L(\rho,\psi)$.
\end{proof}
In the \src{New/Reduction/Zeros/Criterion}{one-sided nonvanishing theorem}, the height $H$ determines the box for both the residue at $\rho$ and the signs on the real axis. The threshold $Y_d$ may depend on the prime chosen for this box.

\subsection{Return to the original modulus}
\begin{proof}[Proof of Theorem~\ref{thm:main}]
Fix a character $\chi\bmod q$, $q\ge1$. On $\Re s\ge1$ classical nonvanishing gives the assertion, with the principal pole at $1$ excepted. Take $\rho$ with $1/2<\Re\rho<1$ and choose $1/2<\beta<\Re\rho$. Put
\[
 Q=18480q,\qquad d_0=Q/2-1=9240q-1 .
\]
Here $8\mid Q$, $d_0\equiv7\bmod 8$, and $(d_0,Q)=1$ since $d_0$ is odd and $Q=2(d_0+1)$. Under Hypothesis~\ref{hyp:scale}, Theorem~\ref{thm:all-y} with $\varepsilon=\beta-1/2$ gives $R_{d,Q}(y)\le y^{\beta-1}$ for all sufficiently large $y$, for each prime $d\equiv d_0\bmod Q$. Apply Theorem~\ref{thm:p1criterion} to the character $\chi_Q$ induced at level $Q$. It gives $L(\rho,\chi_Q)\ne0$. The identity
\begin{equation}
 L(s,\chi_Q)=L(s,\chi)\!
 \prod_{\substack{p\mid Q\\p\nmid q}}(1-\chi(p)p^{-s})\label{eq:level}
\end{equation}
holds by the Euler product for $\Re s>1$ and elsewhere by meromorphic continuation. At $s=\rho$ the left side is nonzero, so $L(\rho,\chi)\ne0$.

Let $\psi$ be the primitive character inducing $\chi$. The additional Euler factors in \eqref{eq:level}, applied from the conductor of $\psi$ to $q$, are nonzero for $\Re s>0$ because $|\psi(p)p^{-s}|<1$. A zero of $L(s,\chi)$ in the open critical strip is therefore a zero of $L(s,\psi)$. The gamma factors are finite and nonzero in the strip, and the primitive functional equation~\cite[\S10.1]{MV} gives $L(1-s,\bar\psi)=0$. Half-plane nonvanishing at $s$ and $1-s$ gives $\Re s\le1/2$ and $1-\Re s\le1/2$, respectively.

For $q=1$ this gives the half-plane assertion for $\zeta$. Suppose $\zeta(s)=0$ outside the negative even integers. Then $s\ne1$ and $s\ne0$, since $\zeta(0)=-1/2$. The factor $\pi^{-s/2}\Gamma(s/2)$ is finite and nonzero at $s$, so the completed zeta function vanishes at both $s$ and $1-s$. If $\Re s<1/2$, then $\Re(1-s)>1/2$ and $1-s\ne1$. The gamma factor is finite and nonzero there, so $\zeta(1-s)=0$, contrary to half-plane nonvanishing. The same assertion rules out $\Re s>1/2$.
\end{proof}

\subsection{The converse on the constructed scales}
\begin{corollary}[Strength of the determinant criterion]\label{cor:scale-equivalence}
Hypothesis~\ref{hyp:scale} is equivalent to GRH for Dirichlet \(L\)-functions.
Under GRH, for fixed permissible \((q_0,d)\) and fixed \(0<\eta<1/100\),
every admissible datum at the scales of \eqref{eq:p1-data-scale} satisfies
\(m|\Tot_D|\le fa(\log b)^5\) at every sufficiently large integer \(x\).
\end{corollary}
\begin{proof}
Theorem~\ref{thm:main} gives the forward implication.
Assume GRH and take \(Q=1\) or \(Q=q_0\). Write
\(\psi(t,\xi)=\sum_{n\le t}\Lambda(n)\xi(n)\). Orthogonality gives
\[
 A_{d,Q}(t)=\frac1{\varphi(Q)}
  \sum_{\xi\bmod Q}\bigl(\psi(t,\xi)-\psi(t,\xi\chi_d)\bigr).
\]
There is one principal character in the first sum. Every twist in the
second sum is nonprincipal by the Chinese remainder theorem, including
when \(Q=1\). The prime number theorem under GRH at the fixed levels
\(Q,Qd\) gives \(A_{d,Q}(t)-t/\varphi(Q)\ll_{d,Q}\sqrt{t}\log^2(2t)\)
\cite[Theorem~13.7]{MV}. Imprimitive Euler factors change the corresponding
primitive sums by \(O_{d,Q}(\log(2t))\).
By \eqref{eq:integralR} and Lemma~\ref{lem:pnt}, \(R_{d,Q}(y)\) equals
the negative tail integral of this centered error against \(dt/t^2\).
Thus \(\lvert R_{d,Q}(y)\rvert\ll_{d,Q} y^{-1/2}\log^2(2y)\).

The fibre of $-1\bmod q_0$ has $A$ members, so its selected subset has $B\le A$. Substitution in \eqref{eq:count}, with $f=A\varphi(q_0)$ and $m<2\sqrt{x}$, gives
\[
 \frac{m|\Tot_D|}{fa}\ll_{d,q_0}\log^2(2x)
       +\sqrt{x/x_1}\log^2(2x_1)=o(L^5).
\]
Indeed \(x_1=\lfloor\exp(x^{\eta/6})\rfloor\) makes the second term tend
to zero, and \(L>x^\eta\). The estimate is uniform over the data at the
fixed parameters. Theorem~\ref{thm:p1-data} supplies a datum at each
sufficiently large integer \(x\), proving Hypothesis~\ref{hyp:scale}.
\end{proof}

\clearpage
\part{The determinant comparison}\label{part:determinant}
\section{The datum, the signed sum and the two lattice bounds}\label{sz:setup}
For a positive multiple $q_0$ of $18480$ and a prime $d\equiv q_0/2-1\pmod{q_0}$,
write $P=(q_0,d)$. The following conditions specify an admissible datum
$D$ for $P$. Theorem~\ref{thm:p1-data} constructs such data at every
sufficiently large integer length. The construction hypotheses on forms
and lifts will be imposed separately for each datum.

\paragraph{Scales and places.}
The positive integers $x,x_1,b$ have logarithms satisfying
\[
 10^{12}C_P\le\log L,\quad L^{100}\le x,\quad x^{10^6}\le x_1,
 \quad x_1^{10^6}\le b,\quad (\log x_1)^6\le L,\qquad L=\log b .
\]
The parameter constant $C_P$ is the fixed constant in
\srcline{New/Reduction/Data/AdmissibleData}{50}{the definition of admissibility}.
In particular $C_P\ge64$ and it bounds
$|c_K|$, where $c_K=\gamma-\Re(L'/L)(1,\chi_d)$ and $\chi_d(n)=(n/d)$.
The length $m$ is a positive power of two in $[\sqrt x,2\sqrt x]$ and the
unitary rank $r$ a positive power of three in $[x^5,3x^5]$.
Take $s=22\log_2 m\le32\log x$,
$\ell_0=3^{\lfloor\log_3s\rfloor+4}$,
$A=3\ell_0 2^s$, $q_1=q_0A$ and $q_2=(q_0d)^{s+6}$, with $q_1x_1\mid b$.
Put $a=xb$, $K=\mathbb Q(\sqrt{-d})$ and $E=K(\zeta_{q_1})$.
Represent its infinite places by the embeddings extending a fixed embedding of
$K$. They form the group $\Phi\simeq(\mathbb Z/q_1)^\times$ and $f=|\Phi|$.
The fibre $\Phi_a$ over $1\bmod q_0$ has size $A$. In the fibre over
$-1$ choose $\Phi_b$ with $|\Phi_b|L\ge A$. Put $k=r+1-m$.
For every prime $p\le k(5r)$ and $0<i<j<\ell_0$, there are no
$\sigma,\tau,\rho\in\Phi_b$ with
$\tau/\sigma\equiv\pm p^i$, $\rho/\sigma\equiv\pm p^j$
modulo the prime-to-$p$ part of $q_1$, with the signs chosen independently.
The place weight $\kappa(\tau)$ equals $b$ on $\Phi_b$ and $a/x_1$ elsewhere.

The distinct level primes $p_0,p_1$ lie in $[b^{10},b^{11}]$ and the
auxiliary prime $N$ has the logarithmic window below. All three are
$1\bmod q_2^2$. Write $M=p_0p_1$, $\log\bar\kappa=f^{-1}\sum_\tau\log\kappa(\tau)$,
and $\mathrm{rd}_E=|\mathrm{disc}(E)|^{1/(2f)}$. There is an integer of $K$
whose product with its conjugate is $M$ and an integer of $E$ with corresponding
product $N$; these choices are part of $D$. The required window is
\[
 \log N_-\le\log N\le\log N_-+\sqrt L,\quad
 \log N_-=m\log\bar\kappa+(m-1)\log(M\mathrm{rd}_E/(2\pi))-15L/2 .
\]

\paragraph{The two scalar quantities.}
For $h\in\mathbb N$ put $j(h)=\log(h!)+h\log(M/\Omega)$,
where $\Omega=2\pi e^{c_K}/\mathrm{rd}_E$. Set $u(t)=t\bmod q_1$ if $t$ is a
unit modulo $q_1$, and $u(t)=1$ otherwise. With $v_+=\max(v,0)$ set
\[
 S(y,\tau)=\sum_{1\le\nu\le a}\Lambda(\nu)(1-\chi_d(\nu))
       (y/\nu-\kappa(u(-\nu)\tau))_+,
 \qquad S_0=a(\log a-\log\bar\kappa+c_1).
\]
Here $c_1=c_{d,1}$ is the modulus-one centering constant in \eqref{eq:R}. Explicitly
$c_1=\Re[(L'/L)(1,\chi_d)-((s-1)\zeta(s))'/((s-1)\zeta(s))|_{s=1}-1]$,
where $(s-1)\zeta(s)$ has its regular value at $1$. Thus $c_1=-c_K-1$.
With the normalization of \eqref{eq:total}, the signed excess
and the common main height are
\begin{equation}\label{sz:height}
 \Tot_D=\sum_{\tau\in\Phi_a}(S(a,\tau)-S_0),\qquad
 H=mk\sum_{\tau\in\Phi}\bigl(j(\kappa(\tau))-S(\kappa(\tau),\tau)\bigr).
\end{equation}
Write $U=faL^3$ and $V=faL^4$. The arguments of $S$ and the sets of places
in these two sums differ. Both determinant estimates use exactly this $H$.

\paragraph{The spaces and maps.}
Let $X$ be the space of algebraic forms on the definite rank-$k$ space
$X_0$ at its prescribed level and Weyl weight. At each of the $f$ chosen places
$w\mid N$ there is a central scaling translation $T_w$ of order dividing
$N-1$. The ramified source is $S=\bigcap_w\ker(\sum_{i=0}^{N-2}T_w^i)$.
The argument of $S(y,\tau)$ distinguishes the prime sum from this space.
Write $Y$ for the principal-level automorphic forms of the prescribed type
on the rank-$r$ space $Y_0$, with the Petersson inner product. The finite
residue group $G$ at the chosen places of $p_0,p_1$ acts through $\rho$.
The central good-type projector $e_0$ has image $Y_T=e_0Y$.
The domain dimension of the normalized maps is $d_S=\dim_\mathbb C S$,
their target dimension is $d_T=\dim_\mathbb C Y_T$, and $d_X=\dim_\mathbb C X$.
For the Poincar\'e map $\Pi:X\longrightarrow Y$ defined by the chosen Bessel
functional and point, the normalization used throughout is
\[
 P_g=c_B c_{\rm ss}^{-1}|G|\,e_0\rho(g)\Pi|_S,\qquad g\in G .
\]
The seesaw gives the common unit $c_{\rm ss}$, and the local divisibility
bounds determine the positive rational $c_B$ in Section~\ref{A:section}.
Centrality implies $P_g=\rho(g)P_1$. At each infinite place $Y_0$ has signature $(r-m/2,m/2)$: the initial
string of length $m-1$ contributes $m/2-1$ hyperbolic pairs and one
negative central line, in addition to $k$ positive directions. The chain $\Theta_\xi$ takes its good forms through $m$ lifts
to algebraic forms on the definite terminal space $Y_m$.

An $m$-fold twist $\xi$ is a weight-zero tame automorphic character of
$U(1)$ with the admissible finite support disjoint from its conjugate
and of norm prime to $MN$. Every power $\xi^j$, $1\le j\le m$, has
nontrivial tame character at every place of that support. Let $N_\xi$ be
the product of the residue norms on the support. It is allowed at $p$ if
$p\nmid N_\xi$. The empty-support trivial twist is denoted by $1$.
For an identification $\iota:\mathbb C\simeq\overline{\mathbb Q}_p$ write
$|z|_\iota=|\iota z|_p$ with $|p|_p=p^{-1}$ and $O_\iota=\{z:|z|_\iota\le1\}$.
Integrality of an algebraic form means that each adelic value, transported
back by the inverse finite local action, lies in the integral Weyl module.
Write $\mathcal I\subset S$ for source forms integral for every $p,\iota$
and $\mathcal L_{\xi,\iota}$ for the corresponding module at the end of the chain.

\paragraph{Index and squared covolume.}
For a finite indexed family $y=(y_i)$ in $Y_T$, the assertion
$\operatorname{IndexLe}(y,\Delta)$ has the following order of quantifiers.
There exist numbers $e_p\ge0$ with $\sum_{p\in F}e_p\log p\le\Delta$ for
every finite set of primes $F$. For every prime $p$ and every $\iota$ above it
there exists one square matrix $u_\iota$ over $O_\iota$ on the index set of
$y$, satisfying
\begin{equation}\label{sz:index}
 |\det u_\iota|_\iota\ge p^{-e_p},\qquad
 \Theta_\xi\!\left(\sum_j (u_\iota)_{ij}y_j\right)\in\mathcal L_{\xi,\iota}
 \quad\text{for all }i\text{ and all }\xi\text{ allowed at }p.
\end{equation}
At a fixed identification the same rows must work for every allowed twist.

Let $v_\sigma(z)$ be the vector obtained by applying
$\sigma\in\operatorname{Aut}(\mathbb C)$ to each of the prescribed
coordinates of $\Theta_1z$ at the terminal class representatives.
The class pairing has a rational positive definite matrix $G_{\rm cl}$
in these coordinates. Set $N_\sigma(z)=v_\sigma(z)^*G_{\rm cl}v_\sigma(z)$.
Then
\begin{equation}\label{sz:covol}
 \operatorname{CovolumeLe}(y,c)\quad\Longleftrightarrow\quad
 \left|\det\bigl(v_\sigma(y_i)^*G_{\rm cl}v_\sigma(y_j)\bigr)_{ij}\right|
       \le e^c\qquad(\sigma\in\operatorname{Aut}(\mathbb C)).
\end{equation}
The Gram matrix is formed \emph{after} acting on coordinates, so it is positive
semidefinite for every field automorphism. This order matters even with rational
$G_{\rm cl}$, since an arbitrary $\sigma$ can fail to commute with complex conjugation.

\paragraph{Statement and dependencies.}
For a fixed admissible $D$, the hypotheses of
Sections~\ref{A:section}, \ref{B:frontier} and \ref{c3:frontier}
imply the pointwise determinant estimate \eqref{eq:det}.
Theorems A and B produce an independent family with small squared
covolume and controlled index. Theorem C applies to that same family.
Section~\ref{sz:comparison} deduces the estimate and gives the
quantifiers required for Hypothesis~\ref{hyp:scale}.
Appendix~\ref{sz:register} records the construction steps that
remain hypotheses.

\section{Short forms and the index of their translates}\label{A:section}
With the conventions of Section~\ref{sz:setup}, let $X_j,Y_j$ be the two towers of hermitian spaces over $E$, $0\le j\le m$. The upper chain acts on the definite groups $U(X_j)$ and restricts at the top to $U(Y_m)$. The lower chain begins on $U(Y_0)$.
\begin{chatAtheorem}[Short forms and index]\label{A:main}
Assume the estimates $\mathsf A_{\rm loc}$ specified below.
The set
\[
 \mathcal S=\{\phi\in\mathcal I:\ N_\sigma(P_g\phi)
       \le \exp(H-m\Tot_D+V)\text{ for every }g,\sigma\}
\]
spans $S$ over $\mathbb C$. For any finite $t\subseteq G$ and any finite indexed family with $y_i=P_{g_i}\phi_i$, $g_i\in t$, $\phi_i\in\mathcal I$, the target lattice satisfies
\begin{equation}\label{eq:short-index}
 \operatorname{IndexLe}(y,2|t|d_SU).
\end{equation}
\end{chatAtheorem}
\noindent\textit{Hypotheses $\mathsf A_{\rm loc}$.}\label{A:inputs}
For use independently of the formal development, the following statements about the prescribed data suffice.
\begin{enumerate}
\item For free-level evaluation use the integral Weyl polynomial model at each place and its unitary Gram norm as specified in Section~\ref{A:source}. The required short independent columns follow from Euclidean transference.
\item For the chosen Bessel functional and point, the orbit expansion, absolute convergence, finite vacuum decomposition and terminal definite pairing in Section~\ref{A:unfold} hold with the coefficient bounds stated there. For the trivial twist the corresponding one-stage kernels satisfy \eqref{A:onelift}, with the row majorants, CM period identity, class counts and integral branching estimates used in deriving \eqref{A:chain}.
\item At small nonsplit primes assume the vacuum product and one-curve inputs in Appendix~\ref{sec:jets}. At ordinary primes assume the standard-datum integral gauge, and at marked primes the spherical transfer identity, diagonal-neighbour gauge and edge estimates applied in Section~\ref{A:local}. The precise averaging and residue estimates at a marked prime are detailed in Section~\ref{A:edgecalc}.
\item At $N$ assume the shifted-lattice elementary divisor bound, the direct-datum gauge $N^{fb}$ and the one-pair gluing and commutation identities for ramified source forms in Section~\ref{A:local}.
\end{enumerate}
Sections~\ref{A:source}--\ref{A:indexsec} derive the two assertions under these hypotheses. The first assertion retains every translate of each short generator. The second permits every integral preimage, with any finite family size. The source concordance identifies the inputs within the four statements that still require separate mathematical proofs.

\subsection{Short vectors at the source}\label{A:source}
For a source form $v$ write $|v(u)|_{\sigma,\mathrm w}^2$ for the squared Weyl norm of the value at a finite adelic point $u$, with $\sigma$ applied to its coordinates. We use this Weyl norm on the coefficient representation at each level. Say $v$ is $B_0$-short if its values are integral in the local Weyl lattices at every finite place and
\[
 |v(u)|_{\sigma,\mathrm w}^2\le B_0\qquad (\sigma\in\operatorname{Aut}(\mathbb C),\ u\in U(X_0)(\mathbb A_f)).
\]
The target norm uses a sum over classes. We take that sum only after projection.

\paragraph{Short bases at a free level.}\label{A:shortbase}
We use the short-basis theorem for forms at a free level, that is, a level with trivial stabilizer at every double coset. In rank $k$ over the CM field $E$, let $D_0$ be a positive integer in the different, and let $J$ be at least $kf$ plus the number of boxes in columns of height below $k$. Then integral forms with pointwise squared norm bounded by $[6D_0(fk)^2]^{6J}$ at all conjugates span the level space. At the initial level $D_0=q_1d$ is admissible, the columns of height below $k$ have at most $fa$ boxes, and the scale estimate reads
\begin{equation}
 4^f[6q_1d(fk)^2]^{6(kf+fa)}\le e^U.
 \label{A:base}
\end{equation}
Consequently there is a spanning set $\mathcal V_0\subset X$ of $(e^U/4^f)$-short forms.

Here is the lattice argument underlying this input. Choose representatives $c$ for the double cosets. At a free level, evaluating a form at the $c$ identifies the space with a direct sum of copies of the coefficient representation. A vector at $c$ extends to a form supported on its coset. Its integral gauge is bounded by the seed gauge, and its squared norm by the seed norm times the square of the stabilizer order, which here equals one.

At $c$ one needs $k$ independent columns in the moved lattice
\[
 L_c=\{y\in E^k:(c^{-1}y)_f\in\widehat{\mathcal O}_E^{\,k}\}.
\]
The short-basis estimate gives a matrix $Y_c\in\mathrm{GL}_k(E)$ with columns in $L_c$ and every conjugate of every entry bounded by $6D_0(fk)^2$. For the transference argument, use $\frac12\operatorname{Tr}_{E/\mathbb Q}\langle\ ,\ \rangle_E$ in the Minkowski space of real dimension $h=2fk$. Locally $c$ is unitary and the standard integral lattice is hermitian self-dual, with conjugate places paired. Thus $L_c$ is hermitian self-dual globally. If $\mathfrak D$ is the different, the Euclidean dual lattice of $L_c$ is $2\mathfrak D^{-1} L_c$: the trace condition for all integral multiples of a vector in $L_c$ is precisely membership of its hermitian pairings in $2\mathfrak D^{-1}$. It follows that the dual of $(2/D_0)L_c$ is $D_0\mathfrak D^{-1} L_c\subset L_c$. The squared Euclidean length of a nonzero $y\in L_c$ is the trace to $\mathbb Q$ from the totally real field of the totally positive algebraic integer $\langle y,y\rangle_E$. It is a positive integer. Thus the dual minimum is at least one.

For a full Euclidean lattice $\Lambda$ of dimension $h$, Banaszczyk's transference theorem~\cite[Theorem~2.2]{banaszczyk} bounds $\mu(\Lambda)\lambda_1(\Lambda^*)$ by $h$, which in particular gives the coarse $h^2$ used here. Apply it to $(2/D_0)L_c$. When a lattice has covering radius at most $r$, its vectors of length at most $3r$ span the real space. Otherwise a point of length $2r$ perpendicular to their span would have a lattice point within $r$, necessarily short and outside the span. Rescaling gives real-spanning vectors in $L_c$ of length at most $3D_0 h^2/2=6D_0(fk)^2$. Their $E$-span has full rank, since an $E$-subspace of smaller rank has real Minkowski dimension less than $2fk$. Select $k$ independent columns and use the same argument for each determinant shift $c^{e_\tau}$.


Weyl evaluation vectors on arrays of roots of unity span the coefficient space by polynomial identity. Acting by $Y_c$ and the determinant factors at the shifted weights preserves their span. The integrality of the translated matrices and the integral coefficients of the Weyl polynomials give local gauge at most one. Put $R_0=6D_0(fk)^2\ge k$. Each entry after moving a root-of-unity array is at most $kR_0$ in modulus and its squared Hilbert--Schmidt norm is at most $k^4R_0^2\le R_0^6$. Thus the polynomial part costs at most $R_0^6$ per box in the squared norm. With the short representatives for the determinant ideals, the determinant factors cost at most $R_0^{6k}$ per place. This gives the squared bound $[6D_0(fk)^2]^{6J}$. Extension by the coset forms and \eqref{A:base} then give $\mathcal V_0$. For the polynomial norm estimate, the Schur operator in singular-value coordinates multiplies the squared norm by at most the product of squared singular values raised to the row lengths. This follows by realizing the polynomial representation inside tensor products of exterior powers. Equivalently one may bound by the squared operator norm to the power of the number of boxes; this also gives $R_0^{6J}$ here.

At the $f$ chosen places $w$ over the level prime $N$, let $T_w$ be the commuting central translations. Put $r_0=N-1$. The level gives $T_w^{r_0}=1$, and
\[
 S_w=\sum_{i=0}^{r_0-1}T_w^i,\qquad
 S=\bigcap_w\ker S_w,\qquad R=\prod_w(1-T_w).
\]

\begin{chatAlemma}\label{A:ram}
The forms $R\psi$ for $\psi\in\mathcal V_0$ span $S$. Each belongs to $\mathcal I$ and satisfies $|R\psi(u)|_{\sigma,\mathrm w}^2\le e^U$ for every $u,\sigma$.
\end{chatAlemma}
\begin{proof}
$S_w(1-T_w)=0$ gives $R(X)\subset S$. Conversely $r_0-S_w=\sum_i(1-T_w^i)$ is divisible by $1-T_w$ in the commutative algebra of translations. On $S$ it acts by $r_0$. Hence $r_0^f S\subset R(X)$ and $R(X)=S$. Applying $R$ to the spanning set gives the assertion about its image.

Each $T_w$ acts by an integral adelic translation and preserves the integral gauge. Thus $(1-T_w)\psi$ is integral. If both terms have pointwise norm at most $\sqrt{B_0}$ after conjugation, their difference has squared norm at most $4B_0$. Iterating over $f$ places multiplies the bound $e^U/4^f$ by $4^f$. For a form in $S$ the resulting finite-place integrality is precisely membership of $\mathcal I$.
\end{proof}

The same operators control the change of dimension. A vector fixed by $T_w$ descends to the level enlarged by that translation. The relative index is at least $N-1$, since its first $N-1$ powers have distinct cosets. At the enlarged free level every class has as many lifts as the relative index. Evaluation on classes therefore gives $\dim\ker(1-T_w)\le d_X/(N-1)$. For two endomorphisms, $B$ maps $\ker(AB)$ to $\ker A$ with kernel contained in $\ker B$. Iteration and rank-nullity now give
\begin{equation}
 d_S\ge\left(1-\frac f{N-1}\right)d_X\ge d_X/2.\label{A:dimension}
\end{equation}
Here $f/(N-1)\le b^{-8}\le1/2$. This factor of two will enter only at the end of the index calculation.

\subsection{The seesaw and the two scalars}\label{A:unfold}
Fix one generator $\phi=R\psi$ from Lemma~\ref{A:ram}. Write $I_\vartheta$ for the lower chain on principal forms and $B_{\vartheta,\Psi}$ for the composite upper lift with chain datum $\Psi$. A twist $\vartheta$ is a weight-zero automorphic character of $U(1)$ with tame part nontrivial precisely on a finite set of places disjoint from its conjugate and of product norm prime to $MN$. In an $m$-fold twist each of its first $m$ positive powers is nontrivial at every place of that set. The lower chain uses the stage characters $\mu_j^0\vartheta^j$. On good target forms and for an $m$-fold twist we write $\Theta_\vartheta$ for this chain at its prescribed terminal level. We use the seesaw identity with the following quantifiers.
\[
 \exists c_{\rm ss}\in\mathbb C^\times\quad\forall\vartheta\quad\exists\Psi_\vartheta
 \quad\forall v\in S,\qquad I_\vartheta(\Pi v)=c_{\rm ss} B_{\vartheta,\Psi_\vartheta}(v).
\]
The datum consists of finitely many translated vacuum tensors with coefficients and level invariance prescribed by the twist. The seesaw theorem unfolds the lower chain against the Poincar\'e series and realizes the resulting tensor as such a combination. We use its identity for the entire source at a fixed twist. Linearity gives, for $z_v=c_{\rm ss}^{-1}\Pi v$,
\begin{equation}
 I_1(z_v)=B_{1,\Psi_1}(v)\qquad(v\in S).\label{A:seesaw}
\end{equation}
Pair the lower chain with any terminal test form. Transposition moves the last lift onto that form. Unfolding the Bessel period then replaces one level integral by a sum over its finite quotient, weighted by the factors from the quotient measures, and an archimedean integral. Repeating this operation leaves a definite pairing at the top. The archimedean integral contributes a common scalar $c_0$ independent of the tame twist. The finite sum gives a tensor $K_\vartheta$ of evaluated upper-chain data. Linearity of both the pairing and the upper lift pulls this sum inside the lift. Equality against every test form identifies the lower chain with $c_0$ times the upper chain for $K_\vartheta$.

For the one-stage integral one can isolate the analytic assumption. Let $R_j$ be the Bessel group of the current grading and let $W_j$ be the equivariant function produced by the preceding stage. If $K_j(r,h)$ denotes the translated theta kernel, the orbit expansion of its Bessel period reads
\[
 \int_{[R_j]}W_j(r)K_j(r,h)\,dr
 =\sum_{\gamma\in R_{j+1}(F)\backslash U(Y_{j+1})(F)}
          W_{j+1}(\gamma h).
\]
Here $W_{j+1}$ is the period of the kernel's constant term on the specified string orbit. The orbit identity and its absolute majorant are inputs for the prescribed grading. To apply it to a rapidly decreasing automorphic test form $z$, reduce $h$ by a rational element. The sum of the absolute kernel terms is bounded by a constant times the norm of its Weil operator at that reduced point. This norm is $O(\operatorname{ht}(h)^{B_*})$ for some fixed exponent $B_*$, while $|z(h)|\operatorname{ht}(h)^{B_*}$ is bounded. The Bessel quotient and the automorphic quotient have finite mass. The product is integrable, so Fubini permits passage through the two integrals and the series. Rational automorphy of $z$ turns the sum over cosets into one integral of $W_{j+1}z$ over $R_{j+1}(F)\backslash U(Y_{j+1})(\mathbb A)$. The transpose of the level-normalized lift contributes the index of its level in the standard level on both sides. This proves the one-stage passage used in the iteration; its growth and orbit assumptions are part of $\mathsf A_{\rm loc}$.

More precisely, put $B=M^{fn_m^2}$. The finite coset expansion writes $K_\vartheta=\sum_x r_x C_0 V_x$ for glued translated vacua $V_x$, with a common scalar $C_0$. The coefficients $B^{2m}r_x$ are integral away from the twist level, and $\sum_x|\sigma r_x|\le B^m$ at the trivial twist. When $c_0 C_0\ne0$ set $u=B^{2m}/(c_0C_0)$. Otherwise take $u=1$ and the zero tensor. In either case $u c_0K_\vartheta$ has integral-away-from-level coefficients, trivial-twist coefficient sum at most $B^{3m}$, and the level invariance inherited from the coset sum. Replace each translate by a supported representative with the same vacuum. This produces the required datum and the common unit $c_{\rm ss}=u^{-1}$. The expansion into ordinary single lifts below is a second expansion, applied to this already normalized glued tensor.

The scalar cancels before we take a norm.\footnote{\ChatAFile{New/UpperBound/Seesaw/SeesawIdentity}{Seesaw identity}.} Expansion of the composite tensor supplies a weak datum $\Psi'$ such that the right side is the restriction to $Y_m$ of a sum of chains of $m$ single lifts from $X_0$ to $X_m$. The expansion is uniform in $v$. At the trivial twist its conjugated coefficients have $\ell^1$ norm at most $e^{x^{60}}$.

For $s\in G$ let $Q_s z=|G|e_0(sz)$, where $e_0$ is the chosen central idempotent. The normalized map of this section is
\[
 P_s v=c_BQ_s z_v,\qquad c_B\in\mathbb Q_{>0}.
\]
Here $c_B$ depends on the finite-place divisibility. The factor $c_{\rm ss}$ has already disappeared in \eqref{A:seesaw}.

\subsection{The finite gains and the height of the chain}
Let $\mathcal P$ be the primes $p<x_1$ nonsplit in $\mathbb Q(\sqrt{-d})$. For an inert $p$ take $I_p$ to be the positive odd exponents and $v_p=1$. For the ramified prime take all positive exponents and $v_d=1/2$. At a split prime set $s_p=0$. Put
\begin{align*}
 s_p(y,\tau)&=v_p\sum_{e\in I_p}(y/p^e-\kappa(u(-p^e)\tau))_+,\\
 \gamma_p&=m\left[\sum_{\tau\in\Phi_a}s_p(a,\tau)
                    +k\sum_{\tau\in\Phi}s_p(\kappa(\tau),\tau)\right],
 & \mathfrak g&=\sum_{p\in\mathcal P}\gamma_p\log p .
\end{align*}
Only finitely many summands are positive. For an inert prime $1-\chi_d(p^e)$ equals two at odd exponents and zero at even ones. For $p=d$ it equals one. Thus regrouping $\Lambda(n)(1-\chi_d(n))$ by prime powers gives
\[
 2\mathfrak g=m\sum_{\tau\in\Phi_a} S(a,\tau)+mk\sum_{\tau\in\Phi}S(\kappa(\tau),\tau).
\]
Here $\gamma_p$ vanishes for $p\ge x_1$, since $y\le a$ and $\kappa\ge a/x_1$. Substituting \eqref{sz:height} yields the identity
\begin{equation}
 H-m\Tot_D+2\mathfrak g=mA S_0+mk\sum_{\tau\in\Phi}j(\kappa(\tau)).\label{A:gain}
\end{equation}
The coefficient of the signed excess is $-m$ after solving for the size.\footnote{\ChatAFile{New/UpperBound/Size/ShortColumns}{Local gains and short generators}.} It is already present in the finite gain of each prime.

We give the deduction of the restricted-chain estimate from the row, kernel, branch and level estimates. Let $J_i$ be the sum of $j$ over the row lengths of the weight diagrams for lift $i$, let $h_0$ be the half-gap in the branch weights, and put
\[
 C_m=m\biggl[A j(a)+k\sum_{\tau\in\Phi}j(\kappa(\tau))\biggr].
\]
\paragraph{A single kernel.}\label{A:kernel}
Let $n=n_i$, $n'=n_{i+1}$, $q_1=q_0\,\mathrm{towerDeg}$ be the cyclotomic modulus, and write $C$ for the fixed parameter bound. If $\ell_{\tau h}$ are the row lengths and $\lambda=a+1$, the logarithm of the Weyl operator norm for one translated vacuum is at most
\begin{equation}
 C(q_1nn'+1)^2+fnn'\log(2\lambda+2)
 +\tfrac12\sum_{\tau,h}\ell_{\tau h}/\lambda+\tfrac12 J_i.\label{A:onelift}
\end{equation}
To see the origin of the factorials, test the Schur jet against a unit vector in the dual Gram norm. A matrix argument with squared singular values $\rho_{\tau h}$ contributes at most the square root of $\prod\rho_{\tau h}^{\ell_{\tau h}}$, times the input and test norms. Put $u=\log(M/\Omega)$. Applying $t^j\le j!e^t$ with $t=\rho_{\tau h}e^{-u-1/\lambda}$ bounds this square root by
\[
 \exp\left(\tfrac12\sum_{\tau,h}
 [\log(\ell_{\tau h}!)+\ell_{\tau h}u+\ell_{\tau h}/\lambda]
 +\tfrac12 e^{-u-1/\lambda}\sum_{\tau,h}\rho_{\tau h}\right).
\]
The trace bound $\sum\rho\le e^u(2\lambda+2)R$ and $e^{1/\lambda}\ge1+1/\lambda$ make the last term at most $\lambda R$ at the Gaussian radius $R$ of the jet estimate. In the Gaussian integral representation for the jet, the residual Gaussian of exponent $-\|y\|^2/(2\lambda+2)$ has normalized mass $(2\lambda+2)^{fnn'}$. The translated theta vacuum contributes the first term of \eqref{A:onelift} by its uniform growth bound. The jet scale has squared modulus $M$ at every conjugate, which gives $u$ in the formula. These vacuum and Schur estimates hold for every conjugate. For one translated vacuum the conjugate coefficient sum is at most one.

The uniform estimate on a moved vacuum in \eqref{A:onelift} can be reduced further. In rank $n_*=nn'$ let $h=2fn_*$ be the real dimension of its archimedean coordinate space and put $D_*=h+1$. The theta matrix-coefficient identity bounds the square of the absolute vacuum value by a sum of $\exp(-\pi|x|^2/\sqrt d)$ over the support lattice in the moved definite space. The Heisenberg translation contributes a phase. For two lattice points the squared distance is the trace to the imaginary quadratic base of the hermitian norm of their difference. The factors of the inverse different in the finite lattice make this trace integral. A real algebraic integer of the imaginary quadratic field is a rational integer, so distinct points are at distance at least one. Disjoint balls of radius $1/2$ then bound the number at distance at most $R$ by $(1+2R)^h$. Completing the square gives
\[
 e^{-\pi R^2/\sqrt d}\le e^{\sqrt d D_*^2}e^{-D_* R},\quad
 \sum_x e^{-D_*|x|}
 \le D_*\int_0^\infty(1+2R)^h e^{-D_* R}\,dR
 \le D_*3^h\le e^{3h}.
\]
Here $1+2R\le3e^R$ and $D_*=h+1$. Thus the logarithm of the vacuum norm is at most $(\sqrt d\,D_*^2+3h)/2\le(2d+6)(q_1n_*+1)^2$. To pass to the conjugate integral vacuum, use the semilinear conjugation of the vacuum to a scalar $a_*$ times a unitary central translate. Evaluating at finite points bounds its squared invariant inner product by $\exp(2C_{\rm vac}n_*q_1^2)$, using the finite-point conjugate theta bound. With the vacuum normalized to squared norm one this gives $|a_*|\le\exp(C_{\rm vac}n_*q_1^2)$. The sum of the two exponents is at most $C(q_1n_*+1)^2$ by the definition of $C$. The matrix-coefficient identity, semilinear conjugation, finite-point theta bound and the Gaussian representation of the Schur jet are the analytic inputs in this deduction.

For an explicit comparison put $X_*=\log x$, $Z=\log L$. The size conditions give $C\le X_*/10^{14}$, $100Z\le X_*$, $q_1\le C e^{11X_*+Z/5}$, $f\le e^{C+11X_*+Z/5}$ and $n,n'\le e^{11X_*/2+12}$. The first term is bounded by $4C^3e^{48+44X_*+2Z/5}\le x^{45}$. The second is bounded by $2e^{C+24+22X_*+6Z/5}\le x^{23}$ using $\log(2a+4)\le2L$. Since $\ell\le2a$, the third is at most $fn\le x^{17}$. Their sum is at most $x^{50}$.

\paragraph{Row lengths.}
Write $n_i$ for the rank of $X_i$. In the diagram at place $\tau$ the comparison row lengths $r_{\tau,l}$ are $a\mathbf1_{\Phi_a}(\tau)$ for $l=0$, $\kappa(\tau)$ for $0<l<k$, and zero thereafter. The diagram estimate bounds the actual length by $r_{\tau,l}+e_l\le2a$, with
\[
 e_l=h_0^2+\mathbf1_{l=0}(b+a/x_1).
\]
Each extra box increases the logarithmic factorial by at most $\log(2a+1)$. Since $2a+1\le b^2$ and $0\le\log(M/\Omega)\le23L$, it adds at most $25L$ to the jet size. Summing over rows and places gives
\[
 J_i\le A j(a)+(k-1)\sum_\tau j(\kappa(\tau))
             +25fL(b+a/x_1+n_i h_0^2).
\]
Put $X_*=\log x$. The admissible bounds $n_i\le n_m\le e^{11X_*/2+12}$, $h_0\le e^{40X_*+25}$ and $L\ge10^{12}X_*$ give $n_i h_0^2\le e^{171X_*/2+62}\le b$. Also $a/x_1\le b$. The total row error is thus at most $75mfLb$. Since $m\le2\sqrt{x}$, $a=xb$, and $L\ge128$, this is at most $faL^3/2=U/2$. Using $j(\kappa(\tau))\ge0$ to replace $k-1$ by $k$ proves
\begin{equation}
 \sum_{i<m}J_i\le C_m+U/2.\label{A:row}
\end{equation}

\paragraph{Choice of the level.}
Let $w=f^{-1}\sum_\tau\log\kappa(\tau)$ and $\Gamma=\log(M\,\mathrm{rd}_E/(2\pi))$, where $\mathrm{rd}_E$ is the root discriminant. The CM period formula with the normalization of $c_1$ gives $\log(M/\Omega)=1+c_1+\Gamma$. The Stirling bound $\log(a!)\le a\log a-a+1+\tfrac12\log a$ and \eqref{sz:height} therefore give
\[
 j(a)-S_0\le a(w+\Gamma)+1+\tfrac12\log a.
\]
The lower endpoint of the admissible window for the level prime is $\log N\ge mw+(m-1)\Gamma-15L/2$. Subtract $a\log N$ after multiplying the preceding inequality by $m$. The terms $ma(w+\Gamma)$ cancel. As $m(1+\tfrac12\log a)\le aL$ and $\Gamma\le23L$, we obtain
\[
 m[j(a)-S_0]-a\log N
 \le m(1+\tfrac12\log a)+a\Gamma+\tfrac{15}2aL\le32aL.
\]
Multiplication by $A\le f$ yields
\begin{equation}
 A\{m[j(a)-S_0]-a\log N\}\le32faL.\label{A:level}
\end{equation}
In particular, by \eqref{A:gain},
\[
 C_m-Aa\log N\le H-m\Tot_D+2\mathfrak g+32faL.
\]
Write $e_\tau$ for the difference of the total weights at $X_m$ and $Y_m$. Summing the branch weights gives $\sum_\tau e_\tau=Aa-fmh_0$. The branch map at the block embedding is an integral matrix times $\prod_\tau\tau(\bar\pi_N)^{-e_\tau}$. The squared modulus of this scalar is $N^{-\sum e_\tau}$ since $\pi_N\bar\pi_N=N$. The squared operator norm of the integral branching matrix is at most $e^{x^{70}}$ by the branch estimate. Restricting the last form at the block embedding $Y_m\hookrightarrow X_m$ therefore multiplies the squared norm bound by at most $\exp(x^{70}-(Aa-fmh_0)\log N)$.

The special feature of the branching matrix is the size of its denominator. Separate the columns at each place into those which keep their height and those which lose their bottom box. If their counts are $u_h,v_h$, the output has $u_h+v_{h+1}$ columns of height $h$, up to the determinant powers at full height. The prescribed interlacing has $u_h=0$ for $h<k$. In the extremal filling of the output diagram the top $k-1$ rows can therefore be wedged on unchanged (with the usual cap at the rank). The signed product of its minors after a row permutation is the product of the two Weyl kernels for the kept and shortened columns. The integral Weyl straightening identity clears this polynomial by the hook number $c$ of the diagram below those rows. If $l$ bounds the number of output boxes in columns of height at least $k$ and below the rank, the hook estimate gives $1\le c\le l^{3l}$. Large strings of low columns contribute no boxes to this estimate. Coefficient extraction in the integral polynomial identity supplies the intertwiner $B$ with highest coefficient $c$ times that product of kernels.

For the norm estimate the two minor identities associated with this construction realize the Gram matrix of the output evaluation vectors as that of ordinary tensor vectors $m(y)$, and give
$\langle v(y),B\xi\rangle=\langle m(y),c E\xi\rangle$, where $E$ embeds the input isometrically in that tensor space. Evaluation vectors at nonsingular arrays span by polynomial identity. Writing $B\xi=\sum_y a_yv(y)$ gives $\|B\xi\|^2=\|\sum a_y m(y)\|^2=\langle\sum a_y m(y),cE\xi\rangle$, hence $\|B\xi\|\le c\|\xi\|$. For these diagrams take $l=fn_m^2\,\mathrm{gap}$: there are at most $n_m$ contributing heights, at most $\mathrm{gap}$ columns at each such height and place, and at most $n_m$ boxes per column. The data scales give $6l\log l\le x^{70}$. This proves the asserted numerical bound from the integral straightening, hook and tensor identities. Those identities are the representation-theoretic input to this use of $\mathsf A_{\rm loc}$.

\begin{chatAproposition}\label{A:chainprop}
Suppose $|v(u)|_{\sigma,\mathrm w}^2\le e^{B_0}$ at every source point for fixed $\sigma$. Under the preceding estimates, the single-kernel and class bounds below give
\begin{equation}
 |B_{1,\Psi_1}(v)(y)|_{\sigma,\mathrm w}^2\le
 \exp(B_0+H-m\Tot_D+2\mathfrak g+U)\qquad(y\in U(Y_m)(\mathbb A_f)).\label{A:chain}
\end{equation}
\end{chatAproposition}
\begin{proof}
At a fixed translation and tuple of classes the $i$th kernel multiplies the norm by at most $\exp(x^{50}+J_i/2)$, uniformly in all points. At each intermediate rank, bounding the maximal-level class number and then multiplying by the congruence index gives at most $\exp(x^{29}(1+\log q))$ classes for a level containing principal level $q$. Indeed, each maximal-level class splits into at most $q^{2f(2n_m)^2}$ classes, by reduction of its matrices modulo $G$. For the prescribed levels this bounds the number by $\exp(x^{30})$. There are therefore at most $\exp(mx^{30})$ tuples and the coefficient sum costs at most $\exp(x^{60})$. The triangle inequality and the uniform input bound give an upper bound for the norm before restriction of
\[
 \exp\!\left(B_0/2+mx^{30}+x^{60}+mx^{50}
                     +\tfrac12\sum_{i<m}J_i\right).
\]
Square this expression, apply the branch estimate, and then use \eqref{A:row}. The logarithm of the bound becomes at most
\[
 B_0+H-m\Tot_D+2\mathfrak g+U/2+
 \bigl[x^{70}+fmh_0\log N+2(mx^{30}+x^{60}+mx^{50})+32faL\bigr].
\]
To bound the bracket, $32faL\le U/4$ since $L^2\ge128$. For the other terms the admissible size bounds, with $\lambda=\log L$, give
\[
 f\le e^{11X_*+\lambda/5+C_P},\quad
 m\le e^{X_*/2+2},\quad h_0\le e^{40X_*+25},\quad
 \log N\le e^{X_*/2+\lambda+32}.
\]
Here $C_P$ bounds the constants of the fixed parameters $(q_0,d)$, $C_P\le\lambda$, $\lambda\le X_*/100$, and $X_*\ge59$. Thus $fmh_0\log N\le e^{52X_*+6\lambda/5+C_P+59}\le x^{54}$. Also $2(mx^{30}+x^{60}+mx^{50})\le6x^{60}\le x^{61}$. Including restriction, their sum is at most $x^{70}+x^{54}+x^{61}\le x^{71}$. Finally $a=xe^L$ and $L\ge10^{12}X_*$ give $x^{71}\le U/4$. The bracket uses at most $U/2$, which together with the row error in \eqref{A:row} accounts for the full $U$.\footnote{\ChatAFile{New/UpperBound/Size/ArchimedeanSize}{Single-kernel, branch and composite estimates}.} This proves \eqref{A:chain}.
\end{proof}
Applying \eqref{A:chain} and \eqref{A:seesaw} to the generator of Lemma~\ref{A:ram} gives the pointwise squared bound $e^{B_1}$ for $I_1(z_\phi)$, where
\[
 B_1=H-m\Tot_D+2\mathfrak g+2U.
\]
One $U$ comes from the generator and one from the chain.
\subsection{Projection, rounding, and the final norm}
The squared norm of a target level form $F$ at $\sigma$ is
\[
 \|F\|_\sigma^2=\sum_{[y]}\frac{|F(y)|_{\sigma,\mathrm w}^2}{|\operatorname{Aut}([y])|},
\]
with the sum over double cosets. The target lattice of Section~\ref{sz:setup} uses this norm on $F=\Theta_1 z$. Compatibility of the lift with good projection and translation expresses the lifted value of $Q_s z$ at $y$ as
\[
 I_1(Q_s z)(y)=\sum_{t\in G}|G|e_0(t)\,I_1(z)(y\,\widetilde{ts}),
\]
where the tilde denotes the lift of the translate in the first block. The central idempotent is a sum of isotypic projectors. The character formula bounds $|G|\,|\sigma(e_0(t))|$ by the sum of the squares of their dimensions: every character value is a sum of roots of unity and the sum of dimension squares over the full regular representation is $|G|$. Thus $|\sigma(e_0(t))|\le1$. For $z=z_\phi$, taking the triangle inequality, squaring and summing over classes therefore gives
\begin{equation}
 N_\sigma(Q_s z_\phi)\le \#\{[y]\}\,|G|^4 e^{B_1}\le e^{x^{99}+B_1}.
 \label{A:proj}
\end{equation}
Here the inverse stabilizer weights are at most one, the class count is at most $e^{x^{30}}$, and the group order bound gives $x^{30}+4\log|G|\le x^{99}$. This is the first passage from a pointwise bound to a class norm.

The local divisibility theorem supplies errors $e_p\ge0$ with $\sum_{p\in F}e_p\log p\le U$ for every finite $F$. Let $\mathcal P'=\mathcal P\cup\{N\}$, and set
\[
 q_p=\lfloor\gamma_p-e_p\rfloor,\qquad c_B=\prod_{p\in\mathcal P'}p^{-q_p}.
\]
Since $N\ge x_1$, $\gamma_N=0$. Rounding costs at most one copy of $\log p$ for each $p\in\mathcal P'$. Consequently
\begin{align}
 \log c_B&\le-\mathfrak g+\sum_{p\in\mathcal P'}e_p\log p+\sum_{p\in\mathcal P'}\log p
          \le-\mathfrak g+U+b,\label{A:floor}
\end{align}
where $\sum_{p\in\mathcal P'}\log p\le\lceil x_1\rceil x_1+\log N\le b$ by admissibility.\footnote{\ChatAFile{New/Reduction/ShortAndEnough}{Projection, rounding and normalization}.}

\begin{chatAproposition}[Theorem A(a)]\label{A:short}
Under $\mathsf A_{\rm loc}$, the following holds.
The forms $\phi\in\mathcal I$ satisfying
\[
 N_\sigma(P_s\phi)\le\exp(H-m\Tot_D+V)\qquad(s\in G,\ \sigma\in\operatorname{Aut}(\mathbb C))
\]
span the ramified source $S$.
\end{chatAproposition}
\begin{proof}
For the fixed generator $\phi=R\psi$, every $\sigma$ fixes $c_B\in\mathbb Q_{>0}$. Linearity of the lift and homogeneity of the squared class norm give $N_\sigma(P_s\phi)=c_B^2 N_\sigma(Q_s z_\phi)$. By \eqref{A:proj} and \eqref{A:floor} the exponent is bounded by
\begin{align}
 x^{99}+B_1+2\log c_B
 &\le x^{99}+H-m\Tot_D+2\mathfrak g+2U+2(-\mathfrak g+U+b)\notag\\
 &=H-m\Tot_D+x^{99}+4U+2b.\label{A:normalize}
\end{align}
For the last size comparison, the admissible inequalities $10^6\log x\le\log x_1$ and $10^6\log x_1\le L$ give $x^{99}\le b$. Also $b\le fa$ and $L\ge128$, so $x^{99}+4U+2b\le fa(3+4L^3)\le faL^4=V$. This bound holds for every translate and conjugate of the chosen generator. Lemma~\ref{A:ram} gives its integrality and gives spanning as $\psi$ ranges over $\mathcal V_0$.
\end{proof}
\subsection{A module on which every prime-to-level twist gains}\label{A:local}
Fix $p$ and $\iota$ over $p$. Write $M_\iota$ for the module of integral forms in $X$. In this subsection $v_\iota(F)\ge t$ means that the moved Weyl gauge of $F$ is at most $p^{-t}$ at every finite adelic argument. The level of a twist $\vartheta$ has norm $Q_\vartheta$. Write
\[
 h_p=\frac{faL^2}{p}
 \begin{cases}m,&p\equiv\pm1\pmod m,\\1,&\text{otherwise},\end{cases}
 \qquad
 \epsilon_p=\begin{cases}2h_p,&p\in\mathcal P,\\2mfb,&p=N,\\0,&\text{otherwise}.\end{cases}
\]
We take these explicit errors as witnesses for the existential divisibility estimate. The published normalization uses a chosen witness. All the deductions about $c_B$ use only its witness properties, so they apply to the published choice as well.

For modules $M',S'$ in the same ambient space, write $[M':S']\mid p^j$ to mean that for every finite family in $M'$ some integral square matrix with determinant dividing $p^j$ has all row combinations in $S'$. This notation makes sense even when $S'$ extends beyond $M'$. Pulling back along a linear map, restricting the starting module, and composing two such clearings preserve this property with the evident additive determinant costs. For restriction to $S$ the row combinations automatically stay in $S$.

\begin{chatAlemma}[Simultaneous local clearing]\label{A:clear}
Under the finite-place inputs of $\mathsf A_{\rm loc}$, the following holds.
There exist $j\in\mathbb N$ and a module $S_\iota$ with
\[
 [M_\iota:S_\iota]\mid p^j,\qquad j\le d_X \epsilon_p.
\]
For every $v\in S_\iota\cap S$, every twist with $p\nmid Q_\vartheta$, and every chain datum $\Psi$ for that twist,
\[
 v_\iota(B_{\vartheta,\Psi}v)\ge\gamma_p-\epsilon_p.
\]
\end{chatAlemma}

\paragraph{The cost of a collision.}\label{A:jets}
Here is the local jet calculation used in the proof. Let $p\in\mathcal P$, fix a stage $i<m$, put $n'=n_{i+1}$, and erase the height-one columns at a set $E_0\subseteq\Phi$. Put $l_\tau=a\mathbf1_{\Phi_a\setminus E_0}(\tau)$. Denote by $c_\tau,b_\tau$ the number of columns and boxes of heights $2$ through $n'-1$ at $\tau$. On writing $\tau'=u(-p^e)\tau$, the contribution compared to the prime sum is
\[
 (l_\tau/p^e-l_{\tau'}-c_{\tau'})_+
 +(b_\tau/p^e-b_{\tau'})_+ -c_\tau/p^e
 \qquad(e\in I_p).
\]
Set $C_p(E_0)=\sum a/p^e$, over $e\in I_p$ and $\tau,\tau'\in\Phi_a\setminus E_0$. The predicted stage valuation $D_{p,i}(E_0)$ includes the gain on the erased boxes. The jet estimate is
\begin{equation}
 D_{p,i}(E_0)\ge \gamma_p/m-C_p(E_0)
 +v_p|\Phi_a\cap E_0|a/x_1-faL^2/(mp).\label{A:collision}
\end{equation}
Here $v_p$ is the uniformizer valuation defined above. The two positive parts record respectively the unpaired long columns and the unpaired boxes of the other columns. The second one is compared with $k(\kappa(\tau)/p^e-\kappa(\tau'))_+$ using
\[
 k\kappa(\tau)\le b_\tau\le k\kappa(\tau)+x^{60},\qquad
 c_\tau\le\kappa(\tau)+x^{60}.
\]
The positive part is 1-Lipschitz. Only exponents below $\lceil\log_p\lceil x_1\rceil\rceil$ can make a positive contribution to the desired weight sum. The comparison costs at most $2u$ per place for $u=a/(mp)$, after allowing the long-long collisions at face value $a/p^e$. When the partner is also long, subtracting its length can indeed exhaust the positive part. Erasing a long column instead leaves a surplus $v_p a/x_1$ per erased long place. Comparing the resulting partner sum to the predicted divisibility costs at most another $2fu$. The factorial digit loss is at most $fn'^3(1+\log_p\max(B,1))\le fu$, where $B$ counts the boxes below height $n'$. Thus the total error is at most $5fu\le faL^2/(mp)$.

For clarity, the analytic input to this count is the normal form for the local vacuum series after rescaling its variables by a scalar $\varpi$ of valuation $1/(p^2-1)$ at an inert prime and $1/(2p-2)$ at $p=d$. Its integral factor multiplies an exponential of Frobenius powers of pairings between partner variables. Explicitly, with $\Gamma_e=\sum_{\tau,h}z_{\tau h}^{p^e}z_{u(-p^e)\tau,h}$, it takes the form
\[
 F(\varpi z)=W(z)\exp\left(\frac{\varpi^{p+1}}{p}
          \sum_{e,j} a_{ej}\Gamma_e^{p^j}\right),
\]
where $W$ and $a_{ej}$ are integral. The exponents are odd in the inert case and include zero in the ramified case. Resolution of the vacuum series gives a product of CM theta factors times an integral series. The normal form for these factors persists under the integral rescaling by the period factor. For a derivative of degree $B$, perform the charged change of directions in each matrix column: on any partner pair touching an erased height-one column, at least one of its two row charges must be positive. This makes the relevant couplings divisible by the local prime in degrees through $B$. The rescaled series then has integral coefficients in degree $B$. Undoing the chain rule costs the sum of the charges of the entries in the filling. In the derivative expansion, coupling one variable at $\tau'$ consumes $p^e$ powers at $\tau$. The minor products enforce the row and column constraints above. The charged-monomial estimate subtracts the allowable paired powers and uses $\operatorname{ord}_p(t!)=(t-s_p(t))/(p-1)$ for the remaining factorials, with $s_p(t)$ the base-$p$ digit sum. This is where the logarithmic term enters. Appendix~\ref{sec:jets} proves the charged-frame and coefficient deductions, starting with the one-curve Frobenius series, its integral period rescaling and the moved-vacuum factorization. At $p=d$ the zero-step partner term subtracts a correction; its contribution at each place is at most $\tfrac12(3a/(dm))\le2a/(dm)$ and is included in the second $2fu$ allowance above. Equation~\eqref{A:collision} uses their version for the moved vacuum with the specified columns erased.

\begin{proof}[Proof of Lemma~\ref{A:clear}]
There are four cases. Away from $N$, expand the glued datum into weak chain data. Coefficients are integral at $p\nmid Q_\vartheta$. One may bound each translated chain and then take the maximum of the gauges. A stage bound propagates by multiplication, and restriction to $Y_m$ preserves the bound away from $N$.

\emph{Outside $\mathcal P\cup\{N\}$.} Here $p\ne d$ and $\gamma_p=0$. The standard datum at each stage has gauge multiplier bounded by a power of $|\overline\pi_M|_p/|\overline\pi_M d|_p=1$. Take $S_\iota=M_\iota$, $j=0$.

\emph{Small nonsplit, nonexceptional.} Put $r=\lceil\log_p m\rceil$ and let $o$ be the order of the class of $-p$ at the base cyclotomic modulus. Exceptionality requires inertness, $p\not\equiv\pm1\pmod m$, and odd $o\le r$. Without erasure, summing \eqref{A:collision} gives valuation at least $\gamma_p-h_p-m C_p(\varnothing)$. Always $C_p(\varnothing)\le 2fa/p$ by the geometric series. For $p\equiv\pm1\pmod m$ this already gives $mC_p(\varnothing)\le h_p$. At the ramified prime there is no collision. At any other inert, nonexceptional prime outside those classes, a collision requires an odd multiple of $o$. If one occurs its exponent is at least $r+1$. The geometric tail is at most $2fa/p^{r+1}$ and $p^r\ge m$. Thus $m C_p(\varnothing)\le h_p$ in every case here. Again the entire module works. Also, for every small nonsplit $p$, multiplication by $p^{j_p}$ with $j_p=\lceil 2mfa/p\rceil$ pays the crude collision cost on all of $M_\iota$.

\emph{Exceptional.} With $\iota$ fixed choose a full mark $\Lambda$ above $p$. Its primes are induced by long places, it is disjoint from its conjugate, and every long place induces a prime in $\Lambda$ or its conjugate. Let $E_\Lambda$ consist of the archimedean places inducing primes of $\Lambda$. Along a colliding odd Frobenius step the two induced primes are conjugates. Hence $C_p(E_\Lambda)=0$. Write
\[
 A_w=a\#\{\tau\in\Phi_a:\tau\text{ induces }w\},\qquad
 Y=\prod_{w\in\Lambda}p^{A_w}T_{r,w},\qquad
 S_\iota=Y^mM_\iota+p^{j_p}M_\iota .
\]
Here $T_{r,w}$ is the rank-$r$ local Hecke operator from the transfer identity. In each stage move one copy of the product from the input to the data. The local transfer expands into integral combinations of neighbouring data after a denominator of exponent $t_p=rf(n_m+r)$ per marked prime. The diagonal neighbour contributes at most $r$ negative valuation units per low column at a marked place. The factors $p^{A_w}$ pay for the long columns there. The log gauge in base $p$ of this stage increases by at most
\[
 |\Lambda|t_p+r\sum_{\tau\in E_\Lambda}c_\tau-D_{p,i}(E_\Lambda).
\]
In this case $v_p=1$. For a marked place $\tau$ the local weight is $a/x_1$ up to the prescribed integer rounding. If $p\le m$, the low-column bound and the admissible scales give $r c_\tau\le a/(mp)+\mathbf1_{\Phi_a}(\tau)a/x_1$. If $p\ge m$, then $r=1$, exceptionality forces $o=1$, and all places inducing a marked prime lie on its long-place orbit. The same bound follows. Summing and absorbing $|\Lambda|rf(n_m+r)$ in the remaining $L^2$ allowance shows that the first two terms together are at most $h_p/m+|\Phi_a\cap E_\Lambda|a/x_1$ by the low-column transfer estimate. Inserting \eqref{A:collision} bounds the increase by $(2h_p-\gamma_p)/m$. After $m$ stages this proves the valuation claim on the first summand of $S_\iota$. The crude collision estimate proves it on the second. The ultrametric inequality proves it on their sum.

For the transfer denominator, let $q_w=z^2$ be the residual cardinality, $t_0$ the rank increment minus one, and $T_{e,i}$ the spherical operator at rank $n_e$ with $i$ uniformizers on the diagonal. The local input for one lift $F_e$ is the polynomial identity
\[
 H_{e+1}(Z)F_e=R(Z)F_e H_e(z^{t_0+1}Z),\quad
 H_e(Z)=\sum_{i=0}^{n_e}(-1)^iq_w^{\binom i2}T_{e,i}Z^i ,
\]
with $R(0)=1$, $R$ integral and independent of the translated vacuum on the support. This is an identity on the spherical source at $w$. Invert $R$ in the integral power series ring and read the coefficient of $Z^i$. It expresses $z^{(t_0+1)i}q_w^{\binom i2}F_eT_{e,i}$ as an integral combination of $T_{e+1,i'}F_e$ for $i'\le i$. For $i\le r$, the $p$-valuation of the prefactor is at most $frt_0$ since $z=p^h$ for $2h\le f$ and $h((t_0+1)i+i(i-1))=h(t_0i+i^2)\le2hrt_0$ for $i\le r\le t_0$. At each subsequent stage no larger index occurs, so after $j$ stages $p^{jfrt_0}$ suffices. Increasing the exponent to $rf(n_m+r)$ incurs an integral scalar since $jt_0\le n_m$. This proves the uniform transfer used above from the displayed spherical polynomial identity.

It remains to bound the index of that sum. For one marked prime the long places inducing $w$ or $\bar w$ form a cycle of length $\nu$ alternating between the ends, with Frobenius power $P_0=p^o$, residue cardinality $q=P_0^\nu$, and $A_w=a\nu/2$. Use the edge elementary-divisor estimate with degree $a$, rank $k$, block size $r$, and exceptional proportion $\theta=1/(m^2f)$. Its size conditions are $0<r\le k/2$, $0<\theta\le1$, $q\ge a^4$ and $a\ge(\nu k\log q/\theta)^5P_0^{3r}$, as verified for the exceptional orbit. In edge coordinates one restricts to the edges at each end and pairs the restrictions. The two end determinant losses together cost at most $c_{\rm div}d_X$. The middle pairing admits an approximate inverse up to at most $\theta d_X$ further generators, at cost $2c_{\rm div}d_X$. Adding $p^{j_p}M_\iota$ pays at most $j_p$ on each further generator. Here
\[
 c_{\rm div}=\nu a\{P_0^{-2r}+r^2/(kP_0)\},\qquad
 3c_{\rm div}\le 12 A_w/(mp).
\]
Thus the exponent per dimension of one scaled operator modulo $p^{j_p}$ is at most $12A_w/(mp)+\theta j_p$. These costs add under products. Indeed, modulo $p^{j_p}M_\iota$, clear first into the image of one operator, lift the coefficients, and clear again before applying that operator. Summing $m$ copies and using $\sum_w A_w\le fa$ and $|\Lambda|\le f$ gives
\[
 m\sum_{w\in\Lambda}\{12A_w/(mp)+\theta j_p\}
 \le 14fa/p+1/m\le h_p.
\]
This proves the asserted colength. For the displayed estimate on $c_{\rm div}$ use $P_0\ge p$, $P_0^{2r}\ge mp$ and $mr^2\le k$. Both terms in braces are at most $1/(mp)$. To identify the input about the middle pairing more precisely, take $w_0^2=p$, let $c_0'$ be the multinomial exponent for rank $k$, and work to precision $w_0^{5\nu c_0'}$. The edge averaging operator acts in the pairing as $p^{\nu c_0'}$ modulo that precision. The space of possible residue defects has dimension times the number of parahoric cosets at most $(1+\theta)$ times the number of monomials. The origin basis accounts for the first copy of that number. At most a proportion $\theta$ of additional generators remains. Section~\ref{A:edgecalc} gives the low-defect count and the averaging argument. It identifies the lattice-point and local representation estimates still needed for the edge factorization.

\emph{The level prime $N$.} Here $\gamma_N=0$. Write $\delta_0$ for the shifted diagonal source lattice and $U_N$ for the sum of translates called the up operator. Use as $S_\iota$ the source forms whose image under $U_N^{m/2}$ is integral at $\delta_0$. If $d_*$ is the summed degree of the last coordinate over the coefficient basis and classes, symmetry of the coordinates gives $k d_*\le d_X fa$. Changing from the dual shifted lattice to the shifted lattice costs at most $d_*$. Indeed, diagonal rescaling multiplies each monomial by its last-coordinate power of $N$, so the determinant exponent is the sum of those degrees. The up operator sends shifted-integral forms into the dual shifted lattice. Pull back the change-of-lattice bound successively. The initial change and the $m/2$ copies cost at most $(1+m/2)d_*\le mfb d_X$. On ramified inputs, gluing the chain in pairs replaces the original input by $U_N^{m/2}v$ on the shifted lattice and contributes a scalar of gauge at most $N^{mfb}$. Each direct lift increases the shifted gauge by at most $N^{fb}$. Final restriction takes that terminal shifted bound to the ordinary target gauge. The total loss is at most $2mfb$, as required.
In more detail each pair of direct stages arises from expansion of the glued kernel into primitive last-column data. The diagonal shifts at the shared odd rank occur with opposite exponents and cancel. The remaining sum over primitive vectors gives their integer cardinality. Moving the resulting up operator to the input gives the product of the commutation scalars at the preceding stages. For pair $i$ the scalar is therefore a product of the gluing scalar, a nonnegative integer, and commutation factors from stages $0,\ldots,2i+1$. The pair's gluing scalar has gauge at most $N^{2fb}$ and all commutation factors have gauge at most one. Multiplication over the $m/2$ pairs gives $N^{mfb}$ on the scalar and $U_N^{m/2}$ on the input. The source ramification is required in the primitive-column identity. The one-pair local Weil action and its scalar estimates are the inputs to this iteration.

\end{proof}

\paragraph{Summing the errors.}\label{A:loss}
Let $K=4faL^2\log x_1$. For ordinary small primes, $\epsilon_p\log p\le K/p$. For $p\equiv\pm1\pmod m$, write $p=mc\pm1$, $1\le c\le x_1$. Then $m/p\le2/c$, giving $\epsilon_p\log p\le K/c$. The level term satisfies $2mfb\log N\le K$ by the upper level bound and $a=xb$. Charge it to $c=1$. For a given $c$, at most four entries $c,mc-1,mc+1,N$ have been charged. Hence for every finite set of primes $F$,
\begin{equation}
 \sum_{p\in F} \epsilon_p\log p\le4K(1+\log x_1)\le faL^3=U .\label{A:budget}
\end{equation}
The last inequality uses $16\log x_1(1+\log x_1)\le L$ from the admissible scales.

\subsection{One matrix for all allowed twists}\label{A:indexsec}
Use the quantifiers of \eqref{sz:index}. The local matrices act on the original indexing of the family, and the integral rows must work for every allowed $m$-fold twist of level norm coprime to $p$.

\begin{chatAproposition}[Theorem A(b)]\label{A:index}
Work under $\mathsf A_{\rm loc}$.
For any finite translate set $t$, any family of $y_i=P_{g_i}\phi_i$ with $g_i\in t$ and $\phi_i\in\mathcal I$ satisfies $\operatorname{IndexLe}(y,2|t|d_SU)$.
\end{chatAproposition}
\begin{proof}
Fix $p,\iota$. Lemma~\ref{A:clear}, restricted to $M_\iota\cap S$, and \eqref{A:budget} prove existence of a common error witness. Use precisely the witness $e_p$ selected in $c_B$, with its corresponding ramified clearing modules. Apply its clearing property to the preimages belonging to a single translate $g$. Its matrix has determinant valuation at most $d_Xe_p$. Each cleared source row $v$ remains ramified. For every allowed twist choose the seesaw datum for that twist. The chain of $z_v$ has valuation at least $\gamma_p-e_p$, and $|c_B|_p=p^{\lfloor\gamma_p-e_p\rfloor}$. Thus the scaled chain is integral for all the twists. At primes outside $\mathcal P\cup\{N\}$ the same formula holds with zero exponent.

Projection and translation preserve this assertion. Indeed the finite group character formula writes each $|G|e_0(t')$ as a sum of character values times integer dimensions, hence as an algebraic integer. The lifted first-block translates preserve integral values, so summing their integral multiples stays in the lattice. By linearity each cleared row of the $y_i$ for this $g$ is integral under every allowed lift. Partition the indices by one witness $g_i$ each and place the matrices for the nonempty parts on the diagonal. Determinant valuations add. Thus $\delta_p=|t|d_X e_p$ works simultaneously at all $\iota$ above $p$. By the selected witness budget and \eqref{A:dimension}, its total logarithmic cost is at most $2|t|d_SU$. The empty family uses the determinant $1$.
\end{proof}
For a family selected from $P_g\mathcal S$, the archimedean estimate applies to the original vectors and the finite-place clearing applies to those same indexed vectors. In particular, if $n$ such vectors are independent, Hadamard bounds the Gram determinant of their untwisted lifts at each conjugate by $\exp(n(H-m\Tot_D+V))$.

\subsection{Residue dimension and averaging at an exceptional prime}\label{A:edgecalc}
Retain the notation of the exceptional case. The orbit length is constrained by $m\le o\nu$ and $\nu\le f$. Hence $q=p^{o\nu}\ge3^m\ge a^4$ at the exceptional primes ($p\ge13$). For the other large-degree condition, $o\le r$ and the tower scale imply $P_0^{3r}\le x_1^{3r^2}\le\sqrt b$. Use $\log q\le o\nu\log x_1$ in the remaining factor. The admissible bounds on $\nu,k,r,f$ give $(\nu k\,o\nu\log x_1/\theta)^5\sqrt b\le a$. This verifies the direction of both large-degree requirements: the residue field dominates the degree, which in turn dominates the precision needed to control the defects.

Here is the numerical part of the residue dimension estimate. In the cycle of $\nu$ letters take polynomials of degree $a$ in each group of $k$ variables. Denote their monomial count by $D_a=\binom{a+k-1}{k-1}^{\nu}$. Put $t$ equal to the defect threshold for the congruence precision, and $E'=4tP_0^{2r}$. Let $s_r$ be the parahoric index, and $S_{\rm res}$ the residue subspace whose translates by the block Levi group have vanishing coefficients at exponents of defect exceeding $3t$. The lattice-point estimate bounds $\dim S_{\rm res}$ by a Levi flag factor times the volume of the defect cone in a product of degree slabs with upper degrees at most $a+E'$. The fraction of a slab in the cone is
\[
 \prod_{j=1}^k \frac{q-1}{q^j-1}.
\]
Its reciprocal counts complete flags. A choice of the distinguished $r$-space and flags in its two blocks contributes at most this many flags. Thus after multiplying by $s_r$, the flag factor and the cone fraction cancel. Bounding each enlarged slab by the volume of a simplex slice gives
\[
 s_r\dim S_{\rm res}
 \le \exp\bigl(\nu(k-1)E'/a\bigr)
       \left(\frac{a^{k-1}}{(k-1)!}\right)^\nu
 \le (1+\theta)D_a.
\]
Indeed, $8\nu ktP_0^{2r}\le\theta a$ gives $2\nu(k-1)E'/a\le\theta\le1$, and $e^z\le1+2z$ for $0\le z\le1/2$. This explains how a proportion $\theta$, rather than a full dimension per edge, reaches the colength bound. The geometric input in this calculation is the enclosure of the admissible lattice points by the enlarged slab together with the stated cone fraction. Its exact statements are in the concordance.

\paragraph{The middle edge pairing.}\label{A:edgeaverage}
Work over the characteristic-zero valuation ring at the marked prime. Write $w^2=p$, $\nu=2h$, $P_0=p^o$, $c_0=k^2o\nu$ (that is, $k^2\log_p q$), and fix an integer precision $c$, eventually $5\nu c_0$. Along the cyclic letters the polynomial has degree $a$ in each group of $k$ variables. Let $t=3kc$. In the large-degree range the Frobenius substitution in one generic matrix has degree at most $D_*=\nu ak(h+2)P_0^h$, with $D_*+k<q$. At each letter choose an integral functional $I_j$ on the generic matrix polynomial ring with three properties. It is invariant under conjugation by the Levi group of the two blocks. It kills monomials whose row degree across the boundary (degree in the first row of the second block minus degree in the last row of the first) exceeds $2tP_0^h$. For a polynomial of degree at most $D_*$ whose Levi conjugates have all those coefficients in an ideal $J$, it equals $p^{c_0}$ times evaluation at the identity modulo $J$.

Here the scale $p^{c_0}$ comes from diagonal weights $p^{c_0}/\binom{d}{\alpha}$ in homogeneous degree $d\le D_*$. They are integral since there are $k^2$ entries and the base-$p$ carry formula bounds the multinomial valuation by $k^2\log_p q$ for $d<q$. On degree $d$ the functional is the trace on the tensor power $(O^k)^{\otimes d}$ of the weighted coefficient operator minus its trace on the submodule generated under the block Schur algebra by tensors of row degree above the threshold. The trace on the full tensor power is the scaled evaluation at the identity by multinomial counting. For a high-weight monomial the image lies in the high-weight submodule, so the two traces agree. Conjugation invariance follows by change of basis in both traces. The remaining assertion that the restricted trace lies in $J$ under the stated coefficient condition uses finite-field interpolation and spanning of the block Schur algebra under $D_*+k\le q$. We retain this integral trace lemma among the local inputs.

Contract the generic Frobenius substitution at letters in a subset $J_0$ using $I_j$, and at its remaining letters by $p^{c_0}$ times identity evaluation. Call the resulting operator $A_{J_0}$. At $J_0=\varnothing$ it equals $p^{\nu c_0}$ times the identity. When all letters except $j$ have been contracted, the weight estimate on the generic substitution says that for a coefficient surviving modulo $w^c$,
\[
 \left|\frac{\text{row weight at }j}{P_0^h}-
    \text{defect of the output monomial at }j+h\right|\le t.
\]
This estimate is the local Frobenius-factor input. Levi invariance at every other letter forces its two root weights to balance, hence only the free copy appears on the left. Output defect $>3t$ forces row weight $>2tP_0^h$. Thus after all letters are contracted, each such coefficient of every Levi translate vanishes modulo $w^c$.

To compare the pairings for $J_0$ and $J_0\cup\{j\}$, first pair the remaining one-matrix polynomial with an all-edges test vector $h'$. If a high-weight monomial survives in a Levi conjugate of the one-matrix coefficient, the same inequality forces defect $>t$. Its pairing with $h'$ is divisible by $w^c$: by transpose adjointness it equals the diagonal pairing weight times that coefficient in the transpose translate of $h'$, and the all-edges depth estimate gives the divisibility. Every high-weight coefficient of the paired polynomial therefore lies in $(w^c)$. The third defining property of $I_j$ shows that insertion of $j$ preserves this pairing modulo $w^c$. Telescoping gives
\[
 \langle A_{\rm all}h,h'\rangle\equiv p^{\nu c_0}\langle h,h'\rangle
 \pmod {w^c}.
\]
The two conclusions explain why the same average yields the residue defect bound and the congruence needed to clear the middle edge pairing.

% Includable section; requires amsmath, amssymb. Labels have prefix B:.
\section{Spreading the source forms and choosing the determinant family}
\label{B:section}

In the notation of \S\ref{sz:setup}, put \(Q=\rho(e_0)\)
for the good-type projector, so that \(Y_T=QY\).
On the full source the maps associated to the residue action are
\[
  G=\prod_{i=0}^1\prod_{w\mid p_i\text{ chosen}}\mathrm{GL}_r(\mathbb F_w),
  \qquad B_g=|G|Q\rho(g)\Pi:X\longrightarrow Y_T.
\]
The restriction of \(B_g\) to \(S\), multiplied by the nonzero
scalar \(c_B c_{\rm ss}^{-1}\) of Theorem~\ref{A:main}, is \(P_g\).

\paragraph{Theorem B (spreading).}
Under hypotheses F0--F7 of \S\ref{B:frontier},
\begin{equation}
\label{B:theorem}
 d_T>0,\qquad
 |t|d_S\le5d_T\log(1+2b^2),\qquad
 \dim\sum_{g\in t}P_g(S)\ge(1-4/b^2)d_T
\end{equation}
for one finite set \(t\subseteq G\).
The arithmetic hypotheses give the two estimates
\eqref{B:near} and \eqref{B:moment} on the same scaled map.

\subsection{From two Hilbert--Schmidt estimates to a new translate}
\label{B:one}
Choose the source coordinates indexed by the pairs consisting of a
source class and a Bessel coordinate (the coordinate isomorphism
includes the central translation by the lattice scale).
In these coordinates \(X\) is a nonzero Euclidean space. Define
\({\Pi_{\rm sc}}=(\operatorname{vol}(K_{\rm pr})\nu)^{-1/2}\Pi\) in these coordinates,
where \(K_{\rm pr}\) is the principal level and \(\nu>0\) is the unipotent
covolume. Both volume and covolume are positive. We use the Hilbert--Schmidt norm
\(\|C\|_{\rm HS}^2=\operatorname{tr}(C^*C)\).
The estimates furnished by the arithmetic argument are
\begin{align}
 \|Q{\Pi_{\rm sc}}\,\operatorname{pr}_Z\|_{\rm HS}^2
 &\ge \dim Z-5b^{-4}d_X\quad(Z\subset X), \label{B:near}\\
 d_T\sum_{g\in G}\|{\Pi_{\rm sc}}^*\rho(g){\Pi_{\rm sc}}\|_{\rm HS}^2
 &\le |G|(1+9b^{-4})d_X^2. \label{B:moment}
\end{align}
Their reduction to more primitive inputs is given below.
For \(Z=X\), \eqref{B:near} gives \(\|Q{\Pi_{\rm sc}}\|_{\rm HS}^2>0\),
since \(b\ge1000\) and \(d_X>0\). Hence \(d_T>0\).
Set \(\epsilon=(8b^2)^{-1}\). The admissible scales imply
\(5b^{-4},9b^{-4}\le\epsilon\) and \(0<\epsilon<1/16\).

Let \(F\subset Y_T\), \(\alpha=\dim F/d_T\), and
\(C_g=\rho(g){\Pi_{\rm sc}}{\Pi_{\rm sc}}^*\rho(g)^*\). By cyclicity of trace,
\[
 \langle C_g,C_h\rangle_{\rm HS}
 =\|{\Pi_{\rm sc}}^*\rho(g^{-1}h){\Pi_{\rm sc}}\|_{\rm HS}^2 .
\]
Summing over both indices and tracing against \(\operatorname{pr}_F\)
gives the identities
\[
 \left\|\sum_g C_g\right\|_{\rm HS}^2
 = |G|\sum_g\|\Pi_{\rm sc}^*\rho(g)\Pi_{\rm sc}\|_{\rm HS}^2,
 \qquad
 \operatorname{tr}(\operatorname{pr}_F C_g)
 = \|\operatorname{pr}_F\rho(g)\Pi_{\rm sc}\|_{\rm HS}^2.
\]
Cauchy--Schwarz in the space of operators and \eqref{B:moment} give
\[
 \left(\frac1{|G|}\sum_g
     \|\operatorname{pr}_F\rho(g){\Pi_{\rm sc}}\|_{\rm HS}^2\right)^2
 \le\alpha(1+\epsilon)d_X^2 .
\]
Since \(2\sqrt{\alpha(1+\epsilon)}\le\alpha+1+\epsilon\),
some \(g\) satisfies
\(\|\operatorname{pr}_F\rho(g){\Pi_{\rm sc}}\|_{\rm HS}^2
\le(1+\alpha+\epsilon)d_X/2\).
Put \(Z=(Q\rho(g){\Pi_{\rm sc}})^{-1}F\).
On \(Z\) the maps \(Q\rho(g){\Pi_{\rm sc}}\) and \(\operatorname{pr}_F\rho(g){\Pi_{\rm sc}}\)
agree: the component orthogonal to \(Y_T\) is orthogonal to \(F\).
Moreover \(Q\) commutes with the unitary \(\rho(g)\).
Thus
\[
 \dim Z-\epsilon d_X
 \le\|Q{\Pi_{\rm sc}}\operatorname{pr}_Z\|_{\rm HS}^2
 =\|\operatorname{pr}_F\rho(g){\Pi_{\rm sc}}\operatorname{pr}_Z\|_{\rm HS}^2
 \le(1+\alpha+\epsilon)d_X/2.
\]
Coordinate change and multiplication of a map by a nonzero scalar
preserve the dimension of the preimage of every subspace. In particular
\(Q\rho(g){\Pi_{\rm sc}}\) and \(B_g\) have preimages of equal dimension.
We have proved, for every \(F\subset Y_T\), the existence of \(g\) with
\begin{equation}
\label{B:spread}
 \dim B_g^{-1}F\le\left(\frac{1+\alpha}{2}
                  +\frac32\epsilon\right)d_X .
\end{equation}
The argument uses the full unprojected moment
\({\Pi_{\rm sc}}^*\rho(g){\Pi_{\rm sc}}\) and the projected inferior estimate for \(Q{\Pi_{\rm sc}}\).

\subsection{Greedy selection with a cardinality bound}
\label{B:greedy}
If \(c=d_T-\dim F>16\epsilon d_T\), apply \eqref{B:spread}
to \(F\). The map to \(Y_T/F\) induced by \(B_g\) has kernel
\(B_g^{-1}F\); therefore
\[
 \dim(F+B_g(X))-\dim F
 \ge (\tfrac12(1-\alpha)-\tfrac32\epsilon)d_X
 \ge \tfrac25(1-\alpha)d_X.
\]
Indeed \(\tfrac12-\tfrac3{32}=\tfrac{13}{32}>\tfrac25\).
Write \(\lambda=2d_X/(5d_T)>0\). Starting with \(F=0\), the first
step is possible, and its gain cannot exceed \(d_T\), so \(\lambda\le1\).
At each step where the codimension exceeds \(16\epsilon d_T\), a
translate reduces it to at most \((1-\lambda)c\le e^{-\lambda}c\).
Such a translate cannot have been selected previously (the gain is
positive). Stop if the threshold is reached, otherwise use at most
\[
 j=\left\lceil\lambda^{-1}\log(1/(16\epsilon))\right\rceil
\]
steps. After \(j\) steps the missing share is at most
\(e^{-j\lambda}\le16\epsilon\); stopping earlier also suffices.
Since \(j\lambda\le1+\log(1/(16\epsilon))\), the set \(t\)
so obtained satisfies simultaneously
\begin{align}
 |t|d_X&\le\tfrac52(1+\log(1/(16\epsilon)))d_T
 \le5d_T\log(1+2b^2),\label{B:count}\\
 \dim\sum_{g\in t} B_g(X)&\ge(1-16\epsilon)d_T
 =(1-2/b^2)d_T.\label{B:full}
\end{align}
For the second inequality of \eqref{B:count}, use
\(1/(16\epsilon)=b^2/2\le1+2b^2\) and
\(\log(1+2b^2)\ge1\). The argument works for the stated dimension
ratios; no implicit rounding of the number of translates is used.

\subsection{The loss from restricting to ramified forms}
\label{B:ramification}
The scalar translations at the \(f\) places dividing \(\pi\) have
order dividing \(N-1\). By the free-level dimension count
in \eqref{A:dimension}, the intersection defining \(S\) has
\begin{equation}
\label{B:codim}
 d_X-d_S\le \frac{f}{N-1}d_X\le b^{-8}d_X,
\end{equation}
using \(f\le b\) and \(N-1\ge b^9\).

Choose a vector-space complement \(X=S\oplus C\). For the
\emph{same} set \(t\), the dimension of
\(\sum_g B_g(X)\) is at most that of \(\sum_g B_g(S)\) plus
\(|t|\dim C\). From \eqref{B:count}, \eqref{B:codim},
\[
 |t|(d_X-d_S)\le 5 b^{-8}\log(1+2b^2)d_T
 \le (2/b^2)d_T .
\]
The last inequality holds, for example, by
\(\log(1+2b^2)\le 2b^2\) and \(b^4\ge5\).
Multiplying the restricted maps by \(c_Bc_{\rm ss}^{-1}\) does not
change their ranges. Therefore
\[
 \sum_{g\in t}P_g(S)=\sum_{g\in t}B_g(S),\qquad
 \dim\sum_{g\in t}P_g(S)\ge(1-4/b^2)d_T.
\]
As \(d_S\le d_X\), \eqref{B:count} also bounds \(|t|d_S\).
This proves \eqref{B:theorem} from \eqref{B:near}--\eqref{B:moment}.

\subsection{Arithmetic origin of the two moment estimates}
\label{B:arithmetic}
We now trace \eqref{B:near} and \eqref{B:moment} further.
Use the filling data of \(D\), write \(M=p_0p_1\), and let
\(R\subset G=G_0\times G_1\) be the reduction of the Bessel subgroup.
Let \(\delta\) be the product of the local Bessel formal degrees and
\(\mu\) the mass of the genus of the standard definite rank-\(r\)
lattice. Put
\[
 I_0=|R|d_X,\qquad I_*=d_X^2/(\delta\mu),\qquad q_g={\Pi_{\rm sc}}^*\rho(g){\Pi_{\rm sc}}.
\]
The class representatives \(h_c\), cutoffs \(\beta_c\) and vector
Bessel function \(W\) are those in the filling data. Write
\(\mathcal A(g;c,c')\) for the matrix
\[
 \sum_{\gamma:\ h_c^{-1}\gamma h_{c'}\in K,\,
             \operatorname{red}(h_c^{-1}\gamma h_{c'})=g}
 \int \beta_c(x)\beta_{c'}(\gamma^{-1}x)
             W(x)W(\gamma^{-1}x)^*\,dx .
\]
Here \(K\) is the lattice stabilizer; the integral is on the
product of real unitary groups in the filling model. The sum is
over rational points. In the source orthonormal basis indexed
by \((c,i)\), the blocks of \(q_g\) are these blocks divided by
\(\nu\), with entrywise conjugation if the inner product is taken
linear in its second variable. In particular
\(\|q_g\|_{\rm HS}^2=\nu^{-2}\sum_{c,c'}\|\mathcal A(g;c,c')\|^2\).

The unfolding giving this identity begins with the Petersson
product of two Poincar\'e series, unfolds the series belonging to
\((c,i)\) to its seed, and evaluates the other series at the
class point \(h_c\) translated by a representative of \(g\).
The finite support of the seed selects exactly the displayed
congruence condition, with factor \(\operatorname{vol}(K_{\rm pr})\).
The real-homeomorphism change of variables produces the product
of the two cutoffs and Bessel functions. Scaling both series
by \((\operatorname{vol}(K_{\rm pr})\nu)^{-1/2}\) gives the factor
\(\nu^{-1}\). F1 is the fitted identification of the period-defined \(\Pi\)
with this series and the normalized seed and block identities.
The absolute summability of the amplitude integrals used here
follows from the majorization in \S\ref{B:cutoff}.

\paragraph{Near the Bessel subgroup.}
Denote by \(r_\psi(g)\) the operator whose blocks come from the
same amplitude sum restricted to the Bessel subgroup; it is a
unitary representation on \(R\), extended by zero off \(R\).
Here is how the action and the trace enter the proof.
For \(t\in R\), a rational Bessel term over the column class and
\(t^{-1}\) specifies a new class and a unitary change of Bessel
coordinates. In the Bessel-only sum on that column, the block
at the new class is this unitary matrix (with the conjugation
fixed by the inner product), and every other block vanishes.
Indeed one moves the term to the identity block, where integrating
the two cutoffs and the Bessel function gives \(\nu\) times the
identity. This identity-block normalization and the change of class
coordinates are among the fitted identities in F1. Composing two
term changes multiplies their matrices. The blocks at inverse group
elements are adjoints by replacing the rational term by its inverse.
Thus \(r_\psi(1)=1\) and \(r_\psi\) is unitary on \(R\).
Applying the same change of variables on the right to the
full sum gives \(q_{gt}=q_g r_\psi(t)\). For \(u\in X\),
\(d=\rho(t){\Pi_{\rm sc}}u-{\Pi_{\rm sc}}r_\psi(t)u\) now satisfies
\({\Pi_{\rm sc}}^*\rho(g)^{-1}d=0\) for every \(g\); hence it is orthogonal to
the span of every \(\rho(g){\Pi_{\rm sc}}X\). It belongs to that very span, so
\(d=0\). This proves \(\rho(t){\Pi_{\rm sc}}={\Pi_{\rm sc}}r_\psi(t)\).
Let \(E_g=q_g-r_\psi(g)\); thus \(E_{sg}=r_\psi(s)E_g\) and
\(E_{gs}=E_gr_\psi(s)\) for \(s\in R\). Write
\(\Omega_i=R G_i\), where \(G_i\) is embedded as one factor.
The quantitative estimate needed is
\begin{equation}
\label{B:nearBessel}
 \sum_{g\in\Omega_i}\|E_g\|_{\rm HS}^2\le b^{-12}I_0
 \qquad(i=0,1).
\end{equation}
For \(g\in R\), \(E_g=({\Pi_{\rm sc}}^*{\Pi_{\rm sc}}-1)r_\psi(g)\), so its squared
Hilbert--Schmidt norm is constant. Restricting \eqref{B:nearBessel}
to \(R\) and dividing by \(|R|>0\) yields
\begin{equation}
\label{B:defect}
 \|{\Pi_{\rm sc}}^*{\Pi_{\rm sc}}-1\|_{\rm HS}^2\le b^{-12}d_X .
\end{equation}
For every subspace \(Z\subset X\),
\[
 \|{\Pi_{\rm sc}}\operatorname{pr}_Z\|_{\rm HS}^2
 =\dim Z+\operatorname{tr}(\operatorname{pr}_Z({\Pi_{\rm sc}}^*{\Pi_{\rm sc}}-1))
 \ge\dim Z-b^{-6}d_X\ge\dim Z-b^{-4}d_X.
\]
Also \(\|q_1\|_{\rm HS}^2\le(1+b^{-6})^2 d_X\le 2d_X\).

To prove \eqref{B:nearBessel}, group the terms by links.
For a subset of reductions \(\Omega\),
let \(\mathcal E(\Omega)\) be the sum of squared cut-amplitude norms, taken
link by link for terms of nonzero frame discriminant. A link
is the submodule of pairs of classes of \(v,\gamma^{-1}v\) in
the quotient of the orthogonal complement of the isotropic frame
span by that span, with \(v\) orthogonal to the frame span and
its \(\gamma\)-translate. Let \(\operatorname{meetNorm}(g)\)
be the product of the norms of places where the frame span
meets its translate nontrivially. The near-pair count gives the following bound. If \(0<z\le\operatorname{meetNorm}(g)\) throughout \(\Omega\),
\[
 \mathcal E(\Omega)/\nu^2\le (b^{2f}/z)^k I_*;
\]
the error from different link classes/zero discriminant is
\(b^{-kf}I_0\) in the normalized cut second moment and the
error of removing the cut is at most \(b^{-40}I_0\) in squared
norm (F2--F4). On a region avoiding nontrivial root-of-unity
scalar cosets the off-Bessel terms equal the remaining terms.
Expanding squares by links, then taking the triangle inequality
in the space of all blocks, gives
\begin{equation}
\label{B:rootbound}
 \sqrt{\sum_{g\in\Omega}\|E_g\|_{\rm HS}^2}
 \le\sqrt{\mathcal E(\Omega)/\nu^2}+ b^{-15}\sqrt{I_0},
\end{equation}
since \(kf\ge40\), and \(2b^{-20}\le b^{-15}\).
For \(\Omega_i\) the reductions match a Bessel element at the
other prime and its conjugate; the scalar-coset avoidance
follows from the first frame vector, and
\(z\ge(M/b)^f\) by the comparability of the level primes.
The filling bounds \(I_*\le(M/b^6)^{kf}I_0\) and \(kf\ge40\)
give
\[
 (b^{2f}/z)^k I_*\le b^{-3kf}I_0\le b^{-30}I_0.
\]
Thus the right side of \eqref{B:rootbound} is at most
\(2b^{-15}\sqrt{I_0}\); squaring and using \(4b^{-30}\le b^{-12}\)
proves \eqref{B:nearBessel}.

\paragraph{The cost of projecting to good types.}
At prime \(p_i=q\), each of the \(f\) fields has \(q\) elements.
The prime-wise good idempotent \(e_i^{\rm pr}\) is the sum over tuples of labels
of degree \(r\) that are allowed at every place and primitive at
at least one place. Allowed means monic squarefree with no root
at zero or at an \(M_0\)-th root of unity in the field, where
\(M_0=q_2\) is the exclusion modulus of \S\ref{sz:setup}.
Primitive means irreducible with a root of order \(q^r-1\).
For each irreducible factor \(p\) of a squarefree label, write
\(d=\deg p\) and choose a root in \(\mathbb F_{q^d}^\times\).
The faithful identification of this cyclic group with its dual
commutes with raising to \(q\), so the degree-\(d\) Frobenius
orbit of roots gives a degree-\(d\) orbit of characters. The
Hall function \(\tau_p\) used in the level projector equals the one-box Green
character of this orbit, in the positive-degree normalization, by
Lemma~\ref{lem:fitted-primary-character}.
Green's theorem~\cite[\S8, Theorem~14 and p.~444]{green}
gives its degree \(\prod_{h=1}^{d-1}(q^h-1)\).
Distinct irreducible factors give distinct orbits. Since the
label is squarefree, parabolic induction of these characters is
the irreducible character indexed by the set of those orbits,
each with partition \((1)\); distinct labels give distinct characters.
The label function is precisely this induction product by construction.
For factor degrees \(d_j\), \(\sum_j d_j=r\), the inducing parabolic
has index \(\prod_{h=1}^r(q^h-1)/\prod_j\prod_{h=1}^{d_j}(q^h-1)\)
by the flag count. Multiplying by the inducing degree yields
\[
 \chi(1)=\frac{\prod_{h=1}^{r}(q^h-1)}{\prod_j(q^{d_j}-1)},
 \qquad \frac{\chi(1)^2}{|\mathrm{GL}_r(\mathbb F_q)|}
 \ge q^{-r}\prod_{h=1}^r(1-q^{-h})\ge w_q,
 \qquad w_q=(1-r/q)q^{-r}.
\]
Thus the positive coefficient at the identity of each character
idempotent is at least \(w_q\).

There are \(q^r\) monic polynomials of degree \(r\); at most
\((1+M_0)q^{r-1}\) have forbidden roots. A nonsquarefree
polynomial is \(p^2h\) with monic \(p,h\) and \(\deg h<r\);
for each possible degree of \(h\) at most \(q^{r-1}\) such
products occur. Hence the weighted count of allowed ones is
at least \(1-\eta\), where \(\eta=(1+M_0+2r)/q\).
There are \(\varphi(q^r-1)/r\) primitive polynomials: a primitive
element cannot lie in a proper subfield and has an orbit of
size \(r\) under Frobenius.
For an integer \(u\) with \(j\) distinct prime divisors, the
\(i\)-th smallest prime divisor is at least \(i+1\), giving
\(\varphi(u)/u\ge\prod_{i=1}^j i/(i+1)=1/(j+1)\);
also \(2^j\le u\). Consequently
\(u\le\varphi(u)(1+2\log u)\). With weight \(w_q\) for each monic polynomial, the count of nonprimitive polynomials is at most
\(\exp[-1/(2r(2r\log q+1))]\).
Taking the product of weighted counts over all \(f\) fields
and subtracting the tuples with no primitive polynomial gives
\[
 (1-e_i^{\rm pr})(1)\le1-(1-\eta)^f+
   \exp[-f/(2r(2r\log q+1))]\le f\eta+b^{-9}\le b^{-8}/2.
\]
The last two inequalities are the scale checks
\(0\le\eta\le1\), \(f\ge 2r(2r\log q+1)\,9\log b\) and
\(f\eta+b^{-9}\le b^{-8}/2\). The tensor product
\(e_0=e_0^{\rm pr}\otimes e_1^{\rm pr}\) acts as \(Q\).
This explains why both good-prime conditions are imposed.

For a complementary idempotent \(x=1-e_i^{\rm pr}\), the convolution
\(f(g)\mapsto\sum_{t\in G_i}x(t)f(gt)\) is an orthogonal projection
on the space of functions on \(G\) with operator-valued entries:
it is the group algebra action of a self-adjoint idempotent.
Apply this to \(E_g\) cut down to \(\Omega_i\). On \(R\)
its value is \(r_\psi(s)\sum_{t\in G_i}x(t)E_t\).
Using \eqref{B:nearBessel} and dividing by \(|R|\) gives
\[
 \left\|\sum_{t\in G_i}x(t)E_t\right\|_{\rm HS}^2\le b^{-12}d_X .
\]
For \(t\in G_i\) the other prime component is the identity.
A diagonal Bessel term is a rational element of \(U_{\rm Bes}\rtimes U(X_0)\)
whose conjugate by its class representative is integral and reduces
to the identity at that other prime. Here \(U_{\rm Bes}\) is the
unipotent radical of the Bessel subgroup. Its middle \(U(X_0)\)-block
inherits integrality and the identity reduction. By freeness at that
prime its rational part is the identity. Thus a diagonal term can
occur only if \(t\) is in the image of the integral unipotents. In
that case translating the identity block by a unipotent representative
multiplies each diagonal entry by its reduced character \(\psi_U(t)\).
Summing the \(d_X\) entries gives \(\operatorname{tr}r_\psi(t)=d_X\psi_U(t)\),
with \(\psi_U\) extended by zero elsewhere.
Centrality replaces \(\psi_U\) by its conjugacy average
\(\psi_U^\natural\) in the trace sum. The local two-frame mass
estimate implies
\begin{equation}\label{B:charmean}
 \sum_{g\in G}|\psi_U^\natural(g)|^2\le2 .
\end{equation}
This is input F5 below. The idempotent's self-adjointness and
idempotence give \(\sum_t|x(t)|^2=x(1)\). By Cauchy--Schwarz the
Bessel trace contribution has absolute value at most
\(d_X\sqrt{(b^{-8}/2)2}=b^{-4}d_X\).
The trace of the remainder has absolute value at most
\(\sqrt{d_X}\,b^{-6}\sqrt{d_X}\le b^{-4}d_X\).
Since \(\rho(x)\) is an orthogonal projection,
\[
 \|\rho(x){\Pi_{\rm sc}}\|_{\rm HS}^2
 =\operatorname{tr}({\Pi_{\rm sc}}^*\rho(x){\Pi_{\rm sc}})\le2b^{-4}d_X.
\]
The two good projectors commute. The squared norm missed by
their product is at most the sum missed by the factors
(orthogonal decomposition with respect to one factor).
Hence \(\|(1-Q){\Pi_{\rm sc}}\|_{\rm HS}^2\le4b^{-4}d_X\).
Subtract this from the inferior estimate following
\eqref{B:defect}, on any subspace \(Z\), to obtain \eqref{B:near}.

\paragraph{The global second moment.}
Set \(\beta(g)=\|q_g\|_{\rm HS}^2\). This nonnegative
function is of positive type: it is the Gram correlation of the
operators \(C_g\) from \S\ref{B:one}. In particular
\(\beta(g)\le\beta(1)\le2d_X\).

We give the finite group argument that converts the local
main term into the global bound. At a chosen place above \(p_0\),
a frame vector is a nonzero vector in the frame span, or in its
dual at the conjugate place. For a set \(T\) of such places,
let \(K_T\subset G\) fix one specified frame vector at each,
and let \(s_T=\sum_{g\in K_T}\beta(g)\). Define
\(D_{\rm fix}=b^{6k}\), \(C_T=D_{\rm fix}^{-|T|}I_*\), and \(x=b^{-4}\).
Call a reduction open if at some chosen place above \(p_0\) all
the relevant leading minors of the frame pairing matrix are
nonzero. The open-link estimate (F6) bounds the normalized
link moment on \(K_T\) intersected with the open reductions by
\((1+x)C_T\).
Since \(I_0\le C_T\) by \(b^{11kf}I_0\le I_*\), and
open reductions avoid scalar cosets, \eqref{B:rootbound} gives
\[
 \sum_{g\in K_T,\,g\text{ open}}\beta(g)
 \le (\sqrt{1+x}+b^{-15})^2 C_T\le(1+4x)C_T.
\]
For the fixer at all \(f\) places, the scalar-coset count is
at most \(M^f|R|\), each term at most \(2d_X\).
The scale bounds imply \(8M^f I_0\le C_T\):
use \(M\le b^{23}\), \(8\le b\), \(23f+1\le5kf\).
Off the scalar cosets, \(z\ge p_0^f\ge b^{10f}\);
the link bound is at most \(b^{-2kf}C_T\le C_T/8\),
and \(b^{-30}I_0\le C_T/8\).
Using \((u+v)^2\le2(u^2+v^2)\) in \eqref{B:rootbound},
the non-scalar contribution is at most \(C_T/2\), while the
scalar contribution is at most \(C_T/2\). Thus \(s_T\le C_T\)
when all places are fixed.

For other fixers a finite-field estimate is needed.
Let \(\Gamma\) be a finite group, let \(\beta:\Gamma\to\mathbb R_{\ge0}\)
be of positive type, let \(q\) be the residue field size and let \(\pi\)
be a representation of \(\Gamma\) in \(\mathrm{GL}_r(\mathbb F_q)\).
Suppose all stabilizers of nonzero columns of either of the
two injective frame maps \(A_0,A_1:\mathbb F_q^j\to\mathbb F_q^r\)
(with the action on the second frame contragredient) have
\(\beta\)-mass at most \(Y_0\).
Then
\begin{equation}
\label{B:minor}
 q\sum_{\det(A_1^t\pi(g)A_0)=0}\beta(g)
 \le 2\sum_{\Gamma}\beta(g)+q^{j+r/2}Y_0 .
\end{equation}
To verify this, take vectors \(u_g\) realizing positive type,
\(\langle u_a,u_b\rangle=\beta(a^{-1}b)\), and the delta and
Fourier orthonormal bases of \(\ell^2(\mathbb F_q^r)\).
The squared norm of
\(\sum_a\delta_{\pi(a)z}\otimes u_a\) is \(|\Gamma|\) times
the stabilizer mass of \(z\). The same identity in the Fourier
basis at \(\pi(a)^{-t}v\) holds for the dual stabilizer.
The cross inner product is
\(|\Gamma|q^{-r/2}\sum_g\psi(v^t\pi(g)z)\beta(g)\)
for a nontrivial additive character \(\psi\), with the Fourier
sign chosen accordingly. Cauchy--Schwarz bounds its real part,
so the weighted character sum has real part at most
\(q^{r/2}Y_0\) for \(v=A_1 y,z=A_0 x'\), both nonzero
parameter vectors. For \(Z_g=A_1^t\pi(g)A_0\), additive
orthogonality gives
\(\sum_{y\in\mathbb F_q^j}\mathbb E_{x'\in\mathbb F_q^j}
  \psi(y^t Z_g x')=|\ker Z_g^t|\), at least \(q\) if
\(Z_g\) is singular. Since \(\beta(g)\ge0\), multiply this pointwise lower bound by
\(\beta(g)\) and sum. The
axes \(y=0\) and \(x'=0\) cost at most twice the total mass;
the remaining pairs cost at most \(q^{j+r/2}Y_0\).
This proves \eqref{B:minor}. The weight in the application is a squared
Hilbert--Schmidt norm, so it satisfies the pointwise sign condition.

The sign condition is necessary in the general statement. Let \(S_3\) act on
\(\mathbb F_2^4\) by permuting the first three coordinates, and take its sign
character as \(\beta\). This one-dimensional character is of positive type.
Every vector and covector stabilizer contains a transposition, because two of
its first three coordinates agree. All their signed masses vanish. Set
\(j=2\), choose \(\operatorname{im} A_0=V=\langle e_1,e_2+e_4\rangle\),
and choose \(A_1\) so that \(\ker A_1^t=(12)V\).
The minor is singular exactly when \(gV\cap(12)V\ne0\).
Comparing the three nonzero vectors of each plane gives
\(g=1,(12),(123),(132)\). Their signs sum to two, whereas the total signed
mass is zero. With \(Y_0=0\), the claimed inequality without the sign
condition would read \(4\le0\).

Choose the set of places and their frame vectors maximizing
\(D_{\rm fix}^{|T|}s_T\); all choices are finite. At a place outside
\(T\), the mass in the current fixer of every additional
vector or covector stabilizer is at most \(s_T/D_{\rm fix}\), by
maximality. If \(T\) is full the preceding all-place bound
already suffices. Otherwise apply \eqref{B:minor} to each
leading minor at such a place; the relevant frame maps are
injective. The scale inequality \(p_0^{j+r/2}\le D_{\rm fix}\) holds for
the minor sizes \(1\le j\le\varrho+1\), and \(p_0\ge b^{10}\).
Each singular-minor mass is therefore at most \(3s_T/b^{10}\).
A non-open reduction has a vanishing relevant minor at this
place; \(\varrho+1\le b\) implies the mass of all non-open
reductions is at most \(xs_T\). Combining with the open estimate,
\[
 s_T\le(1+4x)C_T+xs_T\quad\Longrightarrow\quad
 s_T\le(1+9x)C_T
\]
because \(0\le x\le1/10\) and
\((1-x)(1+9x)-(1+4x)=x(4-9x)\ge0\).
The empty set of places is an admissible competitor, hence
\begin{equation}
\label{B:global}
 \sum_{g\in G}\|q_g\|_{\rm HS}^2
 \le D_{\rm fix}^{|T|}s_T\le(1+9b^{-4})I_* .
\end{equation}

\paragraph{The target dimension.}
By F7 the good forms on \(Y_0\) have the same dimension as those on
the terminal definite comparison space \(W_m\). The precise input is the
cuspidal trace comparison for the two inner forms with matching
Harish--Chandra parameters (F7 below): the rank is at least two,
the number of infinite places is even, and the required parameter
gaps hold. Appendix~\ref{sec:trace} proves the reduction to the trace comparison
and its local inputs for all permitted Hecke operators. Taking the identity
operator there and summing over the good types gives this dimension
equality. Thus \(d_T\) is at most the dimension of all
principal-level forms on the definite \(W_m\).
The level is free: a rational stabilizer element is integral at every place up to conjugacy and congruent to the identity at the chosen places above \(p_0\). The norm of its trace minus \(r\) is at most \((2r)^{2f}<p_0^f\). Divisibility by all the chosen primes forces trace \(r\), hence the identity by archimedean unitarity. Its index
in the maximal level is \(|G|\), so the class-mass formula (or
counting the free fibers of the double coset map with stabilizer
weights) gives \(|G|\mu\) classes.
Each contributes the dimension of the weight representation,
which is \(\delta\) by the Weyl degree formula~\cite[Lecture~24]{fultonharris} and the
parameter match in the admissible Bessel system:
for a definite rank-\(r\) representation of highest weight
\(\lambda\), the local degree is
\(\prod_{i<j}(\lambda_i-\lambda_j+j-i)/(j-i)\);
the discrete-series formal degree here has the same parameter
product at each place. Consequently
\(d_T\le |G|\delta\mu\).
Multiplying \eqref{B:global} by this upper bound cancels the
denominator \(\delta\mu\) in \(I_*\), proving \eqref{B:moment}.


\subsection{The cutoff error and the near-pair sums}
\label{B:cutoff}
The following calculation reduces the analytic cutoff error and
absolute convergence to a single geometric separation estimate, with
the fitted Bessel majorants of F0. Write \(h=\varrho+1\), \(i=r^3f\)
and \(d_0=2hrf\). For an archimedean tuple \(y\) put
\(u(y)=\prod_v\max(1,\ell(y_v)/b^{6h})\), with \(\ell(y_v)\) the
infimum of \(1+\|s y_v\|\) over the Bessel subgroup. Fix the class
\(c\), reduction \(g\), and point \(y\). Give a term leading to
class \(c'\) the weight \(\beta_{c'}(z)\), \(z=\gamma^{-1}y\).
If \(u(z)\le R\), \(R\ge1\), the frames in the support of this
cutoff lie in a ball of radius \(\rho=b^{37h}R\) in the
\(d_0\)-dimensional real space of their coordinates. The box property
and rigidity input F4 state that unequal rational frames have
coordinate distance at least \((N\rho)^{-1}\) in this ball.
A fixed frame has total weight at most one: two terms with the same
frame and reduction differ within a single class by a rational
unipotent stabilizer element, and its translates of the cutoff sum
to one. Translating one frame to the origin, and comparing disjoint
balls of radius \((2N\rho)^{-1}\) with the ball of radius
\(2\rho+(2N\rho)^{-1}\), gives weighted count
\[
 C(R)\le(1+4N\rho^2)^{d_0}
     \le (b^{6h^2}R^2)^{d_0}.
\]
Indeed \(N\le b^{60h}\), \(h\ge23\), and \(134h+1\le6h^2\).
Since \(4d_0<i\), layer integration gives
\[
 \sum_{c',\gamma}\beta_{c'}(z)u(z)^{-i}
   =i\int_1^\infty C(R)R^{-i-1}\,dR
   \le 2b^{12h^3rf}.
\]
For the local full and cut functions, the fitted majorants give
\(\|W_v^{\rm side}(x)\|u_v(x)^{3i}\le w_v(x)\le a_v\) and
\(a_v\int_{U_{\rm Bes}\backslash G_\infty(v)}w_v\le
 b^{400r^2}\dim V_v\). Each \(w_v\) is unipotent invariant.
On the locus where the two kernels differ,
\(u(y)u(z)\ge b\): some local cutoff acts at length
\(>b^{6h+2}/2\). For either kernel use \(2i\) of its spare powers at
each argument and \(i\) more at \(z\). The triangle inequality bounds
the difference by
\[
 2 b^{-2i}\prod_v w_v(y_v)a_v\ u(z)^{-i}.
\]
Multiply by \(\beta_c(y)\beta_{c'}(z)\), sum, and integrate. The
sum over \(c',\gamma\) costs at most \(2b^{12h^3rf}\), and integrating
\(\beta_c\prod_v w_v\) unfolds the unipotent variable with factor
\(\nu\). The row sum of norms of the differences is at most
\[
 4\nu b^{-2r^3 f+12h^3rf+400r^2 f}\prod_v\dim V_v
   \le \tfrac12\nu b^{-k^3f}.
\]
Indeed \(\prod_v\dim V_v\le b^{2kf}\), \(8\le b\), and
\[
 12h^3rf+1+400r^2f+2kf+k^3f\le 2r^3f.
\]
For the latter use \(12h^3+500\le k<r\): the terms other than
\(k^3f\) are at most \(402r^2f\le r^3f\). Without restricting to
the changed locus, the same majorant without the factor
\(2b^{-2i}\) bounds the full kernel and the cut kernel, proving
absolute summability of the integrals. Finally, row squared norms
are bounded by squared row sums, the number of classes is at most
\(I_0\), and \(|G|b^{40}\le4b^{2k^3f}\). The last estimate follows
from \(|G|\le M^{r^2 f}\), \(M\le b^{23}\) and
\(23r^2f+40\le2k^3f\). Summing rows and dividing by \(\nu^2\) gives
the error \(b^{-40}I_0\) used in \eqref{B:rootbound}.

\paragraph{Pairs with equal and unequal links.}\label{B:pairs}
A block squared is bounded by the corresponding sum over ordered
pairs of terms. For a pair in the same nonzero-discriminant link,
its first representative determines both the reduction and the
link, so summing all link squares incurs no extra multiplicity.
For the remaining terms expand the square after partition into the
nonzero-discriminant link blocks and one zero-discriminant block.
Keep the squares of the former and bound the remaining products by
the products of the term sizes, including the entire square of the
zero-discriminant block. The factor from two amplitude majorants
is \(\nu^2\) in either case.

For the arithmetic count let \(d\) be the pair discriminant,
\(j=k-\dim_E(\mathrm{common\ link})\), \(\tau=b^{kf/4}\), and
\(C_p>0\) the pair-count scale of the filling. The geometric
inputs bound the product of term sizes by
\[
 \frac{B}{\tau^{j+2}C_p d^k}
 \quad\text{if}\quad
 (d\eta)^2\le G_\gamma G_{\gamma'},\quad
 D'\ge\max(H_\gamma,H_{\gamma'}),\quad
 \tau^{j+6}D'/\eta^k\le B .
\]
Here \(G_\gamma\) is the level gauge and \(H_\gamma\) the cell-rank
moment, both for the indicated near representative. The
reciprocal-discriminant count is
\(\sum_{\text{pairs of corank }j}d^{-k}\le\tau^{j+2}C_p\).
There are at most \(h^3+1\le\tau\) coranks, so summation gives
\(\tau B\). For equal links one has \(j\le1,\ \eta=z,\ D'=I_*\):
the result is \(\tau^8 I_*/z^k=(b^{2f}/z)^k I_*\).
For the other pairs the required bound on the budget is
\(B=I_0/\tau^{10}\), giving \(\tau B\le b^{-kf}I_0\).
Section~\ref{app:pair-count} proves the weighted lattice count and reduces the
pair count to the one-term multiplicity over a prescribed lattice.
The gauge and cell-rank bounds on individual term sizes are also required.

\paragraph{The coefficient in the open-cell estimate.}
\label{B:openenergy}
For a fixed vector choice at the places in \(T\), take a finite
collection of nonzero open link blocks. The source's open-cell
parametrization assigns to each block a last entry \(t'\), a
representative \(\sigma\) in its floor group (the subgroup used in
the unipotent double-coset parametrization), and integral positions
with the prescribed fixer. It requires
\(|t'|_v^2>N^2\) at every infinite place. With the first
\(\varrho\) entries ranging over the finite entry set, the block
equals the sum of cut amplitudes of the double cosets represented
by \(w(p,t')\sigma\), with \(w(p,t')\) the Weyl representative
for this entry list. The class pair, reduction, and \(\sigma\)
determine the block, since \(\sigma\) determines the link. Thus one
may sum squares without a labeling multiplicity.

Write \(S_{\rm bl}\) for the sum of the block norm squares and
\(R_{k-1}>0\) for the filling's normalizing constant at predecessor
rank. The Siegel estimate after grouping the blocks by \(t'\) is
\[
 R_{k-1}^2 S_{\rm bl}\le(1+b^{-6})C_{\rm S} Z.
\]
Here \(C_{\rm S}\) is the scaled Siegel constant and \(Z\) sums,
over distinct last entries and pairs of heads, the real part of the
product of archimedean band-limit kernels times the local
last-entry density divided by the intermediate group's mass ratio.
The Gram identity at fixed \(t'\) supplies the kernel in this estimate.
For the intermediate group implementing the vector fixers, the
required sieve bound is
\[
 Z\le\left(\prod_{j=0}^{\varrho}C_{{\rm Pois},j}\right)
 \left[
  \left(\prod_{v}\frac{1+b^{-10}}{\deg_v}\right)
  \frac{1+b^{-6}}{b^{6k|T|}}+
  b^{2403(\varrho+1)^9 f}C_{\rm dec}
 \right],
\]
where \(C_{{\rm Pois},j}\) are the stage Poisson constants,
\(\deg_v\) are the local Bessel degrees, and \(C_{\rm dec}\) is
the filling decay constant. In particular the sieve contributes one
factor \(1+b^{-6}\) for the entire fixer set \(T\), including \(T=\varnothing\).
Under the measure and level identifications of \S\ref{spread:normalization}, the scalar calculation there bounds the displayed right side times
\((1+b^{-6})C_{\rm S}\) by
\(R_{k-1}^2(1+b^{-4})b^{-6k|T|}I_*\nu^2\).
Divide by \(R_{k-1}^2\) and take the supremum over finite collections
of nonnegative summands to get F6. Section~\ref{moment:open} gives the rational parametrization.
The fitted Gram/Siegel estimate, remaining local transform estimates and
identification of the measures are needed in \S\ref{moment:full}--\ref{moment:sieve}. The
\srcline{New/Spanning/MainTerm/EnergyEstimate}{631}{formal open-cell bound}
has leading coefficient \(1+b^{-4}\) when \(T=\varnothing\).
For \(T\ne\varnothing\) it has the additional factor \(b^{-6k|T|}\).
The leading coefficient matters for the dimension conclusion. For example,
let the trivial group act on an isometric inclusion of a \(d\)-dimensional
space into one of dimension \(2d\). Every translate has the same image.
The ratio of the second moment of the image projection to its
squared first moment, with trace normalized by \(2d\), equals
\((1/2)/(1/2)^2=2\). A moment estimate with this constant is compatible
with the translates filling only half the target.

\subsection{Arithmetic hypotheses for spreading}
\label{B:frontier}
For clarity, the following are hypotheses of the present
reconstruction until their own proofs are supplied. All filling
statements are for the particular admissible filling associated
to \(D\), with its prescribed local measures, ranks, levels and
Bessel data; they are not claims for arbitrary choices of cutoffs.
Appendix~\ref{sz:register} records which clauses are derived on the constructed scales.
\begin{enumerate}
\item[F0.] The construction of the admissible filling, including
the fitted Bessel systems, their finite-dimensional source
coordinates (nonempty index sets), formal degrees and scale
inequalities used above. In particular the unitary source rank is \(k\),
and the ratios involving \(I_0,I_*\) are part of the filling
admissibility, not consequences of linear algebra.
\item[F1.] The period-to-series identification at the fitted levels,
the normalized identity Bessel block, and the block change of
variables for a rational Bessel term with its unitary coordinate
matrix. These yield the equivariance and the trace calculation in
\S\ref{B:arithmetic}. Absolute convergence of the block integrals
follows from F0 and F4 by \S\ref{B:cutoff}.
\item[F2.] For any \(\Omega\) and positive integer \(z\) with
\(z\le\operatorname{meetNorm}(g)\) throughout \(\Omega\),
\(\mathcal E(\Omega)\le \nu^2(b^{2f}/z)^k I_*\).
This is a source-specific near-pair lattice estimate, not a
general second-moment bound.
\item[F3.] Expansion by links: if \(\mathcal E(\Omega)\le C\nu^2\),
the normalized cut moment of the remaining terms is
at most \(C+b^{-kf}I_0\). The missing proof is the cross-link
near-pair size sum, including zero-discriminant terms, and the
term classification on sets avoiding scalar cosets.
\item[F4.] For \(u(z)\le R\) the cutoff support lies in the
frame box of radius \(\rho=b^{37h}R\); two distinct rational frames
on that support are separated by \((N\rho)^{-1}\), and terms of a
fixed frame and reduction occur in a single unipotent orbit with
cutoff weights summing to at most one, as specified in
\S\ref{B:cutoff}. Section~\ref{moment:rigidity} gives the
integral coordinates, separation and orbit argument in the fitted model.
On the constructed scales that section also constructs the reduced cutoff and proves its partition identity for the levels and scales specified there. Thus F4 follows in the fitted model.
\item[F5.] The conjugacy-average bound \eqref{B:charmean} for the
specific reduced unipotent character. The local mass proposition
with congruence subgroup \(K(M)\) is
\[\int_K\left|\int_K\chi_M(yzy^{-1})\,dy\right|^2dz
\le2\operatorname{vol}(K)^2\operatorname{vol}(K(M)).\]
The function is constant on reduction cosets; integration
over these disjoint cosets followed by cancellation of their
positive volume proves \eqref{B:charmean}. Appendix~\ref{sz:shells} derives this integral bound from explicit
deep-vector frame and shell bounds. The primitive local frame
estimates stated there remain inputs.
\item[F6.] For open reductions in any fixer \(K_T\),
\(\mathcal E(K_T\cap\mathrm{open})\le
 (1+b^{-4})b^{-6k|T|}I_*\nu^2\).
Section~\ref{B:openenergy} reduces this sharp main-term estimate to the displayed Siegel and sieve inputs. The factor
\(1+b^{-4}\) cannot be replaced by an unspecified constant,
since the subsequent dimension argument depends on near-optimal
second moment.
\item[F7.] Use the two-inner-form trace input in \S\ref{sec:trace}
at the identity operator, the equality of the fitted Bessel
formal degree and the compact Weyl degree. These enter the bound on \(d_T\).
Lemma~\ref{lem:fitted-primary-character} supplies the finite character
normalization in the bad-type weight of \S\ref{B:arithmetic}.
\end{enumerate}


\subsection{Selection of one indexed family}
\label{B:selection}
By Theorem~\ref{A:main}, \(\mathcal S\subset\mathcal I\) spans \(S\).
Fix the translate set \(t\) obtained above. Linearity gives
\[
 F_t=\sum_{g\in t}P_g(S)
 =\operatorname{span}_{\mathbb C}\bigcup_{g\in t}P_g(\mathcal S).
\]
Successively choose an element of this union outside the span of the
elements already chosen until their span equals \(F_t\). Its finite
dimension ensures that the procedure ends. Retain a preimage and a
translate for each of the resulting basis vectors:
\begin{equation}
\label{B:witnesses}
 y_i=P_{g_i}\phi_i,\quad g_i\in t,\quad \phi_i\in\mathcal S,\quad
 1\le i\le n,\qquad (1-4/b^2)d_T\le n=\dim F_t\le d_T .
\end{equation}
Since \(1-4/b^2\ge 1/2\) and \(d_T>0\), this gives \(n>0\) and
\(d_T\le2n\). Each chosen preimage belongs to the integral lattice
and satisfies the shortness bound for every translate and conjugation.
Consequently the same witnesses meet both clauses of
Theorem~\ref{A:main} for the family \(y\).

\section{The lower determinant estimate}\label{c3:section}

Let $W$ be the good subspace on $Y_0$ introduced in
\S\ref{sz:setup}. The archimedean factor at $Y_0$
is the prescribed discrete-series type. The algebraic Weyl-valued
model used to measure a lifted vector is on the definite space $Y_m$.
We use the common $H,U,V,L$ and write
\[
 B_0=U/(mk),\qquad
 A_0=\sum_{\tau\in\Phi}
   \bigl(j(\kappa_\tau)-
          S(\kappa_\tau,\tau)\bigr),\qquad H=mk A_0.
\]
With $j$ and $S$ as in \eqref{sz:height},
$k=r+1-m$, and the difference of ranks in a single lift is
$\delta=5r$. The lift of the ordered family will always be the full
chain $\Theta_\eta$. For an $m$-fold twist $\eta$, set
$N_\eta=\prod_{w\in Q(\eta)}q_w$, so $N_1=1$. At a prime $\ell$
the allowed twists satisfy $(\ell,N_\eta)=1$.
Theorem C, below, uses the construction hypotheses in
\S\ref{c3:frontier}.

\paragraph{Local and archimedean structures.}
Use the identifications, integral modules, coordinate vector $v_1(w)$ of $\Theta_1w$, and rational metric $G_{\rm cl}$ of \S\ref{sz:setup}. Write
$\bar\iota=\iota\circ\mathrm{conj}$. The chosen
level makes the stabilizers of the final class representatives trivial.
Indeed, if a rational unitary element $u$ belongs to a conjugate of
the level, then its characteristic polynomial has integral
coefficients (conjugate it into the integral adelic matrices).
For rank $h$, the algebraic integer $h-\operatorname{tr}u$ is
zero only if $u=1$: at a complex place the definite unitary matrix
is diagonalizable with eigenvalues on the unit circle, whose sum
can be $h$ only if they are all one. At all embeddings its modulus
is at most $2h$. The algebraic integer is divisible by every chosen place
above $p_0$, because the conjugated matrix reduces to the identity
there. The admissible congruence $p_0\equiv1$ modulo the cyclotomic conductor
makes $p_0$ completely split in $E$: its Frobenius fixes the roots
of unity generating the ambient cyclotomic field. The $2f$ primes
come in conjugate pairs. Since
$\pi_M\overline{\pi_M}=p_0p_1$ with $p_1\ne p_0$, precisely
one prime of each pair divides $\pi_M$, with multiplicity one.
Thus there are $f$ chosen primes of norm $p_0$. If the trace
difference were nonzero,
its absolute norm would be at least $p_0^f$, but at most
$(2h)^{2f}$. The rank bound $h\le2n_m$, with $n_m=k+m(5r+1)$, and
the level bound $p_0>(8n_m)^2$ make this impossible. The lifted inner product is the finite class sum of Weyl pairings.
Here is the local calculation for one term of this sum.
Let $F$ be the integral embedding of the Weyl module into its
column tensors with orthonormal filling basis. A row of $F$ is the
expansion of the product of minors for that filling in the standard
tableau basis. Standard fillings give linearly independent columns.
Cauchy--Binet gives $F S(a)=M(a)F$, where $S$ is the Schur action
and $M$ the tensor of the exterior-power matrices. Thus the positive
integral Gram matrix $G_W=F^tF$ satisfies
$G_WS(a)=S(a^t)^tG_W$, since $M(a^t)=M(a)^t$. Each determinant
character in the actual weight has the same value on $a^t$ and $a$,
so the relation holds for the full weight action $\pi$.

At a representative $g$, let $a$ be the tuple of its matrices for
$\iota$. The component of $g^{-1}$ for $\bar\iota$ is $a^t$, by
unitarity. For vectors integral at the opposed identifications,
$x=\pi(a^t)\bar\iota(u(g))$ and
$y=\pi(a^{-1})\iota(v(g))$ have integral coordinates by the gauge
definition. Their pairing $x^tG_W y$ equals
$\bar\iota(u(g))^tG_W\iota(v(g))$ by the transpose relation.
It is integral. The class weights are one, so summation gives
\begin{equation}\label{c3:local-pair}
 u\in\mathcal L_{\eta,\bar\iota},\ v\in\mathcal L_{\eta,\iota}
 \quad\Longrightarrow\quad
 |\langle u,v\rangle|_\iota\le1.
\end{equation}
The same class sum at the trivial twist has metric
$G_{\rm cl}=I_{\rm classes}\otimes G_W$, giving the squared conjugate norm
in \eqref{sz:covol}. We use the single matrix from \eqref{sz:index}
at each identification. For injectivity of $\Theta_\eta$ one applies
the sharp-datum injectivity assertion of input K to the elementary
arrows in the lower triangle. The positive level index and Weyl
dimension in the algebraic-to-Petersson norm conversion imply that
a zero algebraic lift has zero weighted chain. Its weights are
nonzero finite group orders, so reversing the injective arrows gives
zero in $W$.

\subsection{A signed height calculus}\label{c3:height}

For a family $u=(u_\iota)$ of complex scalars, write
$\operatorname{ht}(u)\le C$ if there are a multiplier $q\ne0$,
a real $\alpha$ and numbers $h_\ell\ge0$ such that
\begin{equation}\label{c3:heightdef}
 |\sigma q|\le e^\alpha,\qquad |qu_\iota|_\iota\le
 \ell^{h_\ell},\qquad
 \alpha+\sum_{\ell\in F}h_\ell\log\ell\le C.
\end{equation}
All automorphisms, identifications and finite prime sets are included.
This is a predicate, not a conventional nonnegative height function.
In particular $\alpha$ and $C$ may be negative. For products multiply
the multipliers and add both kinds of exponents. This proves the
additivity of the upper bounds, including powers with nonnegative
integer exponent. Multiplication by a family of powers
$N_{\eta_\iota}^{s_\iota}$ with rational exponents changes no finite
bound when the twist is allowed: the positive real root is an
$\iota$-adic unit because some integer power is a power of the
prime-to-$\ell$ integer $N_{\eta_\iota}$.

If $\operatorname{ht}(1)\le C$ then $C\ge0$. Indeed the multiplier
$q$ is algebraic: if it were transcendental, a field isomorphism
$\mathbb C\simeq\overline{\mathbb Q}_2$ could send it to a
transcendental element with arbitrarily large $2$-adic norm, contrary
to the uniform bound. To see the extension assertion, both fields
are algebraically closed of cardinality continuum and transcendence
degree continuum over $\mathbb Q$; extend the isomorphism of rational
function fields using transcendence bases and algebraic closures.
The set of elements algebraic over $\mathbb Q$ is countable, whereas
any nonempty $2$-adic ball contains uncountably many elements (already
in a translate of a ball in $\mathbb Q_2$). Thus the prescribed
transcendental elements exist. Set $K=\mathbb Q(q)$, $d_K=[K:\mathbb
Q]$, and $R=N_{K/\mathbb Q}(q)\ne0$. All embeddings of $K$ at a
finite or infinite place extend to identifications or automorphisms.
Consequently
\[
 \log|R|\le d_K\alpha,\qquad
 \log|R|_\ell\le d_K h_\ell\log\ell.
\]
The ordinary rational product formula, summed over the finite prime
set supporting the numerator and denominator of $R$, yields
$0\le d_K C$. In particular if a nonzero constant $z$ satisfies
$\operatorname{ht}(z)\le C$ and $|\sigma z|\le e^{C'}$
for every $\sigma$, then use multiplier $z$ to bound the height of
$z^{-1}$ by $C'$ and multiply to obtain
\begin{equation}\label{c3:product}
 0\le C+C'.
\end{equation}

\subsection{A regular family of finite-place twists}\label{c3:twists}

We need a single $m$-fold twist $\eta_\iota$ for each local
identification, simultaneously good for all $m$ single lifts, not
one twist for each member of $y$ or one for each $j$. No relation
between $\eta_\iota$ and $\eta_{\bar\iota}$ is required.
A critical pair $p=(\nu,z)$ is a unitary character with its negative
archimedean weight and tame part, and an integer $z$ in the critical
range specified by $D$; write
\[
 R_p(\vartheta,s)=
 \frac{|L^{S}(z,(\nu\vartheta^s)^{-1})|}
 {|L^{S}(z,\nu^{-1})|},\qquad s=\pm1,
\]
where in the denominator $S$ comprises the chosen places and their
conjugates, and in the numerator it additionally comprises $Q(\vartheta)$
and its conjugates. The change of omitted factors is included in the
explicit modulus powers below. For $s=1$ omit $s$.
The positive ratios are viewed as complex scalars, so their
conjugates need not have the same size. A ramified pair has tame part
nontrivial exactly at the chosen places. Let $J_j(\eta)=\eta^j$,
and use the superscript $j$ only for the power of the twist.

The required regularity is as follows:
\begin{enumerate}
\item $(N_{\eta_\iota},\ell)=1$;
\item for each ramified $p$, each $j$ and sign $s$,
$\operatorname{ht}(R_p(\eta_\iota^j,s)^{-1})\le x_1$;
\item for each $j$ there is a nonnegative logarithmic local budget
$\beta_{j,\ell}$ with $\sum_{\ell\in F}\beta_{j,\ell}\le B_0$.
Write $E_{j,\vartheta}$ for the image of the single lift on good
forms on $W_j$ and $P_{j,\vartheta}$ for its orthogonal projection.
There is an operator $\Lambda_\iota$ on $E_{j,\eta_\iota^j}$
with $|\det\Lambda_\iota|_\iota\ge
\exp(-d_T r\beta_{j,\ell})$ such that
$\Lambda_\iota P_{j,\eta_\iota^j}$ takes integral forms
orthogonal to false lifts to integral forms there. Also, for each
main critical pair $p$ of lift $j$, if $\Xi,\Xi'$ are test forms
for its character lift $\phi_\iota$ (integral for every embedding into $\mathbb C_\ell$ and
orthogonal to false lifts in the kernel of the adjoint), then
\[
 |\langle P_{\phi_\iota}\Xi,P_{\phi_\iota}\Xi'\rangle|_\iota
       \le e^{\beta_{j,\ell}}.
\]
Here $P_\phi$ is projection onto $\mathbb C\phi$.
\end{enumerate}

Here is the reduction of regularity to the local arithmetic input.
At fixed $\ell,\iota$, the good Hecke eigensystems on $W_m$ form a
finite set: their commuting normal Hecke operators simultaneously
diagonalize a finite dimensional space. For every $j$ and every such
system one takes the rank-$r$ main summand of its Galois representation;
for each main critical pair take the rank-one summand for the trivial
system. Join both finite collections. The simultaneous twisting
theorem for these collections and the ramified pairs produces $\eta$
whose first $m$ powers have auxiliary modulus away from $\ell$,
short divisible Selmer groups and inverse-ratio norms at most
$\ell^{u_\ell}$, where $G_0$ is the product of the orders of $\mathrm{GL}_{rt}(k_w)$ over the chosen places above $p_0,p_1$, with $t=5r$, and
\[
 u_\ell=(4f+3)r v_\ell(G_0)+2^{-\ell}.
\]
In the source the construction starts at a chosen place away from $\ell$
using its tame inertia generator, constructing lattices in the main
summands and Stickelberger systems for the ramified ratios, and
applying the simultaneous anticyclotomic twisting theorem to all
powers with error $2^{-\ell}$. Identification of the resulting avatar
with an automorphic tame character gives the same modulus for every
power: the avatar is nontrivial at each auxiliary place for each
power. Appendix~\ref{sec:finite} passes from the simultaneous polynomial
datum to a single twist using one root-of-unity tuple for all powers.
The Selmer and Stickelberger construction of the datum remains in R.

The first-order congruence denominator argument on the span of
lift generators applies these Selmer lengths to both the rank-$r$
space and the rank-one line. In more detail, choose a basis $v_i$ from the lift generators.
The first-order congruence calculation gives one operator $G_i$ for
each generator, diagonal on this basis with its $i$-th diagonal
entry of valuation at most the lift-valuation bound. Each $G_i$
after projection preserves the integral module on forms orthogonal
to false lifts. As the residue field of $\overline{\mathbb Q}_\ell$
is infinite, a combination with integral scalar coefficients
can avoid cancellation in every diagonal position: normalize the
entries in each row by one of largest norm, and avoid the finitely
many proper linear hyperplanes defined by their reductions. The
resulting diagonal operator has every entry of valuation at most
the stated bound, and retains integrality by linearity. Multiplying
by $G_0$ clears all stabilizer denominators for pairings. Taking
the product of the eigenvalues gives determinant bound
$\ell^{-\dim(E)r\,\mathrm{clearingBound}(D,\ell)}$ (rank one for
the line). In the rank-one case the denominator bound and the
pairing bound therefore control the line projection itself, with
the class-number factor in the local budget accounting for the
normalization of test forms. Appendix~\ref{c4:central-average} shows that the rank-one
central average changes the pairing by exactly the class factor
$N_0$. Existence of the generator-wise operators remains in R.
Thus one may take
\[
 \beta_{j,\ell}=\bigl(\mathrm{clearingBound}(D,\ell)
                     +v_\ell(N_0)\bigr)\log\ell ,
\]
where $N_0$ is the number of $U(1)$ classes at chosen level. This
choice is actually independent of $j$.

Since $\sum_{\ell\in F}v_\ell(G_0)\log\ell
\le\log G_0$ and
$\sum_{n\ge2}2^{-n}\log n\le\sum_{n\ge2}n2^{-n}<2$,
\[
 \sum_{\ell\in F}u_\ell\log\ell
 \le (4f+3)r\log G_0+2\le x^{999}\le x_1 .
\]
The clearing bound is the sum of the local defect, $(1+2\delta)u_\ell$,
the denominator error, and $v_\ell(G_0)$. Applying the defect envelope and the numerical estimates
in Appendix~\ref{sec:finite} gives
\begin{align*}
 \sum_{\ell\in F}\beta_{j,\ell}
 &\le 4\cdot10^6e^{180} f(m/k)bL
 +(1+2\delta)x^{999}+x^{4000}+\log G_0+x^{30}\\
 &\le B_0.
\end{align*}
The denominator error is a linear combination of prime valuations of
the fixed odd class bound and $G_0$; the source bounds its sum by
$x^{4000}$ and the $N_0$ term by $x^{30}$. Select one such twist for
each identification. With multiplier $1$ the local inverse-ratio
norms imply item 2 by (\ref{c3:heightdef}). The remaining local arithmetic hypotheses in this reduction are
specified under R in \S\ref{c3:frontier}.

\subsection{The height gain in a single lift}\label{c3:single}

Fix $j\in[1,m]$. Let $T_\vartheta$ denote the stage-$j$ theta lift with twist $\vartheta$ restricted to the good forms of the rank-$r$ hermitian space $W_j$ of signature $(k+j-1,m-j)$. All these
sources have dimension $d_T$ by the identity-operator case of the
trace comparison in Appendix~\ref{sec:trace}. Define $Z_j(\eta)=\det(T_{\eta^j}^*T_{\eta^j})$.
The comparison operator is
\[
 U_\vartheta=G_0 h_j^2 g_\vartheta g_\vartheta^* ,
\]
where $g_\vartheta$ is the gluing map from tuples of forms on the
smaller definite space $V^\circ$ and $h_j$ is the hook number of its
negative Weyl shape. Set $A_{\rm gl}=fmbL^2$. Input G gives the seesaw pairing with a tuple $\zeta$ as the sum
over kernel choices $i$ of fixed nonzero $c_i$ times the integral
of $y(g)\,F_i(g)\,\theta^\circ(g,\zeta_i)$, where $F_i$ is the
carrier factor and $\theta^\circ$ the complementary block kernel.
This formula is independent of the twist after multiplication
by $c'_\vartheta$. Testing against tuples supported in one block
shows that a vector killed by the resulting map has zero pairing
with each of these products. The products-spanning input says
precisely that such a good vector is zero. In particular the map
$c_\vartheta' g_\vartheta^*T_\vartheta=\Pi_j$
is independent of $\vartheta$, with $\Pi_j$ injective, and
the quotient $\rho_\vartheta$ of the twisted to untwisted product of
Hecke values satisfies
$\rho_\vartheta=|c_1'|^2/|c_\vartheta'|^2\ne0$.
Hence, for the source inner product,
\[
 T_\vartheta^*U_\vartheta T_\vartheta
 =\rho_\vartheta U^\circ,
 \qquad U^\circ=G_0h_j^2|c_1'|^{-2}\Pi_j^*\Pi_j>0.
\]
Writing $C_\vartheta$ for the compression
$P_{j,\vartheta}U_\vartheta|_{E_{j,\vartheta}}$ and
$F_\vartheta=\det C_\vartheta$, we have
\[
 Z_j(\eta) F_{\eta^j}=\det(U^\circ)\rho_{\eta^j}^{d_T}.
\]
Indeed relative to an orthonormal basis of the image,
$T^*UT=T^*C T$, so determinants multiply, and $\det T^*T=|\det T|^2$.
All factors are nonzero by positivity of $U^\circ$ and injectivity.
Division by the untwisted identity gives
\begin{equation}\label{c3:compression}
 \frac{Z_j(1)}{Z_j(\eta)}
 =\frac{F_{\eta^j}}{F_1}\rho_{\eta^j}^{-d_T}.
\end{equation}

\paragraph{Integral comparison and conjugate norms.}
In the tuple space take the integral tight frame $(f_a)$, at rational
translates of the class representatives integral above $\ell$.
Its frame operator is $h_j^2 I$; conjugated coordinates in the
translated Weyl lattice are integers. The gluing integrality hypothesis takes this
frame to forms integral at both $\iota$ and $\bar\iota$ away from
the modulus. The frame identity gives
\[
 U_\vartheta v=\sum_a G_0
       \langle g_\vartheta f_a,v\rangle\,g_\vartheta f_a .
\]
The factors $G_0$ clear the orders of the class stabilizers, so the
coefficients are integral for integral $v$ by the same local Weyl
pairing as (\ref{c3:local-pair}). The output is integral.
The gluing-period identity also implies that $U_\vartheta$ takes
lifts to forms orthogonal to false lifts.

To construct the frame, let $h=h_j$ and use the filling matrix $F$
just described for the negative shape. Write $G_W=F^tF$.
Let $B_{ij}=b_i(\partial)b_j(0)$ be the integral Fock pairing of
the standard minor polynomials. Its normalization gives
$G_WB=h I$. Indeed substitution by $a$ and coefficient differentiation
are adjoint under transposition of $a$, so $B^tG_W$ commutes with
the Schur representation at every place. The tensor of its
irreducible highest weight modules is irreducible for the product
of the general linear groups. Schur's lemma therefore makes the
operator scalar. Evaluation at the highest polynomial, whose
differentiation pairing with itself is $h$ and whose kernel value
on the standard flag is $1$, fixes that scalar. The last
normalization is the definition of the hook number in the Fock
model. For a tensor basis vector $e_T$ take $a_T=BF^te_T$.
Then $a_T$ is integral and $G_Wa_T=hF^te_T$. Hence
$Fa_T=h\,\operatorname{pr}_{\operatorname{im}F}e_T$.
The projections of an orthonormal basis are a Parseval frame for
the image. Multiplication by $h$ gives frame constant $h^2$ and
$\|a_T\|^2\le h^2$.

If $a_T$ is a frame vector
and $u$ is a rational unitary translate with representation
$\pi(u)$, replace $a_T$ by
$\overline{\pi(u)^{-1}\overline{a_T}}$.
This is a unitary change of frame: the frame operator remains
$h_j^2 I$. Moving the conjugate coordinates back gives the
original integer vector. The Gram identity for the unitary
representation is rational in the coordinates and their
conjugates over the CM field, and remains valid at each
archimedean embedding. Since the original vector is integral,
the moved conjugate norms retain the bound $h_j^2$.
Independently place these vectors in each block of the finite
tuple space, with zero other blocks; summing the frame identity
blockwise gives the tuple tight frame.

At the trivial twist use the tight frame at the unshifted
representatives. Its conjugate squared tuple norms are at most
$h_j^2$. Let $B$ be the set of tuple-coordinate blocks, $I$ the
frame indices; the gluing and counting estimates are
\[
 \|g_1 v\|_\sigma^2\le e^{2A_{\rm gl}}\|v\|_{\sigma,\mathrm{tuple}}^2,\qquad
 |B|(|I|G_0 h_j^2)^2\le e^{2A_{\rm gl}}.
\]
For the cardinality estimate put $N=|B|$ and let $C_f$ be the
number of column fillings for the negative shape. The frame has
$|I|=NC_f$. The class and hook estimates of input G give
$(NG_0)^2\le e^{A_{\rm gl}}$, $C_f,h_j\le e^{x^{999}}$. Since
$N\le e^{A_{\rm gl}/2}$ and $6x^{999}\le A_{\rm gl}/2$,
\[
 N(NC_fG_0h_j^2)^2\le
    \exp(A_{\rm gl}/2+A_{\rm gl}+6x^{999})\le e^{2A_{\rm gl}}.
\]
For $S=e^{2A_{\rm gl}}\|v\|_\sigma^2$, the conjugate coefficient of the
$a$-th frame member is at most $G_0h_j\sqrt S$ by Cauchy--Schwarz.
(The general automorphism changes conjugation to its conjugated
automorphism; the frame bound holds for both.) By triangle inequality
each block of the resulting tuple has norm at most
$|I|G_0 h_j^2\sqrt S$. Summing the squares over $B$ and gluing once
more bounds $\|U_1 v\|_\sigma^2$ by $e^{6A_{\rm gl}}\|v\|_\sigma^2$.
The projection on lifts, applied to a vector orthogonal to the false
lifts, is contractive for every conjugate norm; this uses the
algebraic characterization by the central idempotent and the
Hecke polynomial kernel, not continuity of $\sigma$.
Consequently, for an eigenvector $v\ne0$ of $C_1$ with eigenvalue
$\lambda$, $|\sigma\lambda|^2\|v\|_\sigma^2\le
e^{6A_{\rm gl}}\|v\|_\sigma^2$. The determinant is the product of roots
of the characteristic polynomial with multiplicities, yielding
\[
 |\sigma F_1|\le e^{3d_T A_{\rm gl}}.
\]

At a local identification use $\vartheta=\eta_\iota^j$ and
the clearing operator from regularity. Then $\Lambda_\iota
C_\vartheta$ preserves the integral module of its image. An
eigenvector can be scaled into that module with at least one
transported coordinate of norm one: take the finitely many class
values at representatives integral above $\ell$, take the maximum
of the coordinate norms in the local Weyl model and scale by the
inverse of a coordinate attaining it. Every eigenvalue is therefore
integral. By the product of eigenvalues and the bound on
$\det\Lambda_\iota$,
\[
 |F_{\eta_\iota^j}|_\iota\le e^{d_T r\beta_{j,\ell}}.
\]
In (\ref{c3:heightdef}) take multiplier $F_1$, archimedean budget
$3d_T A_{\rm gl}$, and finite budgets $d_T r\beta_{j,\ell}/\log\ell$.
This proves
\begin{equation}\label{c3:compressionheight}
 \operatorname{ht}(F_{\eta_\iota^j}/F_1)
 \le d_T(3A_{\rm gl}+rB_0).
\end{equation}

\paragraph{The main rank-one ratios.}
The $-A_0$ gain comes from the following two different sides
of a rank-one computation. Put
$J_0=\sum_\tau j(\kappa_\tau)$,
$S_{\Sigma}=\sum_\tau S(\kappa_\tau,\tau)$ so that
$A_0=J_0-S_{\Sigma}$. For a main pair $p$ of lift $j$, let
$\phi_\vartheta$ be its normalized rank-one lift to $U(2z)$
and $\phi_1$ the untwisted lift.

To bound the constant $\|\phi_1\|^2$, compare it to the lift of a
flat pair $q$ of the same rank and empty row. The integral algebraic
part $I_q$ of the flat value at trivial twist is integral everywhere
and equals $G_0\mathrm{denomBound}(q)^2\|\phi_{q,1}\|^2$.
For automorphisms fixing the base field, the single-lift norm formula
gives
\[
 \log|\sigma\|\phi_{p,1}\|^2|
 =C_\sigma+\epsilon_{p,\sigma}
    +\sum_\tau j(l_\tau),\qquad
 l_\tau=p.\mathrm{negWeight}(\tau)-z\ge\kappa_\tau .
\]
Here the constant is common to the two pairs,
$\epsilon_{p,\sigma}\ge-4f$ and
$\epsilon_{q,\sigma}\le4f$ with empty row $l_\tau=0$.
The same moduli of conjugates are realized by automorphisms fixing
the base field. Thus for the multiplier $I_q/\|\phi_1\|^2$
the common constant cancels, the finite budgets are zero and the
archimedean budget is at most
$x^{33}+8f-J_0\le B_0-J_0$. Therefore
\begin{equation}\label{c3:main-arch}
 \operatorname{ht}(\|\phi_1\|^2)\le B_0-J_0 .
\end{equation}
Here one uses monotonicity of $j$ on the
range in question, and bounds the fixed class and denominator
factors by $e^{x^{33}}$. These and the common-constant norm formula
are the numerical hypotheses in input F.

For the finite side let $s_\ell$ be the sum of the local gains at $\ell$ over $\tau\in\Phi$, computed below.
The crucial sharpness estimate is existence of a common nonzero scalar
$c_0$ with $|\sigma c_0|\le e^{B_0}$ and, for every allowed twist and
identification, a test form $\Xi$ satisfying
\begin{equation}\label{c3:test}
 |\langle\Xi,\phi_\vartheta\rangle|_\iota
 \ge |c_0 R_{q_2}(\vartheta)|_\iota \ell^{-s_\ell},
 \qquad 2\sum_{\ell\in F}s_\ell\log\ell\le S_{\Sigma} .
\end{equation}
The prime-sum inequality is just regrouping by prime powers.
In fact put
\[
 g_\tau(v)=(1-\chi_K(v))
       (\kappa_\tau/v-\kappa_{-v\tau})_+ .
\]
The notation $-v\tau$ uses the place of the class of $-v$ if
prime to the cyclotomic conductor and the identity place otherwise.
The weight is at least $a/x_1$, so $g_\tau(v)$ vanishes for
$v\ge x_1$; all terms are nonnegative. To get the local formula, compare the coefficient of each positive
exponent $e$ with the defining sum for the local gain.
At a split prime $1-\chi_K(\ell)^e=0$. At an inert prime this
coefficient is $2$ for odd $e$ and zero for even $e$, and the
local uniformizer weight is $1$. At the ramified prime the
coefficient is $1$ for every $e$, and the local uniformizer
weight is $1/2$. Thus
$2s_\ell=\sum_\tau\sum_{e\ge1}g_\tau(\ell^e)$.
Since $\Lambda(\ell^e)=\log\ell$ and the von Mangoldt
function vanishes at non-prime-powers, summing this equality over
the primes gives exactly $S_{\Sigma}$, and restriction to a finite set
only decreases it. Thus the prime-sum regrouping consumes no error.

Here $q_2$ is an empty-row ramified pair of rank four with the fixed
chosen-level tame character. To specify the origin of the sharper
Taylor input:
the rank $2z$ indices split into blocks of sizes $12,4,4,2z-20$.
The first block gives a Taylor coefficient of weight
$l_\tau\in[\kappa_\tau,\kappa_\tau+x^{42}]$ at a finite indicator
datum; the next carries $q_2$, and the last two carry empty rows
with tame characters whose product is the character selected by
class averaging. The sharp Taylor coefficient is scaled by
$T=M^{12f}c_{\rm vac}\prod_\tau
\mathrm{lossProduct}(l_\tau)^2$ with
$|\sigma T|\le e^{fmb/k}$. Its local valuation upper bound before
scaling is $s_\ell+2\sum_\tau v_\ell(\mathrm{lossProduct}(l_\tau))
-\sum_\tau v_\ell(l_\tau!)$ relative to the integral vacuum.
Multiplication by the factorials and division by $T$ removes both
loss terms, leaving $s_\ell$. The scaled coefficient is $C s(1)$
for a one-variable kernel $s$ on $U(1)$.

Appendix~\ref{jets:sharp} separates the valuation lower bounds
for all coefficients from the upper valuation bound supplied by
one coefficient. The latter direction is needed here.

Class averaging can be made explicit. The rank-one definite
class space is a finite abelian quotient $C$ of the group of finite adelic
points modulo rational points and the chosen level. There
are no stabilizers. Translate the chosen representatives by the
inverse of the representative of the identity class, so that
this representative becomes $1$. The map taking the finite component of a norm-one adele to its
rank-one matrix sends rational points into the rational subgroup,
and integral ideles reducing to $1$ at the chosen places into the
level. Pulling a character of $C$ back along this map therefore
gives a weight-zero tame character with possible conductor only at
the chosen places. At a finite rank-one group element its value is
the original character, since this map composed with the rank-one
adelic determinant is the identity. For a character $\chi$ form the
normalized mean $A_\chi=|C|^{-1}\sum_{g\in C}
\chi(g)s(g_{\rm rep})$. Character orthogonality gives
$\sum_\chi A_\chi=s(1)$; the ultrametric inequality therefore
gives a character with $|A_\chi|_\iota\ge|s(1)|_\iota$.
Call this mean $A_s$.
For the other blocks let $W_\psi$ be the product of the two flat
integral algebraic parts with characters multiplying to the selected
character. The geometric mean $I$ of all such products can be
taken independent of the selected character because each
factorization runs over the entire finite character group. By
multiplicativity some factorization has $|W_\psi|_\iota\ge|I|_\iota$.
All conjugates of $I$ have norm at most $e^{2x^{34}}$ by the
corresponding flat bounds. The kernel pairing identity is
\[
  T I_{q_2} W_\psi(C A_s)
    = N_0^2\langle\Xi,\phi_\vartheta\rangle/R_{q_2}(\vartheta).
\]
Taking norms, using $|N_0|_\iota\le1$, proves~(\ref{c3:test})
with $c_0=T I_{q_2} I$; its archimedean bound is
$e^{fmb/k+3x^{34}}\le e^{B_0}$.
The Taylor sharpness bound, existence of the block test form with
the displayed identity, and the flat integral algebraic estimates
are the substantive inputs of this argument (input F below).

The precise external result here is the
\srcline{New/OneLift/HeckeValues/TwistedValueHeights}{700}{sharp test-form theorem}.
For fixed main pair and rank-four companion, it chooses one nonzero \(c_0\)
before the prime, the identification and the allowed twist. In its proof
the degrees are prescribed by the negative weights and the fixed shift.
The \srcline{New/Derivatives/Sharpness/BoundAttained}{381}{attainment theorem}
supplies the coefficient data; the geometric mean \(I\) used for \(c_0\)
is also fixed before local specialization. The indicator, averaging
character and factorization yielding \(W_\psi\) may depend on the
specialization. The
\srcline{New/OneLift/HeckeValues/TwistedValueHeights}{571}{integral kernel-pairing theorem}
realizes these choices on the required test form. This order of choice
permits the same \(c_0\) and twist at two identifications.

Now take $\vartheta=\eta_\iota^j$, $\phi=\phi_\vartheta$,
$\rho=R_{q_2}(\vartheta)$ and choose another test form
$\Xi'$ at $\bar\iota$ for the \emph{same} $\vartheta$.
Since $\rho$ is real, multiplying the two inequalities gives
\[
 |\langle\Xi,\phi\rangle\langle\phi,\Xi'\rangle|_\iota
 \ge |c_0\bar c_0\rho^2|_\iota\ell^{-2s_\ell}.
\]
Denote the nonzero pairing product by $I_\iota$. Direct projection
onto the line gives
$\langle P_\phi\Xi,P_\phi\Xi'\rangle=I_\iota/\|\phi\|^2$.
By the rank-one norm formula for the Hecke ratio,
\begin{align}
 R_p(\vartheta)^{-1}
 &=N_\eta^{2z-1}\frac{\|\phi_1\|^2}{\|\phi\|^2}\notag\\
 &=N_\eta^{2z-1}\|\phi_1\|^2
 \langle P_\phi\Xi,P_\phi\Xi'\rangle\rho^{-2}
 \frac{\rho^2}{I_\iota}.\label{c3:four-factors}
\end{align}
The first factor after the modulus power costs $B_0-J_0$ by
(\ref{c3:main-arch}); the projection costs $B_0$ by regularity;
the ramified ratio costs $2x_1$. For the last factor use multiplier
$c_0\bar c_0$, archimedean budget $2B_0$ and finite budgets
$2s_\ell$; equivalently the local norm bound is
$\ell^{2s_\ell}$. It costs at most $2B_0+S_{\Sigma}$.
Since $x_1\le B_0$ and the modulus power costs nothing,
\begin{equation}\label{c3:main-height}
 \operatorname{ht}(R_p(\eta_\iota^j)^{-1})
 \le 6B_0-A_0 .
\end{equation}


\paragraph{Counting and combining the three exchanges.}
The three exchange products for the comparison operator are
products of squared critical values with modulus powers.
The unsigned exchange character contributes points $1\le z\le7(m/2)+k/2$,
the two signed flat characters points $1\le z\le m$ and
$1\le z\le m/2$. Reserve for the main-pair estimate the points
$7(m/2)<z\le7(m/2)+k/2$. Since $m$ and $k$ are even, there
are $k/2$ of these, and $7(m/2)+m+m/2=5m$ remaining
ramified points. The ratio of exchange products is therefore
$N_\eta^e(\prod_{a=1}^{k/2}R_{p_a}\prod_{b=1}^{5m}
R_{q_b})^2$, with the prescribed signs on the second product.
This is the partial $L$-value factorization of input G; taking
inverse heights, the squared product doubles the counts, giving
\[
 \operatorname{ht}(\rho_{\eta_\iota^j}^{-1})
 \le k(6B_0-A_0)+10mx_1.
\]
The scale inequality
$3A_{\rm gl}+rB_0+10mx_1+6kB_0\le kLB_0$
and (\ref{c3:compression}), (\ref{c3:compressionheight}) now prove
\begin{equation}\label{c3:single-height}
 \operatorname{ht}\left(\frac{Z_j(1)}{Z_j(\eta_\iota)}\right)
 \le d_T\,k(LB_0-A_0)=d_T\left(\frac{V}{m}-kA_0\right).
\end{equation}


\subsection{From single lifts to the chain}\label{c3:chain}

Put $Z(\eta)=\det(\Theta_\eta^*\Theta_\eta)>0$.
One cannot simply multiply the norms of the single lifts, since
they have different sources; the triangle-of-lifts norm identity
and the equality of Hecke multiplicities are needed. In the source
the $j$-th single-lift Gram operator is a constant $a_j\ne0$
times the common local twist factor $c_Q$ times the Rallis operator
$\mathcal R_j(\eta)$ on its good space. For the unnormalized coset sum write $R$ for the operator with
matrix coefficients $\sum_\gamma K_\eta(\gamma)
\langle x,\gamma y\rangle$; $K_\eta$ is the unramified vacuum
coefficient. For the sharp finite block datum of mass $M_{\rm blk}$
the asserted Petersson Rallis formula of input K (cf.~\cite{gqt} for the general inner product formula) gives the
same sum times $C_jM_{\rm blk}$ with $C_j$ depending only on
the archimedean datum. The latter is fixed under tame weight-zero
twists. Converting to algebraic forms multiplies by the index
of the level and the Weyl dimension. The finite mass relation is
$[K_0:J]M_{\rm blk}=\operatorname{gramFactor}/M_Q$;
the quotient is the chosen-place factor times $c_Q$.
Thus the scalar multiplying $R$ is $a_j c_Q$, where $a_j$
is nonzero and independent of the twist. The paired-series
argument of Appendix~\ref{sec:trace} realizes the two $R$'s
on $W_j$ and $Y_0$ by one Hecke operator and identifies its
two determinants.
For specificity, write $w_n$ for the order of the general linear
group in the $n$-th step of the lower triangle and $M_Q$ for the
number of $Y_1\times\mathrm{RowIndex}$ matrices over the residue
ring of the chosen and twist modulus. The first-arrow Gram formula
in row $i$ is $w_0^{-1}C_i M_Q^{-1}\mathcal R_{i+1}$, transported
to the common starting space, with $C_i$ independent of twist.
Here the Gram pairings are sesquilinear.
The weighted triangular computation works by induction on the row
length $n=m-i$. In the zero row the weighted chain is the identity.
Moving the adjoint of the first step past the $n-1$ subsequent
steps multiplies its scalar by the product of the commutation
constants in those squares. Their moduli, up to a fixed norm
power, are $Q_{i+1,j}$, $i+1<j\le m$; the adjoint law contributes
only scalar factors of modulus one in this sense. Replacing the
paired first-step ascent and descent by its first-arrow Gram
formula inserts $C_i M_Q^{-1}\mathcal R_{i+1}$.
The weights telescope from $w_0w_1^{n-1}/w_{n-1}$
to $w_0w_1^n/w_n$. More explicitly the weighted row transfer
gives the factor $w_0 w_1 w_{n-1}/w_n$ and the first-arrow
pairing contributes $1/w_0$.
Thus at row zero the weighted chain $u_\eta$ has Gram pairing
\[
 \langle u_\eta v,u_\eta v'\rangle
 =\kappa_\eta\frac{w_0w_1^m}{w_m M_Q^m}
    \langle v,\prod_{j=1}^m\mathcal R_j(\eta)v'\rangle,
\]
where the modulus of $\kappa_\eta$ contains exactly one factor
$Q_{ij}$ for each $1\le i<j\le m$. The induction uses the square
commutation and adjoint identities; their arithmetic derivation
is part of input K, not an identity of arbitrary theta lifts.

Conversion from the weighted Petersson model to the final
algebraic model multiplies this Gram form by the Weyl dimension
$h_Y$ and the index of the final level. The precise index
computation is
\[
 \mathrm{ind}(J_\eta)=C_1 N_\eta^u
     \mathrm{gramFactor}(\eta)^m\frac{w_m}{w_0w_1^m}.
\]
The integer $u$ and constant $C_1$ are independent of $\eta$:
at the chosen places all group orders are fixed, and at twist
places the block upper-triangular parabolic has order equal to the
product of the general linear orders of diagonal blocks times
$q_w^{\sum_{a<b}h_a h_b}$ for block sizes $h_a$.
This is obtained by independently choosing each upper off-diagonal
matrix block. Also
\[
 \frac{\mathrm{gramFactor}(\eta)}{M_Q}
 = \mathrm{chosenFactor}\,c_Q,\qquad
 c_Q=N_\eta^{r\delta}\prod_{w\in Q(\eta)}
                      \prod_{s=\delta+1}^{\delta+r}(1-q_w^{-s}).
\]
The product of $1-q_w^{-s}$ is the proportion of rank-$r$
matrices among $(r+\delta)\times r$ matrices: choose successive
independent columns in $q_w^{r+\delta}-q_w^i$ ways for $0\le i<r$.
The precise matching of the parabolic levels and of these factors
is input K. Multiplying the displayed formulas cancels every
$w_n$ and $M_Q^m$, leaving $h_Y C_1 N_\eta^u
\kappa_\eta\mathrm{chosenFactor}^m$ as the scalar multiplying
$c_Q^m$. The resulting operator identity is
\[
 \Theta_\eta^*\Theta_\eta
 =c_\eta c_Q^m\prod_{j=1}^m\mathcal R_j(\eta),\qquad
 |c_\eta|=C N_\eta^{e/2}\prod_{1\le i<j\le m}Q_{ij}(\eta),
\]
where $C>0,e\in\mathbb Z$ are independent of $\eta$. Here
$Q_{ij}$ is the product at $1\le s\le\delta$ of critical ratios
for the character $\nu_j/\nu_i$ at twist $\eta^{j-i}$,
where $\nu_j$ denotes the untwisted character at stage $j$.
This character is a power of the common ratio between consecutive
lift characters and is ramified at the chosen places.
The local triangle identities and the analytic Rallis formula
enter through K. The trace transfer uses the separate hypotheses
of Appendix~\ref{sec:trace}. Once they hold, take
determinants on $W$:
\[
 Z(\eta)=
 \left(c_\eta/\prod_j a_j\right)^{d_T}\prod_j Z_j(\eta).
\]
The common $c_Q^m$ and all Rallis determinants cancel. Since the
Gram determinants are positive, take absolute values, and divide by
the trivial-twist identity (all $Q_{ij}(1)=1$):
\begin{equation}\label{c3:chainfactor}
 \frac{Z(1)}{Z(\eta)}=
 \left(N_\eta^{e/2}\prod_{i<j}Q_{ij}(\eta)\right)^{-d_T}
 \prod_{j=1}^{m}\frac{Z_j(1)}{Z_j(\eta)} .
\end{equation}
Each $Q_{ij}$ contains $\delta$ ramified ratios, and $1\le j-i\le m$,
so their inverse heights cost at most $m^2\delta x_1\le b$ before
raising to the power $d_T$. Summing (\ref{c3:single-height}) over
the same $m$ twists $\eta_\iota^j$, using $mkA_0=H$, gives
\begin{equation}\label{c3:chainheight}
 \operatorname{ht}\left(\frac{Z(1)}{Z(\eta_\iota)}\right)
 \le d_T(b+V-H).
\end{equation}

\subsection{Short auxiliary completion}\label{c3:auxiliary}

We require an auxiliary spanning set in $W$, independent of $y$.
The construction produces a scalar $\gamma(\eta)\ne0$,
$\gamma(1)=1$, with
\begin{equation}\label{c3:gamma}
 |\gamma(\eta)|=N_\eta^{e'/2}
       \prod_{i,j=0}^{m-1}T_{ij}(\eta),
\end{equation}
where $e'\in\mathbb Z$ is fixed, and $T_{ij}(\eta)$ is the
product of the $\delta$ ratios for the character $\nu_{i+1}\nu_{j+1}$
at $\eta^{j+1}$. Its tame character is nontrivial at the chosen
places, so each is a ramified-ratio product. The powers now use
$j+1$, not $j-i$. The spanning vectors $z$ have
\begin{equation}\label{c3:auxbounds}
 \gamma(\eta)\Theta_\eta z\in\mathcal L_{\eta,\iota}
 \ ((N_\eta,\ell)=1),\qquad
 \|\Theta_1z\|_\sigma^2\le\exp(2x_1^2b^2).
\end{equation}

Here is the auxiliary construction reduced to its path-period
identities. For the final grid-level space apply the short-basis estimate
of \S\ref{A:source} with the conductor integer in the different.
The grid's principal level is contained in the chosen level and its
rank $h_0$ is at most $2n_m$, giving trivial stabilizers by the
argument above. The weight has at most $b r^4$ boxes. With
$D_0$ the conductor integer, the numerical specialization is
$6(h_0 f+b r^4)\log(6D_0(fh_0)^2)\le x^{99}b^2$.
Transport through the algebraic realization at the last grid vertex
gives spanning forms all of whose conjugate values have squared
Weyl norm at most $e^{x^{99}b^2}$. There is a fixed onto map $R$
from this space to principal forms on $Y_0$. Pairings at the end
of the chain against a dual Weyl vector $v$ have the expression
\[
  v^t M_{\rm dual}
    \frac{c_\eta}{c_1}\Theta_\eta R\varphi(g)
    =\beta\,\mathcal P_\eta(\varphi,g,v).
\]
Here $\Theta$ in this identity denotes the chain on all principal
forms, $\mathcal P_\eta$ the last-path period (with one factor
a Weyl-valued delta form), and $|c_\eta/c_1|$ is the right side
of (\ref{c3:gamma}). The path data arise by carrying the first
path through all $m^2$ turns; their scalar moduli multiply. The
fixed dual hook number $h>0$ controls recovery from pairings.
There is a single integer $N\ne0$, of size at most
$e^{x^{99}b^2}$, clearing the period denominators for all allowed
twists on integral inputs. The period bound at trivial twist is
\[
 |\sigma\mathcal P_1(\varphi,g,v)|
 \le e^{x^{200}b^2}\|\sigma v\|_{\rm Weyl}
                  \sqrt{B}
\]
when every conjugate value of $\varphi$ has squared norm at most
$B$. These are the path-period and dual-model inputs Aux in
\S\ref{c3:frontier}. Rescale all inputs by a fixed $s\ne0$:
if $\beta\ne0$, take $s=hN/\beta$, else $s=1$.
The rescaled pairings are $hN\mathcal P_\eta$ or zero. The normalized dual model in input Aux has the following two properties. The integral dual-pairing matrix $M_{\rm dual}$ admits an integer
matrix $C$ with $CM_{\rm dual}=hI$ in the normalized dual model.
Testing the pairing against the $i$th row of $C$ gives
$h\zeta_i\in h\mathcal O_\iota$ for the lifted value $\zeta$.
Since $h\ne0$, every coordinate is integral, even at primes
dividing $h$. At infinity the dual model represents
$v\mapsto v^tM_{\rm dual}\sigma\zeta$ by a vector $w$ with
$\|w\|=\|\sigma\zeta\|$. Evaluate the bound at
$v=\sigma^{-1}w$ before conjugation. The resulting inequality is
$\|w\|^2\le hN e^{x^{200}b^2}\sqrt B\,\|w\|$.
The hook estimate $hN e^{x^{200}b^2}\le e^{x^{201}b^2}$ gives
squared Weyl norm at most $e^{2x^{201}b^2}B$. With
$B=e^{x^{99}b^2}$ this is at most $e^{x^{202}b^2}$.
The images $sR\varphi$ span all principal forms since $R$ is onto.
Apply the good projection multiplied by its fixed integral scalar
at the identity translate; it preserves integrality on the lifts,
is surjective onto $W$ and increases the logarithm of the conjugate
squared norm to at most $x^{99}+x^{202}b^2\le2x_1^2b^2$.
This yields (\ref{c3:auxbounds}) with $\gamma=c_\eta/c_1$.

By (\ref{c3:gamma}) and regularity we have, for the same
$\eta_\iota$ as in (\ref{c3:chainheight}),
\begin{equation}\label{c3:gammaheight}
 \operatorname{ht}(|\gamma(\eta_\iota)|^{-2})
 \le 2m^2\delta x_1\le b .
\end{equation}
The nonzero modulus power has zero finite cost; the coefficient 2
here comes from squaring the inverse absolute value, independently
of the index correction.

\subsection{The determinant comparison}\label{c3:compare}

\noindent\textbf{Theorem C (conditional on the geometric and arithmetic inputs
in \S\ref{c3:frontier}).}
Let $y_1,\ldots,y_n$ be independent in $W$,
$n\ge(1-4/b^2)d_T$, with $\operatorname{IndexLe}(y,\Delta)$
and $\operatorname{CovolumeLe}(y,c)$. Then
\begin{equation}\label{c3:theorem}
 d_T(H-3V)\le c+2\Delta.
\end{equation}

\begin{proof}
The case $d_T=0$ follows from $n=0$, the empty Gram determinant
($c\ge0$) and the empty prime-set index bound ($\Delta\ge0$).
Suppose $d_T>0$. Extend the given ordered family to a basis
$w=(y_1,\ldots,y_n,z_1,\ldots,z_{n'})$ by choosing from the
auxiliary spanning set, with $n'=d_T-n$. Such an extension is
obtained by passing to the quotient $W/\operatorname{span}(y)$
and choosing a basis from the images of the spanning set.
The original family is unchanged. Let
\[
 \mathcal G(\eta)=\det(\langle\Theta_\eta w_i,
                            \Theta_\eta w_j\rangle)_{ij}.
\]
A basis change from an orthonormal basis of $W$ shows
$\mathcal G(\eta)=Z(\eta)\det(\langle w_i,w_j\rangle)$: the Gram
matrix is congruent by its change-of-basis matrix.
In particular $\mathcal G(1)\ne0$ and
$\mathcal G(1)/\mathcal G(\eta)=Z(1)/Z(\eta)$.

For fixed $\iota$ at $\ell$, use matrices $U_\iota$ and
$U_{\bar\iota}$ from the index assumption, applying both at the
\emph{same} twist $\eta=\eta_\iota$, allowed at both
identifications. Form the row transformations on the concatenated
family with blocks
$A_\iota=\operatorname{diag}(U_\iota,\gamma(\eta)I_{n'})$
and $A_{\bar\iota}=\operatorname{diag}(U_{\bar\iota},
\gamma(\eta)I_{n'})$. The auxiliary block need not be a matrix
over $\mathcal O_\iota$: it is its lifted vectors that are integral.
Both transformed lifted families are integral, for their
respective identification. By (\ref{c3:local-pair}) all entries
of their mixed pairing matrix have $\iota$-norm at most one.
Its determinant is
\[
 \overline{\det U_{\bar\iota}}\,\det U_\iota\,
 |\gamma(\eta)|^{2n'}\mathcal G(\eta).
\]
The two index determinants have norms each at least
$\ell^{-e_\ell}$ because conjugation carries $\bar\iota$ back
to $\iota$. Using the ultrametric determinant bound gives
\[
 \big|\,|\gamma(\eta)|^{2n'}\mathcal G(\eta)\,\big|_\iota
       \le\ell^{2e_\ell}.
\]
Hence
\begin{equation}\label{c3:local-height}
 \operatorname{ht}(|\gamma(\eta_\iota)|^{2n'}
                    \mathcal G(\eta_\iota))\le2\Delta .
\end{equation}
The two determinant corrections account for the coefficient $2$ on the index.

At infinity, for each $\sigma$ the Gram determinant of the
concatenated conjugate frame is bounded by the product of the
Gram determinants of its two blocks, by Fischer's inequality
(or by taking the exterior product and projecting the added
vectors to the orthogonal complement of the first span).
Hadamard on the second block gives
\[
 \det\operatorname{Gram}(v_\sigma(w))\le
 e^{c+2n'x_1^2b^2}.
\]
To pass to the conjugate of the scalar $\mathcal G(1)$ we must not
assume that $\sigma$ commutes with conjugation. Put
$\sigma'=\mathrm{conj}\circ\sigma\circ\mathrm{conj}$.
Rationality of $G_{\rm cl}$ gives
$\sigma\mathcal G(1)=\det(v_{\sigma'}(w_i)^*G_{\rm cl} v_\sigma(w_j))$.
The determinant Cauchy--Schwarz inequality on the exterior power,
applied to the two wedge products, bounds its absolute value by
the geometric mean of the corresponding two Gram determinants.
Thus
\begin{equation}\label{c3:arch-final}
 |\sigma\mathcal G(1)|\le e^{c+2n'x_1^2b^2}.
\end{equation}

Factor the same nonzero constant, at each $\iota$, as
\[
 \mathcal G(1)=
 (|\gamma(\eta_\iota)|^{2n'}\mathcal G(\eta_\iota))
 (|\gamma(\eta_\iota)|^{-2})^{n'}
 \frac{Z(1)}{Z(\eta_\iota)} .
\]
Apply the height bounds (\ref{c3:local-height}),
(\ref{c3:gammaheight}), (\ref{c3:chainheight}) and then
(\ref{c3:product}) with (\ref{c3:arch-final}):
\[
 0\le c+2\Delta+n'(2x_1^2b^2+b)+d_T(b+V-H).
\]
There are at most $4d_T/b^2$ auxiliary vectors. Since
$4(2x_1^2+b^{-1})\le b$ by the admissible scales,
\[
 n'(2x_1^2b^2+b)\le d_T\,4(2x_1^2+b^{-1})\le d_T b .
\]
The inequality $b\le V$ then absorbs both terms $d_T b$, proving
\eqref{c3:theorem}.
\end{proof}

The family selected in Theorem B enters this comparison through
its original $n$ ordered columns. Its witnesses provide both the
untwisted norm and the uniform local correction in the final
comparison.

\subsection{Elementary scale checks}\label{c3:scales}

The raw scales give
$L^{100}\le x$, $x^{10^6}\le x_1$, $x_1^{10^6}\le b$,
$\log L\ge10^{12}C_P$ and
$C_P\ge64$. In particular $L>e^{200}$,
$x\ge100$, $b\ge x^{10^{12}}$. The source data give
$\sqrt{x}\le m\le2\sqrt{x}$, $x^5\le r\le3x^5$,
$r/2\le k\le r$, $\delta=5r$, $1\le f\le x^{12}$.
Consequently
\[
 B_0=\frac{f x b L^3}{mk}\ge\frac{f bL^3}{6x^{9/2}}
 \ge100f x^{5000},\qquad
 B_0\ge100x_1.
\]
The last inequality follows directly from
$b/x_1\ge x_1^{10^6-1}\ge x^{10^6(10^6-1)}$.
For the leading local defect term, its ratio to $B_0$ is at most
$4(4\cdot10^6e^{180})/L^2<1/2$; the remaining terms of
the regular-family budget satisfy
$(1+2\delta)x^{999}+x^{4000}+\log G_0+x^{30}\le x^{4001}$
using $\log G_0\le x^{32}$ and $x\ge100$.
The same bound on $\log G_0$ gives the twist-error inequality
$(4f+3)r\log G_0+2\le x^{999}$.
For the archimedean multiplier,
$(fmb/k)/B_0=m^2/(xL^3)\le4/L^3$,
hence $fmb/k+3x^{34}\le B_0$, and
$x^{33}+8f\le B_0$ as well. For the one-lift comparison,
divide its required absorption by $kB_0$:
\[
 \frac{3fmbL^2+rB_0+10mx_1+6kB_0}{kB_0}
 \le\frac{12}{L}+2+\frac{40}{L^3}+6<L .
\]
Also $2m^2\delta x_1\le120x^6x_1\le b$,
$4(2x_1^2+b^{-1})\le12x_1^2\le b$ and
$b\le fxbL^4=V$. Finally
$2x^{201}+x^{99}\le x^{202}$ and
$x^{99}+x^{202}b^2\le2x_1^2b^2$.
These estimates use no dimension lower bound and no sign assumption
on the main amount. The group-order and class-number bounds used in these
inequalities are detailed in Appendix~\ref{sec:finite}.
The auxiliary hook and box estimates enter through Aux.

\subsection{Construction hypotheses for Theorem C}\label{c3:frontier}

Theorem C uses the following assertions for the admissible datum.
They identify where its proof uses a construction beyond the
height and determinant arguments.

\begin{description}
\item[I (target and multiplicity).]
The discrete-series spaces on $W_j$ and $Y_0$ satisfy the
cuspidal trace comparison under the hypotheses of
Appendix~\ref{sec:trace}. That appendix deduces both the common
dimension $d_T$ and equality of the determinants of the common
Hecke operator. The integral pairing, coordinate metric and
triviality of the stabilizers used in the determinant calculation
are derived in \S\ref{c3:section}. Injectivity of the chain uses K.

\item[R (arithmetic at the chosen twists).]
For the rank-$r$ main summands of every relevant eigensystem and
the accompanying rank-one lines, assume the Galois representations,
the simultaneous Selmer and Stickelberger polynomial data, and
the tame realization used in \S\ref{c3:twists}. Assume the
crystalline defect is bounded by the stated envelope and the
first-order congruences yield a clearing operator on each generator
with the indicated valuation, including the rank-one pairing
for all test forms. Appendix~\ref{sec:finite} derives the
root-of-unity specialization, explicit defect-envelope sum, geometric
class bounds and resulting $B_0$ budget from its stated arithmetic
inputs. These conclusions supply items 1--3 of \S\ref{c3:twists}.

\item[G (gluing and exchange).]
For each $j$ the period formula with the nonzero $c_i$ and the
products-spanning assertion in \S\ref{c3:single} hold. Its normalized
scalar quotient is the exchange-value ratio
$\rho_\vartheta=|c'_1|^2/|c'_\vartheta|^2$.
Gluing takes the translated integral frame to integral forms at
every place above $\ell$ coprime to the twist. The glued period
is orthogonal to false lifts, and on that subspace the projection
onto lifts is contractive for every conjugate norm. At trivial
twist the gluing inequality has squared constant $e^{2A_{\rm gl}}$ with
$A_{\rm gl}=fmbL^2$; its class, filling and hook estimates are
$(|B|G_0)^2\le e^{A_{\rm gl}}$ and $C_f,h_j\le e^{x^{999}}$.
The three exchange factors are the critical-value products on the
specified integer intervals, with the indicated character, sign and
main-pair weight conditions. These inputs give
\eqref{c3:compression} and \eqref{c3:compressionheight}, and account
for the $k/2$ selected main points and $5m$ remaining points.

\item[F (rank-one values).]
Assume the common-constant logarithmic norm formula, its flat
integral algebraic parts and conjugate denominator bounds in
\S\ref{c3:single}, and the stated rank-one formula for the
Hecke ratio. Use the valuation-attainment estimate of
Appendix~\ref{sec:jets} for the twelve-dimensional block in
the full degree range. Assume the uniform Taylor scaling
$|\sigma T|\le e^{fmb/k}$ and the four-block identity realizes
the specified pairing on a form integral above $\ell$ and
orthogonal to false lifts. The proof of \eqref{c3:test} details the
Fourier averaging, the prime-power formula and the multiplication
of factors at opposed identifications. Together these assumptions
supply the term $-A_0$ in \eqref{c3:main-height}.

\item[K (Rallis and elementary exchanges).]
For the finite sharp data and archimedean edge types of the
single lifts and of the triangle assume the injectivity and
Petersson coset-sum formula in the range
$2\,\mathrm{rank}(\text{source})<\mathrm{rank}(\text{target})$
where it is used. Assume the algebraic realization of each
Gram pairing, including its index and dimension, the local block
mass formula, and the square commutation and adjoint laws for the
triangle, with the $Q_{ij}$ moduli of \S\ref{c3:chain}. The
parabolic index and twist-factor formulas there fix the residual
norm power independently of $\eta$. The determinant calculation
then gives \eqref{c3:chainfactor} for the entire chain.

\item[Aux (auxiliary path).]
Use the shared short-basis bound of \S\ref{A:source} at the grid
level with the box and conductor bounds in \S\ref{c3:auxiliary}.
Assume the onto map $R$ and normalized first-to-last path-period
identity there, with the product modulus \eqref{c3:gamma};
the shifted delta forms test the indicated dual weight.
The common denominator satisfies $1\le N\le e^{x^{99}b^2}$ and
the conjugate period bound is $e^{x^{200}b^2}\|\sigma v\|\sqrt B$.
Assume the normalized dual Weyl pairing has the integral hook
inverse and isometric archimedean duality stated there, with
$hN e^{x^{200}b^2}\le e^{x^{201}b^2}$; the integral good
projection has the squared-norm cost specified there. The
fixed rescaling, local and archimedean dual recovery, and
completion of the original family then give
\eqref{c3:auxbounds}--\eqref{c3:gammaheight}.
\end{description}

On the constructed family, the local Taylor estimate in F follows from the
curve and frame normalizations in \S\ref{sec:unitrow-frame}--\ref{sec:other-witnesses}.
Section~\ref{rankone:main} treats the rank-one character in R.
Section~\ref{pathmath:dual} derives the dual pairing in Aux, and
\S\ref{pathmath:denom} bounds the common denominator under its one-edge and
period inputs. Each statement is for the normalizations specified there.

\section{Comparison on the selected family}\label{sz:comparison}\label{sec:comparison}

\paragraph{The sign of the common height.}
For $0\le y\le a$, the summand in $S(y,\tau)$ vanishes for $\nu>x_1$,
since $\kappa\ge a/x_1$. Every other term is bounded above by
$2y\Lambda(\nu)/\nu$. For every positive integer $h$,
$h/\nu\le2\lfloor h/\nu\rfloor$ when $\nu\le h$. Consequently
\[
 \sum_{\nu\le h}\frac{\Lambda(\nu)}{\nu}
 \le\frac{2}{h}\sum_{\nu\le h}\Lambda(\nu)\lfloor h/\nu\rfloor
 =\frac{2}{h}\log(h!)\le 2\log h.
\]
The equality counts the prime-power divisors of each integer up to $h$.
Thus $0\le S(y,\tau)\le4y\log x_1\le yL$.
Here is a direct lower bound on the archimedean term which suffices for
the cancellation. From $\mathrm{rd}_E\ge1$,
$\log\Omega\le\log(2\pi)+|c_K|\le 2+C_P\le L$ by the first scale
inequality. Since $M\ge b^{20}$, $\log(M/\Omega)\ge19L$.
For an integer $0\le y\le a$, $j(y)\ge19yL$.
Apply this with $y=\kappa(\tau)$ and sum \eqref{sz:height}.
Both $m$ and $k$ are positive by the rank and scale bounds. In particular,
\begin{equation}\label{sz:nonneg}
 H\ge0,\qquad U>0,\quad V>0 .
\end{equation}
This proof needs no estimate on the sign of $\Tot_D$.

\begin{chatAtheorem}[Signed determinant comparison]\label{sz:conclusion}
For an admissible $D$ assume the construction hypotheses of
Theorems A, B and C as specified in
Sections~\ref{A:section}, \ref{B:frontier} and \ref{c3:frontier}.
Then the signed excess of this datum satisfies \eqref{eq:det}.
\end{chatAtheorem}
\begin{proof}
Use the set $t$ and the actual ordered family from \eqref{B:witnesses}.
The vector $y_i=P_{g_i}\phi_i$ has $g_i\in t$ and
$\phi_i\in\mathcal S\subset\mathcal I$, where $\mathcal S$ is the short
set of Theorem~\ref{A:main}. Write $h=H-m\Tot_D+V$.
For a fixed $\sigma$, the matrix in \eqref{sz:covol} is the Gram matrix of
the columns $G_{\rm cl}^{1/2}v_\sigma(y_i)$ in an ordinary Euclidean
complex vector space. Gram--Schmidt gives a determinant equal to the
product of their squared orthogonal residual lengths, hence at most
$\prod_i N_\sigma(y_i)\le e^{nh}$. A dependent family has determinant
zero and satisfies the same inequality. This works for every $\sigma$
even when $h<0$. It proves $\operatorname{CovolumeLe}(y,c)$ with
\begin{equation}\label{sz:upper}
 c=n(H-m\Tot_D+V).
\end{equation}
For the index, the same $\phi_i$ are integral, with the same witnesses
$g_i\in t$. The index estimate \eqref{eq:short-index} therefore gives the budget $2|t|d_SU$.
The bound on $|t|d_S$ in \eqref{B:theorem}, multiplied by $2U>0$, gives
\begin{equation}\label{sz:budget}
 2|t|d_SU\le d_TJ,\qquad J=10U\log(1+2b^2).
\end{equation}
Increasing the budget in \eqref{sz:index} preserves the local matrices,
so $\operatorname{IndexLe}(y,\Delta)$ holds with $\Delta=d_TJ$.

The vectors are independent and $(1-4/b^2)d_T\le n\le d_T$, exactly
the rank assumptions of Theorem C. Its completed basis retains these
original first $n$ columns. Substituting \eqref{sz:upper} and
\eqref{sz:budget} into \eqref{c3:theorem} gives
\[
 d_T(H-3V)\le n(H-m\Tot_D+V)+2d_TJ.
\]
The coefficient of $H$ after rearrangement is $n-d_T\le0$, giving
\begin{equation}\label{sz:signed}
 nm\Tot_D\le(n-d_T)H+(n+3d_T)V+2d_TJ\le d_T(4V+2J).
\end{equation}
Admissibility gives $L\ge128$ and $b\ge3$, whence $4/b^2\le1/2$.
Theorem B gives $d_T>0$, so $0<d_T/2\le n\le d_T\le2n$.
Since $4V+2J\ge0$, division by $n$ bounds $m\Tot_D$ by $2(4V+2J)$.
Finally $1+2b^2\le b^3$ for $b\ge3$, so
$J\le30UL=30V$. It follows that
\[
 m\Tot_D\le2(4+60)V=128faL^4\le faL^5 .
\]
\end{proof}

\begin{corollary}[From the construction estimates to Dirichlet zeros]\label{cor:construction-grh}
Suppose that for each permissible pair $(q_0,d)$ and each $0<\eta<1/100$,
every sufficiently large integer $x$ admits at least one datum from
the construction of Theorem~\ref{thm:p1-data} at $(q_0,d,\eta,x)$
for which the hypotheses of Theorems A, B and C hold simultaneously.
Then every Dirichlet $L$-function is nonzero in $\Re s>1/2$, with the
principal pole at $1$ excepted, and its zeros in $0<\Re s<1$ lie on
$\Re s=1/2$. In particular, every nontrivial zero of $\zeta(s)$ has real
part $1/2$.
\end{corollary}
\begin{proof}
Fix $(q_0,d)$ and $\eta$. At each integer $x$ above the threshold for these
parameters choose one such datum $D$. Theorem~\ref{sz:conclusion} gives
\eqref{eq:det} on this very datum, with the signed sum in
\eqref{eq:total} and the scales of \eqref{eq:p1-data-scale}.
This verifies Hypothesis~\ref{hyp:scale}. Theorem~\ref{thm:main} gives
the stated nonvanishing and the assertion about zeros.
\end{proof}

The pointwise comparison also applies to admissible data outside the
constructed families whenever their construction hypotheses hold.
\subsection{Precision needed for completing the determinant family}\label{sec:precision}
The cost of completing the determinant family depends on the squared norm
of each added column. Keeping the bound from \S\ref{c3:auxiliary}
before its final enlargement gives the following sufficient accuracy for
the two Hilbert--Schmidt estimates.

\begin{proposition}\label{prop:weakerprecision}
Fix the permissible parameters and $0<\eta<1/100$. Work at the scales
of Theorem~\ref{thm:p1-data} for sufficiently large integer $x$.
Suppose $d_X>0$. Retain the shortness and index conclusions of Theorem A,
the restriction bound \eqref{B:codim}, and the pairing and chain-height
estimates used in Theorem C. Assume the auxiliary spanning estimates
of \S\ref{c3:auxiliary}, including the bound $x^{99}+x^{202}b^2$
there for the logarithm of the conjugate squared norm after projection.
Set
\[
                      e=\frac{1}{32b x^{202}} .
\]
In \eqref{B:near} replace $5b^{-4}$ by $e$, and in
\eqref{B:moment} replace $9b^{-4}$ by $e$. If both inequalities hold
for the same normalized map, then $m\Tot_D\le fa(\log b)^5$.
\end{proposition}
\begin{proof}
The projected lower bound with $Z=X$ gives $d_T>0$. If $F$ is a
subspace of the good target and $\alpha=\dim F/d_T$, the trace
average in \S\ref{B:one} gives some translate with projected trace
at most $d_X\sqrt{\alpha(1+e)}\le d_X(1+\alpha+e)/2$.
Applying the lower estimate to the inverse image of $F$ loses
at most $ed_X$, so \eqref{B:spread} holds with $\epsilon=e$.
The greedy argument of \S\ref{B:greedy} stops at codimension at
most $16e d_T$ and uses a set $t$ with
\[
 |t|d_X\le \tfrac52\bigl(1+\log(1/(16e))\bigr)d_T
          \le5\log(1+2b^2)d_T .
\]
Indeed $x^{202}\le b$ eventually because $\log b>x^\eta$.
Restricting to $S$ costs at most $5b^{-8}\log(1+2b^2)d_T$
more dimensions by \eqref{B:codim}. This cost is at most
$d_T/(2b x^{202})$ for large $b$, since $x^{202}\le b$.
The short-vector basis extraction therefore produces an independent
family of size $n$ with $n'=d_T-n\le d_T/(b x^{202})$.

Apply the block determinant argument of \S\ref{c3:compare} to this
family before its last rank estimate. The original columns contribute
$c+2\Delta$. Each added column costs at most $x^{99}+x^{202}b^2$
at infinity by \S\ref{c3:auxiliary} and at most $b$ at the finite
places by \eqref{c3:gammaheight}. Together with \eqref{c3:chainheight}
the height inequality reads
\[
 0\le c+2\Delta+n'(x^{99}+x^{202}b^2+b)+d_T(b+V-H).
\]
The cost of the added columns divided by $d_T$ is at most
$b+x^{-103}/b+x^{-202}\le V$ eventually, since
$V=f x b(\log b)^4$. Also $b\le V$, so $d_T(H-3V)\le c+2\Delta$.
The cardinality bound for $t$ gives the index budget \eqref{sz:budget},
and $n\ge d_T/2$. Substitution of \eqref{sz:upper} and the sign
$H\ge0$ prove \eqref{eq:det} by \eqref{sz:signed}.
\end{proof}

At fixed $\eta$, the bounds $\log b>x^\eta$ and
$x_1=\lfloor\exp(x^{\eta/6})\rfloor$ give $x^{202}=o(x_1^2)$.
Thus the accuracy $(32b x_1^2)^{-1}$ also suffices. More generally,
put $L=\log b$. A normalized error $O(b^{-1}L^{-A})$ in each
estimate suffices for any $A>202/\eta$, since $x^{202}<L^{202/\eta}$.
In particular $O(b^{-1-\delta})$ suffices for every fixed $\delta>0$.
These implications allow constants depending on the fixed parameters
and exponents, with a common bound over the selected data as $x$ varies.


\clearpage
\appendix
\part*{Appendices}
\addcontentsline{toc}{part}{Appendices}
\section{Frobenius couplings, the cost of changing directions, and sharp jets}\label{sec:jets}

Fix a prime $p$ and a field embedding $j:\mathbb C\hookrightarrow\mathbb C_p$ for the jet calculation. Here $\mathbb C_p$ is the completion of $\overline{\mathbb Q}_p$, $v(p)=1$, and $\mathcal O$ denotes its valuation ring. To apply the calculation to the index or height at $(p,\iota)$, use
\[
 j:\mathbb C\xrightarrow{\ \iota\ }\overline{\mathbb Q}_p
       \lhook\joinrel\longrightarrow\mathbb C_p.
\]
The last arrow is isometric. Thus $|j(x)|=|\iota(x)|$ for every complex coefficient $x$, and the lattice coefficient ring is $j^{-1}(\mathcal O)\subset\mathbb C$. The second test form in the height comparison uses $j\circ\mathrm{conj}$ with the same twist. For the nonsplit calculation put
\[
 (\pi,q,i_0)=
 \begin{cases}
 (p,p^2,p+1)&p\text{ inert in } K,\\
 (\sqrt{-d},p,2)&p=d,
 \end{cases}
 \qquad \rho=\frac{v(\pi)}{q-1},\quad
 \gamma=\frac1{p-1}-\rho .
\]
These two local constants differ from the total prime gain in Theorem A.
Write $\tau_e=u(-p^e)\tau$ and let $I^\circ$ consist of the odd exponents at inert primes and all nonnegative exponents at $p=d$.

\subsection{From the CM curve to a matrix of variables}\label{jets:cm}
Let $\theta(0)=1$ be the translated CM theta function and $g=-\log\theta$.
We use the following CM inputs for the normalization of this curve. In the local
series ring, $T=z+\cdots$ is a generator whose degrees are $1\bmod(q-1)$.
The endomorphism $T(\pi z)$ is an integral series in $T$ with linear term
$\pi T$, a unit coefficient $b_1$ at $T^q$, and remaining degrees
$1\bmod(q-1)$. Moreover,
\[
 qg(z)-g(\pi z)\in\pi\mathcal O[[T]],\qquad
 \exp\!\big(\zeta^{-(p+1)}g(\zeta z)-g(z)\big)\in\mathcal O[[T]]
 \quad(\zeta^{q-1}=1).
\]
These inputs concern the integral model and the theta function at its chosen
two-torsion point. The source concordance gives their statements.

Choose $\varpi^{q-1}=-\pi/b_1$ and set $Y(z)=T(\varpi z)/\varpi$.
In the endomorphism series the term of degree $1+j(q-1)$, divided by
$\pi\varpi$ after rescaling, has a factor $\varpi^{j(q-1)}/\pi$.
For $j\ge2$ this belongs to $\pi\mathcal O$. Hence
\[
 \pi^{-1}Y(\pi z)=Y-Y^q+\pi E(Y),\qquad E\in z^2\mathcal O[[z]] .
\]
For $n\ge2$ the coefficient of degree $n$ determines $[z^n]Y$ from lower coefficients by division by $1-\pi^{n-1}$, a unit. Induction proves integrality. Reducing the identity gives $Y-Y^q=z$, hence
$Y\equiv\sum_{j\ge0}z^{q^j}\pmod{\pi}$ by iteration in increasing degree.

Let $g^\flat$ select the degrees $p+1\bmod(q-1)$ of $g$, and put
$s_0=\varpi^{p+1}/p$, $\Xi=-g^\flat(\varpi z)/s_0$.
Averaging over $\zeta^{q-1}=1$ expresses
$\exp(g^\flat-g)^{q-1}$ as the product of the integral exponentials above.
Since $q-1$ is a unit and the constant term is one, successive coefficient comparison gives $\exp(g^\flat-g)\in\mathcal O[[T]]$.
Rescaling in $z$ replaces $T$ by $\varpi Y$, so
$\theta(\varpi z)=W\exp(s_0\Xi)$ with $W$ integral.

Project the logarithmic identity to degrees $i_0\bmod(q-1)$.
The projection of the term in $\mathcal O[[T]]$ is $aT^{i_0}+T^{i_0+q-1}H(T)$ with $a,H$ integral. After scaling and division by $-qs_0$, the term of order $i_0$ is multiplied by the unit $-\pi\varpi^{i_0}/(qs_0)$ and the tail has the extra factor $\varpi^{q-1}$. Thus for some $\lambda\in\mathcal O$,
\[
 (1-\pi^n/q)[z^n]\Xi-\lambda[z^n]Y^{i_0}\in\pi\mathcal O
 \quad(n\ge0).
\]
Below $i_0$ the exponent is zero. For $n>2$, $\pi^n/q\in\pi\mathcal O$.
At $p=d,n=2$, the multiplier is $2$ because $\pi^2=-p$.
Consequently $\Xi$ is integral and its residue is
$\lambda Y^{p+1}$ in the inert case and $\lambda(Y^2-z^2/2)$ in the ramified case.
Here $p$ is odd. In the inert case $Y^{p+1}$ is the product of
$\sum_{j\ {\rm even}}z^{p^j}$ and $\sum_{j\ {\rm odd}}z^{p^j}$ modulo $\pi$.
An even and an odd index have difference $e$ odd and smaller index $j$.
Their monomial is $(z^{p^e}z)^{p^j}$. In the ramified case order the two indices in $Y^2$.
Off-diagonal terms have coefficient $2\lambda$ and diagonal terms have coefficient $\lambda$, except in degree two where the coefficient is $\lambda/2$.
This gives the advertised exponents $I^\circ$, including $e=0$ at $p=d$.

For sharpness one also uses that $\lambda$ is a unit in the inert case. Here is its local test.
Put $x=\wp(z+t)-\wp(t)=z^2+\cdots$ with the normalization of the chosen model,
$p=2m_0+1$ and $h_n=[z^{p-1}]x^n$. The chosen Weierstrass equation gives
$(x')^2=4x^3+a_1^2x^2+4x$ and $x''=6x^2+a_1^2x+2$.
The coefficient of $z^{p-1}$ in $(x^n x')'$ is $p[z^p](x^n x')$.
The latter coefficient before multiplication by $p$ is integral by the integral CM series bound in degrees $<q$. Expanding the derivative gives
\[
 (4n+6)h_{n+2}+(n+1)a_1^2h_{n+1}+(4n+2)h_n\equiv0
 \pmod{\mathfrak m_{\mathcal O}}.
\]
Now $h_0=0$, $h_{m_0}=1$, and $4n+6$ is a unit for $0\le n\le m_0-2$.
Thus $h_1$ is a unit. The identity $g''=g''(0)+x$ gives
$[z^{p+1}]\Xi=-h_1/(p+1)$, so $\lambda$ has unit residue.

For the level-one moved datum the remaining input is a finite resolution
$z=\sum_j m_j(u_j^\dagger z)u_j$ in the chosen embeddings, integral above $p$,
and the factorization of the Taylor series as
$H\prod_j\theta(u_j^\dagger z)^{m_j}$. Here $H$ is integral in the translated curve parameters evaluated on these forms.
The construction in the source first expands the integral vacuum in finitely many Heisenberg translates of one vector. Each rational translated vector has a denominator prime to $p$ and admits the factorization with any multiple of that denominator and the frame denominator. Taking the product of the finitely many denominators gives a common resolution. In its new vectors the denominator is divided out and its square enters the multiplicities, preserving the resolution identity. The theta expansion and the integral quotient for each translated factor enter as hypotheses.

Write $L_j=u_j^\dagger z$. Frobenius on the coefficients gives
$\overline{\tau(u)}^{p^e}\equiv\tau_e(u)\pmod\pi$ for $u$ integral above $p$.
Indeed the cyclotomic generator transforms by the exponent $-p^e$ after conjugation. On the quadratic generator the inert Frobenius and conjugation cancel for $e$ odd, and the ramified generator reduces to zero. This check extends over the integers locally since $p$ is prime to the cyclotomic conductor and to $2$. With coordinate index $a$ the resolution identity therefore gives
\[
 \sum_j m_j L_j^{p^e} L_j
 \equiv\sum_{\tau,a}z_{\tau a}^{p^e}z_{\tau_e a}
 =:\Gamma_e\pmod{\pi}.
\]
Frobenius powers are additive in $\mathcal O/\pi$. Summing the curve exponents with multiplicity gives, modulo $\pi$, an integral combination $\sum_{e,j}\lambda_{ej}\Gamma_e^{p^j}$. Only finitely many terms occur in each degree. The change in the exponent after multiplication by $s_0$ has coefficients of valuation at least $v(s_0\pi)=1/(p-1)$ and zero constant term. Its exponential is integral by the factorial formula. Absorb it and the rescaled $H$ into the integral factor. This proves the multivariate normal form in Section~\ref{A:jets}. Any further integral scalar rescaling changes each coefficient by its homogeneous power. The period factor is a $p$-unit since $p$ does not divide the conductor.

\subsection{The determinant cost for a charged filling}\label{jets:count}
Let $N$ be the number of target rows and $B$ the number of boxes in columns of height $h<N$. Fix the marked places and the nonnegative integral charges $w_{\tau h r}$ for output row $r$. For each $\tau$ let $t_\tau$ count the unmarked height-one columns, $c_\tau$ the remaining columns of heights $2,\ldots,N-1$, and $b_\tau$ their boxes.
Fix an exponent $e\in I^\circ$, write $Q=p^e$, $\tau'=\tau_e$, and set
\[
 K_e(\tau)=\min(t_\tau/Q,t_{\tau'}+c_{\tau'})
     +c_\tau/Q+\min(b_\tau/Q,b_{\tau'}).
\]
We construct matrices on the $h$ given directions at $(\tau,h)$ which make the pairing $\sum_a \beta_a^Q\beta'_a$ divisible by $\pi$. A pair of heights $(h,h')$ is exempt precisely if $h=1$ and either $\tau$ is marked or both $h'=1$ and $\tau'$ is marked. For every exempt pair and each $r$ the hypothesis is $w_{\tau 1 r}>0$ or $w_{\tau' h'r}>0$. Charges at heights absent from the diagram can be set to one.

A coordinate condition on the powered side costs $v(\pi)/Q$, since it can be imposed modulo a $Q$-th root of $\pi$. A coordinate condition on the other side costs $v(\pi)$. Pairing against a fixed first direction on the powered side gives an ordinary linear condition on the other side. For a variable powered side and fixed direction $y'_0$, impose the linear condition
$\sum_a\beta_a(y'_{0,a})^{1/Q}\equiv0\pmod{\pi^{1/Q}}$.
All the roots may be chosen in $\mathcal O$. Raising to the $Q$th power gives the desired pairing congruence modulo $\pi$.

For $h\ne1$ impose on the powered side a pairing condition against the first direction of the other side. This covers $h'=1$ and contributes $c_\tau/Q$.
For $h=1$, choose the cheaper of imposing the single coordinate condition at each unmarked source column and imposing a pairing with the first source direction at each nonexempt partner column. This covers all required pairs with $h=1$ and costs the first minimum. Finally for $h,h'\ne1$ impose all row-coordinate conditions on the cheaper side. These conditions account for the second minimum.

All these conditions can be imposed simultaneously even if a place is its own partner. Each is a linear condition on the row coefficients of a change matrix against an \emph{original} direction or a fixed coordinate. For a single form, select a row whose value has least valuation, subtract integral multiples of it from the other rows to annul their values, then multiply the selected row by the prescribed modulus. The resulting determinant costs at most the valuation of that modulus. A zero form needs no change. Repeat on the current rows for every condition. Integral combinations preserve each earlier divisibility condition. On a height-one side the new direction is just a multiple of the original one, which verifies the imposed first-direction pairings. Multilinearity counts each change determinant once per column of the relevant height. Restrict to $e\le\lfloor\log_p\max(B,1)\rfloor$ since the order of $\Gamma_e$ is at least $p^e+1$. We have proved
\[
 \sum_{\tau,h<N}\kappa(\tau,h)v(\det g_{\tau h})
 \le v(\pi)\sum_{e\ \mathrm{retained}}\sum_\tau K_e(\tau).
\]

Introduce variables $s_{\tau h k r}$ for $k<h$ and $r$ an output row.
In the normal form substitute $z_{\tau r a}=\sum_{h,k}p^{w_{\tau h r}}
\beta_{\tau h k a}s_{\tau h k r}$. Every relevant coupling is divisible by $\pi$.
For a nonexempt pair this follows from the row congruences and Frobenius additivity modulo $\pi$. For an exempt pair the charges contribute a factor $p$. Thus $F(\varpi z(s))$ is integral through degree $B$, by $v(s_0\pi)=1/(p-1)$. The product of column determinants in the $s$-variables defines a homogeneous differential operator $Q_B(\partial)$ of degree $B$ with integral coefficients. On a monomial of degree multi-index $\alpha$ its value at zero uses $\prod_i\alpha_i!$. There are at most $fN^3$ variables, so Legendre's formula gives
\[
 v\big(Q_B(\partial)F(\varpi z(s))|_{s=0}\big)
 \ge \frac{B}{p-1}-fN^3(1+\log_p\max(B,1)).
\]
The chain rule and alternation in the $h$ directions yield the factor $\varpi^B$, one determinant for each column, and $p^{w(i')}$, where $w(i')$ is the sum of charges on the filling $i'$. After division the original minor derivative $D_{i',y}F$ therefore satisfies
\[
 v(D_{i',y}F)+w(i')\ge
 \gamma B-v(\pi)\sum_{e\ \mathrm{retained},\tau}K_e(\tau)
       -fN^3(1+\log_p\max(B,1)).
\]
Each cost is nonnegative so one can include all $e\in I^\circ$ on the right. If $B^\circ=\sum_\tau(t_\tau+b_\tau)$ counts the unerased boxes, then
$\gamma=v(\pi)\sum_{e\in I^\circ,e>0}p^{-e}$ by the geometric series.
Termwise $u-\min(u,v)=(u-v)_+$ gives the exact identity
\begin{align*}
 \gamma B-v(\pi)\sum_{e\in I^\circ,\tau}K_e(\tau)
 &=\gamma(B-B^\circ)-\mathbf1_{p=d}\tfrac12\sum_\tau K_0(\tau)\\
 &\quad+v(\pi)\sum_{\substack{e\in I^\circ\\e>0}}\sum_\tau
 \left[(t_\tau/p^e-t_{\tau_e}-c_{\tau_e})_+
 +(b_\tau/p^e-b_{\tau_e})_+-c_\tau/p^e\right].
\end{align*}
The separate ramified subtraction records the unpowered pairing at $e=0$.
This is the charged-minor bound used in the comparison with the prime sum.

For the Schur entry fix an output filling $i'$ and let $a_i$ be the jets of the entries in that row. Evaluating the product of column minors on $zy$ gives
$D_{i',y}F=\sum_i a_i b_i(y)$ by the action in the standard bitableau basis. Here each $b_i$ is a basis minor product times a common determinant power. The standard basis theorem~\cite[Proposition~4.1 and Theorem~4.11]{brunsvetter} gives independence of these polynomials over the residue field, and multiplication by the common nonzero determinant polynomial preserves it. If some $a_i\ne0$, divide by a coefficient of maximal norm and reduce. The resulting polynomial is nonzero over the infinite residue field, so some integral $y$ gives $|D_{i',y}F|=\max_i|a_i|$. This proves the same estimate for every entry of this row. The all-zero case is immediate.

\subsection{A witness for the sharpness bound}\label{jets:sharp}
The preceding estimate fixes the filling and works for every derivative. The sharpness step of Theorem C constructs a coefficient whose valuation is bounded \emph{above}. The inert calculation uses the unit coefficients $\lambda_{e0}$ of the curve normal form before any nonunit rescaling. The additional inputs concern the twelve-copy least-valuation estimate and its realization by integral theta functions. For clarity, let $\ell_\tau$ lie between the weight and the weight plus $x^{42}$, and put
\[
 h=\lceil6b/k\rceil,\qquad B_\tau=\prod_{i=1}^h ((\ell_\tau-h)_++i),
 \quad S_p=\mathrm{localPrimeSumTotal}(p).
\]
For a tight frame of length $M$, with multiplicities at least $6fb$ and a frame form having unit coefficients at the given prime, the required estimate is the existence of twelve coefficient multi-indices of sum $\ell$ and twelve series $H\Theta$ such that their coefficient product has valuation
\[
 r_p\ \text{or less},\qquad
 r_p=S_p+2\sum_\tau v(B_\tau)-\sum_\tau v(\ell_\tau!).
\]
Here $\Theta$ is the product on the frame and $H$ belongs to the ring generated by its translated curve parameters.
\paragraph{Why an inert coupling yields a small valuation.}
Let $p>2$ be inert and take selected starts $\tau\in\Phi_0$ with odd exponents $e_\tau$.
Require the following uniqueness property. Among pairs $(\sigma,a)$ with
$\sigma\in\Phi_0$, $a$ odd, $a\le e_\tau$ and $\sigma_a=\tau_{e_\tau}$, the only pair is $(\tau,e_\tau)$. Take $j_\tau$ copies of this pairing,
put $J_0=\sum j_\tau$, and write $\nu',\nu''$ for the degrees of its powered and unpowered variables. Thus $|\nu''|=J_0$ and $|\nu'|=\sum p^{e_\tau}j_\tau$.
Use the normal form $F(z)=\Theta(\varpi z)=W\exp(s_0\Omega)$ of the unmodified frame theta product. Here $W(0)=1$ and $v(s_0)=1/(p-1)-1<0$.

There are six integral scalars $\delta_c$ with $\delta_0=1$ and
$\sum\delta_c^2=0$. Indeed any sequence of positive integers tending to $-1$ in $\mathbb Z_p$ can be represented by four integer squares. Compactness in $\mathbb Z_p^4$ yields four squares representing $-1$ in the limit. Append $1$ and $0$. We use the CM product identity that $\prod_c\theta(\delta_c z)$ belongs to the integral curve series ring under this square-sum condition. Its constant term is one. Substitution of the frame forms, multiplication with multiplicities and rescaling show that $U(z)=\prod_c F(\delta_c z)$ has integral coefficients and constant term one. Introduce independent variable blocks $z',z''$ and form
\[
 G=F(z'+z'')\prod_{c\ne0}F(\delta_c z').
\]
Only the first factor depends on $z''$. When differentiating $g_0G$ in an end variable $z''_{\tau_{e_\tau}}$, with $g_0$ integral, divide the result by $s_0G$. The new integral factor reduces to $g_0\,\partial_{\tau_{e_\tau}}\Omega(z'+z'')$ in the residue field: the derivative of $g_0$ and the logarithmic derivative of $W$ both carry $s_0^{-1}$ of positive valuation. Iterating $J_0$ times and then putting $z''=0$ gives $s_0^{J_0}gU$, where modulo the maximal ideal $g$ is the product of these first derivatives with the prescribed multiplicities.

Set the variables outside $\Phi_0$ to zero. In residue characteristic $p$, differentiating $\Gamma_a^{p^j}$ kills the terms with $j>0$ and differentiating the powered side of $\Gamma_a$ gives zero for $a$ odd. The derivative at an end has unique monomial of lowest total degree
$\lambda_{e_\tau,0}z_\tau^{p^{e_\tau}}$ by the uniqueness assumption.
These coefficients are units by the inert calculation above. Thus the coefficient of $z'^{\nu'}$ in $gU$ is a unit. Terms discarded by restricting to $\Phi_0$ cannot contribute to this multi-index. It follows that
\[
 v\big([z'^{\nu'}z''^{\nu''}]G\big)
 =J_0 v(s_0)-v(\nu''!),\qquad
 v(\nu''!)=\sum_\tau v(\nu''_\tau!).
\]
Reversing the initial $\varpi$-scaling subtracts $\rho|\nu'+\nu''|$ from the valuation. The substitutions were integral, so the ultrametric inequality provides a coefficient of the original six-copy product with at most that valuation. In particular the valuation is at most $-J_0$ since
$|\nu'+\nu''|\ge(p+1)J_0$ and $\rho(p+1)=1/(p-1)$.
It is also at most
$\gamma|\nu'+\nu''|-J_0-v(\nu'!)-v(\nu''!)$ by
$v(\nu'!)\le |\nu'|/(p-1)$.
For the prescribed degree vector \(\ell\), the
\srcline{New/Derivatives/Sharpness/Cases}{672}{twelve-copy estimate}
uses a tight frame with a member having a unit coefficient at every place.
At a split prime the local prime sum vanishes. The first six copies may be
taken in degree zero and the unit coefficient estimate on the other six
supplies the factorial bound. At a nonsplit prime \(p>k\) the partner map
is a permutation. The proof uses exponent one and caps the multiplicity
at each source by both its available degree after the frame loss, divided
by \(p\) and rounded down, and the degree at its partner. The
\srcline{New/Derivatives/Sharpness/Cases}{346}{factorial estimate}
accounts for the residual term.

At a nonsplit prime \(p\le k\), use a near exponent on each eligible
place and let \(\nu',\nu''\) be the degrees at the starts and ends.
The remaining six copies supply \(\mu=\ell-\nu'-\nu''\).
The factorial correction from these three vectors is
\[
 \sum_\tau v_p\!\left(\frac{\ell_\tau!}{\nu'_\tau!\nu''_\tau!\mu_\tau!}\right)
 =\sum_\tau
 \frac{s_p(\nu'_\tau)+s_p(\nu''_\tau)+s_p(\mu_\tau)-s_p(\ell_\tau)}{p-1},
\]
where \(s_p\) is the base-\(p\) digit sum. This follows from Legendre's
formula even when one place occurs at both ends of pairs. The
\srcline{New/Derivatives/Sharpness/Cases}{377}{carry bound} and
\srcline{New/Derivatives/Sharpness/Cases}{403}{numerical comparison}
pay for this loss together with the cost of the residual degrees.
At the ramified prime \(p=d\) the near-pair set is empty. The first six
copies then have degree zero. Sections~\ref{sec:unitrow-frame}--\ref{sec:other-witnesses} give the coefficient estimate on the constructed scales, under their indicated local curve and period normalizations.

\paragraph{From the twelve coefficients to one indicator.}
Expanding $H$ and taking a term of maximal norm reduces to monomials in those parameters, of individual degree at most $|\ell|$ since the parameters have zero constant term. In disjoint variable blocks the product coefficient is exactly the product of the selected coefficients. The inequality $3|\ell|\le6fb$ ensures that each local factor $y^a\theta^c=(y\theta^3)^a\theta^{c-3a}$ has nonnegative exponents. Realization of these monomial sections gives an adelic vector with integral finite values and the signed periods of the modulus $(\pi_M)$.

Fix once and for all a nonzero value $c_0$ of the twelve-dimensional integral vacuum. In the finite Fourier expansion of the vector, inversion over the residue classes of $(\pi_M)$ contributes the factor $M^{12f}$. Multiplying the evaluated Fourier identity by the corresponding vacuum value gives integral coefficients of norm at most one, since the characters have prime-adic norm one. Applying the ultrametric inequality and taking a maximizing class $y_0$ gives a coefficient of the vacuum-scaled indicator series with valuation at most $r_p+v(M^{12f}c_0)$. Consequently the single nonzero scalar
\[
 t=M^{12f}c_0\prod_\tau B_\tau^2
\]
works at every prime and identification. Multiplication of the coefficient by $\prod_\tau\ell_\tau!/t$ leaves valuation at most $S_p$. The chosen multi-index and residue class may depend on the prime, identification and modulus. The conjugate bound on $c_0$, together with $B_\tau\le(2b+h)^h$, gives
\[
 \log|\sigma t|\le12f\log M+2fh\log(2b+h)
             +12\,\mathrm{constantBound}\,q_1^2\le fmb/k
\]
under the source's admissible bounds.

The following subsections construct the frame and the twelve coefficients on the scales of Theorem~\ref{thm:p1-data}. The theta realization and the integral Fourier expansion remain necessary for the application to adelic forms. The normal form alone supplies none of the twelve coefficient choices. In particular the charged-filling count cannot be used to discharge the sharpness assumption of Theorem C.

\subsection{Construction of the integral frame with a unit row}
\label{sec:unitrow-frame}

Let $Q=q_1$, $f=\varphi(Q)$, $\zeta=\zeta_Q$,
$C(X)=\Phi_Q(X)$, $\Delta=C'(\zeta)$ and
$\beta=(1+\sqrt{-d})/2$. The integers of $E=K(\zeta)$
are $\mathcal O_K[\zeta]$ since the two cyclotomic conductors
$d$ and $Q$ are coprime (equivalently, the discriminants of
$K$ and $\mathbb Q(\zeta)$ are coprime). The relative different
is $(\Delta)$ by the derivative formula for a power basis.
Differentiating $X^Q-1=C(X)D(X)$ at $\zeta$ shows $\Delta\mid Q$.
Use the fractional ideal $\mathfrak a$ with local factor
$\Delta^{-1}\mathcal O_{E,v}$ at $v\mid\beta$ and
$\mathcal O_{E,v}$ elsewhere. This describes the lattice of the vacuum
in terms of its finite normalization. It is stable under $G=\mathrm{Gal}(E/K)$.
Here the embeddings in $G$ extend the selected embedding of $K$.
Every prime divisor of $Q$ divides $\beta\bar\beta=(d+1)/4$:
for an odd prime this follows from $d\equiv q_0/2-1\pmod{q_0}$,
and for $2$ from $8\mid d+1$. Since $\beta+\bar\beta=1$,
the ideals $(\beta),(\bar\beta)$ are coprime.
In particular $\mathfrak a\bar{\mathfrak a}$ is the inverse different.
For $v\in\mathfrak a$ the number
$N_v=\sum_{\tau\in G}|\tau v|^2=\operatorname{Tr}_{E/K}(v\bar v)$
is therefore integral in $K$. Being real, it is a nonnegative integer.

\paragraph{A second frame length.}
Call an indexed family $v$ of lattice vectors tight of length $H$ when
$\sum_i\sigma(v_i)\overline{\tau(v_i)}=H\delta_{\sigma\tau}$.
The $Q$ powers of $\zeta$ (or of $\zeta^{-1}$) are tight of
length $Q$ by geometric summation. For any $v\in\mathfrak a$, its
indexed family $\zeta^j\omega(v)$, $0\le j<Q$, $\omega\in G$,
is tight of length $QN_v$: the sum in $j$ kills off-diagonal
entries, and the sum in $\omega$ makes every diagonal entry equal.

Take $u=\lceil\log_2 Q\rceil$ and define $b_i,e_i$ by
\[
 C(X)/(X-\zeta)=\sum_{i=0}^{f-1} b_i X^i,\qquad
 e_i=\bar\beta^u b_i/\Delta .
\]
The $b_i$ are integral. At a place away from $\beta$ the quotient
$\bar\beta^u/\Delta$ is integral: $Q\mid(\beta\bar\beta)^u$
by prime factorization, $\Delta\mid Q$, and $\beta$ is a local unit.
At places over $\beta$ the factor $\Delta$ clears the denominator.
Thus all $e_i$ belong to $\mathfrak a$.
Lagrange interpolation gives, for every $\sigma,\tau\in G$,
\[
 \sum_{i=0}^{f-1}\sigma(e_i)\,\overline{\tau(\zeta^{-i})}
   =\bar\beta^u\delta_{\sigma\tau}.
\]
Indeed this is the polynomial $\sigma(C(X)/(X-\zeta))/\sigma(\Delta)$
evaluated at $X=\tau(\zeta)$, times $\bar\beta^u$.

Start with the powers $\zeta^{-j}$ and, for each $i$, the entire indexed
family $\zeta^j\omega(e_i)$. For $0\le i<f$ replace the indexed
pair $\zeta^{-i}, e_i$ by its sum, using the $j=0$, $\omega=1$
copy of $e_i$. The change in the Gram matrix consists of the two
cross terms in the interpolation identity. The new frame has length
\[
 H=Q\left(1+\sum_{i<f}N_{e_i}\right)+\beta^u+\bar\beta^u .
\]
Modulo any prime ideal above a divisor of $Q$, one of
$\beta,\bar\beta$ is zero and the other one, so $H\equiv1$ there.
The expression $H$ is a rational integer, and its Gram matrix is nonnegative.
Thus $H>0$ and $(H,Q)=1$. The indexed member
$y_*=\zeta e_{f-1}=\zeta\bar\beta^u/\Delta$ survives the replacement.

\paragraph{Archimedean size and filling the prescribed length.}
For fixed $q_0,d$, $H$ grows only polynomially in $Q$.
Indeed $Q=q_0 A$ has the same prime divisors as $q_0$, whence
$C(X)=\Phi_{q_0}(X^A)$ (compare primitive roots and degrees).
Its coefficients are bounded in terms of $q_0$, so synthetic division
at any conjugate of $\zeta$ gives $|\tau(b_i)|\ll_{q_0}Q$.
Differentiation yields
$|\tau(\Delta)|=A|\Phi_{q_0}'(\tau(\zeta)^A)|\gg_{q_0}A$.
Also $|\beta|^u\ll_d Q^{O_d(1)}$, which proves the assertion.
On the constructed scales $Q$ is polynomial in $x$ and $\log b>x^\eta$.
Thus $H,Q,f\le b$ eventually. Put $R=\lceil6fb\rceil$.
The level $M\ge b^{20}$ satisfies $M-R(H+Q)\ge QH$ for large $x$.
Choose $j$ with $0\le j<Q$ and $jH\equiv M-R(H+Q)\pmod Q$.
Writing $M-R(H+Q)=jH+iQ$ gives $i\ge0$.
Use the length-$H$ family with multiplicity $j+R$ and the length-$Q$
powers with multiplicity $i+R$. This frame has length exactly $M$,
every multiplicity is at least $6fb$, and it contains both $1$ and $y_*$.

\paragraph{The unit row.}
Let $s_{\rm per}$ denote the period multiplier in the frame forms
$L_y(z)=s_{\rm per}\sum_\tau\overline{\tau(y)}z_\tau$. Its local
normalization is as follows. For the embedding
$\phi_\tau(y)=\iota(\overline{\tau(y)})$ and its place $v_\tau$,
$|\iota(s_{\rm per})|=|\phi_\tau(\Delta)|$ if $v_\tau\mid\beta$, and
$|\iota(s_{\rm per})|=1$ otherwise. These finite absolute values of the
period multiplier are inputs from the normalization of the vacuum.
All $v_\tau$ restrict to the same place on $K$. When they lie away
from $\beta$, take $y=1$; otherwise $\bar\beta$ is a unit there,
so $|\phi_\tau(y_*)|=|\phi_\tau(\Delta)|^{-1}$ at every $\tau$.
This gives one row all of whose coefficients are units, for each
finite prime and embedding. The lattice condition gives integrality
of the other rows. At an inert prime, $p\nmid Q M$ (the level primes
split), so all frame vectors and their conjugates are integral there.
From tightness,
$\sum_i c_i\tau(v_i)L_{v_i}(z)=M s_{\rm per} z_\tau$; the right-hand
scalar is a unit. This supplies the integral resolution used in the
large inert case. Dividing the frame by $\bar\pi_M$ for
$\pi_M\bar\pi_M=M$ also gives the identity resolution for the
normalized inert coupling calculation, with unit rescaling.


\subsection{The Taylor witness at the ramified prime}
\label{sec:ramified-witness}

Write $p=d$, $v(p)=1$, $\mathcal O=\mathcal O_{\mathbb C_p}$,
$\pi=\sqrt{-d}$ and $\rho=\gamma=1/(2(p-1))$.
We use the following local curve data. The translated coordinate $y(z)=z+O(z^2)$
has integral formal addition law, i.e.$y(z+w)\in\mathcal O[[y(z),y(w)]]$. There is an element
$T\in\mathcal O[[y]]$ satisfying the normalized endomorphism identity
of Appendix~A.1. Finally for $\delta_i\in\mathbb Z_p$ with
$\sum_i\delta_i^2=0$ we use the CM theta identity
$\prod_i\theta(\delta_i z)\in\mathcal O[[y]]$, with $\theta(0)=1$.
These are statements about one curve and its normalization.
We give the deduction of a witness at every required multidegree.

\paragraph{Derivatives in one variable.}
Let $m_n$ be the smallest valuation of $g^{(n)}(0)$ as
$g$ ranges over $\mathcal O[[y]]$, with $v(0)=+\infty$.
Only $1,y,\ldots,y^n$ contribute to an $n$th derivative, so the minimum
is attained by one of these powers and is finite since $y^n=z^n+O(z^{n+1})$.
In particular $m_0=0$. For $g\in\mathcal O[[y]]$, the addition law expands
$g(z+w)$ in the products $y(z)^i y(w)^j$ with integral coefficients.
Taking the mixed derivative of orders $a,b$ gives
$m_{a+b}\ge m_a+m_b$.
Put $N=p^j$. The rescaled parameter
$Y(z)=T(\varpi z)/\varpi$, $v(\varpi)=\rho$, satisfies
$Y\equiv\sum_{j\ge0} z^{p^j}\pmod{\pi}$ by the endomorphism recursion in
Appendix~A.1. Hence $v([z^N]T)=-(N-1)\rho$.
Legendre's formula gives $v(N!)=(N-1)/(p-1)$ and therefore
$m_N\le (N-1)\gamma$. Fixing $n>0$, write $N=an+r$ with $0\le r<n$.
Superadditivity gives
\[
 a m_n+\min_{0\le i<n}m_i\le (N-1)\gamma.
\]
Letting $j$ tend to infinity gives $m_n\le n\gamma$ for every $n$.
This argument yields a monomial $g=y^i$, $i\le n$, realizing the bound.

\paragraph{Recovering a multidegree.}
Let $L_j(z)$ be the frame's local linear forms, with
$L_*(z)=\sum_\tau u_\tau z_\tau$ one member such that $v(u_\tau)=0$
for all $\tau$. Write $\Theta=\prod_j\theta(L_j)^{c_j}$, and
$\mathcal A$ for the image of the integral formal power series ring
under substitution of the $y(L_j)$. Choose six $p$-adic integers
$\delta_0=1,\ldots,\delta_5$ with $\sum\delta_i^2=0$ as in Appendix~A.3.
Then $U(z)=\prod_{i=0}^5\Theta(\delta_i z)$ belongs to $\mathcal A$ and
has constant term one. Taking the formal inverse in the generator
variables gives $U^{-1}\in\mathcal A$.
For a multidegree $\ell$, set $n=|\ell|$ and choose $g$ as above.
The coefficient of $z^\ell$ in $g(L_*)$ is
\[
 g^{(n)}(0)\prod_\tau u_\tau^{\ell_\tau}/\prod_\tau\ell_\tau!.
\]
Its valuation is at most $\gamma n-\sum_\tau v(\ell_\tau!)$.
Take $H_0=g(L_*)U^{-1}$ and $H_i=1$ for $1\le i\le5$. Then
$\prod_i(H_i\Theta)(\delta_i z)=g(L_*(z))$.
Its indicated coefficient is a finite sum over
$\alpha_0+\cdots+\alpha_5=\ell$. The factors
$\delta_i^{|\alpha_i|}$ have nonnegative valuation when nonzero.
Consequently at least one product of coefficients of the six unscaled
series has valuation at most $\gamma n-\sum_\tau v(\ell_\tau!)$.
Six further copies contribute their constant terms one.

\paragraph{Comparison with the ramified prime sum.}
Fix the permissible pair and the scale exponent of Theorem~3.1; all estimates
below hold for all sufficiently large integer $x$ at those fixed parameters.
Write $w=a/x_1$, $h=\lceil 6b/k\rceil$ and
$B_\tau=\prod_{i=1}^h((\ell_\tau-h)_++i)$.
The square of $d\equiv q_0/2-1\pmod{q_0}$ is one modulo $q_0$,
and neither $1$ nor $q_0/2-1$ equals $-1$ modulo $q_0$.
Thus no $e\ge1$ has $-p^e\equiv1\pmod{q_0}$.
In particular $\tau$ and $-p^e\tau$ cannot both have weight $b$,
since all such places lie over $-1$. The formula for $S_p$ gives exactly
\[
 \gamma\sum_\tau\kappa_\tau-S_p
 =\tfrac12\sum_\tau\sum_{e\ge1}
       \min(\kappa_\tau p^{-e},\kappa_{-p^e\tau}).
\]
Every summand of the inner sum is at most $w$ and at most $bp^{-e}$.
Set $N=1+\lceil\log_p k\rceil$. Splitting the inner sum at $N$ bounds it
by $Nw+2b/(kp)$. The scale relations give
$Nw=o(b/(kp))$ since $N k p x/x_1\to0$. Also
$x^{42}=o(b/(kp))$ since $k=O(x^5)$ and $\log b>x^\eta$
for a fixed $\eta>0$.
For $\kappa_\tau\le\ell_\tau\le\kappa_\tau+x^{42}$ it follows that
\[
 \gamma |\ell|-S_p\le f\left(\gamma x^{42}+Nw/2+b/(kp)\right)
 \le f\big\lfloor h/p\big\rfloor
 \le\sum_\tau v(B_\tau).
\]
Indeed $h/p\ge6b/(kp)\to\infty$ and every interval of $h$
consecutive integers contains at least $\lfloor h/p\rfloor$ multiples of $p$.
The twelve coefficients just obtained therefore have valuation sum at most
\[
 S_p+\sum_\tau v(B_\tau)-\sum_\tau v(\ell_\tau!).
\]
This gives the required ramified witness with one copy of the loss budget
instead of two. The assumptions used are the stated local CM identities
and the existence of the frame with a unit row. Realizing these coefficient
products by integral adelic test forms is a separate step.

\subsection{The split and inert Taylor witnesses}
\label{sec:other-witnesses}

For these two cases set $\pi=p$ and let $q=p$ at a split prime and
$q=p^2$ at an inert prime. Take $\rho=1/(q-1)$ and
$\gamma=1/(p-1)-\rho$. Use the same one-curve identities as above,
with the following meaning for the endomorphism identity:
$T\in\mathcal O[[y]]$, $T=z+O(z^2)$, and
\[
 T(pz)=pT(z)+b_1T(z)^q+
           \sum_{i\ge2}b_iT(z)^{1+i(q-1)},\qquad
 b_i\in\mathcal O,\quad v(b_1)=0.
\]
For split $p$, this is the height-one CM parameter input, including at
$p=2$. The unscaled ring $\mathcal O[[y]]$ consists of integral
series in the local curve coordinate; its elements need not have integral
coefficients when expanded in the additive variable $z$.

\paragraph{Six copies in every residual degree.}
Choose $\varpi^{q-1}=-p/b_1$. The same coefficient induction as in
Appendix~A.1 makes $Y=T(\varpi z)/\varpi$ integral and gives
$Y-Y^q\equiv z\pmod p$. Indeed after division by $p\varpi$ the higher
terms of the endomorphism acquire factors
$\varpi^{i(q-1)}/p\in p\mathcal O$, and the recursion at degree $m>1$
divides by the unit $1-p^{m-1}$. Reduction yields
$[z^{q^j}]Y\equiv1\pmod p$.
At $N=q^j$, Legendre's formula again gives $v(N!)=(N-1)/(p-1)$.
It follows exactly as above that each multidegree $\nu$ has a six-copy
witness of valuation at most
\begin{equation}\label{eq:jet-six-general}
 \gamma|\nu|-V(\nu),\qquad V(\nu)=\sum_\tau v(\nu_\tau!).
\end{equation}
Here one uses the integral addition law, the six-theta product identity
and a form with all coefficients units. At a split prime $\gamma=0$
and $S_p=0$. Taking six more constant terms proves the entire split case.

\paragraph{Small inert primes.}
Assume $p$ inert and $p\le k$, with the integral normalized frame
and the unit leading coefficients in the inert coupling computation of
Appendix~A.3. In the permissible family $2$ splits since $d\equiv7\pmod8$.
Moreover $p$ is prime to $q_1$: the odd prime divisors of $q_0$
satisfy $d\equiv-1$ modulo each such divisor, so quadratic reciprocity
gives $(p/d)=1$ for them, and $q_1$ has the same prime divisors as $q_0$.
Set $P(\tau)=-p\tau$ on the group of places. For odd $e$ the partner is
$P^e(\tau)$. Call $(\tau,e)$ a short pair if $e$ is positive and odd,
$2e<\ell_0$, and both $\tau$ and $P^e\tau$ belong to $\Phi_b$.
Distinct short pairs have disjoint endpoints, where an endpoint may
mean either end. Two different exponents at one start would give three
members of $\Phi_b$ at the forbidden powers $0,e,e'$. Two pairs ending
at the same point with $e<e'$ would give, from the start of the longer
pair, three at powers $0,e'-e,e'$. When the exponents are equal the
starts coincide since $P$ is a permutation. If an end were a start,
there would be three members at powers $0,e,e+e'$, with
$e+e'<\ell_0$. The same reasoning excludes a loop (take $e'=e$).
Each use of sparseness is valid since $P^i=(-1)^i p^i$ and signs
are allowed independently.

For each start take $j_\tau=\lfloor\ell_\tau/p^e\rfloor$ copies of its
pair; put $j_\tau=0$ elsewhere. Let $\nu'$ record the $p^ej_\tau$
degrees at each start, $\nu''$ the $j_\tau$ at its end. Their sum is
at most $\ell$: the start bound follows by rounding, while
$j_\tau\le 2b/p\le b\le\ell_{P^e\tau}$. Disjointness prevents
double charging. The uniqueness condition of the inert coupling
calculation holds: an odd competing exponent $a\le e$ with the
same endpoint and a selected start would itself be a short pair.
Thus six copies give a witness of valuation at most
$\gamma(|\nu'|+|\nu''|)-\sum j_\tau-V(\nu')-V(\nu'')$.
For the remaining six use \eqref{eq:jet-six-general} at
$\mu=\ell-\nu'-\nu''$. Multiplication in disjoint variable blocks
adds the valuations. The factorial correction at each place is at
most $3(1+\log_2(2b))$: in Legendre's digit-sum formula for the
multinomial coefficient each of its three positive digit-sum terms
divided by $p-1$ is at most the number of digits of $\ell_\tau\le2b$.

To pay for this correction put $w=a/x_1$,
$n_0=\lceil\ell_0/2\rceil$. At an inert prime
$\sum_{e\ge1,\ e\ {\rm odd}}p^{-e}=p/(p^2-1)=\gamma$, so
\[
 \gamma\sum_\tau\kappa_\tau-S_p
 =\sum_\tau\sum_{e\ {\rm odd},\,e>0}
                 \min(\kappa_\tau p^{-e},\kappa_{P^e\tau}).
\]
For $e<n_0$ all terms at a given place are at most $w$, except
possibly its unique short pair, whose contribution is at most
$j_\tau+1$. The remaining geometric tail is at most
$2b p^{-n_0}\le2b/(kp)$. For the last inequality use
$\ell_0>27s$, $2^s=m^{22}\ge x^{11}$, and $k\le 4x^5$,
so $p^{n_0-1}\ge k$ eventually. The total excess per place after
subtracting $j_\tau$, including the factorial correction, is at most
\[
 \gamma x^{42}+3(1+\log_2(2b))+1+n_0 w+2b/(kp).
\]
The first four terms together are $o(b/(kp))$, uniformly for
$p\le k$: use $k^2=O(x^{10})$, $n_0=O(\log x)$, the
exponential growth of $x_1$ in $x^{\eta/6}$ and of $b$ in $x^\eta$.
The display is at most $\lfloor h/p\rfloor\le v(B_\tau)$ for all
sufficiently large $x$. This proves the twelve-copy bound
$S_p+\sum v(B_\tau)-V(\ell)$ uniformly in the small inert primes.

\paragraph{Large inert primes.}
Now $p>k$. We use the integral resolution of the frame at this prime:
each variable, up to a unit $s$, is an integral linear combination
$sz_\tau=\sum_i a_{\tau i}L_i(z)$. For a normalized tight frame
$u_i$ with integral entries and $L_i=s u_i^\dagger z$ this follows
by taking $a_{\tau i}=c_i u_{i,\tau}$ in the chosen embedding.
The series $G_\tau=\sum_i a_{\tau i}y(L_i)$ starts with $sz_\tau$.
Thus $(\prod_\tau G_\tau^{\mu_\tau})\Theta$ has coefficient $s^{|\mu|}$
at multidegree $\mu$. This supplies any residual degree with valuation
zero, using one copy and five constant terms.

Use the exponent-one coupling at every place. Uniqueness in the
coupling calculation follows from bijectivity of $P$. Choose
\[
 i_\tau=\left\lfloor (\ell_\tau-h)_+/p\right\rfloor,\qquad
 j_\tau=\min(i_\tau,\ell_{P\tau}).
\]
As $\ell_\tau\le2b$ and $p>k$, $i_\tau\le h$.
The degree used at a place $t$ is
$p j_t+j_{P^{-1}t}\le(\ell_t-h)_++\min(h,\ell_t)=\ell_t$.
Appendix~A.3 gives valuation at most $-\sum j_\tau$ for the first
six copies. The unit-coefficient construction in the preceding
paragraph fills the remainder. For $r=(\ell_\tau-h)_+$,
$\ell_\tau!\mid r! B_\tau$, and
$v(r!)=i_\tau+v(i_\tau!)$. Also $i_\tau!\mid h!\mid B_\tau$,
the latter by the integrality of $\binom{r+h}{h}$.
Therefore $v(\ell_\tau!)-i_\tau\le 2v(B_\tau)$.
Since $h\ge x^{42}$ eventually, $r\le\kappa_\tau$, and hence
\[
 i_\tau-j_\tau\le(\kappa_\tau/p-\kappa_{P\tau})_+ .
\]
Summing this last bound costs at most the $e=1$ part of $S_p$.
This proves the bound $S_p+2\sum v(B_\tau)-V(\ell)$, uniformly for
all $p>k$ inert.

\paragraph{Inputs to the local deduction.}
Taken together, these computations give the prescribed-degree
coefficient witnesses for every prime on the scales of Theorem~3.1.
The one-curve formal parameter and theta integrality identities must
hold in the specified normalization, and the required frames must exist
with integral coefficients, a unit row and, at inert primes, the
integral resolution and normalized coupling calculation.
Passage to integral sections and the four-block test-form pairing
still require their arithmetic realizations.


\section{The factor two in the finite character moment}\label{sz:shells}
The projection to good types in Section~\ref{B:arithmetic} uses the
unnormalized sum $\sum_{g\in G}|\psi_U^\natural(g)|^2\le2$. The small
constant survives integration over the product of the chosen places and
their conjugates. We give the passage from the estimates for the frame
kernel to this finite sum, with the assumptions on that kernel made explicit.

Let $G_{\rm loc}$ be the unitary group over the semilocal algebra and $K$ its integral subgroup, $K_M$ its
principal congruence subgroup, and $\mu$ the Haar measure of the filling
datum. The character $\chi_M$ is extended by zero off $U K_M$, where
$U$ is the full Bessel unipotent subgroup of $G_{\rm loc}$. On $K$ the support restricts to $(U\cap K)K_M$. It is constant on
right $K_M$-cosets, has modulus at most one, and equals one on $K_M$.
For $K_M\subset Z\subset K$ a compact open intermediate group, set
\[
 a_Z(x)=\int_Z\chi_M(yxy^{-1})\,d\mu(y),\quad
 R_Z=\frac{\mu(Z)\mu(K_M)}{u_F^2},\quad h_F=\varrho+1,
\]
where $u_F>0$ is the unipotent mass of the filling and $h_F$ is its frame
parameter. On $K_M$ the integral equals $\mu(Z)$ by normality.

Here are the precise kernel estimates used on the complement.
In the gauge of scale one, let $P_Z(x)$ be the nonnegative majorant
of the zero-floor frame correlation with empty heads, and let $W_Z(x)$
be its weighted square product. These are the functions
\texttt{formMajorant} and \texttt{weightedSquareProduct} of the filling.
The following four geometric inputs specify the remaining part of F5.
\begin{enumerate}
\item The pairing-kernel estimate gives $|a_Z(x)|^2\le u_F^2 P_Z(x)$.
\item At an active point $y$ and each frame index $i$, the set of stable
products has volume at most the modulus of some element $c$ in it with
$c$ integral and invertible in the semilocal algebra. If every such set
contains a $c$ with integral inverse, the frame reconstruction makes
$(x-1)/M$ an integral matrix. Here the comparison uses the scale-one gauge,
where every bound on matrix size is integrality.
\item If every active point at $x$ has some index with stable-product
volume at most $e<1$, the zero-floor majorant satisfies
$P_Z(x)^2\le e^{2s_F}W_Z(x)$, where $s_F$ is the volume exponent of
the filling.
\item For a measurable $0\le F\le P_Z$ obeying
$F(x)^2\le e^{2s_F}W_Z(x)$ on the positions supporting the majorant,
the dyadic denominator-shell estimate is
\[
 \int F\,d\mu\le e^v R_Z\,2^{10^9h_F^4}
       \qquad(v+10^8h_F^4\le k,\;v\in\mathbb N).
\]
Here the shells group the denominator ideals of the entries of $(x-1)/M$
by the ceiling of the base-two logarithm of their index. The rank $k$
in this assertion is the rank in the filling shape.
\end{enumerate}
These inputs are required for the specific large lattice and its normalized
measure. The following calculation derives the complementary integral from
them and shows which rank margin is spent absorbing its loss.

Suppose $Q\ge2$ and each integral $c$ invertible in the algebra
whose inverse is nonintegral has modulus at most $Q^{-1}$. One can
always take $Q=2$ here. Indeed multiplication by $c$ takes the
integer ring $R$ to the proper finite-index ideal $cR$, and
its modulus equals $[R:cR]^{-1}$. If $x\notin K_M$ and every stable
product set at some active point had volume greater than $Q^{-1}$,
the witnesses $c$ in the second input would all have integral inverse.
Then $(x-1)/M$ would be integral, which in the unimodular integral
model puts $x$ in $K_M$. Thus the third input applies on the complement.
Put $F=\mathbf1_{K_M^c}P_Z$, $B=10^9h_F^4$ and $v=t+B$.
Under $t+2B\le k$ the last input gives
\[
 \int_{K_M^c}P_Z\,d\mu\le Q^{-t}(2/Q)^B R_Z\le Q^{-t}R_Z .
\]
Using the first input there and the exact integrand on $K_M$ yields
\[
 \int |a_Z(x)|^2\,d\mu(x)
 \le\mu(Z)^2\mu(K_M)+Q^{-t}R_Zu_F^2
 =\mu(Z)\mu(K_M)(\mu(Z)+Q^{-t}).
\]

Take $Z=K$, $Q=2$ and $t=h_F$. The filling rank condition is
$h_F+2\cdot10^9 h_F^4\le k$, and the local volume estimate for the
chosen places is $\mu(K)\ge2^{-h_F}$. Hence the parenthesis is at most
$2\mu(K)$. Choose lifts $s_g$ of the finite reduction elements. The
cosets $s_gK_M$ are disjoint. On each, $a_K$ is constant and equals
$\mu(K)\psi_U^\natural(g)$: integration of the conjugated character
is the uniform conjugacy average on the finite quotient. Each coset
has the same positive volume. Summing their integrals and dividing by
$\mu(K)^2\mu(K_M)$ proves \eqref{B:charmean}. Using the whole
product at once gives a single factor two.

\subsection{Integral recovery of a position from its frame}
\label{app:local-frame-recovery}
We prove the stable-product and reconstruction assertions in the second
input of Appendix~\ref{sz:shells}.  The argument takes place over the
entire product of chosen finite places.  Let $A=\prod_v A_v$ be their
product of local fields, $R=\prod_v\mathcal O_v$, and normalize additive
measure by $\operatorname{vol}(R)=1$.  Write $L$ for the coordinate
$R$-lattice of the local Hermitian space and $h$ for its form, linear in
the second argument.  The form is integral and unimodular on $L$.
Write $E_i=\sum_{a<i} A e_a$ for its isotropic frame flag, with
$0\le i\le\varrho$, and $Y_i=E_i^\perp/E_i$.  The Bessel unipotent
group $U$ preserves the flag, fixes $e_i$ modulo $E_i$, and acts
as the identity on $Y_{\varrho}$.  These are the diagonal blocks
of the defining upper unitriangular matrices.  Let $K$ be the
stabilizer of $L$, and $K_M$ its kernel on $L/ML$.  The level $M$
is integral, fixed by the involution and nonzero in each factor.

For $g,g'=gx$ and $w\in E_i^\perp$, call $w'$ a matching lift if
\[
  w'-w\in E_i,\qquad g^{-1}w\in L,\qquad
  g'^{-1}w'-g^{-1}w\in ML .
\]
Define $I_i$ to be the set of values $h(e_i,w)$ for $w$ admitting such
a lift.  At an active point there exist $u,v\in U$ such that
\[
 C=ug\in K,\qquad C'=vg'\in K,\qquad C^{-1}C'\in K_M .
 \tag{*}
\]
Indeed the finite Bessel vector at $y$ has coordinates
$\chi_M(yc^{-1})$ indexed by $c\in Z/K_M$. A nonzero inner product at
$g,g'$ gives one coordinate in which both entries are nonzero. Its
support is $U K_M$, so one can choose $u,v$ with
$ugc^{-1},vg'c^{-1}\in K_M$; now use $c\in K$ and normality.

\begin{proposition}
At each active point $I_i=c_iR$ for some $c_i\in R\cap A^\times$,
and its depth is $[R:I_i]$. If every $I_i$ has index one, then
$x\in K_M$.
\end{proposition}
\begin{proof}
The pairs $(w,w')$ with $w\in E_i^\perp$ and $w'-w\in E_i$
form an $A$-subspace. Intersect it with the inverse image of
$L\oplus L$ under the invertible linear map
\[
 (w,w')\longmapsto (g^{-1}w,M^{-1}(g'^{-1}w'-g^{-1}w)).
\]
The intersection is a full lattice in this subspace: over each DVR
one clears denominators of a basis, and a submodule of a finite free
module is finitely generated. Its image by $(w,w')\mapsto h(e_i,w)$
is a full fractional ideal in every factor, since $(e_i^*,e_i^*)$
has image one for the hyperbolic dual $e_i^*$. Thus $I_i=c_iR$
for a componentwise nonzero generator. Put $f=u^{-1}e_i$ and
$f'=v^{-1}e_i$. These represent $e_i$ modulo $E_i$, and
$g^{-1}f=C^{-1}e_i$, $g'^{-1}f'=C'^{-1}e_i$ both belong to $L$.
For $w\perp E_i$ with $g^{-1}w\in L$,
$h(e_i,w)=h(g^{-1}f,g^{-1}w)\in R$. This gives $I_i\subset R$.
Its inverse volume, the depth by definition, is its index.

Suppose now $I_i=R$ for all $i$. We prove by decreasing induction on
$i$ that every $w\in E_i^\perp$ with $g^{-1}w\in L$ has a matching
lift. For $i=\varrho$ take $w'=v^{-1}u w$. Its class agrees with $w$
by the action on $Y_{\varrho}$. Moreover $g'^{-1}w'=C'^{-1}C g^{-1}w$,
which is congruent to $g^{-1}w$ by $(*)$.

For the induction step choose $z\in E_i^\perp$ and its matching
lift $z'$ with $h(e_i,z)=1$. Put $\alpha=h(e_i,w)\in R$ and
$w_0=w-\alpha z$. Then $w_0\perp E_{i+1}$ and $g^{-1}w_0\in L$.
By induction it has a matching lift $w_0'$ modulo $E_{i+1}$.
Write $w_0'-w_0=\mu e_i+t$ with $t\in E_i$. Then
\[
 \mu=h(z',w_0')-h(z,w_0)\in MR.
\]
Indeed $z'-z\in E_i$, $w_0'\perp E_i$, and $h(z,e_i)=1$.
The two pairings are congruent modulo $M$ on applying $g'^{-1}$ and
$g^{-1}$ respectively, since both pairs of vectors are integral and
congruent modulo $ML$. Now $w'=\alpha z'+w_0'-\mu f'$ has the same
class as $w$ modulo $E_i$, and the subtracted term is zero modulo
$ML$ after applying $g'^{-1}$. This proves the induction.
At $i=0$ the only lift of $w$ is $w$ itself. Taking $w$ in the
columns of $g$ gives $(x^{-1}-1)L\subset ML$. Taking the Hermitian
adjoint and using unimodularity gives $(x-1)L\subset ML$ as well.
Thus $x\in K_M$.
\end{proof}

In particular every active point at $x\notin K_M$ has some $i$ with
$\operatorname{vol}(I_i)\le 1/2$. More precisely the volume is
$\prod_v q_v^{-v(c_{i,v})}$, where $q_v$ is the residue cardinality,
so $1/2$ can be replaced here by $(\min_v q_v)^{-1}$.
This establishes the second geometric input for the stated finite
unimodular model. Converting this depth into a high power of the
volume saving still requires the correlation and integration estimates
in the remaining inputs.

\paragraph{The zero-floor pairing and its normalization.}
Let $\eta$ be Haar measure on $U$, let $u_F=\eta(U\cap K_M)$ and
let $d\dot y$ be quotient measure on $U\backslash G_{\rm loc}$,
normalized by $d\mu=d\eta\,d\dot y$. Write
$V_c(y)=\chi_M(yc^{-1})$ for $K_Mc\in K_M\backslash Z$, and
$H_x(y)=\sum_c\overline{V_c(y)}V_c(yx)$. Then
\[
 a_Z(x)=u_F\int_{U\backslash G_{\rm loc}} H_x(y)\,d\dot y .
\]
In fact, if $u y\in K_Mc$, the two coordinates at $c$ have the
common phase $\chi_M(u^{-1})$. Their product after conjugating the
first equals $\chi_M((uy)x(uy)^{-1})$ by right $K_M$-invariance.
For fixed $y,c$ the set of such $u$ is empty or one right coset
of $U\cap K_M$ within $U$. Its measure in the latter case is $u_F$,
as $U$ is unimodular. If the set is empty then $V_c(y)=0$.
Partitioning $\{u:uy\in Z\}$ by these cosets and applying the
quotient integral formula gives the identity. All coordinates are
bounded with compact support modulo $U$.
In the frame integration coordinates of the filling let $d\theta_0$
be the first floor measure and $d\theta_{>0}$ the measure of the
remaining floors, so $\dot y$ is represented by $y'\sigma$ with
measure $d\theta_{>0}(y')d\theta_0(\sigma)$ on $U$-invariant
integrands. At zero floor the two heads coincide, and the definition
of the majorant therefore reads
\[
  P_Z(x)=
 \left(\int\left|\int H_x(y'\sigma)\,d\theta_{>0}(y')\right|
                  \,d\theta_0(\sigma)\right)^2 .
\]
The triangle inequality now gives $|a_Z(x)|^2\le u_F^2 P_Z(x)$,
the first input of Appendix~\ref{sz:shells}, with coefficient one
for these normalized floor measures.

\paragraph{Counting the denominator ideals.}
Write $n=\varrho+1$ and $H=2^j$ for a dyadic height, $j\ge0$.
The denominator ideal of a matrix $D$ over $A$ is
$J=\{c\in R:cD\text{ is integral}\}$. It is generated by an element
$d\in R\cap A^\times$, componentwise by a uniformizer to the
largest of zero and the pole orders of the matrix entries.
Take $D=(x-1)/M$. Every $x$ belongs to the shell of the least dyadic
$H\ge [R:J]$ and the scale $\bar d d$ makes $(\bar d d)D$ integral.
The scale is invariant under the involution, and its additive modulus
is $[R:J]^{-2}$, since the involution preserves the measure of $R$
and carries ideals to ideals of the same index. Hence this region
costs at most $2 H^2$ in the gauge $2/|\bar d d|$.

There are at most $2H$ such ideals for the chosen places and all
sufficiently large cutoffs. In fact there are at most $4f$ local
factors, since they lie above the two rational primes $p_0,p_1$
and $[E:\mathbb Q]=2f$; any chosen subproduct obeys the same bound.
Each $q_v\ge b^{10}$. Therefore
\[
 \sum_{0<[R:J]<\infty}[R:J]^{-1}
   =\prod_v(1-q_v^{-1})^{-1}
   \le \exp(8f b^{-10})\le 2.
\]
The logarithmic bound uses $-\log(1-s)\le2s$ for $0\le s\le 1/2$.
The last inequality follows from $f\le q_1\le x^{12}$ eventually
and the scales of Theorem~3.1. Every ideal of index at most $H$
contributes at least $1/H$. This gives the required covering count,
indeed with $2H$ in place of $2H^{9n}$.

\paragraph{The finite quotient only needs positions in the integral group.}
The moment in (10.11) is a sum over $K/K_M$, so restrict the outer
integral to $x\in K$. In this region a stronger form of flag recovery
gives the decay needed for summing over $J$. Write
$J(x)=\{c\in R:c(x-1)/M\text{ is integral}\}$ and
$N(x)=[R:J(x)]$. Each chosen place is unramified over its level
prime: the two level primes are prime to the conductor $q_1d$ of a
cyclotomic field containing $E$. Thus $M=p_0p_1$ has valuation one
in each local factor. Since $x$ is integral, the local denominator
ideal is either the whole valuation ring or its maximal ideal. The
latter occurs precisely when the reduction of $x$ at that component
differs from the identity.

Decompose $R$ according to the orbits of its factors under the
involution. For each orbit $o$ contributing to $N(x)$, write $q_o$
for the residue cardinality of any member. This is independent of the
member. The orbit contributes at most $q_o^2$ to $N(x)$.
At every active pair the hypotheses $(*)$ project to each orbit.
The matching conditions defining $I_i$ are componentwise over these
orbits, so $I_i$ is the product of the corresponding ideals $I_{i,o}$.
If $x$ differs from the identity modulo $M$ on $o$, the preceding
proposition applied there gives an $i$ with
$[R_o:I_{i,o}]\ge q_o$. Multiplication over contributing orbits gives
\[
  \prod_{i<\varrho}[R:I_i]\ge\prod_{o\text{ contributing}}q_o
       \ge N(x)^{1/2}.
\]
Hence, for $\varrho>0$, some index has depth at least
$N(x)^{1/(2\varrho)}$. For $\varrho=0$, the induction at the top
quotient already implies $x\in K_M$ wherever the pairing is active.

For precision we isolate the remaining analytic estimate on one shell.
For $H=2^j$ let $S_H$ be the denominator shell introduced above.
The normalized one-frame bound used here is that if $F$ is
nonnegative, is supported in $S_H$, and satisfies $F^2\le t^2 W_Z$
there, then
\[
 \int F\,d\mu\ \le\
    t\,R_Z\,2^{2\cdot10^6 n^3+1287n} H^{10^5 n^2}.
 \tag{**}
\]
It suffices to impose this statement on functions dominated by $P_Z$
and supported in the positions of the frame correlation.
This inequality is an analytic input. So is the third geometric
input in Appendix~\ref{sz:shells}, giving the weighted pointwise
gain $P_Z^2\le e^{2s_F} W_Z$ if every active frame has a depth at least
$e^{-1}$. No estimate for deep vectors at positions outside $K$ is
required below.

Take $Z=K$ and $F=\mathbf 1_{K\cap S_H}P_K$ for $H\ge2$. On this
shell $N(x)\ge\sqrt H$: for $H=2$ this follows from integrality of
$N(x)>1$, and for $H\ge4$ it follows from $N(x)>H/2$.
The previous product estimate permits $e=H^{-1/(4\varrho)}$.
Consequently $(**)$, with $t=e^{s_F}$, gives
\[
 \int_{K\setminus K_M}P_K\,d\mu\le
 R_K\,2^{2\cdot10^6n^3+1287n}
    \sum_{j\ge1}2^{-j(s_F/(4\varrho)-10^5 n^2)} .
\]
The geometric series pays for its prefactor. In fact, if the block
parameter of the large lattice datum is $m_0+1$, the starting
coordinate rank is $k-2m_0-3$, two blocks each of rank $2m_0+4$
are removed, and the exponent subtracts a further $3n+m_0+2$.
Thus $s_F=k-7m_0-3n-13$, with $m_0<9n$. Since
$k\ge n+2\cdot10^9 n^4$, this gives $s_F\ge10^9 n^4$.
Write $A_0=s_F/(4\varrho)-10^5 n^2$.
Using $\varrho<n$ gives
$A_0\ge 2\cdot10^6n^3+1288n+1$.
The sum is at most $2^{1-A_0}$, so the displayed integral is at most
$2^{-n}R_K$. The zero-floor pairing and the exact value on $K_M$
now give
\[
 \int_K |a_K|^2\,d\mu\le
  \mu(K)\mu(K_M)\bigl(\mu(K)+2^{-n}\bigr).
\]
With the normalized volume bound $\mu(K)\ge2^{-n}$ from
Appendix~\ref{sz:shells}, division by $\mu(K)^2\mu(K_M)$ gives (10.11).
This proves the finite-quotient reduction from the weighted pointwise
gain, the one-shell integration bound $(**)$ and the specified
measure normalizations. Other uses of an integral over the entire
local group require their own estimates outside $K$.


\section{Comparing the good spaces on different real forms}
\label{sec:trace}

Let $V_s,V_{s'}$ be two of the rank-$r$ hermitian spaces $Y_0,W_j$
with the prescribed finite frames. These frames identify their groups of
finite adelic points and their principal level $K$. They have a common
doubled Harish--Chandra parameter $l$ and possibly different real
signatures $s,s'$. Write $M_s$ for the space of forms of type $l$ at
level $K$, $e_0$ for the good projector and $M_s^g=e_0M_s$. The Hecke
algebra $\mathcal H$ used here consists of the spherical operators away
from the chosen places, acting simultaneously on the different spaces.
For its elements the desired equality is
\begin{equation}\label{tr:good}
  \operatorname{tr}(T\mid M_s^g)=\operatorname{tr}(T\mid M_{s'}^g)
  \qquad(T\in\mathcal H).
\end{equation}
We derive this under the geometric inputs specified below. The spreading
argument needs $T=1$ at $j=m$. The chain determinant also needs a common
Hecke operator at each $j$.

\subsection{Getting a cuspidal component from a label}\label{tr:label}
At the chosen split places the principal level is the kernel of reduction
onto the product of the finite groups $G_w=\mathrm{GL}_r(k_w)$.
The good idempotent is $e_0=\sum_c e_c$, where $c$ ranges over the good
tuples and $e_c$ is the tensor product of their character idempotents.
Each good tuple has a primitive degree-$r$ label at some $w$ above
the first level prime. Its root in $\mathbb F_{q^r}$,
$q=|k_w|$, generates $\mathbb F_{q^r}^{\times}$. Under the chosen
identification with its character group this gives a character of order
$q^r-1$. Its Frobenius orbit has length $r$, since $0<q^h-1<q^r-1$
for $0<h<r$. The irreducible factor of the label therefore indexes
Green's cuspidal character $\chi$ of $G_w$, corresponding to the degree-$r$
orbit and the one-box partition~\cite[\S8, Theorem~14]{green}.

One can also read off the invariant multiplicity from the polynomial
definition of the label. In the Hall algebra let $*$ be parabolic
induction, $\mathbf{u}$ the constant function, $\mu$ its $*$-inverse,
and $\beta$ the sum of the chosen multiplicative character over
the eigenvalues with multiplicity. Green's function for the irreducible
label is the rank-$r$ part of $C=(\beta\mu)*\mathbf{u}$ up to sign,
with $\beta\mu$ the pointwise product. Block averaging $\delta$ is
compatible with $*$ by the finite Mackey formula (partition a stable
subspace by its intersection with the first block and its image in the
quotient). Since $\delta\mathbf{u}=\mathbf{u}\otimes\mathbf{u}$,
one has $\delta\mu=\mu\otimes\mu$. The eigenvalue sum is constant on
a block extension and equals the sum on its diagonal blocks. Consequently
\[
 \delta(\beta\mu)=(\beta\mu)\otimes\mu+\mu\otimes(\beta\mu),
 \qquad \delta C=C\otimes 1+1\otimes C,
\]
where $1$ is supported in rank zero. Taking two identity diagonal
blocks of positive ranks gives zero average on their radical. This
average of the character is the dimension of the invariant subspace.

Here is the vanishing that the trace comparison requires. Let $N$ be
the unipotent radical of a proper coordinate maximal parabolic, and
let $\pi$ afford $\chi$. The averaging operator
$|N|^{-1}\sum_{u\in N}\pi(u)$ projects onto $\pi^N=0$ by cuspidality.
Multiplication on the left and right and taking the trace gives
$\sum_{u\in N}\chi(aub)=0$ for every $a,b\in G_w$. The coefficients of
the central idempotent are $\chi(1)\overline{\chi(g)}/|G_w|$.
Their sums over these translates also vanish. The map
$Z\mapsto 1+pZ(1-p)$, with $p$ the projection onto the chosen
coordinate subspace, covers $N$ with constant fibre size, so summing
over all matrices $Z$ gives the same vanishing. In $e_c$ the remaining
finite factors are constant during this summation.

Lift the tuple idempotent by a section of reduction and convolve with
$\mathbf 1_K$. On the standard maximal compact this function is the
coefficient function of $e_c$ evaluated at the reduction: precisely the
chosen lift of that reduction contributes. Outside the compact it
vanishes. Right multiplication by the local integral lifts of $N$ stays
in the same compact coset. Thus all the local radical sums that define
cuspidality at $w$ vanish. Convolution on the left by the group-algebra
representative of $T$ preserves these sums. This proves the local
hypothesis for the finite test function of $T e_c$.

\subsection{Comparing traces with two local conditions}\label{tr:geometry}
For clarity, the geometric branch of the argument uses the following
inputs for the kernels and Haar normalizations of these spaces.
They are hypotheses of the present reconstruction.
\begin{itemize}
\item[(i)] At a common compact open level $K'$ for either signature,
the kernel for $l$ has cusp decay of order $10r^5$, and the convergent
pretrace expansion for a level-preserving test is
$\operatorname{vol}(K')\operatorname{tr}=\sum_R I_s(R)$.
Here $I_s(R)$ is the contribution of the characteristic polynomial
$R\in E[X]$ of rational elements. The expansions unfold to class
orbital integrals for irreducible $R$, with absolutely integrable cyclic
triple-kernel integrals also after weighting by a one-step logarithmic
norm increment.
\item[(ii)] For coprime nonconstant local factors, a conjugacy-invariant
distribution supported on the characteristic-polynomial fibre of their
product kills a test function with vanishing radical sums at the split
place. This is the local parabolic descent step.
\item[(iii)] For irreducible $R$ with all roots on the unit circle at
every infinite place, each elliptic class integral factors into a finite
factor and the product of real orbital integrals. Changing the signs
attached to the eigenlines, with total sign product fixed at $1$, gives
a bijection of the corresponding frame classes preserving the finite factor.
For regular $z=(z_i)$ on the compact torus, the real kernels obey
\[
 (-1)^{n_+n_-}\sum_{\sigma\in S_r} O_s(z\circ\sigma)
 =|S_{n_+}\times S_{n_-}|\,\overline{\chi_{l,\mathrm{cpt}}(z)}.
\]
Here $n_+,n_-$ count the two signs, $O_s$ uses the
formal-degree-normalized discrete-series kernel of $l$, and the right side
uses the lowest $K$-type on the compact form at that place.
\end{itemize}
Section~\ref{tr:parabolic-descent} proves (ii) for the split test just constructed. The remaining assumptions concern the indicated kernels on the adelic quotient,
including the finite centralizer measure in (iii).

Suppose the finite test is cuspidal at $w$ and its support at an auxiliary
place consists of integral matrices having zero as a simple root of the
reduced characteristic polynomial. The latter condition excludes a proper
power $R=Q^h$, $h>1$: a monic $Q$ with $Q^h$ integral has integral
coefficients by integral closedness, and its reduction would give zero
with multiplicity divisible by $h$. Each remaining reducible $R$ in the
support has two coprime nonconstant factors over $E$ obtained by
partitioning its primary factors. They stay coprime over the completion.
The functional on the test defined by $I_s(R)$ is conjugation invariant
by change of variables and supported on this fibre. Hence (ii) kills it.

For irreducible $R$, fix a root $z$ with $|z|\ne1$ at one infinite
place, if there is one. Diagonalize a rational representative $\gamma$
there, take its left eigenrow $v$ for $z$, and put
$h(x)=\log\|vx\|$. Then $h(\gamma x)=h(x)+\ell$ with
$\ell=\log|z|$. The centralizer scales this eigenrow, so differences
of $h$ descend under its action. Write $p$ for the convolution
idempotent kernel at infinity. In the unfolded triple integral for the
orbital integral of $p*p*p$ the three successive increments of $h$ from
$x$ to $\gamma x$ sum to $\ell$. Their integrals agree by cyclic
rotation. The integral of the sum of the first two equals that of the
first alone by collapsing the first two kernels using $p*p=p$.
The second integral is therefore zero. Assumption (i) justifies these
integrations and shows that $\ell I_s(R)=0$ after summing the classes.

In the elliptic case, sum (iii) over all orders of the roots at each
infinite place. For fixed signature the eigenline signs of a given
class occur $|S_{n_+}\times S_{n_-}|$ times at each place.
These multiplicities cancel the same factors in the real orbital formula.
The field has an even number $f$ of infinite places, so the total sign
$(-1)^{fn_+n_-}$ is $1$. Every rearrangement of the signature at all
places has total sign product $1$ for the same reason. Matching with
the all-positive signature in (iii) leaves one common finite sum
$C(R)$, and
\[
 I_s(R)=C(R)\prod_{v\mid\infty}\overline{\chi_{l_v,\mathrm{cpt}}(z_v)}
       =I_{s'}(R).
\]
Summing over $R$ and cancelling the positive common volume proves the
two-place trace comparison.

The hypotheses on the data here are quantitative. The field contains
$\mu_{5040}$ because $5040\mid18480\cdot3\mid q_1$ in its cyclotomic
construction. The tower exponent is positive, which gives the parity of
$f$. The admissible ranks give $r\ge2$. At every infinite place the
entries $l_i$ are distinct and congruent to $r-1\pmod2$. For indices of
opposite sign one has
$|l_i-l_j|\ge2(500r^3+10r^5)$ for both signatures. This follows from
the gap in the admissible rows: at $Y_0$ use the initial grid gap; at
$W_j$ the two groups of undoubled parameters lie on opposite sides of
the cut $-\mathrm{liftWeight}_j$ with the prescribed half-gap on each
side. The order $10r^5$ used in (i) is at least $13r^3$. The finite
frames, principal level, and parameter are the same in both applications.

\subsection{Removing the auxiliary local condition}\label{tr:line}
Fix a group-algebra representative $a$ of $T e_c$. Choose a split
place $v$ away from $w,\bar w$, the finite support of the level, and
all nonintegral components of the finitely many elements occurring in
$a$. There are infinitely many split places. At $v$ cut $K$ to the
inverse image of the stabilizer $J$ of the first coordinate line in
$\mathrm{GL}_r(\mathbb F_q)$. Write $E_0$ for averaging over the full
local compact and $A$ for the action of the part of $a$ away from $v$
on this finer level. On the range of $E_0$, $A$ is precisely the old
action of $a$ because the removed components are integral.
Let $P,Q$ denote the unnormalized double-coset operators for
$\operatorname{diag}(\varpi,1,\ldots,1)$ and its inverse at level $J$.
They commute with $A$, and $Q=P^*$ since the inverse double cosets
have the same degree $q^{r-1}$.

Here is the local count. The right cosets for $P$ are indexed by
linear forms $\xi$ with $\xi(e_1)=1$, and those for $Q$ by columns
whose first coordinate is $1$. Representatives are the inverse lift of
the identity with first row replaced by $\xi$, times the diagonal matrix,
and the inverse diagonal matrix times a lift of the identity with first
column replaced by the column, respectively. Multiplying cancels the two
diagonal matrices. For fixed $\xi$ the resulting column $z$ ranges
once over $\xi(z)=1$, since the first row of the row matrix is $\xi$.
The resulting translate of a $J$-fixed vector $u$ depends only on the
line of $z$; write it $u_z$. For $z=e_1$ the count is $q^{r-1}$,
for other multiples of $e_1$ it is zero, and for $z$ off the first
line it is $q^{r-2}$ by two independent affine equations. Since
$\sum_{z\ne0}u_z=(q^r-1)E_0u$, we obtain
\begin{equation}\label{tr:parahoric}
 PP^*u=q^{r-2}\{u+(q^r-1)E_0u\}.
\end{equation}
The positive powers of $P$ are supported on matrices reducing to block
form $\left(\begin{smallmatrix}0&*\\0&B\end{smallmatrix}\right)$ with
$B$ invertible. Products and right multiplication by $J$ preserve that
shape, so zero is a simple root. Removing the place $v$ from $a$ and
cutting there preserve the radical sums at $w$: locally at $v$ all
translates in $a$ are integral and at $w$ the last factor has the
vanishing proved above. Equivalently the removed test is constant along
the $v$ component before cutting, and the cut imposes the independent
$J$ indicator there. Left convolution preserves the vanishing for every
positive power. The two-place result therefore equates
$\operatorname{tr}(A^hP^j)$ on the two fine-level spaces for all $h,j\ge1$.

To recover $\operatorname{tr}(AE_0)$, fix $j$ and decompose into generalized
eigenspaces of $A$. On the $\alpha$-space the trace of $A^hP^j$ equals
$\alpha^h\operatorname{tr} P^j$, since $P$ commutes with the nilpotent
part. For the union of nonzero eigenvalues of the two actions, the
matrix $(\alpha^h)_{1\le h\le N}$ is invertible by Vandermonde, with
$N$ the number of such eigenvalues. Thus the positive power sums of $P$
agree on each nonzero $\alpha$-space. Equation~\eqref{tr:parahoric}
makes $P$ invertible there. Positive power sums determine its multiset
of nonzero eigenvalues with multiplicities, by the same Vandermonde
argument. On this space $B_\alpha$ take determinants in
\eqref{tr:parahoric} to get
\[
  |\det(P\mid B_\alpha)|^2
  =(q^{r-2})^{\dim B_\alpha}(q^r)^{\operatorname{rank}(E_0\mid B_\alpha)} .
\]
Both $P$ and $P^*$ preserve $B_\alpha$, so the restriction still has
the stated adjoint. The common dimension and determinant fix the rank.
Summing $\alpha\operatorname{rank}(E_0\mid B_\alpha)$ gives the desired
old-level trace; $\alpha=0$ contributes zero. This removes the
second-place hypothesis for $T e_c$. Finally
$\operatorname{tr}(T\mid M_s^g)=\operatorname{tr}(Te_0\mid M_s)$
because $T$ preserves the range of the idempotent. Summation over $c$
gives \eqref{tr:good}.

\subsection{The two uses of the comparison}\label{tr:uses}
Taking $T=1$ gives $\dim M_s^g=\dim M_{s'}^g=d_T$ when one of the
spaces is $Y_0$. Applying \eqref{tr:good} to every power of $T$
and using Newton's identities gives equal characteristic polynomials
and equal determinants on the good spaces. For an individual Hecke
character, take the finite union of the joint eigencharacters on the
two principal spaces. For each other character choose a Hecke element
separating it from the selected one and rescale the difference with
its value at the other character to give values zero and one.
Multiplying the separating factors gives a common projector onto the
chosen eigenspace, using the simultaneous eigenspace decomposition
for the commuting normal spherical operators. Its trace is the good
multiplicity. Only the identity and the determinant comparison are
required in the applications just described.

For the latter, one must still identify the operators on the two
spaces with the action of a \emph{single} Hecke element. Fix $j$ and
the tame character. In the convergent expression for the matrix
coefficients of the Rallis operator on principal forms each coset term
is a common scalar (the finite vacuum coefficient with its residue
cardinality normalization) times $\langle z,\pi(\gamma)y\rangle$,
with $\gamma$ supported off the chosen level places. Let $K^0$ be the
standard compact in the off-level group. Average $k\gamma$ over
$k\in K^0/(K^0\cap\gamma K^0\gamma^{-1})$. Compact openness gives a
finite quotient of cosets. This normalized spherical operator has the
same matrix coefficient as $\gamma$ on both spaces: $z$ and $y$ are
fixed there and the compact action is unitary. Thus each pair of
terms lies in the joint evaluation image of $\mathcal H$ in the two
endomorphism spaces, expressed in matrix-coefficient coordinates. That
image is a complex linear subspace of a finite-dimensional space, hence
closed. Taking sums preserves membership. The common residue cardinality
multiplying the normalized series restores the two Rallis operators.
Their determinants on the good spaces therefore coincide by
\eqref{tr:good}. This step uses the Rallis series with its specified
normalization and convergence. The analytic identity relating that
operator to the adjoint times the lift is a separate hypothesis in the
chain argument.

\subsection{Eliminating a split characteristic-polynomial fibre}
\label{tr:parabolic-descent}
Here is a local proof of assertion (ii) of \S\ref{tr:geometry}.
Let $F$ be the local field at the chosen split place, $K=\mathrm{GL}_r(\mathcal O_F)$,
and let $R=R_1R_2$ have coprime monic factors of positive degrees $a,b$
and nonzero constant terms. Write $G=\mathrm{GL}_r(F)$, $M=\mathrm{GL}_a(F)\times
\mathrm{GL}_b(F)$ and $U$ for the upper block radical, identified with
the additive space of $a$ by $b$ matrices. Suppose the compactly supported
locally constant test $f_0$ has
\[
 \int_U f_0(k m u k^{-1})\,du=0\qquad (k\in K,\ m\in M).
\]
For the inflated character idempotent this follows from the radical sums
of \S\ref{tr:label}. Indeed $mu\in K$ forces both diagonal blocks to be
integral and invertible, and then $u\in U(\mathcal O_F)$.
Reduction integrates each residue class with the same measure.
The same conclusion holds for any fixed left and right translates of the
test if the finite radical sums vanish for arbitrary left and right
multipliers: by $G=K M U= M U K$ the outside factors reduce to compact
multipliers and an affine change of the variable in $U$.

\begin{lemma}
Every conjugation-invariant distribution on $G$ supported on the fibre of
$R$ annihilates $f_0$.
\end{lemma}
\begin{proof}
On this closed fibre, the ordered complementary subspaces
$\ker R_1(g),\ker R_2(g)$ have dimensions $a,b$. The projectors onto them
are fixed polynomials in $g$ by B\'ezout's identity. The map to complementary
pairs is therefore continuous and equivariant; the space of pairs is
$G/M$. Over the coordinate pair the fibre is the space $X$ of diagonal
blocks $(A,B)$ with their indicated characteristic polynomials. A test
vanishing on the closed fibre is annihilated by the distribution, since
its compact open support misses that fibre. Tests on the fibre extend
locally constantly with compact support, by a finite compact open cover.
Thus Frobenius descent for locally compact totally disconnected spaces
applies to the distribution on this fibre: $G$ and $M$ are both
unimodular. See Aizenbud, Gourevitch, Rallis and Schiffmann,
\emph{Multiplicity one theorems}, Theorem 2.2~\cite{agrs}.

Concretely the descended distribution on $X$ acts on the test
\[
 (A,B)\longmapsto \int_K\int_U
 f_0(k u\operatorname{diag}(A,B)u^{-1}k^{-1})\,du\,dk .
\]
This is the quotient integral in descent: in the Iwasawa coordinates
$g=kum$ one disintegrates Haar measure as $dk\,du\,dm$, up to a fixed
positive scalar, and absorbs $m$ by invariance of the distribution
on $X$. The displayed integral is compactly supported on $X$.
Indeed the compact support bounds the diagonal and off-diagonal blocks
after conjugation by $K$; the inverse of the map $Z\mapsto ZB-AZ$ is
uniformly bounded there. Its determinant, up to sign, is the nonzero
constant $\operatorname{Res}(R_1,R_2)$, as seen by triangularizing
$A,B$ separately over a splitting field. The same bounds make the
integral locally constant and justify testing the distribution under
the integral.

For fixed diagonal blocks, the upper-right entry of the conjugate by
$u(Z)$ is $ZB-AZ$. Changing variables from $Z$ to $A^{-1}(ZB-AZ)$
expresses the inner integral as a constant times the radical integral
in the hypothesis, hence it vanishes. This argument also covers
nonsemisimple diagonal blocks.
\end{proof}


\section{Simultaneous finite-place specialization and the denominator budget}
\label{sec:finite}\label{c4:finite-repairs}

The regular twist in Section~\ref{c3:twists} must work for every power from
\(1\) to \(m\) at a given identification \(\iota\).
We prove the specialization result with the order of the roots bounded
\emph{before} the auxiliary primes and their polynomials are chosen.
We also evaluate the sum over primes of the explicit crystalline-loss
allowance. Its application to crystalline classes is a separate arithmetic
hypothesis, specified below.

\subsection{A common point for all the polynomials}
\label{c4:root-spec}
Fix a prime \(\ell\). Let \(v\) be the ordinary additive valuation on
\(\overline{\mathbb Q}_\ell\), with \(v(\ell)=1\), \(v(0)=+\infty\);
write \(\mathcal O\) for its valuation ring. Fix a bound \(s_0\) on the
number of variables, a positive rational \(u\), a nonnegative real
\(\beta\), and \(\epsilon>0\). There exists an integer \(h_0\) such
that the following holds. Let \(P_i\) be any finite collection of
polynomials over \(\mathcal O\) in \(s\le s_0\) variables and let
\(x_i\in\mathcal O^s\) satisfy
\[
 v(x_{ij})\ge u,\qquad \sum_i v(P_i(x_i))\le\beta.
\]
There exist \(\ell^{h_0}\)-th roots of unity \(\zeta_j\), common to all
\(i\), such that \(v(\zeta_j-1)\le u\) and
\[
 v\bigl(P_i((\zeta_j-1)_j)\bigr)\le v(P_i(x_i))+\epsilon
 \quad\text{for every }i .
\]
There is no degree restriction on \(P_i\).

For nonzero \(g=\sum a_k Y^k\in\mathcal O[Y]\), put
\(v_w(g)=\min_k(v(a_k)+kw)\). This is additive on products for
\(w\ge0\): in the product take the sum of the smallest indices
attaining the minima in each factor. Only their product attains the
minimum at that index. Factoring \(g\) into linear factors over
\(\overline{\mathbb Q}_\ell\) now gives
\[
 v_w(g)=v_0(g)+\sum_{v(\alpha)\ge0}\min(w,v(\alpha)),
\]
with roots counted with multiplicity. Indeed each nonintegral root
already contributes \(v(\alpha)\) at \(w=0\) and at \(w\ge0\).
For \(v(y)=w\), all factors except those with \(v(\alpha)=w\)
have \(v(y-\alpha)=\min(w,v(\alpha))\).

Write \(\rho_h=1/\varphi(\ell^h)\), \(h\ge1\). Every primitive
\(\ell^h\)-th root has \(v(\zeta-1)=\rho_h\): its conjugates have
the same valuation and their product, up to sign, is
\(\Phi_{\ell^h}(1)=\ell\). Fix \(0<\delta<1\), \(B\ge0\).
Choose \(h_0\) so large that
\[
 B<\delta\sum_{h\in\mathcal H}\frac1{h+1},\qquad
 \mathcal H=\{1\le h\le h_0:\rho_h\le u\}.
\]
This is possible by divergence of the harmonic series. If
\(v_u(g)\le B\), and \(R_h\) is the multiset of roots with valuation
\(\rho_h\), the distinctness of the \(\rho_h\) and the factor formula
give \(\sum_{h\in\mathcal H}|R_h|\rho_h\le v_u(g)\).
Some \(h\) therefore satisfies \(|R_h|\rho_h(h+1)\le\delta\).

For \(\alpha\in R_h\), apply the same factor formula to
\(F(Y)=(Y+1+\alpha)^{\ell^h}-1\). Its roots include
\(\zeta-1-\alpha\) for every primitive \(\ell^h\)-th root \(\zeta\).
The coefficient of \(Y\) has valuation \(h\) because \(1+\alpha\) is
a unit. Hence
\[
 \sum_{\zeta\ \mathrm{primitive}}\min\{1,v(\zeta-1-\alpha)\}
 \le v_1(F)\le h+1.
\]
Averaging over the \(1/\rho_h\) primitive roots yields one \(\zeta\)
for which the sum over \(\alpha\in R_h\) is at most
\(|R_h|\rho_h(h+1)\le\delta<1\). Every truncated valuation in this
last sum equals the untruncated valuation. Returning to the factors
of \(g\), we obtain
\[
 v(g(\zeta-1))\le v_{\rho_h}(g)+\delta.
\]

Take \(\delta=\min(1,\epsilon)/(s_0+2)\) and use \(B=\beta+1\)
in choosing \(h_0\). Move one coordinate at a time. At a coordinate
still to be moved let \(g_i(Y)\) denote \(P_i\) evaluated at its
current coordinates except for this one. Its old coordinate has
valuation at least \(u\). Apply the one-variable result to
\(g=\prod_i g_i\), whose \(v_u\) is at most the sum of the old
valuations. The excesses
\(v(g_i(\zeta-1))-v_{\rho_h}(g_i)\) are nonnegative and sum to at
most \(\delta\). Since each \(v_{\rho_h}(g_i)\) is at most the
old evaluation, each valuation and their sum increase by at most
\(\delta\) at this step. In particular the sum stays below
\(\beta+1\) at all steps. After at most \(s_0\) steps the total
increase for each evaluation is at most \(\epsilon\). The chosen
\(h_0\) depends only on the prescribed initial bounds.

\paragraph{Powers of one character.}
Suppose instead that \(v(P_i(x_i))\le\beta_i<\infty\), and fix
\(m\ge1\). Apply the preceding argument to the polynomials
\[
 P_{i,j}(Y)=P_i\bigl(((1+Y_a)^j-1)_a\bigr),\qquad 1\le j\le m .
\]
For each original coordinate \(x\), the monic polynomial
\((1+Y)^j-(1+x)\) has \(j\) integral roots whose product is
\(\pm x\). At least one has valuation \(\ge u/j\ge u/m\).
Use this preimage for the starting point of \(P_{i,j}\), the radius
\(u/m\), and the bound \(m\sum_i\beta_i\).
The same tuple of roots now satisfies, for every \(1\le j\le m\),
\[
 v(\zeta_a-1)\le u/m,\qquad
 v\bigl(P_i((\zeta_a^j-1)_a)\bigr)\le v(P_i(x_i))+\epsilon .
\]
If \(u<1\), a root of order \(\ell^h\) so chosen cannot have order
dividing \(j\), since otherwise
\(v(\zeta_a-1)=1/\varphi(\ell^h)\ge1/j>u/m\).
Consequently
\(v(\zeta_a^j-1)=\ell^{v(j)}v(\zeta_a-1)\le j\,u/m\le u\),
where \(v(j)=v_\ell(j)\). This gives nontrivial inertia for every
one of the required powers.

\subsection{The arithmetic assertion needed before specialization}
\label{c4:twist-polynomials}
Here is the exact point at which the Selmer construction enters.
Start with a CM field \(E\), a finite set \(S_0\) containing the places
above \(\ell\), an inertial element \(g_0\) above \(S_0\) away from
\(\ell\), and finite collections of Selmer data and Stickelberger
systems as in Section~\ref{c3:twists}. Normalize module length by
\(\operatorname{length}(\mathcal O'/\ell)=1\) for the coefficient
DVR \(\mathcal O'\). The numerical allowances of the data and systems
are denoted by \(\beta_i\). Assume the coefficient DVR and a torsion level \(M_0\) have been chosen together with
\(M_0>\beta_i+1\) for every Selmer datum and residue field size
exceeding \(2\operatorname{rank}(T_i)
\operatorname{length}_{\rm unscaled}(\mathcal O'/\ell^{M_0})+1\).
The following polynomial existence statement is required.
There is \(s_0\ge1\) such that for every rational \(u>0\) with
\(s_0u<1\), and all sufficiently large \(N\), there are
\begin{itemize}
\item an auxiliary set \(Q\) of at most \(s_0/2\) degree-one primes, with
\(Q,\bar Q,S_0\) disjoint in the sense that \(Q\cup\bar Q\) avoids
\(S_0\) and \(Q\cap\bar Q=\varnothing\);
\item characters \(\lambda_a\) of \(\operatorname{Gal}(\bar E/E)\)
with values in \(\mathbb Z/\ell^N\), induced by
\(z\mapsto z/\bar z\) from tame characters of level \(a\), trivial
above \(\ell\), and surjective on inertia at \(a\);
\item integral polynomials \(P_i\) and points \(x_i\) on these
variables with coordinates of valuation at least \(u\) and
\(v(P_i(x_i))\le\beta_i+s_0u\).
\end{itemize}
For \(\zeta_a^{\ell^N}=1\), put
\(\psi=\prod_{a\in Q}\zeta_a^{\lambda_a}\).
For each system, \(P_i((\zeta_a-1)_a)\) is its twisted
group-ring evaluation at \(\psi\). For each datum with this evaluation
of valuation \(<\beta_i+1\), that valuation bounds the normalized
length of its divisible-coefficient Selmer group twisted by \(\psi\).
Sections~\ref{sec:selmer-polynomial-descent} and~\ref{sec:terminal-selmer} give a polynomial construction for abstract Selmer and Stickelberger data, with precise hypotheses on the local conditions, scalar inertia and Euler factors. Application to the critical pairs requires those hypotheses for their actual representations and values.

To deduce the general twisting assertion, it suffices to take
\(0<\epsilon<1\), decreasing the prescribed tolerance if necessary.
Choose \(s_0u\le\epsilon/2\), then \(h_0\) for precision
\(\epsilon/2\), powers \(1,\ldots,m\) and bounds \(\beta_i+1\).
Only now choose \(N\ge h_0\) and the corresponding \(Q\) and
polynomials. Every specialized evaluation at \(\psi^j\) has valuation
\(\le\beta_i+\epsilon<\beta_i+1\). At \(a\), all characters but
\(\lambda_a\) are unramified and a generator of tame inertia has
\(\lambda_a\)-value prime to \(\ell\), so its image under \(\psi^j\)
has the order of \(\zeta_a^j\). The conjugate prime has the same
order by anticyclotomicity. Thus the sum of the valuations of these
inertial values minus one is at most \(2|Q|u\le\epsilon/2\).
This proves the required simultaneous specialization, conditional on
the polynomial existence statement. The application to lifts further
requires the asserted Galois representations, local conditions and
identification of the critical ratios with the Stickelberger systems.

\subsection{The central average for the rank-one line}
\label{c4:central-average}
Fix one rank-one lift datum and let \(\phi\) be the lift of the
constant source form. Put \(P=P_{\mathbb C\phi}\); the case
\(\phi=0\) is immediate. The test-form condition says that a form is
integral at every place above \(\ell\) and orthogonal to the false
lifts lying in the kernel of the normalized descent. The clearing
operator on \(\mathbb C\phi\) applies to forms orthogonal to false
lifts lying in \((\mathbb C\phi)^\perp\). These conditions are
connected by a class sum.

Let \(s_c\) run over representatives of the \(N_0\) chosen-level
classes on \(U(1)\), integral above \(\ell\). To arrange integrality,
write a representative as \(a/\bar a\) as below and use weak
approximation in \(E\) to make \(a\) a unit at all places above
\(\ell\) by multiplication by a rational idele from \(E^\times\).
Let \(T(z)\) be central translation on the target multiplied by
the unitary descent character at \(z\); \(T(z)^*=T(z^{-1})\).
For the lift kernel \(K(g)\) the central action is
\(T(z^{-1})K(g)=K(gz)\). This uses the diagonal scalar action in the
Weil representation (see~\cite{weil} for the representation): its character on the defining block is the
inverse of the factor multiplying the translation. We use here the
specified Weil model and its scalar action. Since the stabilizers
are trivial, \(\phi=\sum_c K(s_c)\), and translating all representatives
permutes the classes. Thus \(T(z)\phi=\phi\) and, for every \(g\),
\(\sum_c T(s_c^{-1})K(g)=\phi\).

Write \(S=\sum_c T(s_c)\). If \(z\perp\phi\), the descent of
\(Sz\) vanishes: its value at \(g\) is
\(\langle K(g),Sz\rangle=\langle\phi,z\rangle=0\), with the
usual kernel realization of the normalized descent. Central
translations commute with each finite Hecke operator and with the
residue-type idempotent because both are sums of translations.
Thus \(S\) carries a false lift in \((\mathbb C\phi)^\perp\)
at a fixed place into a false lift in the descent kernel at the same
place. If \(\Xi'\) is a test form, \(S^*\Xi'\) is therefore
orthogonal to the former false lifts. It is also integral at
\(\iota\): the representatives and their inverses are integral
above \(\ell\), so scalar translation preserves the local Weyl
lattice, and the normalizing character is an \(\ell\)-adic
unit there (the finite-order part is a root of unity and the
algebraic-weight factors are local units).
Finally \(PS^*\Xi'=N_0P\Xi'\), since \(S\phi=N_0\phi\).

Suppose the rank-one clearing estimate with exponent \(\gamma\)
has been established. Its operator on the line is a scalar \(\mu\)
with \(|\iota\mu|\ge e^{-\gamma}\); its pairing clause gives
\(|\iota\langle \Xi,\mu P S^*\Xi'\rangle|\le1\) for any form
\(\Xi\) integral at all places above \(\ell\).
Since \(\langle P\Xi,P\Xi'\rangle=\langle\Xi,P\Xi'\rangle\),
\[
 |\iota\langle P\Xi,P\Xi'\rangle|
 \le e^\gamma\,|\iota N_0|^{-1}
 =\exp\bigl(\gamma+v_\ell(N_0)\log\ell\bigr).
\]
Only one clearing scalar occurs. Applying the two projections
separately would lose this coefficient.

\subsection{Summing the explicit loss}
\label{c4:defect-sum}
Write \(h=r^8\), \(M=p_0p_1\), and let \(\Pi\) be the product of the
positive integers in \((b-30h^3,b+30h^3]\), with \(b\) identified
with its integer value. The local allowance for crystalline defects is
\[
 \lambda_\ell=4\cdot10^6 f(m/k)
     \left[bv_\ell(M)+v_\ell\bigl((2b)!\,\Pi^{h^3}\bigr)\right].
\]
For every positive integer \(A\) and finite prime set \(F\) unique
factorization gives \(\sum_{\ell\in F}v_\ell(A)\log\ell\le\log A\).
The interval has at most \(60h^3\) members. Therefore
\begin{align*}
 \sum_{\ell\in F}\lambda_\ell\log\ell
 &\le 4\cdot10^6 f(m/k)
   \bigl[b\log M+2b\log(2b)+60h^6\log(b+30h^3)\bigr]\\
 &\le 4\cdot10^6 e^{180} f(m/k)bL .
\end{align*}
Indeed \(p_i\le b^{11}\), \(r\le3x^5\le x^6\), and
\(b\ge x^{10^{12}}\), so \(60h^6\le b\), \(30h^3\le b\) and
\(\log(2b)\le2L\). The bracket is already at most \(28bL\).
This computes the sum of the proposed local allowances. The
comparison of crystalline cocycles with these allowances for the
main summands enters the generator-wise congruence estimate.

\paragraph{Field and class counts.}
The elementary field-size bounds can be taken directly from the
parameterization. With tower exponent \(s\ge1\), the root order is
\(q_1=q_0\,3^{\lfloor\log_3s\rfloor+5}2^s\), where \(2^s=m^{22}\),
\(s\le32\log x\), \(m\le2\sqrt{x}\), and \(q_0,d\le\log x\).
Thus \(q_1\le 2^{36}x^{11}(\log x)^2\le x^{23/2}\); the latter
inequality follows from \(\log x\ge100\log L\) and
\(\log L\ge64\cdot10^{12}\).
The field \(E=\mathbb Q(\sqrt{-d},\zeta_{q_1})\) has
degree \(n_E=2f\le2q_1\) and embeds in the cyclotomic field of
conductor \(C=q_1d\le x^{13}\). Here the quadratic field is contained
in the \(d\)-th cyclotomic field by the quadratic Gauss sum
(\(d\equiv3\bmod4\)). The discriminant formula for cyclotomic fields
gives root discriminant at most \(C\); the relative discriminant
formula gives the same upper bound for the subfield \(E\).

Minkowski's ideal-class bound gives an integral ideal of norm
\(k_0\le B=2^{n_E}|\operatorname{disc}E|\) in every class.
An ideal \(I\) of norm \(k_0\) contains \((k_0)\); the ideal
\(I/(k_0)\) in \(\mathcal O_E/(k_0)\) is principal, since this
ring is a product of quotients of DVRs. Hence \(I=(k_0,a)\) and
there are at most \(k_0^{n_E}\) choices of \(a\bmod k_0\). We obtain
\[
 \log h_E\le(n_E+1)n_E\log(2C)\le x^{28}.
\]
Let \(N_0\) be the number of chosen-level classes of the rank-one
unitary group of \(E/E^+\). The map from finite ideles of \(E\)
to the norm-one finite ideles given by \(z\mapsto z/\bar z\) is
onto. Locally this is Hilbert's Theorem~90~\cite[II, Corollary~1.23]{milneCFT} (at a split pair, take
\((z,1)\)). At an unramified nonsplit place one can make an antecedent
a unit by multiplying by a power of a base-field uniformizer,
so one obtains a finite idele. For the quadratic local assertion
one may explicitly take the antecedent \(1+y\) for a norm-one
\(y\ne-1\), and a nonzero element of trace zero for \(y=-1\).
The image of \(E^\times\) is rational and an integral unit idele
congruent to \(1\) modulo \(M\mathcal O_E\) maps into the chosen
level: it reduces to \(1\) at each chosen place above \(p_0,p_1\).
Thus the ray class group modulo \(M\mathcal O_E\) surjects onto
the indicated unitary class quotient. Its order is at most
\(h_E|(\mathcal O_E/M)^\times|\le h_E M^{n_E}\). In particular
\[
 \log N_0\le x^{28}+22 n_E L\le x^{30}.
\]
Finally, at the chosen places the group defining \(G_0\) is a
product of \(\operatorname{GL}_{rt}\) over residue fields,
\(t=5r\). At all chosen places over \(p_i\) the sum of residue
degrees is at most \(n_E\). Counting all matrices gives
\(\log G_0\le (rt)^2 n_E\log M\le x^{32}\), since
\(rt\le45x^{10}\), \(n_E\le2x^{23/2}\), and \(L\le x^{1/100}\).
These estimates also apply to every finite sum of prime valuations
of the three integers with logarithmic weights.

\paragraph{A uniform cohomological integer.}
The first-order estimate uses one positive integer \(O_*\) whose
normalized prime valuations control odd Tate-twist classes at all
torsion levels for \(0\le e<t\), for coefficient DVRs containing the roots of unity of order \(\exp((\mathbb Z/C\mathbb Z)^\times)\), with the unramified classes when
\(e=0\). The required bound is
\(O_*\le 2^{4(C+t+3)^{27}}h_E\); Section~\ref{sec:odd-cocycles-reduction} gives a classical derivation for the specified cocycle groups. With this bound, \(\log O_*\le x^{2000}\)
follows from \(C\le x^{13}\), \(t\le15x^5\), and the preceding
class-number estimate. The denominator-error valuation is exactly
\[
 t v_\ell(O_*)+
 \{8ft+2+(4t+1)h^{22}+h\}v_\ell(G_0).
\]
As \(h\le x^{48}\), the coefficient in braces is at most \(x^{1100}\);
the sum of this error with logarithmic prime weights is therefore
at most \(x^{4000}\). The twist error
\(u_\ell=(4f+3)r v_\ell(G_0)+2^{-\ell}\) has weighted sum
at most \(x^{999}\) by the group-order estimate and
\(\sum_{n\ge1} n2^{-n}=2\).

Adding the defect allowance, \((1+2t)u_\ell\), the denominator
error, \(v_\ell(G_0)\) and \(v_\ell(N_0)\) gives a sum with
logarithmic weights at most \(B_0=fxbL^3/(mk)\).
Indeed the defect term divided by \(B_0\) is at most
\(16\cdot10^6e^{180}/L^2\le 1/2\), by \(m^2\le4x\).
All other terms sum to at most \(x^{4002}\le B_0/2\), since
\(1\le f\), \(k\le r+1\le4x^5\), \(m\le2\sqrt{x}\), and
\(b\ge x^{10^{12}}\). This closes the numerical estimate for
the local denominators with the asserted arithmetic inputs.

\subsection{A uniform integer for odd cyclotomic cocycles}
\label{sec:odd-cocycles-reduction}
Write $C=q_1d$, $Y=C+t+3$ and $G_F=\operatorname{Gal}(\overline{\mathbb Q}/F)$.
Let $\chi_\ell$ be the cyclotomic character and $c$ a fixed complex conjugation.
For a finite extension of $\mathbb Q_\ell$ with integer ring $\mathcal O$ put
$A_n=\mathcal O/\ell^n\mathcal O$.
We use the integral valuation $v_{\mathcal O}$ with $v_{\mathcal O}(\pi)=1$.
The module $Z_e(E,n)$ consists of continuous cocycles for $A_n(-e)$
which vanish on inertia away from $\ell$. For $e>0$ they satisfy
$z(cgc^{-1})=-z(g)$ and vanish on the kernel of $\chi_\ell$ in every
decomposition group above $\ell$. For $e=0$ they vanish on inertia
above $\ell$ as well and there is no parity requirement. These are
cocycles before quotienting by coboundaries. We establish the uniform bound in Appendix D using ordinary
cyclotomic class groups at small primes and the refined Coates--Sinnott
theorem at large primes.

\paragraph{A common integer from Bernoulli polynomials.}
For $1\le e<t$ write $r=e+1$, $H=(\mathbb Z/C\mathbb Z)^\times$,
$D_r=(r+2)!(r+1)$, and put
\[
 b(a)=D_r C^r B_r(a/C),\qquad 0\le a<C.
\]
Here $B_j(X)$ is the Bernoulli polynomial with $B_1(0)=-1/2$.
The polynomial $P(X)=(r+2)!C^{r+1} B_{r+1}(X/C)$ has integral
coefficients. Indeed $(j+1)!B_j(0)$ is integral by induction in the
Bernoulli recursion $(j+1)B_j(0)=-\sum_{i<j}{j+1\choose i}B_i(0)$ for $j\ge1$.
Also $P'(a)=b(a)$ and $P(X+C)-P(X)=D_r C X^r$. Telescoping and using
the integral Taylor expansion modulo $(Cm)^2$ give, for each $m\ge1$,
\[
 D_r\sum_{i=0}^{m-1}(a+iC)^r\equiv m b(a)\pmod{m^2}.
\]
For $u\equiv1\bmod C$ prime to $m$, multiplication by $u$ permutes residues
modulo $Cm$ and preserves every character $\psi$ modulo $C$. Replacing
$ua$ by $ua-Cm\lfloor ua/(Cm)\rfloor$ in this permutation, expanding the
$r$th power modulo $(Cm)^2$, multiplying by $D_r$ and using the last
congruence residue class by residue class gives
\[
 (u^r-1)\sum_{a\bmod C}\psi(a)b(a)
 \equiv D_r r C u^{r-1}
 \sum_{a=0}^{Cm-1}\Big\lfloor\frac{ua}{Cm}\Big\rfloor a^{r-1}\psi(a)
 \pmod m .
\]
The division by $m$ here takes place after a congruence modulo $m^2$
in the ring of character values.

Define $g(a)=2b(a)+(-1)^e{\bf1}_{a=1}+{\bf1}_{a=-1}$ for $a\in H$
and $T=(g(v/w))_{w,v\in H}$, using least nonnegative representatives.
Its eigenvalue on a character $\psi$ is
$\lambda_\psi=2\sum_{a\in H}\psi(a)b(a)+(-1)^e+\psi(-1)$.
The identity $\zeta(1-r,x)=-B_r(x)/r$ expresses the sum as
$-D_r C r L(-e,\psi)$. This value vanishes exactly for
$\psi(-1)=(-1)^e$. Indeed the functional equation for the primitive
inducing character reduces nonvanishing in the other parity to the
Euler product at $1+e>1$, and the extra Euler factors at $-e$ are nonzero.
Thus $\det T\ne0$. The adjugate identity evaluated at the identity
entry of the character vector shows $\lambda_\psi\mid\det T$ in the
ring of character values, including at primes dividing $|H|$.

Set
\[
 X_e=|\det T|\prod_{u\in\{C+1,\,2C+1\}}
              (u^e-1)(u^{e+1}-1).
\]
The two choices of $u$ are coprime; hence at each prime $\ell$ there
is a choice prime to $\ell$ for which both factors jointly divide $X_e$.
If $\psi(-1)=-(-1)^e$, choose $v$ with $X_e=v\lambda_\psi$.
For every $u\equiv1\bmod C$ and every $m\ge1$ coprime to $u$, the resulting congruence is
\[
 (u^{e+1}-1) X_e\equiv
 2w\sum_{a=0}^{Cm-1}\Big\lfloor\frac{ua}{Cm}\Big\rfloor a^e\psi(a)
 \pmod m,\qquad w=vD_r r C u^e .
\]
The multiplier $w$ is independent of $m$.

Here $1\le X_e\le 2^{Y^6}$. In fact the Fourier series of $B_r$ for
$r\ge2$ gives $|B_r(x)|\le 2r!\zeta(r)/(2\pi)^r\le r!$ on $[0,1]$.
As $r+2\le Y$ and $C\le Y$, the entries of $T$ have modulus at most
$4Y^{3Y+1}$. The Leibniz formula bounds its determinant by
$Y^Y(4Y^{3Y+1})^Y$; the four extra factors contribute at most
$(2Y+1)^{4r}$. Their product has binary logarithm below $Y^6$
for $Y\ge 8$, as holds here.

\paragraph{Annihilation by classical Stickelberger reciprocity.}
For each $\psi:H\to\mathcal O^\times$ write $\alpha=\chi_\ell^e\psi$,
inflating $\psi$ to $G_{\mathbb Q}$. Let $V_\psi$ be the continuous
cocycles on $G_{\mathbb Q}$ for $\alpha^{-1}$ with values in $A_n$
which vanish at $c$ and at those elements of $G_{\mathbb Q(\mu_C)}$
lying in inertia away from $\ell$ or in a decomposition group at
$\ell$ with trivial cyclotomic character. Suppose $\alpha(c)=-1$.
We claim $2X_e^2 V_\psi=0$. We may take $n\ge1$.
Set $F=C\ell^n$, $K'=\mathbb Q(\mu_F)$, and choose
$u\in\{1+C,1+2C\}$ coprime to $\ell$.
If $z\in V_\psi$, the restriction $z|_{G_{K'}}$ is a homomorphism:
$\alpha$ is trivial on $G_{K'}$ modulo $\ell^n$. It vanishes on
inertia outside $\ell$. At a place above $\ell$, choose $\tau$ in
local inertia with $\chi_\ell(\tau)=u$. This is possible because the
local $\ell$-power cyclotomic extension is totally ramified with full
group $\mathbb Z_\ell^\times$. Inertia fixes prime-to-$\ell$ roots of
unity and $u\equiv1\bmod C$, hence $\psi(\tau)=1$.
If $h$ belongs to the local decomposition subgroup of $G_{K'}$,
then the commutator $\tau h\tau^{-1}h^{-1}$ has cyclotomic character one
and fixes $\mu_C$. Its $z$-value is zero by the stated local condition;
the cocycle rule evaluates that value as $(u^{-e}-1)z(h)$ in $A_n$.
It follows that $y=(u^e-1)z|_{G_{K'}}$ factors by global reciprocity
through the ordinary class group $\operatorname{Cl}(K')$.
No archimedean modulus occurs because $K'$ is totally imaginary.

Put $\theta_F=\sum_{b\in(\mathbb Z/F)^\times}(b/F)\sigma_b^{-1}$,
where $0<b<F$ and $\sigma_b$ raises roots of unity to the $b$th power.
The classical Stickelberger theorem annihilates $\operatorname{Cl}(K')$
by the integral group-ring element
\[
 (u-\sigma_u)\theta_F=\sum_{b\in(\mathbb Z/F)^\times}
                 \lfloor ub/F\rfloor\,\sigma_b^{-1}
\]
(L.~C.~Washington, \emph{Introduction to Cyclotomic Fields}, 2nd ed.,
Thm.~6.10~\cite{washington}). The extension is abelian over $\mathbb Q$ and $(u,F)=1$,
as required by that theorem. Conjugation by a lift of $\sigma_b^{-1}$
on $G_{K'}$ multiplies the value under $y$ by $b^e\psi(b)$ in $A_n$;
this is the cocycle identity and the fact that the conjugated element
acts trivially. Reciprocity is equivariant for field automorphisms.
Consequently $U_n y=0$, where $U_n$ denotes the sum of
$\lfloor ub/F\rfloor b^e\psi(b)$ over the units modulo $F$.

Write $T_i$ for the full floor sum at modulus $C\ell^i$. Terms divisible
by $\ell$ give $\ell^e\psi(\ell)T_{n-1}$; terms noncoprime to $C$ vanish.
Thus $U_n=T_n-\ell^e\psi(\ell)T_{n-1}$. Apply the Bernoulli congruence
at $m=\ell^n$ and $\ell^{n-1}$ with the same $w$. Since $e\ge1$ it gives
\[
 2wU_n\equiv(1-\ell^e\psi(\ell))(u^{e+1}-1)X_e\pmod{\ell^n}.
\]
The first factor on the right is a unit. Hence
$(u^{e+1}-1)X_e(u^e-1)z|_{G_{K'}}=0$. Since
$(u^{e+1}-1)(u^e-1)\mid X_e$, the cocycle $X_e^2 z$ is zero
on $G_{K'}$. It therefore descends to the abelian group
$\operatorname{Gal}(K'/\mathbb Q)$. Comparing its values on $gc$ and
$cg$, using $z(c)=0$ and $\alpha(c)=-1$, gives $2 X_e^2z=0$.

\paragraph{Integral character descent.}
Suppose $\mathcal O$ contains the roots of unity of order
$\exp H$. Choose lifts $\sigma_a\in G_{\mathbb Q}$ for $a\in H$ with
$\sigma_1=1$, and put $\delta(g)=cgc^{-1}g^{-1}$. All conjugates
of $\delta(g)$ belong to $G_{\mathbb Q(\mu_C)}\subset G_E$ and have
trivial cyclotomic character. For $z\in Z_e(E,n)$ set
\[
 z_\psi(g)=\sum_{a\in H}\alpha(\sigma_a)\,
       z(\sigma_a\delta(g)\sigma_a^{-1}).
\]
The summand does not depend on the choice of lift: replacing $\sigma_a$
by $h\sigma_a$ for $h\in G_{\mathbb Q(\mu_C)}$ multiplies the $z$-value
by $\chi_\ell(h)^{-e}$ and its scalar by $\chi_\ell(h)^e$.
This uses the cocycle identity for conjugation of an element acting
trivially on the coefficient module. Now
$\delta(gg')=\delta(g)\,g\delta(g')g^{-1}$.
Right multiplication permutes the cosets of the lifts, so
$z_\psi(gg')=z_\psi(g)+\alpha(g)^{-1}z_\psi(g')$.
It vanishes at $c$. If $g$ belongs to one of the designated local
sets in $G_{\mathbb Q(\mu_C)}$, both its conjugates by $\sigma_a$
and by $\sigma_a c$ lie in the corresponding local sets of $G_E$,
so their quotient has zero cocycle value. Hence $z_\psi\in V_\psi$.

For $h\in G_E$ orthogonality of the characters gives
\[
 \sum_\psi z_\psi(h)=|H|z(\delta(h))
              =|H|\{z(chc^{-1})-z(h)\}=-2|H|z(h).
\]
If all odd components vanish, the left side is invariant under
$h\mapsto chc^{-1}$ because each even component is zero at $c$ and
has $\alpha(c)=1$. Consequently $4|H|z=0$.
Applying this fact to $2X_e^2 z$ proves in particular that $B_e=16|H|X_e^2$
annihilates $Z_e(E,n)$ for all $\ell$, including $\ell=2$.

\paragraph{A length bound at large primes by global duality.}
Fix $\ell>Y^2$ and $e>0$, and set $j=e+1$, $\varepsilon=(-1)^j$.
First take $\mathcal O=\mathbb Z_\ell$. Let $S_0$ consist of the
$\ell$-adic and archimedean places of $E$.
Write $M_n=\mathbb Z/\ell^n(-e)$ and $M_n^D=\mathbb Z/\ell^n(j)$
for its Cartier dual. Passage to cohomology gives an injection
\[
 Z_e(E,n)\lhook\joinrel\longrightarrow
 \ker\left[H^1(G_{E,S_0},M_n)\longrightarrow
           \bigoplus_{v\mid\ell}H^1(E_v,M_n)\right]^{c=\varepsilon}.
\]
Indeed coboundaries are invariant under precomposition by
$g\mapsto cgc^{-1}$, whereas our cocycles change sign, and $2$
is invertible. The inertia conditions give descent to $G_{E,S_0}$
because the coefficient action is unramified away from $\ell$.
The natural action of $c$ on cohomology includes its action
$(-1)^e$ on coefficients, yielding the displayed parity. At
$v\mid\ell$ a cocycle factors through the image of $\chi_\ell$ by
the local vanishing condition. This image is $\mathbb Z_\ell^\times$
because $E_v$ is unramified over $\mathbb Q_\ell$ ($\ell>C$).
On this abelian group, commutation with an element $b$ for which
$b^{-e}-1$ is a unit makes every cocycle a coboundary: equating
$z(bg)$ and $z(gb)$ gives $(b^{-e}-1)z(g)=(\chi_\ell(g)^{-e}-1)z(b)$.
Such $b$ exists since $0<e<\ell-1$. Complex places contribute zero.

Poitou--Tate duality for finite modules pairs this localization
kernel perfectly with the localization kernel in
$H^2(G_{E,S_0},M_n^D)$; the pairing respects field automorphisms,
so the same sign occurs on the dual. In particular
\[
 \operatorname{length}_{\mathbb Z_\ell}Z_e(E,n)
 \le \operatorname{length}_{\mathbb Z_\ell}
       H^2(G_{E,S_0},\mathbb Z_\ell(j))^{c=\varepsilon}.
\]
Here the integral group surjects onto finite-level $H^2$:
the finite-module groups have $\ell$-cohomological dimension at most
two in this setting, and the finite cohomology groups in the inverse
system satisfy the Mittag--Leffler condition. Applying the sign
idempotent preserves surjectivity.

Let $S$ consist of $\infty$ and the primes of $\mathbb Q$ dividing $C$,
and write $U=\operatorname{Spec}\mathcal O_{E,S}[1/\ell]$.
The integral group just used injects into
$H^2_{\mathrm{\acute et}}(U,\mathbb Z_\ell(j))$. Indeed Galois cohomology
with ramification allowed at $S_0$ computes the cohomology of
$\operatorname{Spec}\mathcal O_E[1/\ell]$ for these coefficients.
By the localization sequence and purity the terms on the left when
the places $v\mid C$ are removed are $H^0(k_v,\mathbb Z_\ell(j-1))$.
They vanish because $|k_v|^{j-1}-1\ne0$.
These uses of finite-module duality and localization are the usual
Poitou--Tate theorem and the localization sequence for an arithmetic
curve (see J.\,S.~Milne, \emph{Arithmetic Duality Theorems}, 2nd ed.,
Parts I and II~\cite{milneADT}).

Apply the refined Coates--Sinnott theorem of C.~Greither and
C.\,D.~Popescu, \emph{An equivariant main conjecture in Iwasawa
theory and applications}, J.~Algebraic Geom., Thm.~6.11~\cite{gp}, to
$E/\mathbb Q,S,\ell,j$. This is an abelian CM extension of the
totally real field $\mathbb Q$, $j\ge2$, $S$ contains every ramified
prime, and $\ell$ is odd. The required $\mu$-invariant vanishes by
the Ferrero--Washington theorem~\cite{fw} since $E(\mu_\ell)$ is abelian over
$\mathbb Q$. Further, $H^1_{\mathrm{\acute et}}(U,\mathbb Z_\ell(j))$
has zero $\ell$-torsion by the coefficient exact sequence: already
$H^0(E,\mathbb Z/\ell(j))=0$ because $0<j<\ell-1$ and local inertia
maps onto $(\mathbb Z/\ell)^\times$. The theorem therefore identifies
the zeroth Fitting ideal of the $\varepsilon$-part of $H^2$ with the
ideal generated on that part by $\Theta_{E/\mathbb Q,S}(1-j)$.
Indeed its sign idempotent for this base field is
$(1+(-1)^jc)/2$.

To read its length, extend scalars to the integers $\mathcal R$ of
an unramified splitting field of $\operatorname{Gal}(E/\mathbb Q)$
over $\mathbb Q_\ell$. Such a field suffices since the group order
divides $\varphi(C)<\ell$. Its group ring is a product of copies of
$\mathcal R$. Base change preserves Fitting ideals, and the normalized
length of the extended module equals its original $\mathbb Z_\ell$-length.
Over each discrete valuation ring the valuation of the Fitting
generator is the length. Thus the required upper bound is the sum of
the valuations of $L_S(-e,\psi^{-1})$ over the characters of
$\operatorname{Gal}(E/\mathbb Q)$ of sign $\varepsilon$. Inflating
them to $H$ gives exactly the $L$-functions at modulus $C$, since
$S$ removes every prime dividing $C$. For the relevant parity their
values equal the eigenvalues $\lambda_{\psi^{-1}}$ of $T$ divided by
$-2D_r C r$, an $\ell$-adic unit. All other eigenvalues of the
integer matrix $T$ are $\ell$-integral. Consequently
\[
 \operatorname{length}_{\mathbb Z_\ell}Z_e(E,n)
 \le v_\ell(\det T)\le v_\ell(X_e).
\]
For general $\mathcal O$ the cocycle equations and local and parity
conditions hold coordinatewise in a $\mathbb Z_\ell$-basis of
$\mathcal O$. Hence $Z_e(E,n;\mathcal O)$ is the scalar extension
of $Z_e(E,n;\mathbb Z_\ell)$, and the length bound becomes
$v_{\mathcal O}(X_e)$, including for ramified coefficient rings.

\paragraph{Lengths at small primes.}
For $\ell\le Y^2$ a uniform detecting set comes from class field theory.
Put $M=C\ell^2$, $K=\mathbb Q(\mu_M)$, $d_K=[K:\mathbb Q]\le M$.
The maximal exponent-$\ell$ abelian quotient of $G_K$ unramified away
from $\ell$ has $\mathbb F_\ell$-rank at most
\[
 \log_2 h_K+\sum_{v\mid\ell}([K_v:\mathbb Q_\ell]+1)
 \ \le\ \log_2 h_K+2M .
\]
Indeed global reciprocity identifies the quotient by the local inertia
images with a quotient of the ideal class group, and local reciprocity
maps the local units onto abelian inertia. Their pro-$\ell$ factors
have $\mathbb Z_\ell$-rank $[K_v:\mathbb Q_\ell]$ and cyclic torsion,
also at $\ell=2$. Choose lifts of generators and adjoin representatives
for $G_E/G_K$, at most $M$ more points.
For $z$ killed by $\ell$, restriction to $G_K$ is a homomorphism
since $K$ contains $\mu_\ell$. Thus vanishing at the points forces
$z=0$. For a general $z$ vanishing there, induct on its exponent,
first applying induction to $\ell z$ and then the exponent-$\ell$ case.
This works for ramified coefficient rings because $\ell z=0$ makes
the twist trivial on every value over $G_K$.

For completeness, Minkowski's bound and $|\operatorname{disc}K|\le M^{d_K}$
give an integral representative of norm $j\le(2M)^{d_K}$ for each
ideal class. Every ideal of norm $j$ contains $j\mathcal O_K$ and
is generated by $j$ and one other element, since the quotient by
$j\mathcal O_K$ is a product of principal local rings. There are at most
$j^{d_K}$ choices. It follows that $\log_2 h_K\le
d_K(d_K+1)\log_2(2M)$. Since $M\le Y^5$, the total number of points
is at most $Y^{20}$ (use $\log_2(2M)\le M$ for $M\ge2$ and $M\ge4$).
The class field theory used here is the idelic reciprocity theorem
and the Hilbert class field theorem; see J.\,S.~Milne,
\emph{Class Field Theory}, chapters V and VIII~\cite{milneCFT}.

At small primes evaluation thus embeds $Z_e(E,n)$ in at most
$Y^{20}$ copies of $A_n[B_e]$. Each copy has length at most
$v_{\mathcal O}(B_e)$. At large primes the preceding stronger
bound $v_{\mathcal O}(X_e)$ is at most $Y^{20}v_{\mathcal O}(B_e)$.
For $e=0$ class field theory identifies
$Z_0(E,n)$ with $\operatorname{Hom}(\operatorname{Cl}(E),A_n)$;
decomposing the class group into cyclic summands bounds its length by
$v_{\mathcal O}(h_E)$. Therefore
\[
 O_*=h_E\left(\prod_{1\le e<t}16\varphi(C)X_e^2\right)^{Y^{20}}
 \quad\text{satisfies}\quad
 \operatorname{length}_{\mathcal O} Z_e(E,n)\le v_{\mathcal O}(O_*)
\]
for every coefficient ring with the stated roots and every torsion level. Since each factor is
at most $2^{4Y^6}$ and there are at most $Y$ factors,
$O_*\le h_E2^{4Y^{27}}$. Dividing by $v_{\mathcal O}(\ell)$
gives the normalized valuation convention used in the denominator
estimate. This calculation makes the integer independent of
$\ell,\mathcal O,e,n$ and retains the factor $h_E$ only once.


\section{Construction hypotheses for the arithmetic comparison}\label{sz:register}
At a fixed admissible datum, Theorem~\ref{sz:conclusion} deduces the signed
inequality from A, B and C. The sharpenings below apply on the constructed
family as $x$ grows with the permissible parameters and scale exponent fixed.

\begin{longtable}{@{}>{\raggedright\arraybackslash}p{.19\linewidth}@{\hspace{.025\linewidth}}p{.755\linewidth}@{}}
\toprule Step & Deductions and remaining requirements\\\midrule\endhead
Local Taylor bound &
Sections~\ref{sec:unitrow-frame}--\ref{sec:other-witnesses} construct the frame and coefficients at split, inert and ramified primes. The local CM identities, the period normalization and the realization by integral test forms enter separately.\\[4pt]
Spreading &
Section~\ref{app:local-frame-recovery} gives a compact-position frame bound. Section~\ref{spread:normalization} computes the scalar in F6. Sections~\ref{moment:rigidity}--\ref{moment:sieve} derive frame separation, a weighted pair count, the primary character formula and the transfer of local errors into both moments. The actual filling still requires the one-term multiplicity, local kernel bounds, orbital integral comparison and compatible measures. Proposition~\ref{prop:weakerprecision} gives an alternative sufficient accuracy.\\[4pt]
Inner forms &
Section~\ref{tr:parabolic-descent} proves local vanishing on a coprime factorization. The global pretrace expansion and the matching of elliptic terms in \S\ref{tr:geometry} remain inputs.\\[4pt]
Twisted lengths &
Sections~\ref{sec:selmer-polynomial-descent}--\ref{sec:terminal-selmer} construct polynomials for the indicated data. Section~\ref{sec:odd-cocycles-reduction} treats the common cohomological integer. Section~\ref{moment:line} verifies the rank-one Selmer conditions for the explicit main line and computes its split vacuum factor. Attaching the actual representations, identifying the local conditions and proving the first-order congruences remain necessary in R.\\[4pt]
Auxiliary forms &
Section~\ref{pathmath:dual} computes the integral pairing correction. The denominator calculation in \S\ref{pathmath:denom} requires the stated kernel and period identities. The required spanning and size of the auxiliary forms remain part of Aux.\\\bottomrule
\end{longtable}
\paragraph{Hypothesis groups on the constructed scales.}
The original register has four groups in $\mathsf A_{\rm loc}$, eight in
F0--F7, and six in I, R, G, F, K, Aux. For the construction at a sufficiently
large cutoff, \S\ref{moment:rigidity} derives F4 from the fitted model and
free principal level. This leaves seventeen other groups to verify on the
selected datum. The groups share inputs, as described below. Several contain substantial written deductions. In
particular the finite character identification in F7 is now
Lemma~\ref{lem:fitted-primary-character}; its inner-form comparison and
formal-degree identification are still required. The simultaneous choice
of datum in Corollary~\ref{cor:construction-grh} applies to all branches.

\subsection*{Where the written proof still depends on construction estimates}
\paragraph{Shortness and integrality.}
The archimedean chain requires the normalized vacuum and period identities,
their growth and expansion bounds, and the integral branching data. At
small nonsplit primes Appendix~\ref{sec:jets} derives a divisibility
estimate from the stated CM and vacuum identities. The marked places and
the level prime also require the transfer, edge and paired-lift inputs
specified in $\mathsf A_{\rm loc}$. These inputs supply both the gain
removed by $c_B$ and the colength of the submodule where the gain holds.

\paragraph{Spreading.}
The restriction to ramified forms uses the loss of at most
$f d_X/(N-1)$ established at the source. The full-source moment still
requires the fitted filling and Bessel block normalizations, the same-link
and unequal-link pair estimates and the remaining
primitive frame estimates of Appendix~\ref{sz:shells}, and the Siegel and
sieve estimates for the open cell. The near-optimal coefficient in F6 is
used in the dimension bound. On the constructed scales the box, separation and orbit statements in F4 follow in the fitted free-level model by \S\ref{moment:rigidity}. The transfer to the good target uses the
degree matching in F7 and the finite character formula of Lemma~\ref{lem:fitted-primary-character} together with Appendix~\ref{sec:trace}.

\paragraph{Determinant height.}
The regular twist uses Galois, Selmer and first-order congruence statements
in R. Appendix~\ref{sec:finite} supplies simultaneous specialization,
the explicit prime sum for the allowance and the passage to test forms
under its stated arithmetic inputs. The compression of the one-lift
operator still requires the normalized period gluing and exchanges in G.
The term $-A_0$ uses the flat value formula, the local curve hypotheses for the sharp witness and
its realization by test forms in F. K supplies the analytic norm formula
and triangular exchanges. Aux supplies the normalized path period, its
surjectivity, and the denominator and size of the grid forms. These are
the inputs to the one-lift height, its telescoping through $m$ lifts and
the extension of the given family at controlled cost.

\paragraph{Three dependencies shared between the branches.}
The short-basis argument is used at two different levels. The source
construction in A and the grid in Aux have different box counts and
require their own conductor and norm estimates. The inner-form comparison
in B at the identity transfers a dimension. In C a determinant transfer
requires the traces of the Hecke operators as well. The common geometric
hypotheses and the extra step for the determinant are in
Appendix~\ref{sec:trace}. Finally, a universal lower bound on the
valuation of a short-form jet in A does not exhibit any coefficient
attaining a small valuation in C. Appendix~\ref{sec:jets} specifies the
additional witness and the arithmetic cases it requires.

The remaining
mathematical work for this exposition lies upstream of \eqref{sz:signed}.
Theorem~\ref{sz:conclusion} applies datum by datum.
For the consequence about Dirichlet zeros, Corollary~\ref{cor:construction-grh}
requires the listed inputs on one constructed datum at each sufficiently
large integer length, for each fixed $(q_0,d,\eta)$.

\section{Normalization of the open-cell moment}

\subsection{The numerical coefficient in the open-cell estimate}
\label{spread:normalization}
We compute the scalar comparison following the Siegel and sieve estimates of
\S\ref{B:openenergy}.  Write \(F=E^+\), \(f=[F:\mathbb Q]\),
\(u=\varrho\), \(n=u+1\), \(r=k+2u+1\), and \(D_E=|\operatorname{disc}E|\).
Let \(m_j\) be the mass of the genus of the standard definite
Hermitian lattice \(\mathcal O_E^j\), with each class weighted by the
reciprocal of the order of its unitary automorphism group.

\paragraph{The consecutive masses.}
The quadratic extension \(E/F\) is unramified at every finite place.
Indeed, complex conjugation acts nontrivially on both
\(\mathbb Q(\zeta_{q_1})\) and \(\mathbb Q(\sqrt{-d})\). Their
conductors are coprime. At any rational prime at least one of the two
fields is unramified, and inertia fixes that field pointwise. Thus the
inertia subgroup in \(E/\mathbb Q\) has trivial intersection with
\(\operatorname{Gal}(E/F)\). In particular
\(a_0:=D_E^{1/(2f)}=|\operatorname{disc}F|^{1/f}\).
Let \(\chi\) be the quadratic Hecke character of \(E/F\).
For \(j\ge1\) one has
\[
 \frac{m_{j+1}}{m_j}
 =\left(a_0^{j+1/2}\frac{j!}{(2\pi)^{j+1}}\right)^f
   L_F(j+1,\chi^{j+1}).
\]
Indeed the standard lattice is self-dual at every finite place, hence
maximal integral there. At a nonsplit finite place the extension is
unramified, and every unit is a norm, including the signed determinant
of the standard form. Thus all local correction factors in the mass
formula are one. Propositions 4.4 and 4.5 and (4.2)--(4.3) of
W.~T. Gan, J.~P. Hanke and J.-K. Yu, \emph{On an exact mass formula of
Shimura}, Duke Math. J. \textbf{107} (2001), 103--133~\cite{ghy}, give
\(m_l=2^{1-fl}\prod_{h=1}^{l}|L_F(1-h,\chi^h)|\).
Taking the ratio leaves only \(2^{-f}|L_F(-j,\chi^{j+1})|\).
The completed Hecke functional equation~\cite{tate} has finite conductor one and
at each real place the gamma factor
\(\Gamma_{\mathbb R}(s+e)=\pi^{-(s+e)/2}\Gamma((s+e)/2)\),
where \(e\equiv j+1\pmod 2\), \(e\in\{0,1\}\).
The quotient of the absolute gamma factors at \(j+1\) and \(-j\)
is \(2j!/(2\pi)^{j+1}\) by the reflection and duplication
formulas. This proves the asserted ratio.

Choose a generator \(\omega\) of the different as in the fitted
measures, and put \(w_v=|\omega|_v\) at each complex place
(one ordinary modulus for each conjugate pair). Then
\(\prod_v w_v=a_0^f\). The mass assigned to the unit sphere in
rank \(i+1\) is
\(s_{i+1,v}=(2\pi)^{i+1}/(w_v^{i+1/2}i!)\).
Put \(P=N^{kf}\prod_{i=k}^{r-1}\prod_v s_{i+1,v}\).
The mass ratio just proved yields the exact identity
\[
 \frac{P m_r}{N^{kf}m_k}
   =\prod_{i=k}^{r-1} L_F(i+1,\chi^{i+1}).
\]
For real \(s\ge3\), the Euler products give
\( |\log L_F(s,\chi^l)|\le f\log\zeta(s)\le 2f\,2^{-s}\).
The last inequality follows by bounding the terms after \(2^{-s}\)
by \(\int_2^\infty t^{-s}dt\). Hence the logarithm of the displayed
product is at most \(2f\,2^{-k}\) in absolute value.
At the constructed scales, \(2f\,2^{-k}+2f/N\le \log(1+b^{-9})\)
eventually: \(f\) is polynomial in \(x\), \(k\asymp x^5\),
\(\log b\le x^\eta+1\), and \(\log N\ge m\log b\) with
\(m\ge\sqrt{x}\).
Using \(-\log(1-N^{-k})\le 2/N\), we obtain
\begin{equation}\label{spread:frame}
 Pm_r\le (1+b^{-9})(N^k-1)^f m_k .
\end{equation}

\paragraph{The factors in the comparison.}
Here are the normalizations required of the filling. Write \(h\) for
its number of level classes, \(e\) for the number of tensor-product
Bessel coordinates, \(H\) for the order of the reduction of the
unitary subgroup on the rank-\(k\) part, and \(J=u(k+u)\).
The reduced unipotent has order \(U=M^{f(2J+u)}\), the reduced
Bessel group has order \(HU\), and \(d_X=he\). The level class
formula is \(h=H(N^k-1)^f m_k\). To obtain it from the fitted
level, the row congruence at each of the \(f\) split level places
has index \(N^k-1\), since \(\mathrm{GL}_k(\mathbb F_N)\) acts
transitively on nonzero rows. The remaining principal level
has index \(H\) by reduction. The mass multiplies by the level
index by orbit--stabilizer; at the fine level the stabilizers
are trivial. This uses the stated identification of the level
and surjectivity onto the respective reduction groups.

Write \(V=(2/a_0)^f\), \(c_j=2\pi^j/(j-1)!\),
\(v_j=\prod_{l=1}^j c_l\), \(c=c_k v_k/v_r\) and
\(\Theta=\pi^{(2u+1)k}v_k/v_r\). Put
\[
 g_v=Mw_v/(2\pi),\quad t_{vi}=g_v^{2(u-i)}/N,\quad
 \Lambda=\prod_{v,i} t_{vi}^{k+2(u-i)},\qquad 0\le i\le u.
\]
In the measures of the filling the scalar quantities take the values
\[
\begin{aligned}
 C_{{\rm Pois},i}&= V/\prod_v t_{vi}^2,&
 R_{k-1}&=c^{f/2}e^{-1}\prod_{v,i}t_{vi}^{k-1+2(u-i)},\\
 c_S&=\left(\frac{\pi^k}{(k-1)!N^k}\right)^f D_E^{1/4}V^k,&
 C_{\rm S}&=(HU)Uh\,c_S,\\
 \nu&=U N^{-kf} V^{-J} D_E^{u/4}/\Lambda,&
 C_{\rm dec}&=\prod_v\Theta
  \left(\frac{g_v^{2u+1}}{MN}\right)^{k-n^7}.
\end{aligned}
\]
In this display \(e^{-1}=1/e\) denotes the reciprocal coordinate
count.  Set \(W=V^n(e/\Lambda)^2/c^f\). Cancellation of one
factor \(t_{vi}\) gives \(\prod_i C_{{\rm Pois},i}=R_{k-1}^2 W\).
The relation between the rank-\(k\) sphere and the frame is
\[
 P c^f D_E^{u/2}=c_S N^{2kf}V^{n+2J}.
\]
Indeed, \(s_{i+1,v}=c_{i+1}2^i/w_v^{i+1/2}\). Multiplication over
\(i=k,\ldots,r-1\) cancels \(v_k/v_r\), with
\(\sum_{i=k}^{r-1}i=(2u+1)(k+u)\). Substituting
\(\prod_v w_v=a_0^f\) proves the relation. It follows that
\[
 \frac{C_{\rm S}W}{\delta I_*\nu^2}
 = \frac{H P m_r}{h}\le 1+b^{-9},
\]
where \(I_*=(he)^2/(\delta m_r)\) and the inequality is
\eqref{spread:frame}. In particular the leading constant survives
the cancellation of all frame scales.

\paragraph{The additive sieve error.}
Write \(s=\sqrt{a_0}\), \(Q=Ms^2/(2\pi)\), and
\(X=Q^{2u+1}/(MN)\). Dividing by \(I_0=HUhe\) and inserting the
constants above gives
\[
 \frac{C_{\rm S}W C_{\rm dec}}{e I_0\nu^2}
 =\left[
  (2/s^2)^{n+2u^2}
  \frac{s^{1-2u}}{2 M^{u(2u+1)}} X^{-n^7}
 \right]^f .
\]
For verification, the quotient before simplification is
\(c_S V^{n+2J} N^{2kf} C_{\rm dec}/(U c^f D_E^{u/2})\).
Use \(\Theta/c=\pi^{(2u+1)k}/c_k\) and
\(C_{\rm dec}=\Theta^f X^{f(k-n^7)}\). The factor \(X^{fk}\)
cancels the remaining factorial and powers involving \(k\).
Here \(Q\ge1\), since each fitted \(g_v\ge1\), and
\(X^{-1}\le MN\le b^{23+60n}\). Since \(s\ge1\), \(u\ge1\),
and \(n+2u^2+(23+60n)n^7\le n^9\) for \(n\ge10^5\), the
quotient is at most \(b^{n^9 f}\). The scale bounds
\(e\le b^{2kf}\), \(b^{11kf}I_0\le I_*\) and \(512000n^9\le k\)
therefore give, for \(t=|T|\le f\),
\[
 C_{\rm S}W b^{2403n^9 f} C_{\rm dec}
    \le b^{2404n^9 f+2kf-11kf}I_*\nu^2
    \le b^{-10-6kt}I_*\nu^2 .
\]
These size bounds are among the fitting estimates. In particular the
last inequality uses \(2404n^9f+10\le 3kf\).

Finally \(f\le b/1000\) eventually, so
\((1+b^{-10})^f\le\exp(b^{-9}/1000)\le 1+b^{-9}\).
After division by \(R_{k-1}^2 b^{-6kt} I_*\nu^2\), the two block
bounds of \S\ref{B:openenergy} now give at most
\[
 (1+z^6)\{(1+z^9)^2(1+z^6)+z^{10}\}\le 1+z^4,
 \qquad z=b^{-1}\le10^{-3}.
\]
To see the last inequality, put \(y=z^4/100\), replace each of
\(z^6,z^9,z^{10}\) by \(y\), and use
\((1+y)^4+y(1+y)\le1+17y\) for \(0\le y\le1\).
The same comparison is valid for the empty fixer. This proves the
scalar estimate from the displayed filling normalizations. After the parametrization in \S\ref{moment:open}, the Siegel estimate for the block sum,
the finite-place and archimedean kernel estimates giving the sieve
bound, and the identification of the fitted measures and level
with those used above remain hypotheses of the application.


\paragraph{Removing the common entries in the sieve.}
There is a separate deduction behind the exponent \(2403n^9f\) in
the bound for \(Z\). Let \(A\) be the finite set of possible
nonzero entries, integral units outside the finite places used in
the sieve. Write \(A^i\) for the possible first \(i\) entries of
a head. Let \(J_i(p,q)\) be its level-\(i\) kernel, with all
unwritten entries padded by one and all entries taken with their
prescribed local scales. For an entry \(y\) of the local product,
put
\[
 F_i(y)=\sum_{p,q\in A^i} J_{i+1}((p,y),(q,y)),\qquad
 D_{i+1}=\sum_{t\in A} F_i(t),\quad D_0=\Re J_0((),()) .
\]
The diagonal kernel on a common entry is positive semidefinite,
so \(F_i\) is real and nonnegative. Assume the local integration
identity \(\int F_i=\Re\sum_{p,q}J_i(p,q)\) and let \(w_i(p)\ge0\)
be weights. More precisely, suppose the following estimates hold
in the measures giving Poisson constant \(P_i=C_{{\rm Pois},i}\):
\begin{enumerate}
\item One-sided Poisson summation applies to \(F_i\), with total
absolute value of its nonzero Fourier coefficients at most
\(\epsilon(\sum_{p\in A^i}w_i(p))^2\) in the dual sum.
\item For heads \(p,q\in A^i\) with distinct last entries,
\( |J_i(p,q)|\le\epsilon w_i(p)w_i(q)\).
\item \((\sum_{p\in A^i}w_i(p))^2\le
  \Gamma\prod_{j<i}P_j\) for \(0\le i<n\).
\end{enumerate}
The first estimate, with the zero frequency evaluated by the
integration identity, bounds \(D_{i+1}\) by
\(P_i(\Re\sum_{p,q}J_i(p,q)+\epsilon(\sum_p w_i(p))^2)\).
In this unrestricted head sum, the equal-last-entry contribution is
\(D_i\); the remaining pairs contribute at most
\(\epsilon(\sum_p w_i(p))^2\) by the second estimate.
At \(i=0\) there are no remaining pairs. Thus
\[
 \frac{D_{i+1}}{\prod_{j\le i}P_j}
 \le \frac{D_i}{\prod_{j<i}P_j}+2\epsilon\Gamma .
\]
The top sum restricted to allowable last entries satisfies \(Z\le D_n\)
by nonnegativity. Iteration gives
\(Z\le(\prod_{j<n}P_j)(D_0+2n\epsilon\Gamma)\).
The fitted kernel bounds needed here are
\[
 D_0\le \frac{(1+b^{-10})^f}{\delta}
           \frac{1+b^{-6}}{b^{6k|T|}},\qquad
 \epsilon=C_{\rm dec}\,b^{2002n^9f},\qquad
 \Gamma=b^{400n^9f}.
\]
Since \(2n\le b^{n^9 f}\), these bounds prove the stated sieve
inequality. They isolate the mass at level zero and the two
different uses of decay. For unequal last entries \(t,s\), the
finite frequency is \((t^{-1}-s^{-1})/M\), a nonzero element
integral outside the chosen places when \(1/M\) is integral there.
A pointwise decay bound at those frequencies supplies item 2.
Section~\ref{moment:sieve} proves the weight bound and sums an explicit
majorant over the full dual lattice. The actual kernels must satisfy
the local domination, integration, support and positivity hypotheses there.

\subsection{Integral coordinates and separation of rational frames}
\label{moment:rigidity}
Let $u=\varrho$, $n=u+1$, and write the ordered frame as
$F=(e_0,\ldots,e_{u-1},z)$, with $f_u=z$. In rational coordinates the Hermitian
space is the orthogonal sum of $u$ hyperbolic planes and $X_0\perp Ez$,
where $X_0$ has identity Gram matrix and $\langle z,z\rangle=-N$.
The $e_i$ pair with their duals $e_i^*$. At a complex place $v$, let
$\Phi_v$ be the chosen embedding and put $s_{va}=g_v^{u-a}/\sqrt N$ for
$0\le a\le u$, where $1\le |g_v|\le b^{23}$ is the local scale.
Changing to the normalized real-place basis multiplies the matrix entry
in row $i$, column $j$ by $s_{vj}/s_{vi}$ on frame rows and columns.
On $X_0$ use scale one and on $e_a^*$ the inverse scale of $e_a$.
The normalized Gram operator $J$ exchanges each hyperbolic pair, is the
identity on $X_0$ and acts by $-1$ on $z$. In particular $\|J\|=1$.

\paragraph{Two families of algebraic integers.}
Let $B_c$ be the finite adelic class basis. In the standard lattice basis
one middle column $x$ of $X_0$ is replaced by $x/\bar\pi+z/N$,
where $\pi\bar\pi=N$; the basis is then moved by the class representative
acting isometrically on $X_0$. Each $B_c$ fixes the frame columns and
leaves the rows along $e_i$ unchanged. Its integral Gram matrix
$J_{\rm lat}=B_c^*J_{\rm rat}B_c$ is independent of the class:
the altered column has squared Hermitian length $1/N-1/N=0$
and pairing $-1$ with $z$. For every rational term $g$ with
$K=B_c^{-1}gB_{c'}$ integral at the finite places, the two families
\[
 Z_{ja}(g)=\langle f_j,g f_a\rangle,\qquad
 A_{ia}(g)=(g f_a)_{e_i}\quad(0\le i<u)
\]
belong to $\mathcal O_E$. Indeed $Z_{ja}(g)$ is the $(f_j,f_a)$ entry
of $J_{\rm lat}K$ and $A_{ia}(g)$ is the corresponding row and
column of $K$. These equalities also hold at places dividing $N$.

\begin{lemma}\label{moment:fixed-frame}
Put $\rho=b^{37n}R$, $R\ge1$. In the normalized bases let $y_v$ be
points with $\|y_v^{-1}f_j\|\le\rho$ and Bessel length at most
$b^{6n+2}R$. Suppose that for a rational term $g$ as above, writing
$g_v$ for its normalized matrix,
\[
 \delta_{va}:=\|y_v^{-1}(g_v-1)f_a\|<(N\rho)^{-1}
 \qquad\text{for all }v,a.
\]
Then $gF=F$.
\end{lemma}
\begin{proof}
With the displayed change of basis,
\[
 |\Phi_v(Z_{ja}(g)-Z_{ja}(1))|
 \le |s_{vj}s_{va}|^{-1}\,\|y_v^{-1}f_j\|\,\delta_{va}
 \le N\rho\,\delta_{va}.
\]
This is Cauchy--Schwarz applied to the normalized Gram operator after
using invariance of the form under $y_v$. An algebraic integer of
modulus strictly below one at each chosen embedding is zero, by its
absolute norm. Hence every pairing agrees with its value for the identity.
Pairing with the $e_i$ and $z$ gives the expected dual and $z$
coordinates of $g f_a$. Unitarity gives the same length as $f_a$.
After those coordinates have been fixed, the difference can have no
$X_0$ component: the remaining contribution to the length is the sum
of its squared absolute coordinates at a complex embedding. Consequently
$(g-1)f_a=\sum_{i<u}d_i e_i$ with $d_i=A_{ia}(g)-A_{ia}(1)\in\mathcal O_E$.

Fix $a$ and descend on $i$. After $d_j=0$ for $j>i$, the $i$th row of
any Bessel unipotent $w$ applied to $(g_v-1)f_a$ equals
$\Phi_v(d_i)s_{va}/s_{vi}$, by unitriangularity on the isotropic flag.
The Bessel subgroup splits as a unipotent times a compact middle
isometry acting isometrically in the normalized Euclidean metric.
Thus the infimum defining the length permits a unipotent $w$ with
$\|w y_v\|\le b^{6n+2}R$ (the additive one in the length absorbs
the approximation to its infimum). It follows that
\[
 |\Phi_v(d_i)|\le b^{23u}\|w y_v\|\delta_{va}
 \le b^{23u+6n+2}R\,\delta_{va}\le\rho\,\delta_{va}<1 .
\]
The norm argument gives $d_i=0$ and closes the induction.
\end{proof}

\paragraph{Support and orbit weights.}
In rational frame coordinates take a reduced cutoff, supported where
the entries above the diagonal in the Gram--Schmidt coefficient matrix,
before archimedean rescaling, belong to a residue domain for
$M\mathcal O_E$ of sup radius $b^{26}$. At indices $a<j$ rescaling
multiplies the entry by $s_{vj}/s_{va}$ of modulus at most one.
The successive orthogonal residuals of $y_v^{-1}F$ have length at
most the Bessel length: the residual norms are unchanged by the
unitriangular action on the flag, and each is bounded by the norm of
the corresponding column after such a move. In normalized coordinates
$\|(wy_v)^{-1}\|=\|wy_v\|$ by $J$-unitarity. Thus each frame vector
has norm at most $n b^{26}\ell(y_v)\le b^{27}\ell(y_v)$.
On the length sublevel in \S\ref{B:cutoff} this is at most $\rho$.
For two contributing terms $p,q$ with the same starting class,
apply Lemma~\ref{moment:fixed-frame} to $p^{-1}q$ at $p^{-1}y$.
The relative term is integral between the two ending class bases since
the original finite matrices are invertible integral. Distinct rational
frames therefore have distance at least $(N\rho)^{-1}$ in the maximum
of the column norms over the embeddings.

If the frames and reductions agree, $p^{-1}q$ fixes $F$ and has
trivial finite reduction in the moved bases. Its component on
$X^\perp/X$ for $X=\langle e_i\rangle$ fixes $z$, hence acts on
$X_0$ by an isometry. The fitted level restricts on this definite
factor to the source level: the middle component on the quotient of
a principal congruence element lies in this level. It follows that
the ending representatives define the same source double coset.
Freeness at that level now makes the rational definite component
the identity. Thus the relative term belongs to the congruence
unipotents fixing the frame. Fixing one term in a fibre, every other
term gives a distinct such unipotent by taking the relative element.
For a cutoff whose translates by the class unipotent group sum to one,
nonnegativity bounds the weight of this fibre by one.

\paragraph{Construction of the cutoff on the required scales.}
We now construct the reduced cutoff just used, for the free level and
class bases of the fitted model. Write $\Gamma_c$ for the rational
Bessel unipotents in the principal congruence subgroup at $c$, and
$\Gamma_c^0$ for those fixing the frame pointwise. We compute the
image of $\Gamma_c$ on $F$. It consists of all upper unitriangular
$n$-by-$n$ matrices with entries in $M\mathcal O_E$ above the diagonal.
In fact $(g-1)/M$ is integral in the moved basis for $g\in\Gamma_c$.
At the split level primes it suffices to impose the congruence on the
chosen member of each pair: unitarity recovers the inverse adjoint
congruence on the conjugate member, as the lattice Gram matrix is
unimodular. At every other place $M$ is a unit. The top-row coordinates
computed above therefore give the inclusion in this unitriangular group.

For surjectivity fix $i<j$ and $a\in\mathcal O_E$. There is a rational
functional $\lambda$ integral on the class lattice, with
$\lambda(f_l)=a\,\mathbf1_{l=j}$ and zero on the other standard
basis vectors outside $X_0$. Indeed the corresponding adelic row of
$aB_c^{-1}$ has the prescribed entries there. Integrality on the
lattice is an open condition on the remaining finite-adelic entries,
so additive strong approximation supplies rational entries on $X_0$.
Write $\lambda=\langle v,\cdot\rangle$. The class lattice is self-dual
(the changed two-column Gram block has determinant $-1$), hence $v$
belongs to it. Put $x=e_i$, so that $\langle x,v\rangle=0$.
Choose $\tau=-M\zeta_3\in\mathcal O_E$, so $\tau+\bar\tau=M$.
Writing $v^\flat=\langle v,\cdot\rangle$, the map
\[
  g=1+M\bigl(xv^\flat-vx^\flat-\tau\langle v,v\rangle x x^\flat\bigr)
\]
preserves the form. In expanding $g^*J_{\rm rat}g$ the squared
term is $M^2\langle v,v\rangle (x^\flat)^*x^\flat$ and the linear
terms cancel it. The formula is integral and congruent to one at the
moved level. It fixes $e_l$ for $l\le i$, adds only multiples of $e_i$
to later isotropic frame vectors, and acts trivially on the middle
quotient; the adjoint action on the dual flag is then forced. Hence
it belongs to $\Gamma_c$. On the frame its matrix is $1+aM E_{ij}$
because $x^\flat$ vanishes on the frame. These elementary matrices
generate the indicated group by successive elimination of the upper
diagonals. The kernel is $\Gamma_c^0$.

The additive lattice $M\mathcal O_E$ in $\mathbb C^f$ has a Borel
domain of representatives of sup radius at most $b^{26}$.
For the fields of Theorem~\ref{thm:p1-data}, the $2f$ elements
$\zeta_{q_1}^a,\sqrt{-d}\,\zeta_{q_1}^a$, $0\le a<f$, are a
rational basis of algebraic integers by coprimality of the two
conductors. Their Minkowski images are a real basis. Rounding the
coordinates in this basis gives, for the sublattice generated by
their $M$-multiples, covering radius at most
$Mf(1+\sqrt d)/2$. The larger lattice has no greater covering
radius. Since $M\le b^{22}$, $f\le q_1\le b$ and $d$ is fixed,
the claim holds eventually. Choose a nearest representative in
the sup norm, breaking ties lexicographically among the finitely
many minimizing representatives. The resulting domain is Borel
by discreteness and countability.

Let $T(y)$ be the upper unitriangular matrix expressing the frame in
its Gram--Schmidt residuals, conjugated back by the diagonal scales.
When $y$ is replaced by $\gamma y$, $\gamma\in\Gamma_c$, this matrix
is right multiplied by the inverse frame matrix of $\gamma$. This
follows since adding earlier columns to a column preserves its
orthogonal residual; diagonal conjugation restores the original
rational coefficients. Reduce the entries of $T(y)$ into the chosen
additive domain in order of increasing superdiagonal. Right
multiplying by $1+M a E_{ij}$ changes entry $(i,j)$ by $Ma$,
affecting only more distant superdiagonals in that column. Existence
and uniqueness of each additive representative now give precisely one
coset of $\Gamma_c^0$ that takes $y$ into the reduced set.

Finally $\Gamma_c^0$ is discrete at infinity. A compact archimedean
set allows only finitely many rational matrices satisfying the fixed
finite-place lattice conditions, since their entries have bounded
denominators and bounded conjugates. There is a Borel fundamental
set for its free action by left translations: cover the ambient Lie
group by countably many open sets each disjoint from its distinct
translates and discard the orbits covered by earlier sets. Multiply
the indicator of this fundamental set by the reduced-set indicator.
Its sum over $\Gamma_c$ is exactly one, since reduction selected
one coset and the first indicator selects one element in that coset.
The same conclusion holds with any nonnegative measurable
$\Gamma_c^0$ cutoff whose translates sum to one. Thus the support
and weight properties hold for the class cutoffs used in the
amplitudes, regardless of that last choice.

On the constructed scales, F4 therefore follows for the fitted free
level, integral class bases, archimedean scales and reduced-cutoff
definition above. The existence of the actual filling with these
data (F0) and its identification with the period map (F1) are
separate requirements. The finite and archimedean basis computations
and the reduced-cutoff construction correspond to
\lineSrc{New/Spanning/Unfolding/Terms}{618}{frame integrality},
\lineSrc{New/Spanning/Unfolding/Packing}{286}{frame separation}, and
\lineSrc{New/Spanning/Unfolding/ReducedPoints}{637}{the box partition}.

\subsection{The reciprocal-discriminant sum at the level prime}
\label{app:pair-count}
Write $h=\varrho+1$, $H_0=b^{h^3 f}$ and $\tau=b^{kf/4}$.
We separate a lattice estimate in \S\ref{B:pairs} from the multiplicity
with which a near term can lie over a lattice. Let $\Lambda$ be any lattice
in the genus of the standard definite Hermitian lattice of rank $k$ over
$\mathcal O_E$. At finite places it is self-dual. For a rank-$s$ lattice
$A$ in its rational span, write $A^\#$ for its Hermitian dual there and
$d(A)=[A^\#:A]^{1/2}$ when $A$ is integral. These are cardinality
indices. This is the positive square root of the finite index of the
Gram determinant, the discriminant convention for the pair count. In
particular $d(B)=i\,d(A)$ for $B\subset A$ of index $i$.

\begin{lemma}\label{lem:pair-row-count}
Let $s+j=k$. For each row coset $y$ at the level prime let $t_y$ be its
integral row on $\Lambda$ modulo the ideal $\mathfrak n$ of norm $N^f$.
Suppose the $t_y$ are distinct. For each $y$ take all rank-$s$ sublattices
$B\subset\Lambda$ with $t_y(B)=0$ and $d(B)\le H_0$. At the scales
of the filling their combined weight satisfies
\[
 \sum_y\sum_B d(B)^{-k}
 \le (N^f)^j H_0^{2\varrho+6}
 (200fk\delta_E)^{f(k+1)j}\,2^k .
\]
Here $\delta_E$ is the root discriminant and $(\varrho+1)2f\le2^{2\varrho}$
is the field-degree bound for the filling.
\end{lemma}
\begin{proof}
Put $V=EB$, $A=\Lambda\cap V$, $d=d(A)$, and $i=[A:B]$.
The number of all rows modulo $\mathfrak n$ annihilating $B$ is at most
$i(N^f)^j$. Indeed, locally at $\mathfrak p^m\Vert\mathfrak n$ extend a
basis of the saturated summand $A_{\mathfrak p}$ to an ambient basis,
and diagonalize $B_{\mathfrak p}$ inside it with elementary-divisor
exponents $a_1,\ldots,a_s$. For $q=\operatorname{N}\mathfrak p$
there are $q^{mj+\sum_a\min(a_a,m)}$ choices of row. Multiplication over
the primes gives the assertion, also for a lattice in a moved basis.
Here the rows for the actual row cosets are the distinguished rows of
$y^{-1}$ in coordinates on the unimodular lattice. Congruent inverse
rows put $y^{-1}y'$ in the row stabilizer and give the same coset.

Put $\mu=2\varrho+4$. Since $di\le H_0$ and $i,d\ge1$,
\[
 i(di)^{-k}\le H_0^{\mu+2} d^{-(k+1)} i^{-(s+\mu)}.
\]
For fixed $V$, the sum of $i^{-(s+\mu)}$ over its finite-index sublattices
is at most $2^k$. At a finite place the series for sublattices in a
free module of rank $s$ is $\prod_{a=0}^{s-1}(1-q^{a-z})^{-1}$:
triangular bases over the valuation ring give one diagonal exponent
at each pivot and residues above the pivots. The Euler product at
$z=s+\mu$ is bounded by $\zeta_E(\mu)^s$.
Indeed $\zeta_E(\mu)\le\zeta(\mu)^{2f}$, and
$\log\zeta(\mu)\le 2^{-\mu}(1+2/(\mu-1))$ by integral comparison.
The displayed field-degree bound makes $2f\log\zeta(\mu)<\log 2$.

For clarity, the remaining sum over subspaces satisfies
\[
 \sum_{\operatorname{codim} V=j} d(\Lambda\cap V)^{-(k+1)}
 \le (200fk\delta_E)^{f(k+1)j}.
\]
The case $j=0$ follows from self-duality. For $j>0$ let
$A'=\Lambda\cap V^\perp$. Saturation and self-duality identify
$A^\#/A$ and $(A')^\#/A'$ in cardinality: restriction of the pairing
maps $\Lambda$ onto $A^\#$ locally, since $A$ is a direct summand
locally, and its kernel is $A'$. Thus both quotients have the cardinality
of $\Lambda/(A+A')$. In the Euclidean norm summing the squared lengths
over the $f$ complex places, the covolume of $A'$ in its real span is
$(2^{-f}\sqrt{|D_E|})^j d(A)$. To see the normalization, take a
pseudo-basis of $A'$. The determinant of its complex Gram matrix at
each chosen embedding enters the real covolume to the first power.
The coefficient ideals enter by their absolute norms. Their product
with the Gram factors is $[(A')^\#:A']^{1/2}$, by taking determinants
of the Hermitian pairing locally. The remaining factors are the
Minkowski covolumes of $j$ copies of $\mathcal O_E$.

Let $m=2fj$ and let $\lambda_a$ be the Euclidean successive minima of
$A'$. Minkowski's second theorem, using the cube of side $2/\sqrt m$
inside the unit ball, gives $\prod_{a=1}^m\lambda_a\le
m^{m/2}\operatorname{covol}(A')$. One can choose $E$-independent vectors
$X_1,\ldots,X_j$ with $\|X_{a+1}\|\le\lambda_{2fa+1}$ for $0\le a<j$:
the preceding $E$-span has real dimension $2fa$. Therefore
\[
 (\prod_a \|X_a\|)^{2f}\le (8fj\delta_E)^{fj}d(A).
\]
The chosen tuple determines $V$ by orthogonality. Every nonzero
$x\in\Lambda$ has $\|x\|\ge1$ since the trace from $E^+$ of its Hermitian
norm is a positive integer. The points are thus $1$-separated. Put
$D=2fk$ and $p=2f(k+1)$. Disjoint balls of radius $1/2$ give
$\#\{0<\|x\|\le R\}\le(1+2R)^D$ for $R\ge1$.
Integration against $pR^{-p-1}\,dR$ bounds $\sum_{x\ne0}\|x\|^{-p}$
by $3^D p/(p-D)=3^D(k+1)\le5^p$
(the last inequality follows from $k+1\le(5/3)^{2k}$).
Sum the product bound over all tuples and use $j\le k$ and
$8\cdot5^2=200$. This proves the subspace sum and, on multiplying
the preceding bounds, the lemma.
\end{proof}

\paragraph{Class and pair multiplicities.}
For a free source level the classes have equal quotient measure.
Suppose its row quotient has size $Y\ge(N^k/b)^f$, and the translated
class lattice is identified by a rational unitary isometry with the
kernel of the corresponding row in $\Lambda$ modulo $\mathfrak n$.
Integrate the uniform row bound of Lemma~\ref{lem:pair-row-count}
over the compact adelic quotient with probability measure. Right
translation preserves the measure, so division by $Y$ bounds the
average over classes. At the fitted scales
\[
 b^f H_0^{2\varrho+6}2^k\le\tau,\qquad
 N^f(200fk\delta_E)^{f(k+1)}\le\tau .
\]
For the first comparison use $k\le f$, $b\ge2$ and
$2+h^3(2\varrho+6)\le k/4$. For the second use $N\le b^{60h}$,
$(200fk\delta_E)^{1000}\le b$ and $60h+(k+1)/1000\le k/4$.
Consequently the sum of $d(B)^{-k}$ over classes and their eligible
lattices is at most $\#\{c\}\tau^{j+1}/N^{kf}$.

Finally suppose each near pair determines a common label and a rank-$(k-j)$
common lattice of discriminant at most $H_0$ contained in both term lattices, with pair discriminant
equal to the discriminant of the common lattice. The single-term estimate
needed is that, at fixed first class, label and such lattice, there are at
most $b^{2000h^6(j+1)f}$ possible terms, counting the second class as
part of the term. Thus the number of ordered pairs in a fibre is at most
the square of this bound. For $j\le h^3$ and $k\ge512000 h^9$ that square
is at most $\tau$. Summing over labels and multiplying by the preceding
class sum gives $C_p\tau^{j+2}$, where
$C_p=\#\{\text{labels}\}\#\{c\}/N^{kf}$.
This proves the reciprocal-discriminant reduction in \S\ref{B:pairs}
given the indicated fitting data and single-term multiplicity. The use of this
bounded-discriminant count for all the near pairs requires the geometric
bound $d\le H_0$ in addition to $j\le h^3$; the source derives these together
in the \lineSrc{New/Spanning/RemainingTerms/Sum}{90}{gauge comparison}. The local
cell-rank and gauge estimates that bound the weights of individual pairs
are separate inputs. The
\lineSrc{New/Spanning/RemainingTerms/LatticesOfCorank}{373}{row count},
\lineSrc{New/Spanning/RemainingTerms/SiegelLemma}{630}{subspace sum} and
\lineSrc{New/Spanning/RemainingTerms/Count}{187}{pair estimate}
specify the same normalizations.

\subsection{The primary Hall factor in the level projector}

Let $k=\mathbb F_q$. In this subsection $P$ denotes a monic irreducible
polynomial over $k$, with $P(0)\ne0$, of degree $d\ge1$.
Fix an injective character $T:\overline{k}^{\times}\to\mathbb C^{\times}$
and a root $\alpha$ of $P$ in $k_d=\mathbb F_{q^d}$. Write
$T(\alpha)=\exp(2\pi i a/(q^d-1))$ and put $\eta=T^a$.
On $k_d^\times$ the result depends only on $a\bmod (q^d-1)$.
The assignment $\alpha\mapsto\eta|_{k_d^\times}$ is an isomorphism
from a cyclic group to its dual and commutes with the $q$-power map.
Thus $\eta|_{k_d^\times}$ is regular, meaning that its Frobenius orbit
has size $d$.

Here is the Hall formula defining the primary factor. For
$g\in\operatorname{GL}_d(k)$, let $V=k^d$ be the $k[X]$-module with
$X$ acting as $g$. Write $\mu(0,W)$ for the M\"obius function of its
submodule lattice. If $\beta_\eta(h)$ sums $\eta$ on the eigenvalues
of $h$, with algebraic multiplicities, then
\[
 \tau_P(g)=(-1)^d\sum_{W\subseteq V\,\text{a submodule}}
       \beta_\eta(g|_W)\mu(0,W).
\]
Indeed the Hall product with the constant function one sums over invariant
subspaces, and its Hall inverse evaluated on a module gives
the M\"obius function by the defining recursion. These are the
\lineSrc{Lib/LocalGL/Green/HallAlgebra}{549}{eigenvalue sum and Hall product}
and the \lineSrc{Lib/LocalGL/Green/TauFunctions}{133}{rank and sign convention}
used in the projector.

\begin{lemma}\label{lem:fitted-primary-character}
The function $\tau_P$ on $\operatorname{GL}_d(k)$ is the cuspidal
character attached to the orbit of $\eta|_{k_d^\times}$. Its value
at the identity is $\prod_{j=1}^{d-1}(q^j-1)$.
\end{lemma}
\begin{proof}
Only semisimple $W$ contribute. Indeed, for a nonsemisimple module $W$
its socle, the sum of all its simple submodules, is nonzero and proper.
Induction on length gives vanishing for each proper nonsemisimple
submodule. The sum of the M\"obius values on the other proper submodules
equals the sum over the entire submodule lattice of the socle, hence is
zero. The M\"obius recursion gives $\mu(0,W)=0$.

For each distinct monic irreducible factor $H$ of the characteristic
polynomial, let $s_H=\deg H$, $Q_H=q^{s_H}$, and $m_H$ be the multiplicity
of $k[X]/(H)$ in the socle of $V$. Each nonzero semisimple submodule is
specified by subspaces of $\mathbb F_{Q_H}^{m_H}$. Write $n_H$ for their
dimensions and $B_H=\sum_{H(z)=0}\eta(z)$, summing over distinct roots.
Then $\beta_\eta(g|_W)=\sum_H n_H B_H$. For a vector space of
dimension $n$ over $\mathbb F_Q$, its M\"obius value is
$(-1)^n Q^{\binom n2}$. This follows from the M\"obius recursion and
the finite $q$-binomial identity
\[
 \sum_{n=0}^{m}{m\brack n}_Q(-1)^n Q^{\binom n2}u^n
       =\prod_{j=0}^{m-1}(1-Q^j u),
\]
where the Gaussian coefficient counts the $n$-dimensional subspaces.
M\"obius values multiply on products of lattices. It follows that the
unsigned Hall sum before multiplication by $(-1)^d$ is
\[
 \left.
 \left(\sum_H B_H u_H\frac{\partial}{\partial u_H}\right)
 \prod_H\prod_{j=0}^{m_H-1}(1-Q_H^j u_H)
 \right|_{u_H=1\ \mathrm{for\ all}\ H}.
\]
With two or more distinct factors every term vanishes. With one factor $H$
the value of $\tau_P(g)$ is
\[
       (-1)^{d-1}
       \left(\sum_{H(z)=0}\eta(z)\right)
              \prod_{j=1}^{m_H-1}(1-q^{s_H j}).
\]
In this case $s_H\mid d$, so the roots occurring in the sum lie in $k_d$.
Moreover $m_H=\dim_{k_d}\ker(g-z)$ after extending scalars, as is seen
on each cyclic primary summand. These values on all conjugacy classes agree
with Green's cuspidal character formula for the regular character
$\eta|_{k_d^\times}$, in the form stated by D.~Prasad,
\emph{The space of degenerate Whittaker models for general linear groups
over a finite field}, Int. Math. Res. Not. (2000), no.~11,
\href{https://repository.ias.ac.in/36158/1/36158.pdf}{Theorem 2, pp.~581--582}~\cite{prasad}.
At $g=1$ take $H=X-1$, $m_H=d$ and $B_H=1$. The two factors of
$(-1)^{d-1}$ cancel.
\end{proof}

For a squarefree label the fitted character is, by definition, the Hall
product of these primary factors, hence their parabolic induction.
Distinct irreducible polynomials give distinct cuspidal labels (their
degrees differ or their Frobenius orbits differ). Thus the squarefree-label
argument and flag count of \S10.4 apply directly to the fitted character
with the positive sign. This discharges the finite primary-character
identification within F7. The comparison of inner forms and the matching
of archimedean formal degrees in F7 concern different objects.

\subsection{A variable error bound for the good projector}

Let $G=G_0\times G_1$ be the residue group at the two level primes. Regard
each factor as a subgroup of $G$. Write $X$ for the complex Hilbert space
with orthonormal basis indexed by the fitted classes and Bessel coordinates,
$d_X=\dim X>0$, and $Y$ for the principal-level target with its Petersson
metric. Write $C$ for the isomorphism from source forms to these coordinates.
If $J$ is the principal level, $\nu>0$ the fitted unipotent covolume and
$\Pi$ the Poincar\'e map, fix the single operator
\[
       P=(\operatorname{vol}(J)\nu)^{-1/2}\Pi C^{-1},\qquad
       q_g=P^*\rho(g)P ,
\]
where $\rho$ is the unitary residue-group action on $Y$. Both moment
inequalities use this $P$. The unfolding needed here identifies each block
of $q_g$ with the conjugate of its full amplitude sum divided by $\nu$.
At the level of scalars this follows from the fitted Petersson identity
with factor $\operatorname{vol}(J)$ for the unscaled map, since the inner
product uses two copies of the scaling factor. The period-to-series identity
and the Petersson identity for the actual fitted functions are hypotheses
of this identification. The source gives the
\lineSrc{New/Spanning/SecondMoment/SpanningTheorem}{153}{coordinate map and scaling} and
the \lineSrc{New/Spanning/SecondMoment/SpanningTheorem}{168}{unfolded block identity}.

Let $R\subseteq G$ be the reduced Bessel subgroup and $r$ its unitary
representation on $X$, extended by zero off $R$. Assume the fitted
equivariance $\rho(s)P=Pr(s)$ for $s\in R$. Set $E_g=q_g-r(g)$.
For $s\in R$ we then have
$E_{sg}=r(s)E_g$ and $E_{gs}=E_g r(s)$. Define nonnegative errors by
\[
 h_R^2=\frac{1}{|R|d_X}\sum_{s\in R}\|E_s\|_{\mathrm{HS}}^2,\qquad
 h_i^2=\frac{1}{|R|d_X}\sum_{g\in R G_i}\|E_g\|_{\mathrm{HS}}^2
 \quad (i=0,1).
\]
For $t\in G_i$, assume the diagonal Bessel trace identity
$\operatorname{tr}r(t)=d_X\psi_i(t)$, with $\psi_i$ the reduced
unipotent character extended by zero outside its support. Write
\[
  \overline\psi_i^{\,\mathrm{cl}}(t)
     = |G_i|^{-1}\sum_{v\in G_i}\psi_i(vtv^{-1}),\qquad
  K_i=\sum_{t\in G_i}|\overline\psi_i^{\,\mathrm{cl}}(t)|^2 .
\]
The character average is normalized; the sum defining $K_i$ is unnormalized.
Conjugating $t$ by $G$ gives the same average because the other factor
commutes with $t$. Hence the full-group bound (10.11) implies $K_i\le2$.

Write $e_i$ for the self-adjoint good central idempotent in $\mathbb C[G_i]$ and
$x_i=1-e_i=\sum_{t\in G_i}x_i(t)t$. Put $a_i=x_i(1)$ and
$Q=\rho(e_0)\rho(e_1)$. The coefficients of $e_0\otimes e_1$ at
$(g_0,g_1)$ are the products of the two sums over allowed label tuples
having a primitive member at their respective primes. Expansion sums over
precisely the good label families appearing in the actual target projector.
This proves equality of the two projectors, coefficient by coefficient, for
the same label idempotents. See the
\lineSrc{New/Spanning/SecondMoment/SpanningTheorem}{197}{tensor identity}.

\begin{proposition}\label{prop:variable-projector-loss}
With the identities and conventions above, every subspace $Z\subseteq X$
satisfies
\[
 \|QP\operatorname{pr}_Z\|_{\mathrm{HS}}^2
 \ \ge\ \dim Z-d_X\left(h_R+
     \sum_{i=0}^1\bigl(h_i+\sqrt{a_i K_i}\bigr)\right).
\]
\end{proposition}
\begin{proof}
For $s\in R$, $E_s=(P^*P-1)r(s)$, so
$\|P^*P-1\|_{\mathrm{HS}}=h_R\sqrt{d_X}$. Pairing with
$\operatorname{pr}_Z$ in the Hilbert--Schmidt inner product gives
$\|P\operatorname{pr}_Z\|_{\mathrm{HS}}^2\ge\dim Z-h_R d_X$.

Fix $i$ and put $x=x_i$. On $\ell^2(G;\operatorname{End}(X))$, with
counting measure and Hilbert--Schmidt fibres, set
$(T_x F)(g)=\sum_{t\in G_i}x(t)F(gt)$. The right translations are unitary
and form a representation, so $T_x$ is an orthogonal projection.
Apply it to $F(g)=\mathbf1_{R G_i}(g)E_g$. For each $s\in R$,
$(T_xF)(s)=r(s)\sum_t x(t)E_t$. Consequently
\[
  |R|\Big\|\sum_t x(t)E_t\Big\|_{\mathrm{HS}}^2
     \le \|T_x F\|_2^2\le\|F\|_2^2=|R|d_Xh_i^2 .
\]
The absolute trace of $\sum_t x(t)E_t$ is
therefore at most $d_X h_i$. For the Bessel part, centrality of $x$ permits
averaging $\psi_i$ over conjugacy classes. The identity coefficient of
$x^*x=x$ gives $\sum_t|x(t)|^2=a_i$. Thus
\[
  \left|\sum_t x(t)\operatorname{tr}r(t)\right|
        \le d_X\sqrt{a_i K_i}.
\]
Adding the two contributions bounds
$\|\rho(x_i)P\|_{\mathrm{HS}}^2
 =\operatorname{tr}(P^*\rho(x_i)P)$ by $d_X(h_i+\sqrt{a_iK_i})$.
Since the two projections commute, $1-Q\le\rho(x_0)+\rho(x_1)$ as
positive operators. Restricting the source to $Z$ can only decrease the
squared Hilbert--Schmidt norm of the discarded part. Subtracting the two
losses from the first bound proves the assertion.
\end{proof}

For example, the arithmetic bounds $h_i\le b^{-6}$, $a_i\le b^{-8}/2$
and $K_i\le2$ imply (10.2), since $h_R\le h_i$ and
$3b^{-6}+2b^{-4}\le5b^{-4}$ for $b\ge1$.
On the constructed scales, if the former bounds for $h_i,a_i$ hold, it
suffices instead that $K_i=O(b^{6-\delta})$ for a fixed $\delta>0$,
with uniform implied constants at the fixed parameters. Indeed the total
loss is $O(b^{-6}+b^{-1-\delta/2})=o((32b x^{202})^{-1})$ because
$\log x=o(\log b)$. This relaxes the character input specifically for
the projected inequality. The full moment still has its own hypotheses.

For comparison, the fitted source invokes the
\lineSrc{New/Spanning/Unfolding/Moments}{446}{off-Bessel square sum},
the \lineSrc{New/Spanning/Unfolding/Moments}{751}{diagonal trace}, the
\lineSrc{New/Spanning/Unfolding/Moments}{991}{character mass} and the
\lineSrc{New/Spanning/SecondMoment/SpanningTheorem}{221}{bad idempotent mass}.
They give the
\lineSrc{New/Spanning/SecondMoment/SpanningTheorem}{257}{projected estimate}
for $QP$. The
\lineSrc{New/Spanning/SecondMoment/SpanningTheorem}{270}{full estimate}
uses the same $q_g$ and relabels the product group by its chosen places,
a bijection preserving the unnormalized sum.

\subsection{Rational parametrization of the open blocks}\label{moment:open}

Write the Hermitian space as $X\oplus H\oplus X^*$ in a Witt basis,
where $X$ has ordered basis $e_1,\ldots,e_u$, $X^*$ has its dual
isotropic basis, and $H=X_0\perp Ez$ with $X_0$ standard positive
definite and $h(z,z)=-N$. Use the convention that $h$ is linear in the
second argument. Let $U$ be the unipotent radical of the stabilizer of
the complete flag in $X$. For $g$ write $D_j$ for the leading minor
of order $j$ of the matrix of pairings of $(e_1,\ldots,e_u,z)$
with its $g$-image, and $D_0=1$.

\paragraph{Rational normal form.}
If $D_j\ne0$ for $1\le j\le u$, there is a unique expression
$g=v w_p\sigma v'^{-1}$ with $v,v'\in U(E)$, $p_i\in E^\times$,
and $\sigma$ fixing the hyperbolic pairs and acting on $H$ by an
isometry. Here $w_p e_i=p_i e_i^*$, $w_p e_i^*=\bar p_i^{-1}e_i$,
and $w_p$ is the identity on $H$. Indeed, in the first hyperbolic
pair put $p=h(e_1,g e_1)$. The isotropic column $g e_1/p$ has the
form $e_1^*+y+c e_1$ in its orthogonal splitting, where
$c+\bar c=-h(y,y)$. There is a unipotent isometry $v_1$ fixing
$e_1$, sending $e_1^*$ to this column and sending a middle vector
$h_0$ to $h_0-h(y,h_0)e_1$. Hence $(v_1 w_1(p))^{-1}g$ fixes $e_1$.
Any isometry fixing $e_1$ factors as the extension of its induced
isometry on $e_1^\perp/E e_1$ times a unipotent of this same form:
after removing the induced isometry, the image of $e_1^*$ determines
$y,c$, and preservation of the pairings determines the other columns.
Iterate on the quotient. Left and right flag unipotents preserve each
leading minor (they act on the first $j$ frame vectors by
unitriangular operations, also when the last vector is $z$). Thus the
successive pivots are $p_i=D_i/D_{i-1}$ and the recursion can run to
completion. For uniqueness, two expressions have the same pivots.
The middle isometry is also determined: the relation on $X^\perp/X$
consisting of the pairs of classes of $a,g^{-1}a$ with
$a\in X^\perp\cap g X^\perp$ is invariant under the two unipotents,
and for $w_p\sigma$ it is the graph of $\sigma|_H$, with the second
coordinate as input. Finally $w_p\sigma$ conjugates $U$ to the
opposite flag unipotent; their intersection is trivial, by block
triangularity in opposite orders. This determines both unipotents.

\paragraph{The last entry and the blocks.}
For $t=h(z,\sigma z)$ write $\sigma z=\omega-(t/N)z$ with
$\omega\in X_0$. Preservation of its length gives
\[
        t\bar t-N^2=N h(\omega,\omega).
\]
The frame determinant and penultimate minor of $w_p\sigma$ are
$t\prod p_i$ and $\prod p_i$. Thus the nonzero-discriminant
condition forces $\omega\ne0$, and $|t|_v^2>N^2$ at every
archimedean place (ordinary complex absolute value). It also gives
$t=D_{u+1}/D_u$.

For a contributing cut term, both archimedean arguments have Bessel
length at most $b^{6n+2}$, with $n=u+1$. Move each independently to
its reduced box by the class congruence unipotents of
\S\ref{moment:rigidity}. Each frame column now has norm at most
$b^{37n}$. The pairings of the two frames, converted back by their
scales, have modulus at most $Nb^{74n}\le b^{134n}$ since
$N\le b^{60n}$; the corresponding rational pairings are algebraic
integers by the class-basis calculation. Their leading minors are
unchanged by the two upper unitriangular frame operations. Expansion
of a $j$-by-$j$ determinant for $j\le n\le b$ bounds its squared
modulus at every embedding by $n^{2n}b^{268n^2}\le b^{270n^2}$.
Only finitely many algebraic integers satisfy this bound. Put $A$
equal to the set of ratios of the nonzero members of this finite set. Then all
$p_i$ and $t$ lie in $A$. For fixed classes $c,c'$, reduction $s$ and
link $\Lambda$, every contributing term has the same $\sigma$ by
the graph description. Assuming absolute summability of the cut amplitudes, summing first over $U(E)\times U(E)$ and then
over the finitely many heads $p\in A^u$ gives exactly the cut block:
enforce the original class and reduction restrictions on each term.
The rational normal form proves injectivity of this parametrization
and absolute summability justifies the regrouping. Across blocks,
$(c,c',s,\sigma)$ determines $\Lambda$ and hence the block.
Consequently summing the squares of these block sums incurs no
multiplicity. It suffices to bound every finite collection and then
take the supremum.

For an integral representative in the moved frames the same elimination
works outside a finite conjugation-stable set containing the support
of the pivots, their inverses, the level, and the frame denominators:
each step divides only by its pivot or its conjugate and uses the
integral Hermitian pairing. Class changes on $X_0$ fix the hyperbolic
pairs and $z$, so the middle factor after elimination is the
class-translated middle isometry. This explains the integrality
condition on the floor-group factors in the finite-place sum. Absolute convergence and the bounds on the resulting orbital integrals enter separately.

The rational elimination and link identification above correspond to
\lineSrc{Lib/DiscreteSeries/Bruhat/BigCell}{289}{rational elimination},
\lineSrc{Lib/DiscreteSeries/Bruhat/BigCell}{374}{uniqueness}, and
\lineSrc{New/Spanning/Unfolding/Link}{309}{the graph of the middle isometry}.
The block reindexing occurs in
\lineSrc{New/Spanning/MainTerm/GlobalToLocal}{383}{the sum over heads}.

\subsection{From open-cell energy to the full moment}\label{moment:full}

The following transfer makes the dependence on the leading open-cell
error explicit. It also identifies a sufficient replacement for F6
while keeping the other errors of the manuscript. Let
$P=(\operatorname{vol}(J)\nu)^{-1/2}P_{\mathrm{raw}}$ in the
orthonormal fitted source coordinates, where $J$ is the principal
level and $\nu$ the unipotent covolume. For the unitary residue-group
action $\rho$ set
\[
 \beta(s)=\|P^*\rho(s)P\|_{\mathrm{HS}}^2,\qquad
 I_*=d_X^2/(\deg_{\mathrm{Bess}}\,m_r).
\]
Here $m_r$ is the definite genus mass. The period-to-series
identification for this $P$ says that its translated Gram blocks are
the full amplitude sums divided by $\nu$, up to the chosen convention
for complex conjugation. In particular all the sums below refer to this
single $P$. For a set $T$ of chosen places above the first level prime
and one nonzero frame vector at each place in $T$, let $K_T$ be their
simultaneous fixer and put $A_T=b^{-6k|T|}I_*$. Write
$M_T=\sum_{s\in K_T}\beta(s)$. The vector choices are included when
maximizing over $T$. Positivity of type follows from
\[
  \beta(s^{-1}t)=\langle \rho(s)PP^*\rho(s)^*,\rho(t)PP^*\rho(t)^*
                    \rangle_{\mathrm{HS}},
\]
by cyclicity of the trace.

Retain the cutoff and unequal-link error, the identity-moment bound,
the full-fixer bound and the nonopen estimate used in the original
moment argument. More precisely, at every fixer suppose the
square root of the full open mass is at most the square root of the
nonzero-discriminant link energy divided by $\nu$, plus
$b^{-15}\sqrt{I_0}$, where $I_0=|R|d_X\le A_T$ is the identity-moment scale. At the fixer with all places selected suppose $M_T\le A_T$.
At a maximizer of $b^{6k|T|}M_T$ with a place still free, assume the
nonopen mass is at most $b^{-4}M_T$. The manuscript proves the nonopen estimate by its finite-field minor
inequality, since maximality makes the mass of each additional vector
stabilizer at most $b^{-6k}M_T$. Its all-place bound uses the separate
scalar-coset and near-pair estimates. Suppose the open link energy is bounded, at every
fixer, by $(1+\delta)\nu^2 A_T$ for some $\delta\ge0$. Then
\[
    \sum_s\beta(s)\le
    \frac{(\sqrt{1+\delta}+b^{-15})^2}{1-b^{-4}} I_*.
\]
Indeed, if the maximizing fixer contains all places, comparison
with the empty fixer proves the claim directly. Otherwise the open
mass is at most $(\sqrt{1+\delta}+b^{-15})^2 A_T$ by the square-root
estimate, and $M_T$ is bounded by this plus $b^{-4}M_T$.
Divide by $1-b^{-4}$ and compare with the empty fixer.

At the constructed scales $e=(32 b x^{202})^{-1}$ satisfies
$b^{-4}=o(e)$, since $\log x=o(\log b)$. Thus
$\delta\le e/2$ suffices eventually: subtracting 1 from the displayed
factor gives $(\delta+2b^{-15}\sqrt{1+\delta}+b^{-30}+b^{-4})/
(1-b^{-4})$. In particular a uniform open-link error
$O(b^{-1-\epsilon})$ for any fixed $\epsilon>0$ suffices, provided
the other inputs just specified hold. Finally
$d_T\le |G|\deg_{\mathrm{Bess}} m_r$ converts this into
the second of the two normalized inequalities of Proposition~\ref{prop:weakerprecision}.
This propagates the relaxed target one level upstream and preserves
the required $b^{-6k|T|}$ saving at every fixer. It does not estimate
the open energy or the projected moment.

For the stronger coefficients of the pinned source, the chain is
\lineSrc{New/Spanning/Unfolding/FillingAdmissible}{743}{admissibility of the actual filling},
\lineSrc{New/Spanning/SecondMoment/SpanningTheorem}{168}{the fitted Gram blocks},
\lineSrc{New/Spanning/MainTerm/EnergyEstimate}{631}{open link energy},
\lineSrc{New/Spanning/Unfolding/SquareExpansion}{559}{the cutoff and cross-link error},
\lineSrc{New/Spanning/Unfolding/MomentBounds}{264}{open mass},
\lineSrc{New/Spanning/Unfolding/MomentBounds}{441}{maximizing fixer},
and \lineSrc{New/Spanning/SecondMoment/SpanningTheorem}{288}{target-dimension comparison}.
The source obtains $1+b^{-4}$ for link energy, $1+4b^{-4}$ for open
mass, and $1+9b^{-4}$ for full mass. The threshold for the projected
estimate must still be verified separately on $P$.

\paragraph{The outstanding estimate on the actual filling.}
For each finite collection of distinct blocks as above in one fixer,
write $B_j$ for its sum of cut double-coset amplitudes over heads.
The unresolved analytic assertion at this point is
\[
   \sum_j\|B_j\|^2
       \le (1+e/2)b^{-6k|T|} I_*\nu^2.
\]
Its quantifiers include every allowed choice of fixed vectors. The
source's stronger argument groups by $t$; the Gram matrix of the
scaled archimedean orbital vectors and the finite last-entry density
then give a bound for $R_{k-1}^2\sum_j\|B_j\|^2$. The next inequality
bounds the real part of the sum over last entries and pairs of heads,
with that density divided by the intermediate group's mass ratio.
Compare \lineSrc{New/Spanning/MainTerm/EnergyEstimate}{422}{archimedean Gram identification},
\lineSrc{New/Spanning/MainTerm/SiegelTopBound}{846}{the fixed-tail Siegel estimate},
\lineSrc{New/Spanning/MainTerm/FiniteSieveSystem}{280}{intermediate group},
\lineSrc{New/Spanning/MainTerm/GlobalSieve}{1145}{global sieve}, and
\lineSrc{New/Spanning/MainTerm/EnergyEstimate}{561}{assembly of the fixed-fixer estimate}.
The manuscript's Appendix F.1 computes the numerical comparison
conditional on the fitted measures and inequalities. The rational
normal form and the error transfer above supply two intermediate
arguments; an independent derivation of the fitted orbital and sieve
bounds is still required.

\subsection{The Fourier sums in the open-cell sieve}\label{moment:sieve}
\subsection*{The weight of the heads in the open-cell sieve}
Write $n=\varrho+1$, $f=[E:\mathbb Q]/2$, and let $v$ run over one embedding of the CM field $E$ from each conjugate pair, with ordinary complex modulus. Let $\Sigma$ be the finite set of places used for integration; it may contain places beyond those at which vectors are fixed. The constants from the filling used below are
\[
 b\ge1000,\quad n\ge 10^5,\quad a_0=|\operatorname{disc}E|^{1/(2f)}\le b,\qquad
 1\le g_v\le b^{23},\qquad 1\le N\le b^{60n}.
\]
Put $t_{va}=g_v^{2(\varrho-a)}/N$ for $0\le a\le\varrho$ and $P_a=(2/a_0)^f/\prod_v t_{va}^2$. An entry list is a finite set $A\subset E^\times$ integral off $\Sigma$. For a head $p\in A^i$, put $d_j=p_1\cdots p_j$ and $c_{vj}=\prod_{a=0}^{j-1}t_{va}$ for $1\le j\le i$. For an element $d\ne0$ integral at $\Sigma$, set $I_\Sigma(d)=\prod_{w\in \Sigma} q_w^{v_w(d)}$, with the additive valuations normalized by a uniformizer. Set $R=b^{60n^2}$.
The weight $w_i(p)$ is zero unless all $d_j$ are integral at $\Sigma$ and $|c_{vj}v(d_j)|\le R$ at every $v$. On that support it equals
\[
 \prod_{j=1}^i\left(I_\Sigma(d_j)^{1/2}\prod_v\max\{1,|c_{vj}v(d_j)|^{-1}\}\right)^{n^6}.
\]
\begin{lemma}
For $0\le i\le\varrho$ and every entry list,
\[
 \left(\sum_{p\in A^i}w_i(p)\right)^2
 \le b^{400n^9f}\prod_{a<i}P_a .
\]
\end{lemma}
\begin{proof}
On the support, $d_j\in\mathcal O_E$, and $c_{vj}^{-1}\le N^j\le R$, hence $|v(d_j)|\le R^2$. Since $I_\Sigma(d_j)$ divides $|N_{E/\mathbb Q}d_j|$, the factor before the power $n^6$ is at most
\[
 \prod_v\max\{|v(d_j)|,c_{vj}^{-1}\}\le R^{2f}.
\]
Thus each weight is at most $b^{120n^8fi}$. The head is recovered by successive quotients of its partial products. Distinct integers in the Minkowski embedding have sup distance at least one, since the norm of their difference is a nonzero integer. Disjoint open balls of radius $1/2$ for that norm show that the number of integers of norm at most $R^2$ in the sup metric is at most $(1+2R^2)^{2f}$. Thus
\[
 \sum_p w_i(p)\le (1+2R^2)^{2fi} b^{120n^8fi}\le b^{130n^9f}.
\]
Indeed $1+2R^2\le b^{120n^2+1}$ and $2n(120n^2+1)\le10n^9$. Also $t_{va}\le b^{46\varrho}$ yields $\prod_{a<i}P_a\ge b^{-(1+92\varrho)if}\ge b^{-100n^2f}$. Squaring and using $260n^9+100n^2\le400n^9$ proves the claim.
\end{proof}
This proves item 3 of the sieve on \S\ref{spread:normalization} for its specified weights and scales.

\subsection*{Summation of the dual-frequency majorant}
For $y\in E^\times$ integral off $\Sigma$, extend $I_\Sigma(y)=\prod_{w\in \Sigma}q_w^{v_w(y)}$ to rational values, and put $D(y)=N(\{a\in\mathcal O_E:ay\in\mathcal O_E\})$. Then $D(y)=\prod_{w\in \Sigma}q_w^{\max(0,-v_w(y))}$ and $I_\Sigma(y)\le |N_{E/\mathbb Q}y|$. For $d\ge0$ set $h_d(s)=s^d\min(1,s)^6$.
\begin{lemma}
Suppose $c_v\in\mathbb C^\times$, $|c_v|^{-1}\le x_v\le B$, $B\ge1$. For every finite set $X$ of nonzero $\Sigma$-integers,
\[
 \sum_{y\in X}\frac{I_\Sigma(y)^{d/2}}{D(y)^3}
       \prod_v h_d(|c_vv(y)|^{-1})
 \ \le\ \prod_v x_v^d (bB^2)^{3f}.
\]
\end{lemma}
\begin{proof}
The norm bound cancels the factors $|v(y)|^{-d}$. Since $\min(1,|c_v v(y)|^{-1})\le B/\max(1,|v(y)|)$, it suffices to sum $D(y)^{-3}\prod_v\max(1,|v(y)|)^{-6}$ and bound the result by $b^{3f}$.
Fix the denominator ideal $\mathfrak d$, of norm $D$, and partition by $2^{m_v}\le\max(1,|v(y)|)<2^{m_v+1}$, $m_v\ge0$. All these points lie in $\mathfrak d^{-1}$. Divide the $v$-coordinate by $L_v=2^{m_v+1}$. Distinct scaled points have sup distance at least $s=(D\prod_v L_v^2)^{-1/(2f)}$, since the norm of a nonzero difference in $\mathfrak d^{-1}$ is at least $1/D$. Packing disjoint balls of radius $s/2$ in the sup ball of radius $1+s/2$ bounds the count by $(1+2/s)^{2f}\le 3^{2f}D\prod_v L_v^2$. Summing the dyadic boxes costs at most $(36\sum_{m\ge0}2^{-4m})^fD$ before the factor $D^{-3}$. Finally $\sum_{\mathfrak d}N\mathfrak d^{-2}=\zeta_E(2)\le\zeta(2)^{2f}<4^f$ by the Euler products. The resulting constant $(36\cdot16/15\cdot4)^f$ is below $b^{3f}$.
\end{proof}
For the fitted frequency scales write $d=k-n^7$, let $\Theta$ be the archimedean volume factor in Appendix F.1, and assume $|\lambda_v|^{-1}\le g_v/M$. At step $a$ set $c_v=\lambda_v/t_{va}$, $x_v=g_v^{2\varrho+1}/(MN)$, and $B=b^{23(2\varrho+1)}$. Then $|c_v|^{-1}\le x_v\le B$. With $C_{\rm dec}=\prod_v\Theta x_v^d$ and $(bB^2)^3\le b^{n^9}$ the lemma gives
\[
 \sum_{y\in X}\!
 \left[\prod_v\Theta b^{2000n^9} h_d(|c_vv(y)|^{-1})\right]
 \frac{I_\Sigma(y)^{d/2}}{D(y)^3} b^{n^9 f}
 \le C_{\rm dec} b^{2002n^9 f}.
\]
The same inequality holds for the whole nonzero dual set by monotone convergence. This supplies the global summation of the pointwise decay majorant, including its numerical exponent. Applying it to the actual transforms still requires the archimedean and finite-place kernel bounds by this majorant, along with the scaled Poisson formula. It also bounds a single off-diagonal frequency meeting these conditions.

\subsection*{The one-sided Poisson step}
Let $X_\Sigma=\prod_{v\mid\infty}\mathbb C\times\prod_{w\in \Sigma}E_w$, with ordinary area in each complex coordinate and finite Haar measures giving $\mathcal O_w$ volume one. Let $\iota_r(y)=((r_v v(y))_v,(y)_w)$ with $r_v>0$. Put $\mathcal O_{E,\Sigma}=\{y\in E:v_w(y)\ge0\ (w\notin \Sigma)\}$. Fix a generator $\omega$ of the different, as in the fitted data. Use the complex character $z\mapsto\exp(z-\bar z)$ and at $w\in \Sigma$ the finite components of the global trace character with infinite component $u\mapsto\exp(-2\pi i\operatorname{Tr}_{E_\infty/\mathbb R}u)$ evaluated on $-x_w/\omega$, with the signs chosen so the resulting dual characters on $X_\Sigma$ annihilate $\iota_r(\mathcal O_{E,\Sigma})$ at frequencies
\[
 \xi(y)=\big((2\pi i\,v(y)/(r_v v(\omega)))_v,(y)_w\big),\qquad y\in\mathcal O_{E,\Sigma}.
\]
Here the archimedean factor of this dual character is $\exp(2\pi i\sum_v2\operatorname{Re}(v(y)z_v/(r_vv(\omega))))$. Denote integration against the character by $\widehat F(\xi(y))$.
\begin{lemma}
Suppose $F\ge0$ is integrable on $X_\Sigma$ and lower semicontinuous at the points $\iota_r(t)$ of a finite set $V\subset\mathcal O_{E,\Sigma}$. If $\sum_{y\ne0}|\widehat F(\xi(y))|\le L$, where $y$ runs over $\mathcal O_{E,\Sigma}$, then
\[
 \sum_{t\in V}F(\iota_r(t))
 \le \frac{(2/a_0)^f}{\prod_v r_v^2}
 \left(\int_{X_\Sigma}F+L\right).
\]
\end{lemma}
\begin{proof}
The diagonal $\Sigma$-integers form a lattice before scaling. Indeed $\mathcal O_{E,\Sigma}$ represents every coset in $\prod_{w\in \Sigma} E_w/\mathcal O_w$ by the Chinese remainder theorem for fractional ideals, and the residual lattice in the archimedean coordinates is $\mathcal O_E$. Thus its covolume after scaling is $v_L=2^{-f}\sqrt{|\operatorname{disc}E|}\prod_vr_v^2$. The displayed characters give its full dual: in the standard adelic pairing the annihilator of diagonal $E$ is $E$. A character surviving averaging on integral components outside $\Sigma$ is indexed by $y/\omega$ for $y$ integral off $\Sigma$, because the trace dual of $\mathcal O_w$ is the inverse different. The archimedean dilation divides its frequency by $r_v$.
On the quotient by the scaled lattice, periodize $F$ by its nonnegative translates. Tonelli gives an integrable periodic function with Fourier coefficients $v_L^{-1}\widehat F(\xi(y))$ when normalized quotient probability measure is used (up to reversing the character index). Absolute summability gives a continuous Fourier series $S$ equal almost everywhere to that periodization, with $S(0)\le v_L^{-1}(\int F+L)$. For every $\varepsilon>0$, lower semicontinuity at the finitely many points provides a neighborhood of zero on which $\sum_{t\in V} F(x+\iota_r(t))\ge\sum_{t\in V} F(\iota_r(t))-\varepsilon$. This finite sum is bounded by the periodization there. Arbitrarily small neighborhoods contain points of almost-everywhere equality with $S$; continuity at zero now gives the stated inequality.
\end{proof}
Taking $r_v=t_{va}$ gives the precise factor $P_a$. The scale parameter of the preceding frequency estimate is $\lambda_v=2\pi i/v(\omega)$, whose inverse modulus equals $g_v/M$ in the fitted normalization. Thus local absolute Fourier domination by the preceding majorant times $(\sum_p w_i(p))^2$ implies item 1 of the sieve for an integrable nonnegative density lower semicontinuous at the sampled points.

For the application, the density representative must have the stated
lower semicontinuity at the particular lattice samples. Measurability
of the cutoff in \S\ref{moment:rigidity} alone permits changing a value
at a sample without changing any Fourier coefficient.

\subsection*{Oscillation at one complex floor}
In the floor basis, the Hermitian matrix $J$ has an orthonormal summand. Choose disjoint pairs $B_1,B_2$ of axes in it, leaving $d$ orthonormal axes $I$ for Fourier transformation. The remaining axes, including the hyperbolic ones, may have an indefinite form. Write $m$ for the total complex coordinate dimension in this chart. Use ordinary area measure on $\mathbb C$, and write $\psi(z)=\exp(z-\bar z)$.
The quadric measure below uses the level coordinate $\operatorname{Im}\beta$ in the convention $|z_1|^2+|z_2|^2=2\operatorname{Im}\beta$ for a positive two-coordinate block, with the contribution of the other coordinates subtracted from the level.

\begin{lemma}
Let $\mu$ and $\mu'$ be the measures on two fixed levels of the Hermitian form in these coordinates, obtained by integrating Lebesgue measure in the coordinates outside $B_1$, respectively $B_2$, and the level measure on the sphere specified below in the remaining pair. For $R,R'>0$ write $Q_R=\{y:\max_j|y_j|\le R\}$. For $\xi\ne0$ and $h\in L^2(\mu)$ supported in $Q_R$, one has
\[
 \int_{Q_{R'}}\!\left|\int \psi(\xi v^*Jy) h(y)\,d\mu(y)\right|^2\!d\mu'(v)
 \ \le\ C(R)C(R')\left(\frac{\pi^2}{|\xi|^2}\right)^d\!\int |h|^2\,d\mu.
\]
Here $C(R)=4\pi^2 R^2(\pi R^2)^{m-d-2}$.
\end{lemma}
\begin{proof}
For squared radius $s>0$ in $\mathbb C^2$, the ball volume is $\pi^2s^2/2$. Differentiating against the level variable $s/2$ gives fibre mass $2\pi^2s$. Give it the uniform probability distribution on that sphere times this mass, and give the zero and negative radii mass zero. A sphere meeting a two-coordinate box of radius $R$ has $s\le 2R^2$, so each restricted fibre has mass at most $4\pi^2R^2$, independently of all outside coordinates and the level.

Write $y=(a,u,z)$ in the source chart, with $u=y_I$, $z=y_{B_1}$ and $a$ the remaining coordinates. Write $v=(a',u',z')$ with $u'=v_I$ and $z'=v_{B_2}$. Since the rows and columns on $I$ and $B_2$ agree with those of the identity and the two solved blocks are disjoint, the phase decomposes into
\[
 v^*Jy=(u')^*u+(z')^*a_{B_2}+L_{a'}(y).
\]
Here $a_{B_2}$ denotes the indicated coordinates of $a$. With $a'$ fixed, absorb $\psi(\xi L_{a'}(y))$ into $h$. Integrate over $z$ first against its fibre measure to obtain $H_{a'}(a,u)$. Cauchy--Schwarz in $z$ bounds $\int |H_{a'}|^2\,da\,du$ by $4\pi^2R^2\int |h|^2d\mu$, uniformly in $a'$. The transform in $u$ with phase $\psi(\xi (u')^*u)$ has squared $L^2$ norm multiplier $(\pi^2/|\xi|^2)^d$. Indeed its real frequency map is a scaled orthogonal map with absolute determinant $(|\xi|/\pi)^{2d}$ in the $2\pi$ Fourier convention. Cauchy--Schwarz in $a$ costs $(\pi R^2)^{m-d-2}$ and removes the phase depending on $z'$. The resulting bound for the squared output, before integration in $u'$, is uniform in $z'$. The output fibre and the integral in $a'$ cost the two factors in $C(R')$. Plancherel in $u$ now gives the inequality.
\end{proof}

For the common-entry Fourier term of the cut Bessel kernel, the phase at floor $j$ is $\psi(\xi\langle v,y\rangle)$ with $y=\sigma^{-1}f_j$ the orbit point of the frame vector. The partial correlation on the next floor descends to $y$ when it is invariant under the stabilizer of $f_j$, and the floor quotient measure must push forward to the quadric measure just used. Under those identifications and a bound $R$ for the support in this basis, it can be taken as $h$ in the lemma. A fixed translation in $v$ multiplies $h$ by a unit phase and preserves the estimate. Thus the factor $|\xi|^{-d}$ in operator norm comes directly from $d$ free complex coordinates, with two disjoint pairs used to handle the two quadric equations. To obtain the full fitted floor estimate one still needs the actual support radii in each position, the bounds for the remaining correlation norms, and their integrated shell estimates. The other local phase arising from unequal last entries additionally requires the Heisenberg equivariance identity.

The head weights, frequency decay sum and common-entry Poisson step correspond to the
\lineSrc{New/Spanning/MainTerm/GlobalSieve}{1103}{head-weight bound},
\lineSrc{New/Spanning/MainTerm/GlobalSieve}{818}{frequency sum}, and
\lineSrc{New/Spanning/MainTerm/GlobalSieve}{416}{scaled Poisson inequality}.
For the passage from local Plancherel to the full correlation bound see the
\lineSrc{New/Spanning/Frames/RealBandLimits}{648}{archimedean floor estimate}
and the \lineSrc{New/Spanning/MainTerm/FiniteSieveSystem}{116}{finite-place density estimate}.

\section{Polynomials for the twisted Selmer groups}

\subsection{Auxiliary primes, variation of local conditions, and a polynomial for the length}
\label{sec:selmer-polynomial-descent}
Let $O$ be the integer ring of a finite extension of $\mathbb Q_\ell$, with uniformizer
$\pi$ and residue field $k$. Write $v(\ell)=1$, and use $\lambda_O$ for ordinary
composition length. The normalized length is $\lambda_O/v_\pi(\ell)$.
We first explain how the auxiliary primes for several lattices can be chosen by a
single descent, and then construct the polynomials governing the finite-level lengths.
Throughout this subsection the finite-level modules consist of cocycles satisfying
the specified linear relations. Passing to a quotient by coboundaries in the
middle of the argument would change the lengths.

\begin{lemma}[Varying a local condition]\label{lem:selmer-tilt}
Let $C$ have finite length over $O$, and let $A,B:C\to(O/\pi^h)^d$ be linear.
There is a set of at most $d$ elements of $k$ such that for every $a\in O$
outside those residue classes,
\[
 \lambda_O\ker(aA-B)\ \leq \lambda_O\{z\in\ker A:Bz\in A(C)\}.
\]
\end{lemma}
\begin{proof}
Put $C'=\sum_{u\in O}\ker(uA-B)$. On each summand $Bz=A(uz)$, so
$B(C')\subseteq A(C')$. Let $L$ be the inverse image of $A(C')$ in $O^d$.
It is free of rank at most $d$. Write $p:L\twoheadrightarrow A(C')$ for reduction.
By freeness lift $p$ to $s:L\to C'$ through $A$, then lift $Bs$ through $p$
to an endomorphism $U$ of $L$. Thus $(aA-B)s=p(a-U)$. Away from the roots
in $k$ of the characteristic polynomial of $U\bmod\pi$, $a-U$ is invertible.
Consequently $(aA-B)(C')$ contains $A(C')$. Subtracting the image lengths
from $\lambda_O C'$ bounds the length of the kernel on $C'$ by
$\lambda_O(C'\cap\ker A)$. The entire kernel of $aA-B$ belongs to $C'$,
and $B(C')\subseteq A(C)$. This proves the assertion.
\end{proof}

Fix a torsion level $\ell^M$, and put $e=v_\pi(\ell^M)$.
For a finite family of lattices of ranks $n_i$, write $Z_{i,h}(J)$ for the
cocycles modulo $\pi^h$ unramified outside the designated places and $J\cup\bar J$,
satisfying the original relations and the following new conditions.
At each $w\in J$ fix a tame inertia generator $\gamma_w$, a Frobenius
$\phi_w$ and a slope $s_{i,w}\in O$. Require
\[
 z(\gamma_w)=s_{i,w}z(\phi_w),\qquad
 z(\gamma_w^{-c})=s_{i,w}z(\phi_w^{-c}),\qquad
 g^{-c}=(cgc^{-1})^{-1}
\]
for the chosen complex conjugation $c$. Take the auxiliary primes and conjugates
off the old ramification, distinct from one another, and prime to $\ell$.
At a prospective new prime $w$, let $C_{i,h}$ allow ramification at $w,\bar w$
without imposing either new relation. The maps $A,B$ evaluate a cocycle on the
two indicated inertia and Frobenius elements, respectively. Here $d=2n_i$.

The equality $\ker A=Z_{i,h}(J)$ on $C_{i,h}$ follows directly. At either new
place the original lattice is unramified, so the cocycle restricts to a
homomorphism on inertia. Its image is an $\ell$-group. Wild inertia and the
prime-to-$\ell$ part of tame inertia map trivially, and the $\ell$-part of
tame inertia is procyclic. Vanishing on the chosen generator therefore
removes all new ramification. For a unit $a$, imposing slope $a^{-1}$ gives
the kernel of $aA-B$.

\begin{proposition}[Simultaneous descent]\label{prop:selmer-tilt-descent}
Suppose $|k|>2n_i e+1$ for every $i$. Suppose at some stage that an allowable
new prime for datum $i$ and a level $1\leq h\leq e$ satisfy
\[
 B(Z_{i,h}(J))\not\subseteq A(C_{i,h}).
\]
Here allowable means that the auxiliary data extend to the new prime with
arbitrary slope for $i$ and zero slope for every other datum.
One can add this prime and decrease
$D(J)=\sum_i\sum_{j=1}^e\lambda_O Z_{i,j}(J)$ by at least one.
Starting at a singleton of slope zero, this gives a set of size at most
$1+\sum_{i,j}\lambda_O Z_{i,j}(\varnothing)$ at which
$B(Z_{i,j}(J))\subseteq A(C_{i,j})$, for every allowable prime, index and level.
\end{proposition}
\begin{proof}
For the chosen $i$ take the union of the exceptional sets of
Lemma~\ref{lem:selmer-tilt} for all $e$ levels. At most $2n_i e$ residues
are excluded. Choose a nonzero remaining residue and a unit lift $a$.
At every level the new length is at most that of the old cocycles with
$Bz\in A(C_{i,j})$. This is a proper submodule at the indicated level,
so its length drops there. For other indices choose slope zero, which
simply removes the new ramification. All lengths are finite. For instance
each continuous cocycle at finite level, together with the given residual
representation, gives a homomorphism to a fixed finite affine group,
unramified off a fixed finite set enlarged by the ramification of the representation.
The finiteness theorem for number fields of bounded degree unramified
outside a fixed set gives finitely many possibilities
(J.~S.~Milne, \emph{Algebraic Number Theory}, Theorem~8.42~\cite{milneANT}).
Each step spends one unit of a nonnegative integer. At the initial prime
slope zero gives precisely the original cocycles. The bound follows, and
is independent of the order $\ell^N$ of the auxiliary characters.
\end{proof}

\paragraph{A single matrix for all coefficient specializations.}
Let $G=G_E$. Fix a lattice $O^n$ with continuous action $\rho$, a finite
ramification set $S$, linear homogeneous relations on cocycle values with
coefficients in $O$, and locally constant characters
$h_j:G\to\mathbb Z/\ell^N$ for $j\in J$. With exponents read in
$[0,\ell^N)$ put
\[
 R_X(g)=\rho(g)\prod_{j\in J}(1+X_j)^{h_j(g)}.
\]
For an integral point $x$ over any finite coefficient extension $O'$,
let $Z(x)$ be the locally constant solutions in $(O'/\ell^M)^n$ of
$z(gh)=z(g)+R_x(g)z(h)$, the given relations and vanishing on inertia
outside $S$. This definition makes sense for arbitrary $x$.
At $x_j=\zeta_j-1$ with $\zeta_j^{\ell^N}=1$ the matrices give the
desired twisted action.

There is a polynomial matrix $A(X)$ with finitely many columns such that
$Z(x)=\ker(A(x)\bmod\ell^M)$ in the coordinates described below.
To see that one set of columns suffices, let $G'$ be the intersection of
the kernels of $\rho\bmod\ell^M$ and all $h_j$. Every solution restricts
on $G'$ to a homomorphism. Let $H$ be the common kernel of all continuous
homomorphisms $G'\to\mathbb Z/\ell^M$ unramified outside the places above
$S$. There are finitely many by the same bounded-degree finiteness theorem
over the fixed field of $G'$. Replace $H$ by its normal core in $G$.
Every such $z$ vanishes on $H$: composing its restriction with additive
characters into $\mathbb Z/\ell^M$ detects each nonzero element of its
finite abelian target of exponent dividing $\ell^M$. The equation for
$z(gh)$ now shows $z$ factors as a function on $G/H$. Use $n[G:H]$
unknowns, one $n$-vector for each coset. Write all the equations in these
unknowns as rows over $O[X_j:j\in J]$, retaining all original pairs
$g,h$ when writing the cocycle equation. Conversely any solution of this
matrix gives a locally constant function satisfying every equation.
The polynomial ring is Noetherian, so a finite subset of the rows
generates their row module and suffices over every specialization.

\begin{lemma}[One polynomial sharp at a chosen point]\label{lem:selmer-polynomial}
For every integral point $x_0$ over a finite coefficient extension there is
a polynomial $P\in O[X_j:j\in J]$ such that
\[
 \frac{\lambda_{O'}Z(x)}{v_{\pi'}(\ell)}\leq v(P(x))
\]
over every finite coefficient extension $O'$ with uniformizer $\pi'$,
with equality at $x_0$.
\end{lemma}
\begin{proof}
If $A$ has $d$ columns, append the $d$ rows of $\ell^M I_d$.
At each specialization the minimum valuation of its maximal minors is
the normalized length of $Z(x)$. Indeed, put $A(x)$ in Smith normal
form over $O'$. Writing $t_j$ for its elementary-divisor valuations in
uniformizer units, with $t_j=\infty$ for zero positions, the kernel
modulo $\ell^M$ has length $\sum_{j=1}^d\min(t_j,v_{\pi'}(\ell^M))$.
The row lattice of the augmented matrix is the row lattice of $A(x)$
plus $\ell^M O'^d$, so its index has the very same length. Its index
ideal is generated by the maximal minors. Select a minor attaining the
minimum at $x_0$. Its valuation bounds the minimum at every other
point. Division by $v_{\pi'}(\ell)$ gives the assertion.
\end{proof}

For coefficients $K'/O'$ ($K'=\operatorname{Frac}(O')$) use the same
cocycle and linear relations. The $\ell^M$-torsion of this solution
module is exactly the finite-level module at a character point, via
$u\bmod\ell^M\mapsto u/\ell^M\bmod O'$. If its normalized length is
strictly below $M$, every solution is already killed by $\ell^M$.
Indeed the image of a continuous cocycle in its discrete target is finite
by compactness, so each solution is torsion. An element of order
greater than $\pi'^{v_{\pi'}(\ell^M)}$ would generate a cyclic
submodule whose $\ell^M$-torsion alone has normalized length $M$.
The polynomial therefore bounds the divisible-coefficient length
whenever $v(P(x))<M$.

\paragraph{Where the arithmetic enters.}
The descent proves its terminal inclusion without producing the
inequality $\lambda_O Z_{i,e}(J)\leq G_i+L_i$ from that inclusion.
Here $G_i$ is the supremum of lengths of equivariant maps from the
torsion lattice to its conjugate Tate dual and $L_i$ is the sum of
local excess lengths away from $\ell$. Section~\ref{sec:terminal-selmer} proves the implication using global duality and the reciprocity argument for these local conditions. The construction of the allowable prime characters
with the required values at one another's Frobenius elements is also
needed here. Section~\ref{sec:terminal-selmer} constructs the primes and proves the terminal bound under the stated hypotheses on the underlying Selmer data. Lemma~\ref{lem:selmer-polynomial} supplies the
polynomial once a point of small length has been found. A polynomial
sharp at the untwisted point need not be small there. The next
calculation explains how to produce a small point from the tied bound.

\paragraph{Releasing the local conditions along a line.}
Take $J$ and its distinguished prime $w_0$ as above, write $q=\ell^N$, and
assume the following local specifications. Every $w\in J$ has norm $1\bmod q$.
The anticyclotomic $h_j$ have values $h_j(\gamma_w)=\delta_{jw}$ and
$h_j(\phi_w)=\delta_{jw_0}$ for $w\ne w_0$ (the inertia prescription applies
also to $w_0$). The old $\rho$ is unramified at the new places.
For each nonzero slope at $w$, $\rho(\phi_w)$ and
$\rho(\phi_w^{-c})$ act trivially modulo $\ell^M$.
Fix one datum of rank $n$, and suppose that after releasing the zero slope
at $w_0$ the normalized untwisted length, with all the other ties retained,
is at most $\beta$. This last release costs at most
\[
 \sum_{g\in(\phi_{w_0},\phi_{w_0}^{-c})}v(\det(\rho(g)-1)),
\]
where a sum over a displayed pair counts both entries. Indeed the untwisted
tame relation gives $(\rho(g)-1)z(\gamma)=0$ modulo $\ell^M$ when $N\geq M$:
Frobenius raises tame inertia to the norm and that norm is $1\bmod\ell^M$.
Evaluation therefore embeds the quotient by the two restored ties into the
two kernels. Each kernel has normalized length at most its determinant
valuation by Smith normal form. If the Frobenius matrices agree with
prescribed comparison matrices modulo $\ell^M$, the same kernel bound
uses the comparison matrices.

For $0<u\leq1$ rational, $Nu\geq M$, take the line
$X_{w_0}=Y$ and $X_w=s_wY$ for $w\ne w_0$. Let $v(y)=u$ in a finite
coefficient extension. For any $\alpha\equiv1\pmod y$, the sum of
$\alpha^j$ for $0\leq j<\ell^N$ has valuation at least $Nu$.
Indeed, factor the sum into $N$ geometric sums of length $\ell$.
Each is congruent to $\ell$ modulo $y$, and $v(\ell)=1\geq u$.
It follows that all $(1+X_w)^q$ specialize to $1\bmod\ell^M$.
The point twist is a representation modulo $\ell^M$. Let $z$ be a solution
for that representation with all auxiliary ties removed. For
$w\ne w_0$ use the pair $\phi,\gamma$ at $w$ or its conjugate inverse.
In the affine representation $g\mapsto(t\mapsto R_x(g)t+z(g))$,
the inertia image has exponent dividing $q$: the linear part is scalar
$\alpha\equiv1\pmod y$ and the translation of its $q$th power is
$(1+\alpha+\cdots+\alpha^{q-1})z(g)=0$. It therefore factors through
the cyclic tame $\ell$-quotient. Frobenius commutes with this inertia
image since the norm is $1\bmod q$; so does its inverse. Comparing the
translations gives
\[
 ((1+y)\rho(\phi)-1)z(\gamma)=s_w y z(\phi).
\]
The hypotheses on the slope give
$s_w((1+y)\rho(\phi)-1)=s_w y I$ modulo $\ell^M$.
Thus the discrepancy $z(\gamma)-s_w z(\phi)$ is killed by
$(1+y)\rho(\phi)-1$. Evaluating all discrepancies embeds the quotient by
the retained ties into their respective kernels. With
\[
 H(Y)=\prod_{w\ne w_0}\,
       \prod_{g\in(\phi_w,\phi_w^{-c})}\det((1+Y)\rho(g)-1),
\]
the normalized cost of removing these ties is at most $v(H(y))$.

Apply Lemma~\ref{lem:selmer-polynomial} to the retained ties at the origin,
and restrict its polynomial to this line, obtaining $F(Y)$ with
$v(F(0))\leq\beta$. For any integral polynomial $f=\sum a_jY^j$ define
$v_u(f)=\min_j(v(a_j)+ju)$. Finitely many nonzero polynomials attain their
respective $v_u$ at the same point of valuation $u$: write $y=t\xi$ with
$v(t)=u$, divide each $f(tX)$ by a coefficient of minimal valuation, and
choose the nonzero residue of $\xi$ outside the roots of the resulting
nonzero residue polynomials. A finite unramified extension supplies this
residue. Since $H$ has unit leading coefficient and degree
$2n(|J|-1)$, at this point
\[
 v(F(y))+v(H(y))\leq\beta+2n(|J|-1)u .
\]
This bounds the untied normalized length at $x=(y,(s_wy)_{w\ne w_0})$.
Applying Lemma~\ref{lem:selmer-polynomial} at $x$ to the untied module now
gives the desired polynomial, with $v(P(x))$ bounded by the same quantity
and all coordinates of valuation at least $u$.

\paragraph{The group-ring polynomial.}
For an integral Stickelberger element $\Theta_J=\sum_b c_b[b]$, put
$P_J(X)=\sum_b c_b\prod_j(1+X_j)^{h_j(b)}$, using the nonnegative
representatives of the exponents. Here the characters descend to the ray
class group in which $b$ ranges. Substitution of roots of unity recovers exactly
the character evaluation. Suppose the one-prime norm relation multiplies this
evaluation by $(1-\alpha q_w^{-d})(1-\alpha^{-1}q_w^{-d})$ for each
$d$ in a fixed list of $r$ exponents, with $q_w$ the prime norm and
$\alpha=\vartheta(\phi_w)\psi(\phi_w)$ for a character of the lower level.
Assume the coefficients and $\vartheta$ are integral with $\vartheta$ unit
valued. Put $\beta'=v(\boldsymbol{1}(\Theta_\varnothing))+
2r v(\vartheta(g_0)-1)<\infty$ with $\vartheta(g_0)\ne1$.
Choose $M'>\beta'+1$, and require $Nu\geq M'$ and
$v(\vartheta(\phi_{w_0})-\vartheta(g_0))\geq M'$. At $w_0$ and the trivial
character the Euler factors are congruent modulo valuation $M'$ to
$1-\vartheta(g_0)^{\pm1}$ since $q_{w_0}\equiv1\pmod{\ell^N}$.
If $r>0$ these differences have valuation $<M'$; if $r=0$ there is no factor.
Hence the one-variable evaluation $F_0(Y)$ of $\Theta_{\{w_0\}}$ using
$h_{w_0}$ has $v(F_0(0))\leq\beta'$.

Set $Z=1+Y$ and $D=r(|J|-1)$. On the line $X_{w_0}=Y$, $X_w=0$ otherwise,
denote $P_J$ by $P_*(Y)$. At every $Z$ of order dividing $\ell^N$, successive
applications of the norm relation give
\[
 Z^D P_*(Y)= F_0(Y)
 \prod_{w\ne w_0}\prod_d
 (1-A_wq_w^{-d} Z)(Z-A_w^{-1}q_w^{-d}),
 \qquad A_w=\vartheta(\phi_w).
\]
Indeed $h_{w_0}(\phi_w)=1$ at these primes. Subtracting the right side
from the left yields an integral polynomial divisible by $Z^{\ell^N}-1$,
by monic division and the $\ell^N$ distinct roots in characteristic zero.
The product on the right has degree $2D$ and unit leading coefficient.
At a point $v(y)=u$ generic for this product and $F_0$, its value times
$F_0(y)$ has valuation at most $\beta'+2Du$. For $2Du<1$ this
is strictly below $M'$. Meanwhile $v((1+y)^{\ell^N}-1)\geq Nu$ by the
geometric sum estimate above. The congruence therefore gives
$v(P_*(y))\leq\beta'+2Du$, since $1+y$ is a unit.
This supplies the required witness on the same auxiliary set for each system.

Thus, given allowable prime data and the terminal Selmer length estimate,
one constant $s_0\geq2|J|\max(1,\max_i n_i,\max r)$ controls the number
of variables and all witness losses. Proposition~\ref{prop:selmer-tilt-descent}
makes its choice uniform in $N$ when those prime data exist for every $N$.
One may enlarge $N$ after fixing $u$ to satisfy both torsion thresholds.
The individual witnesses may differ. The common roots-of-unity specialization
in the preceding argument applies to the polynomials so obtained.
The construction of the Galois representations and the identification of the
actual critical-value ratios with group-ring elements satisfying the displayed
norm relation remain separate arithmetic obligations.


\subsection{Arithmetic at the terminal set of auxiliary primes}
\label{sec:terminal-selmer}
We give the length estimate used in \S\ref{sec:selmer-polynomial-descent}
for a Selmer datum satisfying the following hypotheses. Write $T=O^n$,
$G=G_E$ for a CM field $E$, and $S$ for a finite set containing the primes over $\ell$, outside which $\rho$ is unramified. Fix the torsion level $M$ and take $N>M$. Fix the common element $g_0$ in inertia at a bad place away from $\ell$ and a common open neighborhood imposing the prescribed congruences there. At each $v\mid\ell$
specify a finite free lattice $Z_v$ of local integral cocycles containing
the coboundaries. Impose on the torsion cocycles all the homogeneous
linear relations among evaluations valid on $Z_v$. Require
\[
 \sum_{v\mid\ell}\operatorname{rank} Z_v
 \leq (\#\{v\mid\ell\}+r_2(E))n .
\]
Assume $\rho(G)$ spans $M_n(O)$ over $O$. Suppose there is
$\sigma$ in inertia at a prime of $S$ away from $\ell$ such that
$\rho(\sigma)=\rho(\sigma^{-c})=\xi I$, with $\xi^2-1$ a unit.
Add the normalization $z(\sigma)=0$ to the global cocycles.
For $W=T/\pi^j$, let $W^*=\operatorname{Hom}(W,\mu_{\ell^\infty})$
be the Tate dual with the contragredient action and let $W^{*c}$ carry
the action pulled back by $g\mapsto g^c=cgc^{-1}$. In additive
notation for characters, scalars in $O$ act by precomposition.
Define $G_T=\sup_j\lambda_O\operatorname{Hom}_G(W,W^{*c})$.
For $v\nmid\ell$ let $b_v$ be the supremum over $j$ of the positive
part of $\lambda_O C_v(W)-\lambda_O W$, where $C_v(W)$ consists of
restrictions to $G_{E_v}$ of all continuous global cocycles with values
in $W$ and unrestricted ramification. Put $L_T=\sum_{v\in S,\ v\nmid\ell} b_v$.
Suprema over all nonzero scalar ideals give the same definitions.
The argument presumes finiteness of these terms as in the datum.

\paragraph{Producing allowable primes.}
Put $q=\ell^N$. The following construction also explains the choice of
the first prime $w_0$. Fix finitely many conjugation-stable places $B$
containing the places over $\ell$, at which the new character is to be
locally trivial. Let $D$ be a finite set disjoint from $B$ and let
$g$ belong to the normal subgroup generated by inertia at $D$. For any open
neighborhood of $g$, one can choose a degree-one prime $w\ne\bar w$ off a
prescribed finite exceptional set, a Frobenius lift in that neighborhood
and an anticyclotomic character $h:G\to\mathbb Z/q$ trivial locally at $B$,
unramified off $w,\bar w$, and onto on inertia at $w$. Its tame generator
may be chosen to have image $1$ and its Frobenius lift to have image $0$.

Indeed enlarge $B$ to a finite set $B'$ disjoint from $D$ and generating
the ideal class group: choose representatives prime to the given excluded
set and take their prime divisors. The extension of $E$ formed by adjoining $\mu_q$
and $q$th roots of generators of the finitely generated $B'$-unit group
is finite Galois and unramified at $D$, since these generators are units
there and $q$ is invertible. Its Galois kernel contains $g$. Apply
Chebotarev~\cite[Theorem~8.31]{milneANT} to the compositum with a finite Galois quotient specifying
the neighborhood: one can choose a place above $w$ with the prescribed
Frobenius element in that quotient. The set of possible primes has
positive Dirichlet density, hence contains degree-one primes since
the primes of higher residue degree have convergent norm-reciprocal sum.
Avoid primes ramified over $\mathbb Q$. A degree-one unramified prime
cannot be fixed by $c$, since its decomposition group in $E/E^+$ is trivial.
We also exclude $B'$ and any other specified finite set.
The residue field at $w$ contains $\mu_q$ and the $B'$-units reduce
to $q$th powers. A surjective discrete logarithm from its multiplicative
group to $\mathbb Z/q$ therefore kills every $B'$-unit.
On finite ideles which are units off $B'$ take the logarithm of
the residue at $w$. This map descends to a character of all finite idele
classes: principal ideles and this subgroup generate, and their
intersection consists of principal $B'$-units. Its kernel is open.
Precompose with $a\mapsto a/\bar a$ and use global Artin reciprocity.
The result has the claimed ramification and is trivial on the full local
multiplicative groups at $B$ because $B$ is conjugation-stable and
$w,\bar w\notin B$. It surjects on inertia since $w\ne\bar w$.
Multiply the chosen Frobenius lift by a power of its tame generator to
give character value zero. This preserves its image in the prescribed
finite quotient since $w$ can be chosen unramified there.

Start with $B=\{v\mid\ell\}$, and use $g_0$ in the prescribed bad
inertia away from $\ell$. Take $w_0$ with Frobenius close enough to
$g_0$ for all the finitely many congruences on the data and the
Stickelberger systems. Exclude the rational primes below the whole prescribed bad set when choosing $w_0$, and exclude the rational prime below $w_0$ in every subsequent choice. Take all slopes zero there. For extensions of
$J$, put $B=\{v\mid\ell\}\cup(J\setminus\{w_0\})\cup
\overline{(J\setminus\{w_0\})}$. For the datum being shortened let
$H$ be the subgroup generated by $[g,\sigma]$ and $[g,\sigma^c]$,
$g\in G$, using $[a,b]=aba^{-1}b^{-1}$. It is normal, since
$a[g,b]a^{-1}=[ag,b][a,b]^{-1}$, and stable under $c$.
Its image under $\rho$ and every abelian character is trivial. For
any $h\in H$, take as target $g=\gamma_0 h$, where $\gamma_0$ is the
tame generator at $w_0$. This lies in the normal subgroup generated by
inertia at $w_0$ and the primes for $\sigma,\sigma^c$, disjoint from $B$.
At this target the old characters take values $1$ for $h_{w_0}$ and $0$
for the others. Both $g,g^c$ act trivially modulo $\ell^M$ and on
$\mu_q$. These are open conditions and so can be imposed on the new
Frobenius. The new character is trivial at the old later primes, even
on their Frobenius elements, by local triviality at $B$. We can keep
the new prime and its conjugate away from $S_0$, all old primes and
the rational prime below $w_0$. This extends all the local specifications
of \S\ref{sec:selmer-polynomial-descent}, for arbitrary slope on the
chosen datum. In particular, allowable Frobenii approximate every
$\gamma_0 h$ with $h\in H$ arbitrarily closely in $G$. The old characters are held constant and the actions on $T/\ell^M$, its conjugate and $\mu_q$ held trivial by shrinking the neighborhood.

\paragraph{Comparing a Selmer group with its dual.}
Fix $j\le e$. Define local subgroups $\mathcal L_v\subset H^1(E_v,W)$
as follows. At $v\mid\ell$ use cocycles satisfying the relations of
$Z_v$, modulo coboundaries. At $v\in S$ away from $\ell$ use the image
of $C_v(W)$. At a later auxiliary prime and its conjugate use the
cocycles satisfying the tie, modulo coboundaries. Elsewhere use
unramified classes, including at $w_0,\bar w_0$. The local conditions
contain coboundaries before quotienting: at a later prime inertia acts
trivially, and a nonzero slope requires $\rho(\phi)=1$ on $W$.
Every global class has a unique representative vanishing at $\sigma$,
because $\xi-1$ is invertible. Consequently the normalized tied cocycles
are exactly $H^1_{\mathcal L}(E,W)$. Use the orthogonal complements
$\mathcal L_v^\perp$ for the local Tate pairings to define its dual
$Y=H^1_{\mathcal L^\perp}(E,W^*)$. The finite-module Greenberg--Wiles
formula (A. Smith, \emph{Galois cohomology and the Greenberg--Wiles
formula}, Theorem 1.8~\cite{smith}) gives
\[
 \lambda_O Z_{i,j}(J)-\lambda_O Y
 = \sum_{v\text{ finite}}
       (\lambda_O\mathcal L_v-\lambda_O H^0(E_v,W))
       -r_2(E)\lambda_O W .
\]
Here $H^0(E,W)$ and $H^0(E,W^*)$ vanish since $\sigma$ acts by $\xi$
and $\xi^{-1}$ respectively, the cyclotomic character being trivial at
$\sigma$. Taking logarithms of the cardinality formula to base $|k|$
gives these lengths. At an unramified place the local difference is
zero. At any other finite place, if $D_v$ is the specified local
cocycle group, the difference equals $\lambda_O D_v-\lambda_O W$,
by the exact sequence defining coboundaries. At a tied place a cocycle
is determined by its values on inertia and Frobenius, and the former
value is determined by the latter. Thus $\lambda_O D_v\leq\lambda_O W$.
At a place in $S$ away from $\ell$, the excess costs at most $b_v$.

For $v\mid\ell$ let $r_v=\operatorname{rank} Z_v$.
All scalar coordinate evaluations on $Z_v$ generate a free
submodule of $Z_v^\vee$ of rank at most $r_v$. Each local solution to
the defining linear relations gives a homomorphism from this submodule to
$O/\pi^j$, by sending a coordinate evaluation to the corresponding
coordinate of the solution. Every dependence on the evaluations is a
relation by definition. This map on solutions is injective, so
$\lambda_O D_v\le jr_v$. The stated bound on the sum of ranks offsets the
possible archimedean loss. We obtain
\[
                    \lambda_O Z_{i,j}(J)\leq\lambda_O Y+L_T .
\]
The choice of $C_v(W)$ at old bad places matters: using all local
cocycles there could enlarge the local error in the definition of $L_T$.

\paragraph{Controlling the dual at a terminal set.}
Identify the underlying constant group of roots of unity of order dividing
$\ell^M$ with the corresponding subgroup of $\mathbb Q/\mathbb Z$
using a compatible primitive root; all Galois actions retain their
cyclotomic factors. Suppose $z\in Z_{i,j}(J)$ has some value with a unit
coordinate modulo $\pi^j$. On $H$, $z$ and each dual cocycle $y$
are homomorphisms. The span of $z(H)$ is $W$: it is stable under
$\rho(G)$, and $z([g,\sigma])=(1-\xi)z(g)$ since $z(\sigma)=0$.
A vector with a unit coordinate generates $W$ under the full matrix
algebra. The same spanning assertion holds for $z(H^c)$.

For $t\in O$ define the biadditive function
\[
 B_t(h,h')= y(h')(t z(h))-y(h^c)(t z(h'^c)),\qquad h,h'\in H .
\]
We show $B_t(h,h)=0$. The restrictions of $y$ to inertia at $w_0$ and
$\bar w_0$ are zero: the dual of the unramified condition there is
unramified, and inertia acts trivially. The same holds for $z$ by the
zero tie. Their values at $\gamma_0 h$ and $(\gamma_0h)^c$ therefore
coincide with those at $h,h^c$. Choose a new auxiliary prime as in the
preceding construction, with Frobenius $\phi$ sufficiently close to
$\gamma_0h$ that these values of $y,z$ and their conjugates agree.
They are locally constant. Both $\phi$ and $\phi^c$ act trivially
on $W,W^*$. The terminal inclusion produces a cocycle $z'$ allowed to
ramify at this new pair with
\[
 z'(\gamma)=z(\phi),\qquad z'(\gamma^{-c})=z(\phi^{-c})=-z(\phi^c).
\]
The sum of the local invariants of the global cup product of $tz'$ and
$y$ vanishes. At old places the classes are orthogonal by construction;
away from the new pair and the old places both are unramified. At the
new pair the dual class is unramified. For trivial action the local
pairing of a cocycle and an unramified dual class evaluates the latter
on Frobenius against the value of the former on normalized tame inertia
(up to a common sign from the cup-product order). In detail, cup
product followed by the tame residue contracts the restriction to
$\operatorname{Hom}(I^{\rm tame},W)$ with the unramified Frobenius
value; the normalization sends the tame generator acting on a
$q$th root of a uniformizer by the chosen primitive $q$th root to
$1$. For a cyclic summand of order $m\mid\ell^M$, identify it locally with $\mu_m$ using the compatible root. A Kummer class of $a\in E_w^\times$ has tame value $v_w(a)\bmod m$ in this normalization: the Kummer extension of a unit is unramified since $m$ is prime to the residue characteristic. Pairing with an unramified character evaluates that character on the local Artin symbol of $a$, hence on $v_w(a)$ times Frobenius. Additivity over the cyclic summands gives the assertion for $W$. At the conjugate place the normalized generator is the inverse
conjugate because $c$ inverts roots of unity. Replacing the two normalized
generators by $\gamma,\gamma^{-c}$ multiplies both evaluations by the
same unit modulo $\ell^M$. Thus the vanishing reads
$y(\phi)(t z(\phi))-y(\phi^c)(t z(\phi^c))=0$, as claimed.

Polarizing gives $B_t(h,h')+B_t(h',h)=0$. Conjugate its first argument
by $\sigma$. The two transformation rules are
\[
 B_t(\sigma h\sigma^{-1},h')=B_{\xi t}(h,h'),\qquad
 B_t(h',\sigma h\sigma^{-1})=B_{\xi^{-1} t}(h',h).
\]
Indeed $\sigma$ and $\sigma^c$ have opposite scalar actions on $W$,
and the dual actions invert these scalars. Substituting $\xi t$ in
the polarized identity with conjugated argument and subtracting the
original identity gives $B_{(\xi^2-1)t}(h,h')=0$. Thus $B_t=0$ for all $t$.

The assignment $z(h)\mapsto y(h^c)$ extends uniquely to an
$O$-linear map $\eta_y:W\to W^{*c}$. To see well-definedness, take a
relation $\sum_j a_j z(h_j)=0$. Evaluating
$\sum_j a_j y(h_j^c)$ on $t z(h'^c)$ gives zero for all $h',t$, by
$B_{a_jt}(h_j,h')=0$. These vectors span $W$ additively, including
all their scalar multiples. Conjugating $h$ by $g\in G$ gives
$\eta_y(\rho(g)z(h))=g^c\cdot y(h^c)$ since $H$ acts trivially
on both modules. Hence $\eta_y$ is equivariant. Coboundaries vanish
on $H$, so this is a linear map on dual classes. If $\eta_y=0$,
$y$ vanishes on $H$. Normalize $y(\sigma)=0$ by its unique
coboundary since $\xi^{-1}-1$ is a unit. Then
$y([g,\sigma])=(1-\xi^{-1})y(g)$ shows $y=0$.
Consequently $\lambda_O Y\leq G_T$ whenever such a $z$ exists.

Finally let $j$ be the largest exponent of the cyclic submodules of
the tied module at level $e$ (so these cyclic orders are $\pi^j$).
If $j=0$ the module is zero. Otherwise every value of every such cocycle
is in $\pi^{e-j}W_e$, for $W_e=T/\pi^e$. Identify $T/\pi^j$ with
this submodule by multiplication by $\pi^{e-j}$. Dividing values in
this identification embeds the top-level module in $Z_{i,j}(J)$:
all equations are homogeneous over $O$ and all inertia restrictions
are preserved. A cocycle of maximal order yields a unit coordinate
somewhere. The terminal inclusion was obtained for every level.
Applying the preceding argument at $j$ now gives
$\lambda_O Z_{i,e}(J)\leq G_T+L_T$ for every datum.
This closes the terminal estimate for the indicated Selmer data.

\paragraph{The numerical bound from representation data.}
Write $f=r_2(E)$ and $G_0$ for the order of the product of finite general
linear groups in the determinant datum. Suppose $S$ consists of the places
over $\ell p_0p_1$. At each $w\in S$ away from $\ell$ assume there is an
inertial $u_w$ with $v(\det(\rho(u_w)-1))\le n v_\ell(G_0)$. The local
Euler characteristic formula at residue characteristic different from $\ell$,
together with Tate duality~\cite[Part~I]{milneADT}, gives
\[
 \lambda_O Z^1(E_w,W)-\lambda_O W=\lambda_O H^0(E_w,W^*).
\]
Indeed $\lambda H^1=\lambda H^0+\lambda H^2$ and
$\lambda B^1=\lambda W-\lambda H^0$. The group of dual invariants is
killed by $\rho^*(u_w)-1$. Since inertia fixes the $\ell$-power roots of
unity here, its kernel has length at most $v_\pi(\det(\rho(u_w)-1))$
by duality and Smith normal form. This bounds the restricted cocycles
as well. There are at most $4f$ such places. Assume the same determinant
bound for each of $g_0,g_0^{-c}$, giving two further terms for the first
prime.

For the global term suppose there is $g\in G$ acting by a scalar unit
$\alpha$, with $g^{-c}$ acting by another scalar unit $\alpha'$.
Each equivariant map $W\to W^{*c}$ takes values killed by
$\delta=\alpha-\chi_{\rm cyc}(g)\alpha'$.
Its value at a fixed standard basis vector determines the map, since
the full matrix algebra is spanned by the action. The $\delta$-torsion
of $W^*$ has length at most $n v_\pi(\delta)$, hence this also bounds
the length of the space of such maps. In particular, if for some
$h\in\{0,1\}$ one has $\chi_{\rm cyc}(g)=p_h^n$ and
$\alpha'=\chi_{\rm cyc}(g)^{t-1}\alpha$ with $t\ge1$, the normalized
global term is at most $n v_\ell(p_h^{nt}-1)$. When $nt$ is at most
the dimension of the general linear block at $p_h$, the integer
$p_h^{nt}-1$ divides its order: the order contains the factor
$p_h^{d nt}-1$ for residue degree $d\ge1$. Thus the normalized sum of the
global, local, and first-prime terms is at most
$(4f+3)n v_\ell(G_0)$. This argument needs the scalar and inertial
determinant information on the actual representations; it does not
construct an attached representation from a Hecke eigensystem.

For the critical ratios take the specified tame generator $g_0$ at the
chosen $w\nmid\ell$ over a level prime. Suppose its Stickelberger
character $z=\vartheta(g_0)$ satisfies $z\ne1$ and $z^{q_w-1}=1$,
and its list of Euler exponents has size at most two. Then
$(z-1)\mid(q_w-1)$ in the valuation ring: reducing
$0=1+z+\cdots+z^{q_w-2}$ modulo $z-1$ gives the divisibility.
In particular $\beta'-v(\boldsymbol{1}(\Theta_\varnothing))
\le4v_\ell(q_w-1)$. Since the chosen general linear block is over
the field of size $q_w$, $q_w-1\mid G_0$. Assume the interpolation
identity at each character $\psi$ used in the argument is
\[
 v(\psi(\Theta_J))=v(\boldsymbol{1}(\Theta_\varnothing))
                        +d\,v(\iota(R_\psi)),\qquad d\in\{1,2\},
\]
where $R_\psi$ is the critical ratio with the stated Euler factors removed.
With the simultaneous specialization error $2^{-\ell}$ this gives
$d\,v(\iota(R_\psi))\le4v_\ell(G_0)+2^{-\ell}$. A negative valuation
already satisfies the required upper bound. For $1\le n\le r$ and $f\ge1$,
both this estimate and the normalized Selmer length estimate give the common
allowance
\[
                     u_\ell=(4f+3)r v_\ell(G_0)+2^{-\ell}
\]
for all the indicated powers. In particular the inverse norm of the ratio is
at most $\ell^{u_\ell}$. The interpolation identity and the prescribed
inertial and scalar properties are the application-specific inputs in
this deduction.

\section{The rank-one character and the auxiliary path}

\subsection{The Galois line for a main critical pair}
\label{rankone:main}
Let $(\nu,z)$ be one of the main pairs of \S11, with source $U(1)$, target $U(2z)$ and $\nu$ of weight $-c_v$ at the complex place $v$ of $E$. The constant source system has Galois character $1$, inertia polynomial $X-1$ at every chosen place, and Hodge--Tate weight zero. The character needed after normalization is different. We describe it explicitly. Use the Artin map sending a uniformizer to geometric Frobenius in the unramified quotient, and count the cyclotomic character $\epsilon_\ell$ with Hodge--Tate weight $-1$.

For a full idele $x$ write $|x|_{\mathbb A_E}$ for the product of the normalized absolute values (squared modulus at a complex place). Set
\[
 \Theta(x)=\nu(x/\bar x)^{-1}|x|_{\mathbb A_E}^{1-z}.
\]
It is trivial on $E^\times$ by the product formula and automorphy of $\nu$. Let $\sigma$ run over the embeddings $E\hookrightarrow\mathbb C$, and let $s_\sigma=+1$ on the chosen CM type and $-1$ on its conjugate. The infinite component is $\prod_\sigma\sigma(x_\infty)^{a_\sigma}$, with
\[
 a_\sigma=1-z+s_\sigma c_v,\qquad v=\{\sigma,\bar\sigma\}.
\]
Indeed $\nu^{-1}$ evaluates on $x_v/\bar x_v$ with exponent $c_v$, and $|x_v|_v^{1-z}$ contributes $1-z$ to both exponents. The half-rank shift is $(1-(2z-1))/2=1-z$.

Fix $\iota:\overline{\mathbb Q}\hookrightarrow\overline{\mathbb Q}_\ell$, identifying algebraic values in $\mathbb C$ through a fixed archimedean embedding. If $x$ is a finite idele, put
\[
 t_\ell(x)=\iota(\Theta_f(x))
        \prod_{\sigma} (\iota\sigma)(x_{w(\iota\sigma)})^{a_\sigma},
\]
where $w(\iota\sigma)$ is the place over $\ell$ induced by $\iota\sigma$, and the latter map on a local component means its continuous extension. The smooth finite part has algebraic values in a fixed number field. To see this, take a modulus killing its unit character, choose finitely many representatives of the ray class group, and raise each representative to the order of its class. Its finite character value then has a power equal to a product of conjugates of a field element. These finitely many algebraic values together with the finite unit image and the normal closure of $E$ suffice. Thus $t_\ell$ takes values in a finite extension of $\mathbb Q_\ell$. It is continuous, and on a diagonal $u\in E^\times$ the two factors cancel because $\Theta_f(u)=\prod_\sigma\sigma(u)^{-a_\sigma}$. Global class field theory gives a character $T_\ell:G_E\to\overline{\mathbb Q}_\ell^\times$ whose pullback is $t_\ell$. In detail the finite idele class quotient by the closure of $E^\times$ is compact (the ideal class group is finite and finite local units form a compact group); over a totally imaginary field it is the global reciprocity quotient. Compactness makes the valuation of $T_\ell$ identically zero, so its coefficient ring of integers supplies a stable rank-one lattice.

At $w\nmid\ell$ the local character under reciprocity is $\iota\Theta_{f,w}$, including its tame unit character. At an unramified place $T_\ell(\operatorname{Frob}^{\rm geom}_w)=\iota\Theta_{f,w}(\varpi_w)$. At $w\mid\ell$ its restriction to local units is the product of the finite character $\iota\Theta_{f,w}$ and $\prod_{\sigma\mid w}(\iota\sigma)^{a_\sigma}$. On a sufficiently small open subgroup the finite factor is trivial. The rank-one local character theorem (Tate--Lubin--Tate) gives the labelled Hodge--Tate weight
\[
   -a_\sigma=z-1-s_\sigma c_v .
\]
With geometric reciprocity the character whose pullback on units is the identity embedding is the Lubin--Tate character at that embedding; it has weight $-1$ there and zero at the other embeddings. This fixes the sign. The local character is potentially crystalline and is crystalline when its finite factor is unramified. A general reference for the unit criterion and labelled weights is B.~Conrad, \emph{Lifting global representations with local properties}, Appendix B, especially Proposition B.4. Moreover
\[
 T_\ell^c T_\ell=\epsilon_\ell^{2(1-z)}.
\]
Indeed complex conjugation inverts the $\nu$-factor, and the avatar of the idele norm is $\epsilon_\ell$, as one checks on the dense set of unramified geometric Frobenius elements. A finite order twist $\eta$ on $U(1)$ multiplies $T_\ell$ by the avatar of $x\mapsto\eta(x/\bar x)^{-1}$, so all labelled weights are unchanged.

The line, its normalization and its local invariants are thus available before constructing a representation for the lifted form. To identify it in that representation one still needs the unramified theta-lift Hecke formula and the attachment of a representation to the target system. More precisely, if at almost all places the normalized target trace at geometric Frobenius equals the trace of the twisted line plus the sum of $\epsilon_\ell^{-i}$ for $0\le i<2z-1$, Chebotarev density and continuity extend that equality to $G_E$. Equality of traces in characteristic zero identifies semisimplifications. This derivation has no implication for integral extension classes between the constituents, which enter the congruence estimate in R.


\subsection{The normalized pairing at the endpoint of the auxiliary path}
\label{pathmath:dual}
We give the two assertions about the dual Weyl pairing used to recover an integral form from its period pairings in \S11.5. These assertions include primes dividing the hook number.

Fix a finite ordered collection of column heights $0<s<n$ at each of the infinite places; regard determinant powers separately. Over $\mathbb Z$ take the tensor product $T$ of the exterior powers $\bigwedge^s\mathbb Z^n$, with the tensor wedge basis orthonormal. For a tuple of square matrices $X$, write $e(X)$ for the tensor whose factor for a column of height $s$ is the wedge of the first $s$ columns at that place. The coordinates of $e(X)$ are products of flag minors. Let $P$ be the integral span of these polynomials, and give $W=P^\vee$ its evaluation realization in $T$. More explicitly, the coordinate map $T^\vee\twoheadrightarrow P$ sends a linear functional $a$ to $a(e(X))$. The group $P$ is free of finite rank as a subgroup of the homogeneous polynomials. Hence the dual inclusion $W\hookrightarrow T$ is split over $\mathbb Z$. In particular $W=T\cap W_{\mathbb Q}$. Over $\mathbb C$, the evaluations span $W_{\mathbb C}$, since a polynomial vanishing on all square matrices is zero. Invertible matrices suffice by density. The Weyl metric here is precisely the restriction of the tensor metric; its kernel on real frames is
\[
 K(Y,X)=e(Y)^t e(X)=\prod_{\tau,s}\det((Y_\tau^t X_\tau)_{[1,s],[1,s]}),
\]
with a factor for each column. This identity follows by Cauchy--Binet.

Replace each height $s$ by $n-s$ to obtain $T',W'$. Apply the oriented complementary-wedge map in each factor, and reorder tensor factors. The resulting signed permutation of the tensor bases preserves their integral lattices and their metrics. It maps $W_{\mathbb Q}$ onto $W'_{\mathbb Q}$: for any invertible $g$ at a place, the image of its first $s$ columns wedged together is $\det(g)$ times the wedge of the last $n-s$ columns of $g^{-t}$, up to the fixed ordering sign. This is Jacobi's complementary-minor identity, obtained by wedging a basis minor with its complementary wedge and evaluating the resulting top form. Reversal of the column order of the frame makes these the first columns of a new frame. The inverse complementary map gives the reverse inclusion. Saturation now gives an integral isometric isomorphism $J:W\to W'$. Its sign is fixed by sending the highest unit frame to the prescribed opposite unit frame. At each place with $c$ columns the determinant multiplier is $\det(g)^c$; for dual highest weights the two determinant exponents sum to $-c$. Thus the determinant powers cancel in the invariant pairing. This construction applies place by place to their tensor product.

Choose an integral polynomial basis $p=(p_i)$ of the complementary shape and the dual basis of $W'$. Let $G$ be the restricted tensor Gram matrix in that dual basis. Thus $K(Y,X)=p(Y)^tG p(X)$ for the complementary shape. Put
\[
 H_{ij}=p_i(\partial)p_j(0),\qquad
 f(X)=K(I,X),\qquad h=f(\partial)f(0).
\]
Then $H$ is a symmetric integral matrix. We claim that $HG=hI$. To see this, give the complex polynomial space the Fischer Hermitian form in which monomials $X^\alpha$ are orthogonal with squared norm $\alpha!=\prod_i\alpha_i!$. This form is invariant under unitary row substitutions: multiplication by a linear variable is adjoint to its conjugate differential operator, and unitary substitution preserves those adjunctions and the norm of $1$. In consequence the metric on $P_{\mathbb C}$ with matrix $H$ and the dual of the restricted tensor metric with matrix $G^{-1}$ are both invariant. The span of the frame tensors is irreducible: the tensor of highest wedges has the sum of the factor highest weights, and that weight space is one-dimensional. Complete reducibility and highest-weight theory therefore place the generating highest tensor in a single irreducible summand. Its orbit spans the frame tensors. This argument also applies to the product over places. Schur's lemma gives $H=cG^{-1}$. At the unit frame $v=p(I)$, the tensor $e(I)$ has squared norm $v^tGv=1$. The coefficient vector of $f$ is $Gv$, so
\[
 h=(Gv)^t H(Gv)=c.
\]
In particular $h$ is a positive integer. For the endpoint pairing take $A=GJ$. The integral matrix $C=J^{-1}H$ satisfies $CA=hI$. For $v\in W_{\mathbb C}$ the functional $u\mapsto u^tAv$ equals the Hermitian pairing with $\overline{Jv}$ in $W'$, which has norm $\|v\|$. This proves both the inverse and the isometry in the specified normalization.

There is also a dimension-free bound for the integer just obtained. For $D$ boxes in the complementary shape, expanding $f$ gives at most $\prod_s s!\le D!$ signed monomials, counted with multiplicity, each of total degree $D$. Each pair contributes at most $D!$ in absolute value to its Fischer product. Consequently $h\le(D!)^3$. When $D\le br^4$, $r\le3x^5$, and $L=\log b\le x^{1/100}$, this gives $\log h\le3br^4\log(br^4)\le x^{22}b$ for all sufficiently large $x$. Hence any period denominator with $\log N\le x^{99}b^2$ satisfies $hN\exp(x^{200}b^2)\le\exp(x^{201}b^2)$.

\subsection{A uniform integer clearing the final path}
\label{pathmath:denom}
The following reduction makes the uniformity in the twist explicit. Let $Q$ be the product of its auxiliary places, $q_Q$ the product of their absolute norms, and $t$ a finite idele generating the grid modulus. At the endpoint write $u=\overline{t\alpha}/(t\alpha)$, where $t$, $q_1\alpha$, $\alpha^{-1}$, and $q_Q M^b/t$ are integral finite adeles. These four integrality conditions are the hypotheses on the grid modulus and its fixed fractional generator used below. With $A_0=M^b q_1$ one has
\[
 q_QA_0u=(q_Q M^b/t)\,\overline{tq_1\alpha}\,\alpha^{-1}
 \quad\text{and}\quad q_Q A_0u^{-1}=\overline{q_QA_0u}.
\]
Both right sides are integral. Away from $q_Q$, $A_0u$ and $A_0u^{-1}$ are therefore integral at every embedding. In particular, either sign of an integral determinant exponent $l$ costs at most the common denominator $A_0^{m|l|}$ at the block point $\mathrm{diag}(g,u^m)$.

For precision the remaining local input is the following. In the gauges given by the translated Weyl lattices, each one-edge kernel on the last path carries an integral vector into $d_0^{-1}$ times the integral lattice for every prime away from $q_Q$, with
\[
 d_0=(M^b)^{2f(2n_m)^2}(M^b d)^{br^4}.
\]
Here $d$ is the original quadratic prime. Assume the middle branching matrix is integral and equivariant with one-dimensional exponents $l_\tau$ satisfying $|l_\tau|\le br^4+\lceil20b\rceil$, and the normalized path period expands into a finite sum of these branch pairings with algebraic-integer coefficients. This last condition includes its class-volume normalization. Define
\[
 N_0=(M^bq_1)^{mf(br^4+\lceil20b\rceil)}d_0^{2m}.
\]
For integral starting values, induction on the $m$ edges in each half makes each endpoint vector integral after multiplication by $d_0^m$ in its translated lattice. The middle pairing of the translated vectors is integral since its matrix is integral. Returning to the original frames costs only the determinant characters of $u^m$; the formula above clears their product. Bilinearity gives the factor $d_0^{2m}$, and summation preserves integrality. Thus this one $N_0$, independent of $Q$, clears all such periods away from $q_Q$.

On the admissible scales $m\le2\sqrt{x}$, $n_m=k+m(5r+1)\le x^6$, $r\le3x^5$, $q_1,f\le x^{12}$, $d\le\log x$ eventually for fixed parameters, and $\log M\le22L$, $L\le x^{1/100}$. Substitution gives $\log(M^bq_1),\log(M^bd)\le 24b x^{1/100}$ for sufficiently large $x$. It follows that $\log d_0\le x^{26}b^2$ and $mf(br^4+\lceil20b\rceil)\log(M^bq_1)\le x^{35}b^2$, hence $\log N_0\le x^{36}b^2$. This suffices for the $x^{99}b^2$ budget. The normalized endpoint Weyl duality above is unconditional in the indicated Weyl model. Applying the denominator formula to the arithmetic path still requires the one-edge integral estimate, the stated branching transformation and the finite integral expansion. It also leaves the construction of the surjection onto principal forms, the period identity through the turns, and the archimedean period bound to be established.

\paragraph{The denominator of one matrix indicator.}
Here is the local calculation behind $d_0$. Fix a prime $p\nmid q_Q$, write $N=M^b$, and let $a,a'$ be the ranks at the ends of an edge. Set $s=aa'$ and $\mathcal A=(\mathcal O_E/\mathfrak m)^s$, where $Nq_Q\in\mathfrak m$. In the Schr\"odinger realization at the places of $\mathfrak m$, use the integral vacuum as the characteristic function of the integral lattice (with its chosen vacuum scalar). Projection onto one residue matrix $S$ on this lattice has the finite Fourier expansion
\[
 {\bf1}_{x=S}=\frac1{|\mathcal A|}\sum_{\xi\in\widehat{\mathcal A}}\xi(-S)\xi(x).
\]
Indeed the sum of a nontrivial character over the dual group is zero, and at $x=S$ every term equals one. The order $|\mathcal A|$ divides $(Nq_Q)^{2fs}$ because $\mathcal O_E/(Nq_Q)$ surjects onto $\mathcal O_E/\mathfrak m$ and $[E:\mathbb Q]=2f$. Thus $N^{2fs}/|\mathcal A|$ is $p$-integral, as are the roots of unity in the numerators.

For clarity, the local model needed to interpret this equality is the following. Use the self-dual finite additive pairing for the standard vacuum. In its split Schr\"odinger coordinates a character of the residue variables acts by multiplication on the vacuum, equivalently by a Heisenberg operator in the opposite isotropic direction. In coordinates with generator $t$ for $\mathfrak m$ its parameter is $\overline{t^{-1}z}$ on the conjugate support, for an integral lift $z$ of the residue vector, extended by zero elsewhere. The character identification uses the primitive additive character for this modulus; choosing the opposite sign replaces $z$ by $-z$. Multiplication by $N$ makes every such parameter integral above $p$, since $Nq_Q/t$ is integral and $q_Q$ is a $p$-adic unit. Consequently the edge indicator multiplied by $N^{2fs}$ is a $p$-integral combination of integral-vacuum translates at translation level $N$. This calculation accounts for the vacuum scalar once in the indicator; inserting a second copy alters the normalization.

The geometric arithmetic input can now be stated on individual vacuum translates. For any such translate $v$, denote by $F_v$ its normalized Taylor series at the basic CM period scale, before multiplication of the variables by $Nd$. Require
\[
 \partial^\alpha[F_v(Nd\,Z)]_{Z=0}\in\mathcal O_p
 \qquad (|\alpha|=D')
\]
for the degrees $D'$ of the occurring Schur entries, where $\mathcal O_p$ is the valuation ring at the chosen embedding. Require also covariance in the translated Weyl gauges: after changing source and target to integral frames above $p$, their change on standard data preserves the $\mathcal O_p$-span just described. The usual rational change of frame gives this reduction whenever weak approximation on both unitary groups, rational equivariance of the jet kernel, and preservation of this standard-data span by local integral elements hold. These are separate local-model inputs.

In these frames the entry $P(\partial)F_v(cZ)|_{Z=0}$ for an integral homogeneous Schur polynomial $P$ of degree $D'$ costs the denominator $(Nd)^{D'}$ if $c$ is $p$-integral. This follows by multiplying by $(Nd)^{D'}$, using homogeneity, and expanding $P$ in integral monomials in the derivatives. The same assertion holds for integral linear combinations of translates. The sup norm of the output of an integral matrix on a vector is bounded by the sup norm of that vector, by the ultrametric inequality. Together with the Fourier calculation this gives the one-edge denominator $N^{2fs}(Nd)^{D'}$. Since $s\le(2n_m)^2$ and $D'\le br^4$, this divides $d_0$. The result reduces the asserted one-edge bound to scaled CM derivatives of vacuum translates and their local covariance; the elementary Fourier expansion alone does not establish those arithmetic properties.

\subsection{The Selmer datum for a main critical line}\label{moment:line}

Fix a main pair $(\nu,z)$ at stage $1\leq j\leq m$, and put $t_0=2z-1$.
Write $p_0,p_1$ for the two distinct level primes and $f=[E:\mathbb Q]/2$.
The tame character of $\nu$ at each chosen place is the $j$th power of
the prescribed residue character of order $q_2=(q_0d)^{s+6}$.
There is one chosen place from each conjugate pair above each $p_h$,
and $p_h\equiv 1\pmod {q_2^2}$. These are the level and main-pair
conditions of the construction. In particular $p_h$ splits completely in $E$,
whose conductor divides $q_1d\mid q_2$. Enlarge the coefficient field
at $\ell$ to contain the character values and all labelled embeddings.
Let $\mathcal O$ be its integer ring, $\pi$ a uniformizer, and $T$ the
stable line of \S\ref{rankone:main}, with character $\tau$.
Its weights at a pair of complex embeddings are $z-1-c_v$ and $z-1+c_v$,
and $\tau^c\tau=\epsilon_\ell^{1-t_0}$.
The only places at which the line can ramify are those above
$\ell p_0p_1$.

\paragraph{Local rank at $\ell$.}
For $v\mid\ell$, let $Z_v$ be the continuous $T$-valued cocycles whose
rational cohomology class belongs to the Bloch--Kato group
$H^1_f(E_v,V)$, $V=T[1/\ell]$. This specifies the integral condition
in G.2 by imposing, at every torsion level, the homogeneous evaluation
relations valid on $Z_v$. Coboundaries lie in $Z_v$. The module is
torsion-free and finite over $\mathcal O$: the coboundaries are an
image of $T$ and $H^1(E_v,T)$ is finitely generated by local Galois
cohomology. Rational cocycles on the compact decomposition group
have bounded image, so tensoring $Z_v$ with the coefficient field
recovers all rational cocycles with finite class. The character is
potentially crystalline, hence de Rham. The Bloch--Kato dimension formula
(\href{https://virtualmath1.stanford.edu/~conrad/BSDseminar/refs/BKTamagawa.pdf}{\emph{L-functions and Tamagawa numbers of motives}, Corollary 3.8.4}~\cite{bk})
gives $\dim H^1_f=\dim H^0+\dim(D_{\rm dR}/\mathrm{Fil}^0)$.
Adding the coboundary dimension $1-\dim H^0$ gives
\[
  \operatorname{rank}_{\mathcal O} Z_v
  =1+\#\{\sigma\mid v : \operatorname{HT}_{\sigma}(V)<0\}.
\]
Since the two weights at a conjugate pair sum to $2z-2\geq0$, at most
one is negative. Summing yields the bound $\#\{v\mid\ell\}+f$
required in G.2.

\paragraph{An inertia scalar and the old local terms.}
Write $Q=q_2$. At a chosen $w\nmid\ell$ over $p_h$ the value $\xi$
of $\tau$ on a tame generator has order $Q/\gcd(Q,j)$ by the
rank-one formula under local reciprocity. Inversion of a generator or
character preserves that order. Choose $a\in\{2,5\}$ different from
$\ell$. Since $a\mid q_0$, we have $a^s\mid Q$ and
$a^s\ge 2^s=m^{22}>2m$ on the constructed scales. If the residual
order of $\xi$ divided $2$, then $Q\mid 2j\ell^b$ for some $b$,
which would imply $a^s\mid 2j$. Thus $\xi^2-1$ and $\xi-1$ are
units. Inertia at $w$ has trivial cyclotomic character, so
$\tau(\sigma^{-c})=\tau(\sigma)=\xi$ for the chosen generator $\sigma$;
here $\sigma^{-c}=(c\sigma c^{-1})^{-1}$. This also handles the
conjugate place. At least one of $p_0,p_1$ differs from $\ell$,
so such a $\sigma$ always exists. The matrix algebra condition in
rank one holds since the identity is in the image.

For $W=T/\pi^u T$ the local Tate dual has no invariants at any such
$w$: inertia acts on it by $\xi^{-1}$ on a generator with $\xi-1$
a unit. Local duality and the Euler characteristic at $w\nmid\ell$
then give $\operatorname{length} Z^1(E_w,W)=\operatorname{length}W$.
Thus the restricted old local cocycles of G.2 have zero positive excess.
The two determinant losses for releasing the first auxiliary prime also
vanish if its Frobenius and inverse conjugate have the prescribed
residual inertia values. This is an open congruence condition in G.2.

\paragraph{The global term.}
The conjugate Tate dual of $W$ has character
$\epsilon_\ell(\tau^c)^{-1}=\tau\epsilon_\ell^{t_0}$.
As an $\mathcal O$-module the Tate dual is free of rank one over
$\mathcal O/\pi^u$: the trace-dual fractional ideal of $\mathcal O$
is principal. Choose $p_h\ne\ell$. For an arithmetic Frobenius lift
at a degree-one place over $p_h$ the cyclotomic value is $p_h$;
ramification of $T$ there does not affect the displayed character
quotient. Every equivariant map from $W$ to its conjugate Tate dual
therefore has image killed by $p_h^{t_0}-1$. Its value on a basis vector
determines the map. Uniformly in $u$, the length of the space of
equivariant maps, divided by $v_\pi(\ell)$, is at most
$v_\ell(p_h^{t_0}-1)$. The group-order budget covers this integer:
$t_0=2z-1\le 2r-1\le r t$ for $t=5r$, and the factor
$p_h^{t_0}-1$ occurs in the order of the chosen
$\mathrm{GL}_{rt}(\mathbb F_{p_h})$.
We have thus verified the rank, scalar, and length hypotheses in G.2
for the explicit untwisted line with the stated local conditions.
The residue-field size condition in G.1 can be met by an unramified
coefficient extension at each fixed torsion level (its ramification
index is unchanged).

\paragraph{Split Hecke action on the lift of the constant.}
Put $N=2z$ and take a split $w\nmid\ell$ away from the level
and the conductor of a tame twist $\eta$. Write $q=\#k_w$.
Use the split Schr\"odinger model with local vacuum
$\mathbf1_{\mathcal O_w^{N\times1}}$. For a splitting character $a$
of rank $d$ its Frobenius factor is
\[
 c_d(a)=q^{d/2}a(\varpi_w/\bar\varpi_w)\upsilon(\varpi_w)^{-d}
\]
with $\varpi_w$ interpreted as a local uniformizer idele, and
$\upsilon$ the unitary splitting Hecke character. Write
$\alpha=c_1(a_W)$ and $\beta=c_N(a_V)$ for the source and target
factors of the pair. Conjugation of the local variable on the source
side corresponds to the adjoint Hecke action.

Normalize the Hecke polynomial in rank $d$ by
$P_d(X)=\sum_i(-X)^i q^{\binom{i}{2}}T_i$ for the unaveraged
double-coset sums. The cosets in $T_i$ are indexed by
$i$-planes $U\subset k_w^d$: one can take a representative
$\gamma_U$ integral with elementary divisors $\varpi_w$ on $i$ axes
and $1$ elsewhere and with reduction kernel $U$.
For a matrix $y$ of width $d$, $y\gamma_U^{-1}$ is integral exactly
when $y$ is integral and $\bar y$ kills $U$. Indeed, in adapted
integral bases the inverse divides the coordinates of $U$ by
$\varpi_w$. If $y\gamma_U^{-1}$ is integral, multiplication by
$\gamma_U$ gives integrality of $y$.

On the target vacuum the transpose coset therefore gives a multiplier
$\alpha^i$ times the indicator of integrality of $y\gamma_U^{-1}$,
where $y$ is the transpose of the local column. On the dual source
vacuum the inverse coset gives the same expression with $y$ the
column and multiplier $\beta^i$. These are the split Weil
change-of-variables formulas: the determinant has valuation $i$,
its unramified splitting-character value is the $i$th power of the
uniformizer value, and the inverse square-root Jacobian contributes
$q^{di/2}$ with $d$ the rank of the opposite space. Thus for integral
$y$ of residual rank $a$ and width $d$, the resulting weighted sum,
with $Z$ incorporating the character and Jacobian, is
\[
 \sum_{i=0}^{d-a}{d-a\brack i}_q(-Z)^i q^{\binom{i}{2}}
 = (Z;q)_{d-a},\qquad (Z;q)_h=\prod_{i=0}^{h-1}(1-q^iZ).
\]
Here the Gaussian binomial counts the $i$-planes in $\ker\bar y$;
the finite $q$-binomial identity follows by expanding the product.
For a nonintegral matrix every term vanishes. The column and its
transpose have the same rank $a\in\{0,1\}$. Consequently
\[
 (Z;q)_{N-a}=(Z;q)_{t_0}(q^{t_0}Z;q)_{1-a}.
\]
On the target take $Z=\alpha X$. The polynomial argument on the
source is then $q^{t_0}\alpha X/\beta$.

For the global passage, suppose the chosen section of the theta
kernel has this vacuum component at $w$, and that the global lift
pairs the source form and the kernel equivariantly under their finite
translations. In the usual theta integral this equivariance follows
by change of integration variable with invariant measure, provided
the kernel and the chosen section have the specified transformation
law. It moves a target Hecke sum onto the vacuum and then moves the
adjoint source sum onto the source form. The constant source on $U(1)$
has Hecke polynomial $1-X$, since the sole rank-one uniformizer
translation fixes it. Hence its lift has reciprocal Hecke polynomial
\[
 (\alpha X;q)_{t_0}\bigl(1-q^{t_0}\alpha\beta^{-1}X\bigr).
\]
Multiplying the global lift by the source level index, as in the
normalization of the manuscript's character lift, preserves this
identity. A nonzero lift is required to read its eigenvalues.

The negative linear coefficient, after the target normalizing twist
by $\alpha^{-1}$ at geometric Frobenius, is
$q^{t_0}/\beta+\sum_{a=0}^{t_0-1}q^a$.
To check the normalization, write $u=\nu\eta$. In even rank the
splitting convention is $a_V=u\,\upsilon_{U(1)}^{z}$, with
$\upsilon_{U(1)}(x/\bar x)=\upsilon(x)^2$. Thus
$\beta=q^{z}u(\varpi_w/\bar\varpi_w)$.
Consequently the first term is $q^{z-1}u(\varpi_w/\bar\varpi_w)^{-1}$,
the uniformizer value of $u(x/\bar x)^{-1}|x|^{1-z}$ in \S\ref{rankone:main}.
The other terms are the values of $\epsilon_\ell^{-a}$ on geometric
Frobenius. If a continuous representation attached to the target
has these Hecke polynomials, on the density-one set of degree-one
unramified primes of $E$ this gives the normalized trace
\[
 \eta^{-1}_{\ell}\tau+\sum_{a=0}^{2z-2}\epsilon_\ell^{-a},
\]
where $\eta_\ell$ is the avatar of $x\mapsto\eta(x/\bar x)$.
Degree-one primes away from ramification are split over $E^+$;
the remaining primes of absolute residue degree at least two have
convergent reciprocal-norm sum. Chebotarev~\cite[Theorem~8.31]{milneANT} and continuity extend this
trace equality to $G_E$, and characteristic-zero character theory
identifies semisimplifications.

The local computation uses the displayed split action. For its
application to the manuscript's lift one must identify the actual
unramified section and the equivariant global pairing with these
normalizations, and supply the attached target representation.
The integral extension classes and generator-wise clearing operator
require separate first-order congruences. For the higher-rank branch
one also needs ramified label compatibility and the full integral
matrix span for a lattice in the main summand.

\paragraph{Source correspondence.}
At revision \texttt{72fead1}, the calculations can be compared with
the
\ifdefined\href\href{https://github.com/lsernau/grh-lean/blob/72fead102f663d5c7a407964a13dc0cf9d977207/grh/src/New/Congruences/Selmer/MainSummands.lean\#L294}{Selmer datum for a main summand}\else Selmer datum for a main summand\fi,
the
\ifdefined\href\href{https://github.com/lsernau/grh-lean/blob/72fead102f663d5c7a407964a13dc0cf9d977207/grh/src/Lib/Weil/Unramified/HeckeTransfer.lean\#L413}{split vacuum Hecke calculation}\else split vacuum Hecke calculation\fi,
the
\ifdefined\href\href{https://github.com/lsernau/grh-lean/blob/72fead102f663d5c7a407964a13dc0cf9d977207/grh/src/New/OneLift/HeckeValues/CharacterLiftNorm.lean\#L215}{specialization to the constant lift}\else specialization to the constant lift\fi,
and the
\ifdefined\href\href{https://github.com/lsernau/grh-lean/blob/72fead102f663d5c7a407964a13dc0cf9d977207/grh/src/New/Congruences/GaloisRepresentations/RowsFromEigensystems.lean\#L340}{normalized Frobenius trace identity}\else normalized Frobenius trace identity\fi.


\section*{Sources and preparation}
Levent Alp\"oge developed the original argument with Claude \cite{alpo}.
Luke Sernau, colleagues at Anthropic, and Claude produced the readable
edition used here. The repository credits Claude with writing the formal
text and accompanying documents, and the Anthropic collaborators with
directing the rewriting \cite{edition}. Jimmie Adriazola's earlier manuscript
reconstructs the original development \cite{ja}. Zhang, Slijep\v{c}evi\'c and
Rubinstein give further accounts \cite{zhang,slijepcevic,rubinstein}.

Kimi assisted with the preliminary reconstruction. OpenAI Codex drafted the
analytic exposition and compared it with the indicated source declarations.
ChatGPT drafted, compared and revised the three sections on short forms,
spreading and the determinant lower bound, and their synthesis. The
short-form exposition also draws on an earlier working draft.

The arithmetic working draft includes additional AI-assisted reconstructions of the
local frame, moment, cohomology, and period arguments. The large-prime bound
in \S\ref{sec:odd-cocycles-reduction} uses the refined Coates--Sinnott theorem~\cite{gp};
\S\ref{spread:normalization} uses the definite unitary mass formula~\cite{ghy}. Hyperlinks to
the formal source document the origin of the construction. The arithmetic
claims of this draft are limited to the deductions from the stated hypotheses.

\Needspace{13\baselineskip}
\printbibliography[heading=bibintoc,title={References}]

\end{document}
